% GENERATED from data/values by sync/generate_latex_catalog.py; do not edit.
\section{Proofs of catalogue values}\label{app:value-proofs}
Each statement and proof in this appendix is generated from the structured database value record; the database is the editable source of truth.

\subsection{Effective Range of Trivial concepts}\label{val:effective_range_trivial_concepts}
\noindent\textbf{Statement.} The Effective Range of Trivial concepts is $n$.
\begin{proof}
The two concepts $\emptyset$ and $[n]$ disagree on every coordinate, so every coordinate lies in the effective range.
\end{proof}

\subsection{Log Shattering of Trivial concepts}\label{val:log_shattering_trivial_concepts}
\noindent\textbf{Statement.} The Log Shattering of Trivial concepts is $\log(n+1)$.
\begin{proof}
The shattered sets are exactly the empty set and the $n$ singleton subsets: a two-point set would require four traces, but the class has only two. Thus there are $n+1$ shattered sets.
\end{proof}

\subsection{Membership and Proper Equivalence Queries Complexity of Trivial concepts}\label{val:membership_and_proper_equivalence_queries_complexity_trivial_concepts}
\noindent\textbf{Statement.} The Membership and Proper Equivalence Queries Complexity of Trivial concepts is $1$.
\begin{proof}
One membership query at any coordinate distinguishes the two constant concepts. At least one query is needed for a two-concept class.
\end{proof}

\subsection{Membership Query Complexity of Trivial concepts}\label{val:membership_query_complexity_trivial_concepts}
\noindent\textbf{Statement.} The Membership Query Complexity of Trivial concepts is $1$.
\begin{proof}
One membership query at any coordinate distinguishes the empty and full concepts.
\end{proof}

\subsection{Projected Membership Queries Complexity of Trivial concepts}\label{val:projected_membership_queries_complexity_trivial_concepts}
\noindent\textbf{Statement.} The Projected Membership Queries Complexity of Trivial concepts is $1$.
\begin{proof}
Every nonempty projection remains a pair of complementary constant concepts, and one membership query distinguishes that pair.
\end{proof}

\subsection{Recursive Teaching Dimension of Addressing}\label{val:recursive_teaching_dimension_addressing}
\noindent\textbf{Statement.} The Recursive Teaching Dimension of Addressing is $1\quad(n\ge2)$.
\begin{proof}
Every addressing concept contains a unique address coordinate in $[n]$: its positive label there is inconsistent with every other addressing concept.  Hence all concepts have one-point teaching sets in the initial class and can be removed in a single recursive-teaching layer.  A nontrivial class cannot have recursive teaching dimension zero, so the value is exactly one.
\end{proof}
\noindent\textbf{Reference.} \cite{maass1992lower}.

\subsection{Recursive Teaching Dimension of Tagged Singletons}\label{val:recursive_teaching_dimension_tagged_singletons}
\noindent\textbf{Statement.} The Recursive Teaching Dimension of Tagged Singletons is $1$.
\begin{proof}
In the first layer, every tagged pair $\{i,n+f(i)\}$ is uniquely taught by its positive ordinary coordinate $i$.  After removing those pairs, every tag singleton is uniquely taught by its positive tag coordinate; the empty concept is then alone.  Thus there is a recursive teaching plan of width one.  Since the initial class is nontrivial, width zero is impossible, so $\mathrm{RTD}=1$.
\end{proof}
\noindent\textbf{Reference.} \cite{maass1992lower}.

\subsection{Recursive Teaching Dimension of Half-intervals}\label{val:recursive_teaching_dimension_halfintervals}
\noindent\textbf{Statement.} The Recursive Teaching Dimension of Half-intervals is $1$.
\begin{proof}
The half-intervals form a chain. At every recursive stage, an endpoint of the remaining chain is identified by one labelled coordinate, so the minimum teaching-set size is one.
\end{proof}

\subsection{Recursive Teaching Dimension of Full Cube}\label{val:recursive_teaching_dimension_full_cube}
\noindent\textbf{Statement.} The Recursive Teaching Dimension of Full Cube is $n$.
\begin{proof}
In any nontrivial subfamily of the cube, identifying a chosen concept requires fixing every coordinate on which another remaining concept can differ. For the full cube this is $n$, and this bound is attained at the initial stage.
\end{proof}

\subsection{Recursive Teaching Dimension of Singletons plus empty set}\label{val:recursive_teaching_dimension_singletons_plus_empty_set}
\noindent\textbf{Statement.} The Recursive Teaching Dimension of Singletons plus empty set is $1$.
\begin{proof}
Each singleton can be removed using its unique positive coordinate as a one-point teaching set; the empty concept is then the only remaining concept. Thus every recursive round has minimum teaching-set size one.
\end{proof}

\subsection{Recursive Teaching Dimension of Singletons}\label{val:recursive_teaching_dimension_singletons}
\noindent\textbf{Statement.} The Recursive Teaching Dimension of Singletons is $1$.
\begin{proof}
Each singleton is uniquely specified by one positive example. Removing singleton concepts in any order therefore has teaching-set size one at every nontrivial stage.
\end{proof}

\subsection{Recursive Teaching Dimension of Trivial concepts}\label{val:recursive_teaching_dimension_trivial_concepts}
\noindent\textbf{Statement.} The Recursive Teaching Dimension of Trivial concepts is $1$.
\begin{proof}
Any one coordinate distinguishes $\emptyset$ from $[n]$, so either concept can be removed with a teaching set of size one. A nonempty two-concept class cannot have teaching dimension zero.
\end{proof}

\subsection{Star number of Majority}\label{val:star_number_majority}
\noindent\textbf{Statement.} The Star number of Majority is $n$.
\begin{proof}
Choose a base subset whose size is at least two away from the majority threshold. Flipping any one of its $n$ base coordinates then leaves the majority bit unchanged, giving all $n$ one-coordinate neighbours. No star can use more base coordinates.
\end{proof}
\noindent\textbf{Reference.} \cite{maass1992lower}.

\subsection{Star number of Addressing}\label{val:star_number_addressing}
\noindent\textbf{Statement.} The Star number of Addressing is $n-1$.
\begin{proof}
Project to all primary coordinates except one, say $p$.  The addressing concept whose primary coordinate is $p$ becomes the empty trace, while every other concept becomes its distinct primary singleton; these traces form an $(n-1)$-star.  Conversely, a $d$-star has $d+1$ distinct traces (its centre and leaves), each realized by a distinct original concept.  Since the addressing class has $n$ concepts, $d\le n-1$.  Thus its star number is exactly $n-1$.
\end{proof}
\noindent\textbf{Reference.} \cite{maass1992lower}.

\subsection{Teaching Dimension of Trivial concepts}\label{val:teaching_dimension_trivial_concepts}
\noindent\textbf{Statement.} The Teaching Dimension of Trivial concepts is $1$.
\begin{proof}
At any coordinate, $\emptyset$ has label $0$ and $[n]$ has label $1$, so one example teaches either concept. No empty teaching set distinguishes the two concepts.
\end{proof}

\subsection{Star number of Tagged Singletons}\label{val:star_number_tagged_singletons}
\noindent\textbf{Statement.} The Star number of Tagged Singletons is $n$.
\begin{proof}
Project to all $n$ ordinary coordinates.  The empty concept gives the star centre, and each tagged singleton $\{i,n+f(i)\}$ restricts to the leaf $\{i\}$, producing an $n$-star.  Conversely, in a star witness each selected tag coordinate either contributes its tag singleton as a leaf or prevents every ordinary coordinate carrying that tag from being a one-coordinate leaf.  Taking all ordinary coordinates is optimal, so no star has more than $n$ leaves.
\end{proof}
\noindent\textbf{Reference.} \cite{maass1992lower}.

\subsection{Star number of Full Cube}\label{val:star_number_full_cube}
\noindent\textbf{Statement.} The Star number of Full Cube is $n$.
\begin{proof}
Around any cube vertex, flipping each of the $n$ coordinates gives all $n$ neighbors, hence an $n$-star. No star can use more coordinates than the domain.
\end{proof}

\subsection{Star number of Singletons plus empty set}\label{val:star_number_singletons_plus_empty_set}
\noindent\textbf{Statement.} The Star number of Singletons plus empty set is $n$.
\begin{proof}
The empty concept is adjacent in the 1-inclusion graph to each of the $n$ singleton concepts, giving an $n$-star. There are only $n$ coordinates.
\end{proof}

\subsection{Star number of Singletons}\label{val:star_number_singletons}
\noindent\textbf{Statement.} The Star number of Singletons is $n-1$.
\begin{proof}
Restrict to any $n-1$ coordinates. The singleton outside that set becomes the empty trace, while the remaining singletons form its $n-1$ neighbours, giving an $(n-1)$-star. No restriction can yield a larger star.
\end{proof}

\subsection{Star number of Trivial concepts}\label{val:star_number_trivial_concepts}
\noindent\textbf{Statement.} The Star number of Trivial concepts is $1$.
\begin{proof}
Restrict to any one coordinate: the two complementary constant concepts form a 1-star. No restriction can contain a 2-star because it has only these two traces.
\end{proof}

\subsection{co-VC dimension of Addressing}\label{val:covc_dimension_addressing}
\noindent\textbf{Statement.} The co-VC dimension of Addressing is $n$.
\begin{proof}
On the n private address coordinates, the all-zero labeling is absent, but every labeling differing from it at exactly one coordinate is realized. This is a hollow star of size n. Conversely, a minimally inconsistent sample of size k has, for each of its k deleted examples, a concept consistent with the remaining sample. These k concepts are distinct: a concept witnessing two deletions would satisfy the full inconsistent sample. Hence k cannot exceed the n concepts. Thus co-VC is exactly n, also at n=1.
\end{proof}
\noindent\textbf{Reference.} \cite{maass1992lower}.

\subsection{co-VC dimension of Tagged Singletons}\label{val:covc_dimension_tagged_singletons}
\noindent\textbf{Statement.} The co-VC dimension of Tagged Singletons is $2$.
\begin{proof}
For an ordinary coordinate $i$ and its tag $t=n+f(i)$, the trace $\{i\}$ is absent on $\{i,t\}$, while its two one-coordinate neighbours $\emptyset$ and $\{i,t\}$ are present.  This is a hollow two-star.  A missing singleton trace can have only its own tag as a doubleton neighbour; a missing doubleton trace cannot have a three-coordinate neighbour because every class member has size at most two; and the empty trace is always present.  Hence no hollow three-star occurs, so $\mathrm{coVC}=2$.
\end{proof}
\noindent\textbf{Reference.} \cite{maass1992lower}.

\subsection{co-VC dimension of Half-intervals}\label{val:covc_dimension_halfintervals}
\noindent\textbf{Statement.} The co-VC dimension of Half-intervals is $2$.
\begin{proof}
On two ordered coordinates $a<b$, the restricted traces include the empty trace and the full two-point trace while omitting the trace containing only $b$, yielding a hollow 2-star. A chain cannot realize a hollow star of larger dimension.
\end{proof}

\subsection{co-VC dimension of Full Cube}\label{val:covc_dimension_full_cube}
\noindent\textbf{Statement.} The co-VC dimension of Full Cube is $0$.
\begin{proof}
Every restriction of the full cube is itself a full cube, so it has no missing trace. Hence it contains no hollow star of positive dimension.
\end{proof}

\subsection{co-VC dimension of Singletons plus empty set}\label{val:covc_dimension_singletons_plus_empty_set}
\noindent\textbf{Statement.} The co-VC dimension of Singletons plus empty set is $2$.
\begin{proof}
On any two coordinates, the restrictions contain the empty set and both singletons but omit their doubleton, a hollow 2-star. A larger hollow star would require a missing set whose one-coordinate neighbours are all in the class, which is impossible because the class contains only sets of size at most one.
\end{proof}

\subsection{co-VC dimension of Singletons}\label{val:covc_dimension_singletons}
\noindent\textbf{Statement.} The co-VC dimension of Singletons is $n$.
\begin{proof}
On the full domain, the empty trace is missing while all $n$ of its one-coordinate neighbours, the singleton concepts, are present. This is a hollow $n$-star.
\end{proof}

\subsection{co-VC dimension of Trivial concepts}\label{val:covc_dimension_trivial_concepts}
\noindent\textbf{Statement.} The co-VC dimension of Trivial concepts is $2$.
\begin{proof}
On any two coordinates, the two traces are the empty and full two-point traces. Either of the other two traces is missing while both of its one-coordinate neighbours are present, giving a hollow 2-star. No larger hollow star is possible with only two traces.
\end{proof}

\subsection{Largest Strongly Shattered Set of Addressing}\label{val:largest_strongly_shattered_set_addressing}
\noindent\textbf{Statement.} The Largest Strongly Shattered Set of Addressing is $0$.
\begin{proof}
Two distinct addressed concepts have different base singleton coordinates and different address words, hence differ on at least two coordinates. Thus the 1-inclusion graph has no edge, so no nonempty set is strongly shattered.
\end{proof}
\noindent\textbf{Reference.} \cite{maass1992lower}.

\subsection{Largest Strongly Shattered Set of Half-intervals}\label{val:largest_strongly_shattered_set_halfintervals}
\noindent\textbf{Statement.} The Largest Strongly Shattered Set of Half-intervals is $1$.
\begin{proof}
Every coordinate except the first is strongly shattered by two consecutive initial segments. No two-dimensional cube occurs in a chain of initial segments.
\end{proof}

\subsection{Star number of Half-intervals}\label{val:star_number_halfintervals}
\noindent\textbf{Statement.} The Star number of Half-intervals is $2$.
\begin{proof}
The 1-inclusion graph of the chain is a path: an internal initial segment has its two consecutive neighbours, and no concept has more.
\end{proof}

\subsection{Projected Recursive Teaching Dimension of Trivial concepts}\label{val:monotone_recursive_teaching_dimension_trivial_concepts}
\noindent\textbf{Statement.} The Projected Recursive Teaching Dimension of Trivial concepts is $1$.
\begin{proof}
Every nontrivial projection is again a pair of complementary constant concepts and has recursive teaching dimension one; trivial projections have value zero.
\end{proof}

\subsection{Minimum Teaching Set Size of Trivial concepts}\label{val:minimum_teaching_set_size_trivial_concepts}
\noindent\textbf{Statement.} The Minimum Teaching Set Size of Trivial concepts is $1$.
\begin{proof}
A single coordinate distinguishes either concept, while the empty sample is consistent with both concepts.
\end{proof}

\subsection{Membership and Equivalence Queries Complexity of Trivial concepts}\label{val:membership_and_equivalence_queries_complexity_trivial_concepts}
\noindent\textbf{Statement.} The Membership and Equivalence Queries Complexity of Trivial concepts is $1$.
\begin{proof}
A single membership query distinguishes the two concepts, so allowing equivalence queries cannot require more.
\end{proof}

\subsection{Oriented Diameter of Singletons plus empty set}\label{val:oriented_diameter_singletons_plus_empty_set}
\noindent\textbf{Statement.} The Oriented Diameter of Singletons plus empty set is $1$.
\begin{proof}
The empty set and any singleton form an oriented pair on one coordinate. No oriented pair can differ on two coordinates, since the class contains no nonempty set of size two.
\end{proof}

\subsection{Oriented Diameter of Singletons}\label{val:oriented_diameter_singletons}
\noindent\textbf{Statement.} The Oriented Diameter of Singletons is $0$ for $n=1$; $1$ for $n\ge2$.
\begin{proof}
For $n=1$ the sole concept has no coordinate with both labels, so the oriented diameter is zero. For $n\ge2$, projecting to one coordinate gives both its empty and singleton traces, so the value is at least one. A trace has at most one positive coordinate, excluding a larger empty/full pair.
\end{proof}

\subsection{Maximum Degree of Tagged Singletons}\label{val:maximum_degree_tagged_singletons}
\noindent\textbf{Statement.} The Maximum Degree of Tagged Singletons is $\left\lfloor\frac{n}{2}\right\rfloor+\lceil\log n\rceil+2$.
\begin{proof}
The $1$-inclusion graph has an edge from $\emptyset$ to every tag singleton, and an edge from the tag singleton $\{n+\ell\}$ to each tagged pair carrying tag $\ell$.  For $\ell=1$, the definition of $f$ assigns exactly $\lfloor n/2\rfloor+\lceil\log n\rceil+1$ ordinary coordinates to that tag.  Its degree is therefore $\lfloor n/2\rfloor+\lceil\log n\rceil+2$.  Later tag groups are smaller, every tagged pair has degree one, and the empty concept has only $\lceil\log n\rceil$ neighbours, proving this is the maximum.
\end{proof}
\noindent\textbf{Reference.} \cite{maass1992lower}.

\subsection{Maximum Degree of Addressing}\label{val:maximum_degree_addressing}
\noindent\textbf{Statement.} The Maximum Degree of Addressing is $0$.
\begin{proof}
Distinct concepts have different addressed base coordinates and different address words, so their symmetric difference has size at least two. Therefore the 1-inclusion graph has no edge and maximum degree zero.
\end{proof}
\noindent\textbf{Reference.} \cite{maass1992lower}.

\subsection{Maximum Degree of Majority}\label{val:maximum_degree_majority}
\noindent\textbf{Statement.} The Maximum Degree of Majority is $n$.
\begin{proof}
At a base subset sufficiently far from the majority threshold, toggling any of the $n$ base coordinates leaves the majority bit fixed, so the corresponding concepts are 1-inclusion neighbours. The domain has only $n+1$ coordinates, and flipping the final majority coordinate alone never stays in the graph, giving maximum degree $n$.
\end{proof}
\noindent\textbf{Reference.} \cite{maass1992lower}.

\subsection{Maximum Degree of Half-intervals}\label{val:maximum_degree_halfintervals}
\noindent\textbf{Statement.} The Maximum Degree of Half-intervals is $2$.
\begin{proof}
Consecutive initial segments differ on exactly one coordinate, so the 1-inclusion graph is a path and its maximum degree is two.
\end{proof}

\subsection{Maximum Degree of Singletons plus empty set}\label{val:maximum_degree_singletons_plus_empty_set}
\noindent\textbf{Statement.} The Maximum Degree of Singletons plus empty set is $n$.
\begin{proof}
In the 1-inclusion graph, the empty concept is adjacent to every one of the $n$ singletons and no other vertex has greater degree.
\end{proof}

\subsection{Maximum Degree of Full Cube}\label{val:maximum_degree_full_cube}
\noindent\textbf{Statement.} The Maximum Degree of Full Cube is $n$.
\begin{proof}
Every vertex of the $n$-cube has one $1$-inclusion neighbor for each coordinate.
\end{proof}

\subsection{Maximum Degree of Singletons}\label{val:maximum_degree_singletons}
\noindent\textbf{Statement.} The Maximum Degree of Singletons is $0$.
\begin{proof}
Distinct singleton concepts differ on exactly two coordinates, so the $1$-inclusion graph has no edge.
\end{proof}

\subsection{Maximum Degree of Trivial concepts}\label{val:maximum_degree_trivial_concepts}
\noindent\textbf{Statement.} The Maximum Degree of Trivial concepts is $0$.
\begin{proof}
The two vertices $\emptyset$ and $[n]$ differ in $n$ coordinates, so the $1$-inclusion graph has no edge (for $n>1$).
\end{proof}

\subsection{Log Shattering of Singletons}\label{val:log_shattering_singletons}
\noindent\textbf{Statement.} The Log Shattering of Singletons is $0$ for $n=1$; $\log_2(n+1)$ for $n\ge2$.
\begin{proof}
For $n=1$, only the empty coordinate set is shattered, so log shattering is zero. For $n\ge2$, each coordinate takes both labels and is shattered. No pair is shattered because no concept has two positive coordinates. Thus there are exactly $n+1$ shattered sets. A proper pair projection also has the zero trace, but still lacks the doubleton.
\end{proof}

\subsection{Order Sample Compression of Singletons}\label{val:order_sample_compression_singletons}
\noindent\textbf{Statement.} The Order Sample Compression of Singletons is $1$.
\begin{proof}
The order-compression construction for VC-dimension-one classes gives an order sample compression scheme of size one.  Size zero is impossible for a class with two distinct concepts, since an empty compression set has one fixed reconstruction and cannot be consistent with samples distinguishing those concepts.  Hence the value is exactly one.
\end{proof}
\noindent\textbf{Reference.} \cite{darnstadt2016order}.

\subsection{Proper Ordered Sample Compression of Singletons}\label{val:proper_ordered_sample_compression_singletons}
\noindent\textbf{Statement.} The Proper Ordered Sample Compression of Singletons is $n-1$.
\begin{proof}
Darnstädt--Doliwa--Simon--Zilles order the singleton hypotheses.  For the sample containing negative labels on all but one coordinate, the remaining singleton is the first compatible proper hypothesis only after all $n-1$ negatives are retained; any proper subsample instead reconstructs a later singleton inconsistent with one omitted negative label.  Thus every proper order scheme needs $n-1$.  Keeping those negatives gives the matching scheme.
\end{proof}
\noindent\textbf{Reference.} \cite{darnstadt2016order}.

\subsection{Effective VC Radius of Singletons}\label{val:effective_vc_radius_singletons}
\noindent\textbf{Statement.} The Effective VC Radius of Singletons is $1$.
\begin{proof}
The effective range is all of $[n]$. Every singleton is shattered, while no two-point subset is shattered.
\end{proof}

\subsection{Yang dimension of Singletons}\label{val:yang_dimension_singletons}
\noindent\textbf{Statement.} The Yang dimension of Singletons is $n-1$.
\begin{proof}
For the delta-function class, distinct nonempty subfamilies have distinct least containing cubes, so the Taylor resolution of Yang's canonical ideal is minimal. Its top syzygy has homological degree $n-1$, which is therefore the Yang dimension.
\end{proof}
\noindent\textbf{Reference.} \cite{yang2018homological}; see the full arXiv version, Section "Delta Functions", which constructs the minimal $(n-1)$-simplex resolution of the delta-function class..

\subsection{Partial Equivalence Queries Complexity of Singletons}\label{val:partial_equivalence_queries_complexity_singletons}
\noindent\textbf{Statement.} The Partial Equivalence Queries Complexity of Singletons is $1$.
\begin{proof}
Ask the partial-equivalence query that labels every coordinate $0$.  Every singleton target disagrees at its unique positive coordinate, which the counterexample reveals; one query identifies the target.
\end{proof}

\subsection{Membership and Equivalence Queries Complexity of Singletons}\label{val:membership_and_equivalence_queries_complexity_singletons}
\noindent\textbf{Statement.} The Membership and Equivalence Queries Complexity of Singletons is $1$.
\begin{proof}
The improper equivalence query for the all-zero labeling is allowed.  Its counterexample against a singleton target is that target's unique positive coordinate, so one query identifies it.
\end{proof}

\subsection{Minimum Teaching Set Size of Singletons}\label{val:minimum_teaching_set_size_singletons}
\noindent\textbf{Statement.} The Minimum Teaching Set Size of Singletons is $1$.
\begin{proof}
For the target $\{i\}$, the positive label at coordinate $i$ is consistent with no other singleton. An empty sample is consistent with every singleton.
\end{proof}

\subsection{Projected Recursive Teaching Dimension of Singletons}\label{val:monotone_recursive_teaching_dimension_singletons}
\noindent\textbf{Statement.} The Projected Recursive Teaching Dimension of Singletons is $1$.
\begin{proof}
After any projection, remove the remaining singleton concepts first, each with one positive example; the possible empty concept then remains alone. Thus every projection has recursive teaching dimension at most one.
\end{proof}

\subsection{Teaching Dimension of Singletons}\label{val:teaching_dimension_singletons}
\noindent\textbf{Statement.} The Teaching Dimension of Singletons is $1$.
\begin{proof}
The positive label at the unique point of a singleton identifies it. No target is identified by the empty sample.
\end{proof}

\subsection{Membership and Proper Equivalence Queries Complexity of Singletons}\label{val:membership_and_proper_equivalence_queries_complexity_singletons}
\noindent\textbf{Statement.} The Membership and Proper Equivalence Queries Complexity of Singletons is $n-1$.
\begin{proof}
Query membership at $n-1$ distinct coordinates.  A positive answer identifies its singleton target, and if all answers are negative the only remaining singleton is identified.  Conversely, an adversary can answer $0$ to any membership query, eliminating only that singleton.  Against a proper equivalence query proposing $\{i\}$, it returns $i$ whenever possible, again eliminating only $\{i\}$.  Thus no allowed query removes more than one candidate in the worst case, and $n-1$ queries are necessary and sufficient.
\end{proof}
\noindent\textbf{Reference.} \cite{angluin1988queries}.

\subsection{Projected Membership Queries Complexity of Singletons}\label{val:projected_membership_queries_complexity_singletons}
\noindent\textbf{Statement.} The Projected Membership Queries Complexity of Singletons is $n-1$.
\begin{proof}
A projection to $k<n$ coordinates is the class consisting of the empty trace and its $k$ singleton traces, so identifying the empty trace requires $k$ negative membership answers; this is also sufficient.  The full projection is the original singleton class and needs $n-1$ queries.  Maximizing over projections therefore gives $n-1$.
\end{proof}
\noindent\textbf{Reference.} \cite{angluin1988queries}.

\subsection{Membership Query Complexity of Singletons}\label{val:membership_query_complexity_singletons}
\noindent\textbf{Statement.} The Membership Query Complexity of Singletons is $n-1$.
\begin{proof}
Query coordinates until a positive answer occurs.  If the first $n-1$ answers are negative, the remaining coordinate is forced to be the target singleton, so $n-1$ queries suffice.  Conversely, after fewer than $n-1$ negative answers at least two singleton concepts remain consistent.
\end{proof}

\subsection{Order Sample Compression of Trivial concepts}\label{val:order_sample_compression_trivial_concepts}
\noindent\textbf{Statement.} The Order Sample Compression of Trivial concepts is $1$.
\begin{proof}
The order-compression construction for VC-dimension-one classes gives size at most one.  A sample containing one labeled coordinate distinguishes the two constants, while size zero has only one reconstruction and cannot represent both labels.  Thus the order sample compression value is one.
\end{proof}
\noindent\textbf{Reference.} \cite{darnstadt2016order}.

\subsection{Proper Ordered Sample Compression of Trivial concepts}\label{val:proper_ordered_sample_compression_trivial_concepts}
\noindent\textbf{Statement.} The Proper Ordered Sample Compression of Trivial concepts is $1$.
\begin{proof}
Order the two constant hypotheses.  A one-point labeled sample selects the least compatible constant and is itself a teaching set for that choice; the decompression is therefore proper.  As no size-zero proper scheme can distinguish the two constants, the proper ordered sample compression value is exactly one.
\end{proof}
\noindent\textbf{Reference.} \cite{darnstadt2016order}.

\subsection{Projected VC Radius of Trivial concepts}\label{val:projected_vc_radius_trivial_concepts}
\noindent\textbf{Statement.} The Projected VC Radius of Trivial concepts is $1$.
\begin{proof}
Every nonempty projection is again the two complementary constant concepts, which shatter one point but no two points.
\end{proof}

\subsection{Effective VC Radius of Trivial concepts}\label{val:effective_vc_radius_trivial_concepts}
\noindent\textbf{Statement.} The Effective VC Radius of Trivial concepts is $1$.
\begin{proof}
The effective range is all of $[n]$. Every singleton is shattered, whereas no two-point subset is shattered because only two traces are available.
\end{proof}

\subsection{Yang dimension of Trivial concepts}\label{val:yang_dimension_trivial_concepts}
\noindent\textbf{Statement.} The Yang dimension of Trivial concepts is $1$.
\begin{proof}
The canonical suboplex of the two constant functions is the disjoint union of the simplex labelled all $0$ and the simplex labelled all $1$.  It has nonzero reduced homology only in dimension $0$, and Yang's homological dimension is one plus the highest nonzero homology degree.  Therefore $\mathrm{Y}=1$.
\end{proof}
\noindent\textbf{Reference.} \cite{yang2018homological}.

\subsection{Partial Equivalence Queries Complexity of Trivial concepts}\label{val:partial_equivalence_queries_complexity_trivial_concepts}
\noindent\textbf{Statement.} The Partial Equivalence Queries Complexity of Trivial concepts is $1$.
\begin{proof}
A one-point partial query is a membership query and distinguishes the two constant concepts.
\end{proof}

\subsection{Largest Strongly Shattered Set of Full Cube}\label{val:largest_strongly_shattered_set_full_cube}
\noindent\textbf{Statement.} The Largest Strongly Shattered Set of Full Cube is $n$.
\begin{proof}
The full domain is strongly shattered: the whole $n$-cube occurs in the class. No shattered set can have more than $n$ coordinates.
\end{proof}

\subsection{Largest Strongly Shattered Set of Singletons plus empty set}\label{val:largest_strongly_shattered_set_singletons_plus_empty_set}
\noindent\textbf{Statement.} The Largest Strongly Shattered Set of Singletons plus empty set is $1$.
\begin{proof}
Any singleton coordinate is strongly shattered by the empty set and that singleton. A strongly shattered set of size two would require a two-dimensional cube, hence a set of size two, which is absent.
\end{proof}

\subsection{Largest Strongly Shattered Set of Singletons}\label{val:largest_strongly_shattered_set_singletons}
\noindent\textbf{Statement.} The Largest Strongly Shattered Set of Singletons is $0$.
\begin{proof}
A nonempty strongly shattered set would give a $1$-inclusion edge, but the singleton family has none. The empty set is strongly shattered.
\end{proof}

\subsection{Largest Strongly Shattered Set of Trivial concepts}\label{val:largest_strongly_shattered_set_trivial_concepts}
\noindent\textbf{Statement.} The Largest Strongly Shattered Set of Trivial concepts is $0$.
\begin{proof}
A nonempty strongly shattered set would yield a $1$-inclusion edge. The class has none, while the empty set is strongly shattered.
\end{proof}

\subsection{Threshold Dimension of Singletons}\label{val:threshold_dimension_singletons}
\noindent\textbf{Statement.} The Threshold Dimension of Singletons is $0$ for $n=1$; $1$ for $n\ge2$.
\begin{proof}
For $n=1$ no coordinate has both labels, so the threshold dimension is zero. For $n\ge2$ a one-coordinate projection contains both the empty and full traces. Any threshold chain on two or more coordinates would include a doubleton trace, impossible for singleton concepts. Thus the value is one.
\end{proof}

\subsection{Size of Addressing}\label{val:size_addressing}
\noindent\textbf{Statement.} The Size of Addressing is $n$.
\begin{proof}
There is exactly one concept for each index $i\in[n]$. Distinct indices give distinct private positive address coordinates, so these n concepts are distinct, irrespective of their bit encoding.
\end{proof}
\noindent\textbf{Reference.} \cite{maass1992lower}.

\subsection{Size of Majority}\label{val:size_majority}
\noindent\textbf{Statement.} The Size of Majority is $2^n$.
\begin{proof}
Every subset of the first $n$ coordinates determines exactly one concept: the last coordinate is then fixed by whether that subset has strict majority. Thus the class has $2^n$ concepts.
\end{proof}
\noindent\textbf{Reference.} \cite{maass1992lower}.

\subsection{Size of Tagged Singletons}\label{val:size_tagged_singletons}
\noindent\textbf{Statement.} The Size of Tagged Singletons is $n+\lceil\log n\rceil +1$.
\begin{proof}
The definition lists one empty concept, $\lceil\log n\rceil$ tag singleton concepts, and one tagged two-point concept for each of the $n$ base coordinates. These concepts are distinct.
\end{proof}
\noindent\textbf{Reference.} \cite{maass1992lower}.

\subsection{Size of Half-intervals}\label{val:size_halfintervals}
\noindent\textbf{Statement.} The Size of Half-intervals is $n$.
\begin{proof}
There is one initial segment for each of the $n$ possible right endpoints.
\end{proof}

\subsection{Size of Singletons plus empty set}\label{val:size_singletons_plus_empty_set}
\noindent\textbf{Statement.} The Size of Singletons plus empty set is $n+1$.
\begin{proof}
The class consists of the $n$ singleton subsets of $[n]$ together with the empty subset.
\end{proof}

\subsection{Size of Full Cube}\label{val:size_full_cube}
\noindent\textbf{Statement.} The Size of Full Cube is $2^n$.
\begin{proof}
Every subset of the $n$-element domain is a concept, giving $2^n$ concepts.
\end{proof}

\subsection{Size of Singletons}\label{val:size_singletons}
\noindent\textbf{Statement.} The Size of Singletons is $n$.
\begin{proof}
There is one singleton concept for each of the $n$ coordinates.
\end{proof}

\subsection{Path Dimension of Trivial concepts}\label{val:path_dimension_trivial_concepts}
\noindent\textbf{Statement.} The Path Dimension of Trivial concepts is $1$.
\begin{proof}
Restricting to any one coordinate yields the empty and full one-point traces, a path of length one. A path of length two would require an intermediate trace that is not present.
\end{proof}

\subsection{Threshold Dimension of Trivial concepts}\label{val:threshold_dimension_trivial_concepts}
\noindent\textbf{Statement.} The Threshold Dimension of Trivial concepts is $1$.
\begin{proof}
On any one coordinate the complementary constant traces form a threshold chain of length one. No two-coordinate threshold chain can be formed from only the two constant traces.
\end{proof}

\subsection{Hamming Radius of Trivial concepts}\label{val:hamming_radius_trivial_concepts}
\noindent\textbf{Statement.} The Hamming Radius of Trivial concepts is $\left\lceil \frac{n}{2}\right\rceil$.
\begin{proof}
For a center $T\subseteq[n]$, the two distances are $|T|$ and $n-|T|$. Their maximum is minimized by a balanced $T$, giving $\lceil n/2\rceil$.
\end{proof}

\subsection{Diameter of Trivial concepts}\label{val:diameter_trivial_concepts}
\noindent\textbf{Statement.} The Diameter of Trivial concepts is $n$.
\begin{proof}
The only distinct pair is $\emptyset,[n]$, whose symmetric difference has size $n$.
\end{proof}

\subsection{Oriented Diameter of Trivial concepts}\label{val:oriented_diameter_trivial_concepts}
\noindent\textbf{Statement.} The Oriented Diameter of Trivial concepts is $n$.
\begin{proof}
Restricting to all $n$ coordinates contains the pair $\emptyset,[n]$, which realizes the defining oriented pair on a set of size $n$.
\end{proof}

\subsection{Size of Trivial concepts}\label{val:size_trivial_concepts}
\noindent\textbf{Statement.} The Size of Trivial concepts is $2$.
\begin{proof}
By definition the class consists exactly of $\emptyset$ and $[n]$.
\end{proof}

\subsection{VC Dimension of Tagged Singletons}\label{val:vc_dimension_tagged_singletons}
\noindent\textbf{Statement.} The VC Dimension of Tagged Singletons is $1$.
\begin{proof}
Every coordinate has both labels among the class, so VC dimension is at least one.  No two ordinary coordinates can both be positive in one concept; an ordinary coordinate can be positive only together with its unique tag; and no two tag coordinates are simultaneously positive.  Thus every two-coordinate projection omits at least one of the four patterns, so VC dimension is exactly one.
\end{proof}
\noindent\textbf{Reference.} \cite{maass1992lower}.

\subsection{Hamming Radius of Tagged Singletons}\label{val:hamming_radius_tagged_singletons}
\noindent\textbf{Statement.} The Hamming Radius of Tagged Singletons is $2$.
\begin{proof}
The empty set is a Hamming-ball center of radius two because every concept has size at most two.  Conversely, a radius-one center must be empty or a singleton in order to be within one of $\emptyset$.  Among the at least two tag singleton concepts, one is then at distance two from that center.  Thus the minimum enclosing radius is exactly two.
\end{proof}
\noindent\textbf{Reference.} \cite{maass1992lower}.

\subsection{Diameter of Tagged Singletons}\label{val:diameter_tagged_singletons}
\noindent\textbf{Statement.} The Diameter of Tagged Singletons is $4$.
\begin{proof}
Every concept has support of size at most two, so the symmetric difference of two concepts has size at most four.  Two tagged singleton concepts with distinct ordinary coordinates and distinct tags are disjoint two-element sets, and hence have symmetric difference four.  Thus the diameter is exactly four.
\end{proof}
\noindent\textbf{Reference.} \cite{maass1992lower}.

\subsection{Oriented Diameter of Tagged Singletons}\label{val:oriented_diameter_tagged_singletons}
\noindent\textbf{Statement.} The Oriented Diameter of Tagged Singletons is $2$.
\begin{proof}
For an ordinary coordinate $i$ and its tag $n+f(i)$, the restrictions of $\emptyset$ and $\{i,n+f(i)\}$ give the oriented pair $\emptyset,\{i,n+f(i)\}$ on two coordinates.  Conversely, the full trace in an oriented pair must be the restriction of a class member containing all of the selected coordinates; every member has size at most two.  Therefore the oriented diameter is exactly two.
\end{proof}
\noindent\textbf{Reference.} \cite{maass1992lower}.

\subsection{Effective Range of Tagged Singletons}\label{val:effective_range_tagged_singletons}
\noindent\textbf{Statement.} The Effective Range of Tagged Singletons is $n+\lceil \log n \rceil$.
\begin{proof}
For each ordinary coordinate $i\in[n]$, the tagged singleton containing $i$ differs from $\emptyset$ at $i$.  For each tag coordinate $n+\ell$, the singleton $\{n+\ell\}$ differs from $\emptyset$ there.  Hence the effective range is the full domain of $n+\lceil\log n\rceil$ coordinates.
\end{proof}
\noindent\textbf{Reference.} \cite{maass1992lower}.

\subsection{Effective Range of Addressing}\label{val:effective_range_addressing}
\noindent\textbf{Statement.} The Effective Range of Addressing is $\begin{cases}0&n=1,\\ n+\lceil\log_2 n\rceil&n\ge2.\end{cases}$.
\begin{proof}
Put $r=\lceil\log_2n\rceil$. For $n\ge2$, each of the n private coordinates takes value one at its own target and zero at the others. Every bit coordinate also varies: if one were constant, the r-bit encoding could realize at most $2^{r-1}<n$ distinct words. Thus all n+r coordinates are effective. At n=1 there is one concept and no bit coordinate, so no coordinate varies.
\end{proof}
\noindent\textbf{Reference.} \cite{maass1992lower}.

\subsection{Effective Range of Majority}\label{val:effective_range_majority}
\noindent\textbf{Statement.} The Effective Range of Majority is $n+1$.
\begin{proof}
Each of the first $n$ coordinates can be toggled between two concepts. The majority coordinate also changes by comparing subsets of sizes $\lfloor n/2\rfloor$ and $\lfloor n/2\rfloor+1$. Hence every one of the $n+1$ coordinates is effective.
\end{proof}
\noindent\textbf{Reference.} \cite{maass1992lower}.

\subsection{Effective Range of Half-intervals}\label{val:effective_range_halfintervals}
\noindent\textbf{Statement.} The Effective Range of Half-intervals is $n-1$.
\begin{proof}
Coordinate 1 belongs to every initial segment. For each coordinate $i>1$, consecutive initial segments $[i-1]$ and $[i]$ disagree at $i$. Hence exactly the $n-1$ coordinates from 2 through $n$ are effective.
\end{proof}

\subsection{Effective Range of Singletons plus empty set}\label{val:effective_range_singletons_plus_empty_set}
\noindent\textbf{Statement.} The Effective Range of Singletons plus empty set is $n$.
\begin{proof}
For every coordinate $i$, the empty concept and the singleton at $i$ disagree at $i$. Hence every coordinate is effective.
\end{proof}

\subsection{Effective Range of Full Cube}\label{val:effective_range_full_cube}
\noindent\textbf{Statement.} The Effective Range of Full Cube is $n$.
\begin{proof}
For every coordinate, the cube contains two concepts that differ exactly on that coordinate, so the effective range is the full domain.
\end{proof}

\subsection{VC Dimension of Singletons plus empty set}\label{val:vc_dimension_singletons_plus_empty_set}
\noindent\textbf{Statement.} The VC Dimension of Singletons plus empty set is $1$.
\begin{proof}
Every one-point set is shattered by the empty concept and its singleton. No two-point set is shattered because the required doubleton restriction is absent.
\end{proof}

\subsection{VC Dimension of Full Cube}\label{val:vc_dimension_full_cube}
\noindent\textbf{Statement.} The VC Dimension of Full Cube is $n$.
\begin{proof}
All $2^n$ labelings of the full domain occur in the class, so it is shattered. The domain has only $n$ coordinates.
\end{proof}

\subsection{VC Dimension of Singletons}\label{val:vc_dimension_singletons}
\noindent\textbf{Statement.} The VC Dimension of Singletons is $0$ for $n=1$ and $1$ for $n\ge2$.
\begin{proof}
For $n=1$, the class consists only of the all-one concept, so no nonempty set is shattered and the VC dimension is zero. For $n\ge2$, every one-point set is shattered: its singleton labels it one, and any other singleton labels it zero. No pair is shattered, since no concept labels both coordinates one. For $n\ge3$, the zero pattern on a pair also occurs, using a singleton outside the pair.
\end{proof}

\subsection{VC Dimension of Trivial concepts}\label{val:vc_dimension_trivial_concepts}
\noindent\textbf{Statement.} The VC Dimension of Trivial concepts is $1$.
\begin{proof}
Every singleton is shattered by the restrictions of $\emptyset$ and $[n]$. No two-point set is shattered because the class has only two concepts.
\end{proof}

\subsection{Effective Range of Singletons}\label{val:effective_range_singletons}
\noindent\textbf{Statement.} The Effective Range of Singletons is $0$ for $n=1$; $n$ for $n\ge2$.
\begin{proof}
For $n=1$ there is only one concept and no varying coordinate. For $n\ge2$, at each coordinate $i$ the concept $\{i\}$ and any other singleton have different labels. Hence every one of the $n$ coordinates belongs to the effective range.
\end{proof}

\subsection{Threshold Dimension of Full Cube}\label{val:threshold_dimension_full_cube}
\noindent\textbf{Statement.} The Threshold Dimension of Full Cube is $n$.
\begin{proof}
For any ordering of the $n$ coordinates, every initial segment belongs to the cube, realizing a threshold chain of length $n$.
\end{proof}

\subsection{Oriented Diameter of Full Cube}\label{val:oriented_diameter_full_cube}
\noindent\textbf{Statement.} The Oriented Diameter of Full Cube is $n$.
\begin{proof}
The full cube contains both $\emptyset$ and $[n]$, realizing the defining oriented pair on all $n$ coordinates.
\end{proof}

\subsection{Threshold Dimension of Singletons plus empty set}\label{val:threshold_dimension_singletons_plus_empty_set}
\noindent\textbf{Statement.} The Threshold Dimension of Singletons plus empty set is $1$.
\begin{proof}
The empty set and any singleton give a threshold chain of length one. A chain of length two would contain a doubleton, which the class does not contain.
\end{proof}

\subsection{Threshold Dimension of Tagged Singletons}\label{val:threshold_dimension_tagged_singletons}
\noindent\textbf{Statement.} The Threshold Dimension of Tagged Singletons is $2$.
\begin{proof}
On $S=\{n+f(i),i\}$, in the order $n+f(i),i$, the restrictions of $\emptyset$, the tag singleton $\{n+f(i)\}$, and the tagged singleton $\{i,n+f(i)\}$ form the three threshold traces.  A threshold witness of dimension three would need a class member whose restriction contains all three selected coordinates, but every class member has size at most two.  Hence the threshold dimension is two.
\end{proof}
\noindent\textbf{Reference.} \cite{maass1992lower}.

\subsection{Positive Recursive Teaching Dimension of the Positive-Teaching Triangle}\label{val:positive_recursive_teaching_dimension_positive_teaching_triangle}
\noindent\textbf{Statement.} The Positive Recursive Teaching Dimension of Positive-Teaching Triangle is $2$.
\begin{proof}
In the full class, every positive one-coordinate sample is shared by another concept, so no target has a positive teaching set of size one. Its two positive coordinates teach each target, giving a first removal of cost two; the remaining two targets then have one-coordinate positive teaching sets. Thus positive RTD is exactly two. The recurrence is verified by the dedicated positive-teaching-triangle checker.
\end{proof}
\noindent\textbf{Reference.} Direct calculation; verified by the dedicated positive-teaching-triangle checker.

\subsection{Preference-Based Teaching Dimension of the Positive-Teaching Triangle}\label{val:preferencebased_teaching_dimension_positive_teaching_triangle}
\noindent\textbf{Statement.} The Preference-Based Teaching Dimension of Positive-Teaching Triangle is $1$.
\begin{proof}
Put 011 first in the preference order, so it is taught by the empty sample. Among the remaining targets, 110 and 101 are distinguished by one coordinate. Thus a preference-based plan of width one exists; width zero is impossible for a non-singleton class. The exhaustive preference-order calculation is in sync/verify\_positive\_teaching\_triangle.py.
\end{proof}
\noindent\textbf{Reference.} Direct calculation; verified by sync/verify\_positive\_teaching\_triangle.py..

\subsection{Recursive Teaching Dimension of the Positive-Teaching Triangle}\label{val:recursive_teaching_dimension_positive_teaching_triangle}
\noindent\textbf{Statement.} The Recursive Teaching Dimension of Positive-Teaching Triangle is $1$.
\begin{proof}
The class admits a recursive teaching plan of width one: first remove 011 using its label 0 at $x_1$, then distinguish 110 and 101 by one coordinate. No non-singleton class has recursive teaching dimension zero. Equivalently, finite-class preference-based teaching dimension is one; sync/verify\_positive\_teaching\_triangle.py exhausts the corresponding preference orders.
\end{proof}
\noindent\textbf{Reference.} Direct calculation; verified by sync/verify\_positive\_teaching\_triangle.py..

\subsection{VC Dimension of the Extremal-Extension Separation Class}\label{val:vc_dimension_extremal_extension_gap}
\noindent\textbf{Statement.} The VC Dimension of Extremal-Extension Separation Class is $2$.
\begin{proof}
A two-coordinate projection is shattered, while no three-coordinate projection has all eight traces. The exhaustive calculation over all coordinate subsets is in sync/verify\_extremal\_extension\_gap.py.
\end{proof}
\noindent\textbf{Reference.} Direct calculation; verified by sync/verify\_extremal\_extension\_gap.py..

\subsection{Extremal-Extension VC Dimension of the Extremal-Extension Separation Class}\label{val:extremal_extension_vc_dimension_extremal_extension_gap}
\noindent\textbf{Statement.} The Extremal-Extension VC Dimension of Extremal-Extension Separation Class is $3$.
\begin{proof}
Enumerate all finite families on the same four-coordinate domain that contain the class. Among the shattering-extremal superfamilies, the least VC dimension is three. The complete same-domain enumeration is in sync/verify\_extremal\_extension\_gap.py.
\end{proof}
\noindent\textbf{Reference.} Direct calculation; verified by sync/verify\_extremal\_extension\_gap.py..

\subsection{Self-Directed Queries Complexity of the Ben-David--Eiron Family}\label{val:selfdirected_queries_complexity_ben_david_eiron_selfdirected}
\noindent\textbf{Statement.} The Self-Directed Queries Complexity of Ben-David--Eiron Fixed-VC Self-Directed Family is $\ge n$.
\begin{proof}
Lemma 1 of Ben-David--Eiron proves recursively that after every queried coordinate and either returned label, the remaining version space contains a copy of $\mathcal C_{n-1}^d$. Hence every self-directed learner can be forced to make at least $n$ mistakes.
\end{proof}
\noindent\textbf{Reference.} \cite{ben1998self}, Lemma 1..

\subsection{Recursive Teaching Dimension of the Ben-David--Eiron Family}\label{val:recursive_teaching_dimension_ben_david_eiron_selfdirected}
\noindent\textbf{Statement.} The Recursive Teaching Dimension of Ben-David--Eiron Fixed-VC Self-Directed Family is $\le 39.3752d^2-3.6330d$.
\begin{proof}
Ben-David--Eiron give $\mathrm{VC}(\mathcal C_n^d)=d$. The explicit quadratic theorem of Hu--Wu--Li--Wang therefore bounds $\mathrm{RTD}(\mathcal C_n^d)$ by $39.3752d^2-3.6330d$, uniformly in $n$. Here $d$ is fixed and $n$ is the family variable.
\end{proof}
\noindent\textbf{Reference.} \cite{ben1998self,hu2017quadratic}..

\subsection{Self-Directed Queries Complexity of Axis-Parallel Boxes}\label{val:selfdirected_queries_complexity_axis_parallel_boxes}
\noindent\textbf{Statement.} The Self-Directed Queries Complexity of Axis-Parallel Boxes is $2$.
\begin{proof}
The learner approaches the target box from two opposite corners of $X_n^d$, predicting negative labels.  The first mistakes identify the two opposite target corners and hence the box.  Conversely, both corners can be forced as mistakes, so the optimum is exactly two.
\end{proof}
\noindent\textbf{Reference.} Theorem 4 of \cite{goldman1994power}..

\subsection{Partial Equivalence Queries Complexity of Axis-Parallel Boxes}\label{val:partial_equivalence_queries_complexity_axis_parallel_boxes}
\noindent\textbf{Statement.} The Partial Equivalence Queries Complexity of Axis-Parallel Boxes is $\Theta(d\log n)$.
\begin{proof}
For every fixed dimension $d\ge1$, the partial-equivalence-query complexity of $\mathrm{BOX}_n^d$ is bounded above and below by constant multiples of $d\log n$.  In particular it diverges with $n$ whereas the self-directed complexity remains two.
\end{proof}
\noindent\textbf{Reference.} \cite{maass1992lower,goldman1994power}; the latter explicitly invokes the matching box bound from the former in its proof of the strict model separation..

\subsection{Partial Equivalence Queries Complexity of Addressing}\label{val:partial_equivalence_queries_complexity_addressing}
\noindent\textbf{Statement.} The Partial Equivalence Queries Complexity of Addressing is $1\quad(n\ge2)$.
\begin{proof}
Query the partial hypothesis that labels every address point negative and is undefined on the bit points.  For target concept $i$, its unique positive address point $i$ is the only possible counterexample, so the reply identifies the target in one query.  A nontrivial class cannot be identified with zero queries.
\end{proof}
\noindent\textbf{Reference.} Directly from the addressing-class definition; this is also the $\mathrm{LC}_{\mathrm{partial}}(\mathrm{ADDRESSING}_n)=1$ entry in \cite{maass1992lower}..

\subsection{VC Dimension of Addressing}\label{val:vc_dimension_addressing}
\noindent\textbf{Statement.} The VC Dimension of Addressing is $\lfloor\log n\rfloor\quad(n=2^r,\ r\ge1)$.
\begin{proof}
Restrict here to $n=2^r$ with all r-bit words present, so $\lfloor\log_2n\rfloor=\log_2n=r$. The address coordinates realize every binary pattern, so they form a shattered set of size $\lfloor\log n\rfloor$. Since the class has exactly $2^{\lfloor\log n\rfloor}$ concepts, no larger set can be shattered.
\end{proof}
\noindent\textbf{Reference.} \cite{maass1992lower}.

\subsection{Membership and Equivalence Queries Complexity of Addressing}\label{val:membership_and_equivalence_queries_complexity_addressing}
\noindent\textbf{Statement.} The Membership and Equivalence Queries Complexity of Addressing is $\Theta(\log n)\quad(n=2^r,\ r\ge1)$.
\begin{proof}
Restrict here to $n=2^r$ with all r-bit words present, so $\lfloor\log_2n\rfloor=\log_2n=r$. For the addressing family, Maass--Turán's table of matching bounds gives $\mathrm{LC}_{\mathrm{ARB\text{-}MEMB}}(\mathrm{ADDRESSING}_n)=\Theta(\log n)$, exactly the database's model with arbitrary equivalence and membership queries.
\end{proof}
\noindent\textbf{Reference.} \cite{maass1992lower}, Table 1..

\subsection{Membership and Proper Equivalence Queries Complexity of Addressing}\label{val:membership_and_proper_equivalence_queries_complexity_addressing}
\noindent\textbf{Statement.} The Membership and Proper Equivalence Queries Complexity of Addressing is $\Theta(\log n)$.
\begin{proof}
The $r=\lceil\log_2 n\rceil$ membership queries at the bit points reveal the injective code of the target, so $r$ queries suffice.  For the lower bound, against every failed proper equivalence query the adversary returns the proposed hypothesis's address point, which excludes only that hypothesis.  At any fixed transcript of $m$ membership answers and $p$ such failures, at most $2^m(p+1)\leq2^{m+p}$ targets can have been distinguished.  Thus identifying all $n$ targets requires $m+p\geq\lceil\log_2 n\rceil$ queries.
\end{proof}
\noindent\textbf{Reference.} Direct calculation from the addressing construction of \cite{maass1992lower}.

\subsection{Stable Unlabeled Complexity of Half-Intervals}\label{val:stable_unlabeled_complexity_halfintervals}
\noindent\textbf{Statement.} The Stable Unlabeled Complexity of Half-intervals is $1$.
\begin{proof}
For a realizable sample from $\mathcal I_n$, retain its largest positive point if it has one and reconstruct the corresponding initial segment.  If it has no positive point, retain nothing and reconstruct $[1]$; every negative sample point is then at least $2$, since every member of $\mathcal I_n$ contains $1$.  The scheme is sample-consistent and stable: deleting an unretained example never changes the largest positive point or the fact that there are no positives.  Size zero cannot compress a nonconstant class, so the optimum is one.
\end{proof}
\noindent\textbf{Reference.} Direct calculation for the half-interval class; compare the stable-compression framework in \cite{hanneke2021stable}.

\subsection{Projected Distinguishing Range of Half-Intervals}\label{val:projected_distinguishing_range_halfintervals}
\noindent\textbf{Statement.} The Projected Distinguishing Range of Half-intervals is $n-1$.
\begin{proof}
The full-domain projection is the n-concept chain of initial segments. A set of k queried coordinates induces at most k+1 distinct traces on this chain, so at least n-1 coordinates are needed; querying coordinates $2,\ldots,n$ attains this bound. No projection can have more than n concepts, and every m-concept binary class has distinguishing range at most m-1 by greedily refining its trace partition. Hence the maximum over projections is exactly n-1.
\end{proof}
\noindent\textbf{Reference.} Direct calculation from the half-interval construction.

\subsection{Projected Membership Queries Complexity of Half-Intervals}\label{val:projected_membership_queries_complexity_halfintervals}
\noindent\textbf{Statement.} The Projected Membership Queries Complexity of Half-intervals is $\lceil\log_2 n\rceil$.
\begin{proof}
Every coordinate projection of the half-interval class is again a chain of at most n distinct initial-segment traces, and binary search identifies its target in at most $\lceil\log_2 n\rceil$ membership queries. The full-domain projection is the original n-concept chain, for which every membership-query decision tree has at least n leaves and therefore depth at least $\lceil\log_2 n\rceil$.
\end{proof}
\noindent\textbf{Reference.} Direct calculation; the ordinary half-interval membership-query value is recorded from \cite{angluin1988queries}.

\subsection{Projected Distinguishing Range of Trivial Concepts}\label{val:projected_distinguishing_range_trivial_concepts}
\noindent\textbf{Statement.} The Projected Distinguishing Range of Trivial concepts is $1$.
\begin{proof}
Every nonempty projection of the two trivial concepts consists of the all-zero and all-one traces, which are distinguished by any one retained coordinate. The empty projection has one trace. Thus the maximum distinguishing range over projections is one.
\end{proof}
\noindent\textbf{Reference.} Direct calculation.

\subsection{Projected Distinguishing Range of Full Cube}\label{val:projected_distinguishing_range_full_cube}
\noindent\textbf{Statement.} The Projected Distinguishing Range of Full Cube is $n$.
\begin{proof}
A projection of the n-dimensional full cube onto k coordinates is the k-dimensional full cube. Every one of its k coordinates is needed to distinguish the two traces that differ only there, and all k coordinates suffice. Maximizing over $0\leq k\leq n$ gives n.
\end{proof}
\noindent\textbf{Reference.} Direct calculation.

\subsection{Distinguishing Range of the Distinguishing-Range Projection Gap}\label{val:distinguishing_range_distinguishing_range_projection_gap}
\noindent\textbf{Statement.} The Distinguishing Range of Distinguishing-Range Projection Gap is $2$.
\begin{proof}
Coordinates three and four give the four distinct traces 10, 11, 01, and 00, so two suffice. One binary coordinate yields at most two traces and cannot distinguish four concepts.
\end{proof}
\noindent\textbf{Reference.} Direct calculation.

\subsection{Projected Distinguishing Range of the Distinguishing-Range Projection Gap}\label{val:projected_distinguishing_range_distinguishing_range_projection_gap}
\noindent\textbf{Statement.} The Projected Distinguishing Range of Distinguishing-Range Projection Gap is $3$.
\begin{proof}
Projection to coordinates one, two, and four gives 000, 001, 011, and 100. Every one of its three coordinates is needed: deleting any coordinate merges a pair of these traces. Thus this projection has distinguishing range three. No projection can have distinguishing range above three because the class has four concepts.
\end{proof}
\noindent\textbf{Reference.} Direct calculation.

\subsection{Projected Distinguishing Range of Singletons}\label{val:projected_distinguishing_range_singletons}
\noindent\textbf{Statement.} The Projected Distinguishing Range of Singletons is $n-1$.
\begin{proof}
On the full domain, omitting two coordinates $i,j$ makes the singleton concepts $\{i\}$ and $\{j\}$ have the same empty trace, so every distinguishing coordinate set has size at least $n-1$. Retaining all but one coordinate gives distinct traces for all n singleton concepts, so this bound is attained. A projection to k<n coordinates consists of its k singleton traces together with the empty trace, and every retained coordinate is needed to distinguish its singleton from the empty trace; its distinguishing range is therefore k\leq n-1. The full projection has the already calculated value n-1.
\end{proof}

\subsection{Littlestone Dimension of Addressing}\label{val:littlestone_dimension_addressing}
\noindent\textbf{Statement.} The Littlestone Dimension of Addressing is $\lfloor\log n\rfloor\quad(n=2^r,\ r\ge1)$.
\begin{proof}
Restrict here to $n=2^r$ with all r-bit words present, so $\lfloor\log_2n\rfloor=\log_2n=r$. Querying the $\lfloor\log n\rfloor$ address coordinates in sequence yields a complete mistake tree: every address bit pattern labels one concept. A deeper tree would require more than the $2^{\lfloor\log n\rfloor}$ concepts in the class.
\end{proof}
\noindent\textbf{Reference.} \cite{maass1992lower}.

\subsection{Repetition-Free Recursive Teaching Dimension of Full Cube}\label{val:repetitionfree_recursive_teaching_dimension_full_cube}
\noindent\textbf{Statement.} The Repetition-Free Recursive Teaching Dimension of Full Cube is $n$.
\begin{proof}
In the full cube, whichever concept is taught first must be distinguished from every one-bit neighbour, so its teaching set contains all n coordinates; hence $\mathrm{rfRTD}\ge n$.  Conversely, induct on n: split the cube into its two (n-1)-dimensional fibres at one coordinate.  Teach the 1-fibre first, adding that coordinate to every teaching set in a repetition-free plan for the fibre, then teach the 0-fibre with its repetition-free plan.  The added coordinate distinguishes the two blocks' coordinate sets, and it also separates each first-block concept from every still-present 0-fibre concept.  The order is at most n.  Thus $\mathrm{rfRTD}=n$.
\end{proof}
\noindent\textbf{Reference.} \cite{doliwa2014recursive}.

\subsection{Projected Distinguishing Range of the Eluder--Range Separation Class}\label{val:projected_distinguishing_range_eluder_range_separation_class}
\noindent\textbf{Statement.} The Projected Distinguishing Range of Eluder--Range Separation Class is 5.
\begin{proof}
The displayed class contains pairs differing only at each of the five coordinates: $(01111,11111)$, $(00111,01111)$, $(10000,10100)$, $(10000,10010)$, and $(01000,01001)$ respectively for $x_1,\ldots,x_5$. Hence every distinguishing coordinate set contains all five coordinates, so $\mathrm R(\mathcal D_{\mathrm{ER}})=5$. The full-domain projection is permitted in the maximum defining $\mathrm R^p$, while no projection has more than five coordinates; therefore $\mathrm R^p(\mathcal D_{\mathrm{ER}})=5$.
\end{proof}

\subsection{Combinatorial Eluder Dimension of the Eluder--Range Separation Class}\label{val:combinatorial_eluder_dimension_eluder_range_separation_class}
\noindent\textbf{Statement.} The Combinatorial Eluder Dimension of Eluder--Range Separation Class is 4.
\begin{proof}
With base $01111$, query $x_5,x_4,x_2,x_1$ and use respectively $01000,01001,00111,11111$ as witnesses; each witness agrees with the base at every earlier queried coordinate and disagrees at its own coordinate. Thus $\mathrm{Edim}\geq4$. For the reverse inequality, a length-five witness would query all five coordinates. Directly applying the defining compatibility condition to each of the eight possible bases gives maximum lengths $4,4,4,4,4,4,4,4$ (in the displayed order $00111,01000,01001,01111,10000,10010,10100,11111$): at every compatible partial sequence no fifth unused coordinate has a witness. Hence $\mathrm{Edim}=4$.
\end{proof}

\subsection{Proper Stable Sample Compression of Half-intervals}\label{val:proper_stable_sample_compression_halfintervals}
\noindent\textbf{Statement.} The Proper Stable Sample Compression of Half-intervals is $1$.
\begin{proof}
For a realizable labeled sample of the initial-segment class, if it contains a negative example, retain its least negative coordinate $j$ and reconstruct $[j-1]$; realizability gives $j\geq2$. If it contains no negative example, retain nothing and reconstruct $[n]$. The reconstruction belongs to the half-interval class and is consistent with the whole sample. The selected least negative point, when present, is unchanged on deleting unselected examples; the empty case remains empty under every subsample. Thus this is a proper stable unlabeled scheme of size one. Size zero is impossible for a nontrivial class.
\end{proof}

\subsection{Hitting Size of Common-Core Singletons}\label{val:hitting_size_common_core_singletons}
\noindent\textbf{Statement.} The Hitting Size of Common-Core Singletons is $1$.
\begin{proof}
The common coordinate $0$ belongs to every member of $\mathcal C^{\cap}_n$, so it is a hitting set of size one. The class has nonempty members, so the empty set cannot be a hitting set.
\end{proof}

\subsection{co-VC Dimension of Common-Core Singletons}\label{val:covc_dimension_common_core_singletons}
\noindent\textbf{Statement.} The co-VC dimension of Common-Core Singletons is $n$.
\begin{proof}
Projecting away the common coordinate gives the $n$ singleton class. Its absent all-zero trace has all $n$ one-coordinate flips present, giving a hollow $n$-star and hence $\mathrm{coVC}\geq n$. Conversely every projection of $\mathcal C^{\cap}_n$ has at most $n$ distinct concepts, while a hollow $d$-star has $d$ distinct realized leaves. Therefore no projection has co-VC dimension above $n$.
\end{proof}

\subsection{Largest Strongly Shattered Set of Common-Core Singletons}\label{val:largest_strongly_shattered_set_common_core_singletons}
\noindent\textbf{Statement.} The Largest Strongly Shattered Set of Common-Core Singletons is $0$.
\begin{proof}
Any two distinct members $\{0,i\}$ and $\{0,j\}$ differ in the two private coordinates $i,j$. Thus the one-inclusion graph has no edge, so the class contains no positive-dimensional cube. Its largest strongly shattered set consequently has size zero.
\end{proof}

\subsection{Minimum Degree of Common-Core Singletons}\label{val:minimum_degree_common_core_singletons}
\noindent\textbf{Statement.} The Minimum Degree of Common-Core Singletons is $0$.
\begin{proof}
Distinct concepts $\{0,i\}$ and $\{0,j\}$ differ at both private coordinates $i,j$, so no two concepts are adjacent in the one-inclusion graph. Every vertex has degree zero, hence so does the minimum degree.
\end{proof}

\subsection{Effective VC Radius of Universal Class}\label{val:effective_vc_radius_universal_class}
\noindent\textbf{Statement.} The Effective VC Radius of Universal Class is $1$ for $n\geq2$.
\begin{proof}
For $n\geq2$, the effective coordinates of $\mathcal U_n$ are exactly the nonempty proper subsets $A\subset[n]$: its column is the membership vector of $A$ across $u_1,\ldots,u_n$, so it is nonconstant precisely in that case. Every such singleton coordinate is shattered. But any effective $A$ and its complement have only the two joint patterns $01$ and $10$ across the concepts, so that pair is not shattered. Thus the effective VC radius is one.
\end{proof}

\subsection{Largest Strongly Shattered Set of Universal Class}\label{val:largest_strongly_shattered_set_universal_class}
\noindent\textbf{Statement.} The Largest Strongly Shattered Set of Universal Class is $0$.
\begin{proof}
For distinct $u_i,u_j$, the incidence functions differ on exactly the subsets containing exactly one of $i,j$, of which there are $2^{n-1}$. Thus the one-inclusion graph has no edge, so the class contains no positive-dimensional cube and its largest strongly shattered set has size zero.
\end{proof}

\subsection{Littlestone Dimension of Tagged Singletons}\label{val:littlestone_dimension_tagged_singletons}
\noindent\textbf{Statement.} The Littlestone Dimension of Tagged Singletons is $2$.
\begin{proof}
Choose a tag coordinate $t$.  On the $t=1$ branch, querying an ordinary coordinate assigned tag $t$ separates its tag singleton from its tagged pair; on the $t=0$ branch, querying a distinct tag coordinate separates $\emptyset$ from that tag singleton.  This shatters a depth-two tree.  At an ordinary coordinate, the positive branch is a singleton class, while at a tag coordinate the positive branch is a star and has Littlestone dimension one.  Thus no root has both branches of Littlestone dimension at least two, proving the value is exactly two.
\end{proof}
\noindent\textbf{Reference.} \cite{maass1992lower}.

\subsection{Hitting Size of Universal Class}\label{val:hitting_size_universal_class}
\noindent\textbf{Statement.} The Hitting Size of Universal Class is $1$.
\begin{proof}
The domain point $[n]$ belongs to every universal-class concept $u_i$, hence it is a hitting set of size one. Since each $u_i$ is nonempty, the empty set cannot be a hitting set.
\end{proof}

\subsection{Hitting Size of Common-Core Cube}\label{val:hitting_size_common_core_cube}
\noindent\textbf{Statement.} The Hitting Size of Common-Core Cube is $1$.
\begin{proof}
Every concept contains coordinate $0$, so $\{0\}$ is a hitting set. Since the class has nonempty concepts, no empty hitting set exists.
\end{proof}

\subsection{Largest Strongly Shattered Set of Common-Core Cube}\label{val:largest_strongly_shattered_set_common_core_cube}
\noindent\textbf{Statement.} The Largest Strongly Shattered Set of Common-Core Cube is $n$.
\begin{proof}
Varying $A\subseteq[n]$ realizes every labeling of the $n$ private coordinates while the common coordinate is fixed. Thus the class contains an $n$-cube. No contained cube can have dimension above the $n$ varying coordinates.
\end{proof}

\subsection{Effective VC Radius of Common-Core Cube}\label{val:effective_vc_radius_common_core_cube}
\noindent\textbf{Statement.} The Effective VC Radius of Common-Core Cube is $n$.
\begin{proof}
The common coordinate is constant and hence absent from the effective range; the $n$ private coordinates form the effective range. Every labeling of them occurs by varying $A\subseteq[n]$, so the entire effective range is shattered. Thus the effective VC radius is $n$.
\end{proof}

\subsection{Minimum Degree of Common-Core Cube}\label{val:minimum_degree_common_core_cube}
\noindent\textbf{Statement.} The Minimum Degree of Common-Core Cube is $n$.
\begin{proof}
Two concepts $\{0\}\cup A$ and $\{0\}\cup B$ are one-inclusion neighbours exactly when $A,B$ differ at one private coordinate. Thus the one-inclusion graph is the $n$-dimensional cube, which is $n$-regular.
\end{proof}

\subsection{Shattering-Cover Projection VC Dimension of Common-Core Cube}\label{val:shattering_cover_projection_vc_dimension_common_core_cube}
\noindent\textbf{Statement.} The Shattering-Cover Projection VC Dimension of Common-Core Cube is $n$.
\begin{proof}
Every projection that retains all $2^n$ distinct concepts must retain all $n$ private coordinates, since deleting one identifies the two concepts differing only there. On those coordinates the projection is the full $n$-cube, whose VC dimension is $n$ and whose shattered-set family has size $2^n$. Thus every admissible projection has VC dimension at least $n$, and the private-coordinate projection attains it.
\end{proof}

\subsection{Subset Teaching Dimension of Common-Core Cube}\label{val:subset_teaching_dimension_common_core_cube}
\noindent\textbf{Statement.} The Subset Teaching Dimension of Common-Core Cube is $n$.
\begin{proof}
Deleting the constant common coordinate identifies this class with the full $n$-cube and preserves all minimum teaching samples on the varying coordinates, hence preserves the subset-teaching recursion. The established full-cube calculation therefore gives $\mathrm{STD}=n$.
\end{proof}

\subsection{Log Size Ratio of Universal Class}\label{val:log_size_ratio_universal_class}
\noindent\textbf{Statement.} The Log Size Ratio of Universal Class is $\frac{\log(n-1)}{\log(2^n+1)}$ for $n\geq2$.
\begin{proof}
The universal class has $n$ concepts and domain $2^{[n]}$, of cardinality $2^n$. Substitution into the definition gives $\log(n-1)/\log(2^n+1)$.
\end{proof}

\subsection{Log Size Ratio of Common-Core Cube}\label{val:log_size_ratio_common_core_cube}
\noindent\textbf{Statement.} The Log Size Ratio of Common-Core Cube is $\frac{\log(2^n-1)}{\log(n+2)}$.
\begin{proof}
The common-core cube has $2^n$ concepts on its $n+1$ point domain. Substitution into the definition gives $\log(2^n-1)/\log(n+2)$.
\end{proof}

\subsection{Log Size Ratio of Even-Parity Class}\label{val:log_size_ratio_even_parity_class}
\noindent\textbf{Statement.} The Log Size Ratio of Even-Parity Class is $\frac{\log(2^{n-1}-1)}{\log(n+1)}$ for $n\geq2$.
\begin{proof}
Exactly half of the $2^n$ subsets of $[n]$ have even cardinality, so $|\mathcal E_n|=2^{n-1}$. Substitution into the definition gives the displayed value.
\end{proof}

\subsection{Littlestone Dimension of Half-intervals}\label{val:littlestone_dimension_halfintervals}
\noindent\textbf{Statement.} The Littlestone Dimension of Half-intervals is $\lfloor\log_2 n\rfloor$.
\begin{proof}
A balanced binary-search tree on the ordered endpoints is shattered: at a queried coordinate, the target initial segment lies on exactly one side of that coordinate, and each branch retains an interval of feasible endpoints.  This gives depth $\lfloor\log_2 n\rfloor$.  Conversely, a depth-$d$ shattered Littlestone tree has $2^d$ leaves realized by distinct concepts, while the class has only $n$ initial segments.  Hence $d\le\log_2 n$ and the value is logarithmic.
\end{proof}
\noindent\textbf{Reference.} \cite{littlestone1988learning}.

\subsection{Minimum Degree of Even-Parity Class}\label{val:minimum_degree_even_parity_class}
\noindent\textbf{Statement.} The Minimum Degree of Even-Parity Class is $0$.
\begin{proof}
Flipping one coordinate changes parity. Hence no two even subsets differ in exactly one coordinate, so the $1$-inclusion graph is edgeless and has minimum degree zero.
\end{proof}

\subsection{Largest Strongly Shattered Set of Even-Parity Class}\label{val:largest_strongly_shattered_set_even_parity_class}
\noindent\textbf{Statement.} The Largest Strongly Shattered Set of Even-Parity Class is $0$.
\begin{proof}
A nonempty strongly shattered set contains a one-dimensional subcube, hence an edge of the $1$-inclusion graph. The even-parity class has no such edge because flipping one coordinate changes parity. The empty set is strongly shattered, so the value is zero.
\end{proof}

\subsection{Effective VC Radius of Even-Parity Class}\label{val:effective_vc_radius_even_parity_class}
\noindent\textbf{Statement.} The Effective VC Radius of Even-Parity Class is $n-1$ for $n\geq2$.
\begin{proof}
Every coordinate is effective. For every $(n-1)$-subset of coordinates, either parity choice of the omitted coordinate extends each trace, so that projection is the full $(n-1)$-cube. The full $n$-set is not shattered because only even patterns occur. Hence the effective VC radius is $n-1$.
\end{proof}

\subsection{Effective VC Radius of Cube with Apex}\label{val:effective_vc_radius_cube_with_apex}
\noindent\textbf{Statement.} The Effective VC Radius of Cube with Apex is $1$.
\begin{proof}
Every coordinate is effective and every singleton is shattered. For each $i\in[n]$, the trace on $\{\star,i\}$ realizes $00,01,10$ but not $11$, because the apex is zero at $i$ and every base-cube concept is zero at $\star$. Thus not every effective pair is shattered, so the value is one.
\end{proof}

\subsection{Largest Strongly Shattered Set of Cube with Apex}\label{val:largest_strongly_shattered_set_cube_with_apex}
\noindent\textbf{Statement.} The Largest Strongly Shattered Set of Cube with Apex is $n$.
\begin{proof}
The base subfamily is a full $n$-cube with $\star$ fixed to zero, so $[n]$ is strongly shattered. A strongly shattered $(n+1)$-set would require a full $(n+1)$-cube and hence $2^{n+1}$ concepts, but $\mathcal A_n$ has only $2^n+1$. Thus the value is exactly $n$.
\end{proof}

\subsection{Projected No-Clashing Teaching Dimension of Full Cube}\label{val:projected_noclashing_teaching_dimension_full_cube}
\noindent\textbf{Statement.} The Projected No-Clashing Teaching Dimension of Full Cube is $\lceil n/2\rceil$.
\begin{proof}
Every coordinate projection of the full n-cube is a full m-cube for some \(m\le n\), whose no-clashing teaching dimension is \(\lceil m/2\rceil\) by Theorem 19 of Kirkpatrick--Simon--Zilles. The maximum occurs for the full projection.
\end{proof}
\noindent\textbf{Reference.} \cite{KSZ2019}.

\subsection{Projected No-Clashing Teaching Dimension of Trivial Concepts}\label{val:projected_noclashing_teaching_dimension_trivial_concepts}
\noindent\textbf{Statement.} The Projected No-Clashing Teaching Dimension of Trivial concepts is $1$.
\begin{proof}
Every nonempty coordinate projection remains the two constant concepts and has no-clashing teaching dimension one; the empty projection has one trace and value zero. Thus the maximum projected value is one.
\end{proof}

\subsection{Stable Labeled Sample Compression of Full Cube}\label{val:stable_labeled_sample_compression_full_cube}
\noindent\textbf{Statement.} The Stable Labeled Sample Compression of Full Cube is $\lceil n/2\rceil$.
\begin{proof}
The database-owned proper stable labeled full-cube value proof gives an explicit stable scheme of width $\lceil n/2\rceil$, so the same upper bound holds without imposing properness. The stable-compression theorem gives the lower bound by projected NCTD, which is $\lceil n/2\rceil$ on the full cube. Thus equality holds. The former assertion that labeled compression dominates VC without a constant was invalid; that exact lower bound applies to the unlabeled model.
\end{proof}
\noindent\textbf{Reference.} The explicit block certificate in the proper stable labeled full-cube value proof; survey theorem thm:stable-comp-low-NCTD-star and \cite{KSZ2019}..

\subsection{No-Clashing Teaching Dimension of Quadratic-Residue Tournament Classes}\label{val:noclashing_teaching_dimension_quadratic_residue_tournament_classes}
\noindent\textbf{Statement.} The No-Clashing Teaching Dimension of Quadratic-Residue Tournament Classes is $1$.
\begin{proof}
For any tournament-induced class $\mathcal C(G)=\{C_a\}$, assign $C_a$ the one-point sample on coordinate $a$.  For distinct vertices $a,b$, the tournament edge between them makes one assigned sample inconsistent with the other concept, so the map is no-clashing. A class with more than one concept has no-clashing teaching dimension at least one.
\end{proof}
\noindent\textbf{Reference.} \cite{simon2023tournaments}.

\subsection{Recursive Teaching Dimension of Quadratic-Residue Tournament Classes}\label{val:recursive_teaching_dimension_quadratic_residue_tournament_classes}
\noindent\textbf{Statement.} The Recursive Teaching Dimension of Quadratic-Residue Tournament Classes is $\Theta(\log_2 p)$.
\begin{proof}
Theorem 15 of Simon proves that every quadratic-residue tournament $G_p$ has minimum teaching dimension at least $\lfloor\tfrac12\log_2p-\log_2\log_2p\rfloor$. Recursive teaching dimension is at least the minimum teaching dimension of its initial class, giving the same lower bound. For the matching order upper bound, every remaining subclass has at most $p$ concepts. Repeatedly conditioning on a smaller nonempty fibre of a nonconstant coordinate identifies some concept within at most $\lfloor\log_2p\rfloor$ steps. Use such a teacher at each recursive removal.
\end{proof}
\noindent\textbf{Reference.} \cite{simon2023tournaments}.

\subsection{Littlestone Dimension of Full Cube}\label{val:littlestone_dimension_full_cube}
\noindent\textbf{Statement.} The Littlestone Dimension of Full Cube is $n$.
\begin{proof}
Querying distinct coordinates along every root-to-leaf path gives a shattered tree of depth $n$. A deeper tree would require more than $n$ independent coordinates.
\end{proof}

\subsection{No-Clashing Teaching Dimension of the Projected No-Clashing Teaching Separation Class}\label{val:noclashing_teaching_dimension_projected_noclashing_teaching_gap}
\noindent\textbf{Statement.} The No-Clashing Teaching Dimension of Projected No-Clashing Teaching Separation Class is $1$.
\begin{proof}
Assign respectively to $0000,0001,0010,0011,0100,1101$ the one-example samples $(2,0),(3,1),(3,0),(2,1),(1,1),(0,1)$. Every pair avoids a clash under this map, proving the upper bound; a class with more than one concept cannot have NCTD zero. The calculation is also checked by sync/verify\_projected\_noclashing\_gap.py.
\end{proof}

\subsection{Projected No-Clashing Teaching Dimension of the Projected No-Clashing Teaching Separation Class}\label{val:projected_noclashing_teaching_dimension_projected_noclashing_gap}
\noindent\textbf{Statement.} The Projected No-Clashing Teaching Dimension of Projected No-Clashing Teaching Separation Class is $2$.
\begin{proof}
The trace on coordinates $\{1,2,3\}$ is $\{000,001,010,011,100,101\}$ and has no width-one no-clashing teacher, while it has one of width two. Exhaustive enumeration of all teacher maps also checks that every coordinate trace has NCTD at most two; hence the projected maximum is exactly two. See sync/verify\_projected\_noclashing\_gap.py.
\end{proof}

\subsection{Positive No-Clashing Teaching Dimension of the Projected Positive No-Clashing Teaching Separation Class}\label{val:positive_noclashing_teaching_dimension_projected_positive_noclashing_gap}
\noindent\textbf{Statement.} The Positive No-Clashing Teaching Dimension of Projected Positive No-Clashing Teaching Separation Class is $1$.
\begin{proof}
Assign $001,010,111$ respectively the positive one-example samples $(2,1),(1,1),(0,1)$. Each pair is prevented from clashing by one of these samples. Width zero is impossible for a non-singleton class. See sync/verify\_projected\_positive\_noclashing\_gap.py.
\end{proof}

\subsection{Projected Positive No-Clashing Teaching Dimension of the Projected Positive No-Clashing Teaching Separation Class}\label{val:projected_positive_noclashing_teaching_dimension_projected_positive_noclashing_gap}
\noindent\textbf{Statement.} The Projected Positive No-Clashing Teaching Dimension of Projected Positive No-Clashing Teaching Separation Class is $2$.
\begin{proof}
Projection to coordinates $\{1,2\}$ gives $\{01,10,11\}$, which has positive no-clashing teaching dimension two. Exhaustive enumeration verifies that no projection has a larger value. See sync/verify\_projected\_positive\_noclashing\_gap.py.
\end{proof}

\subsection{Projected No-Clashing Teaching Dimension of the Projected Positive No-Clashing Teaching Separation Class}\label{val:projected_noclashing_teaching_dimension_projected_positive_noclashing_gap}
\noindent\textbf{Statement.} The Projected No-Clashing Teaching Dimension of Projected Positive No-Clashing Teaching Separation Class is $1$.
\begin{proof}
Every coordinate projection has ordinary no-clashing teaching dimension at most one; exhaustive enumeration of all traces verifies this, while the identity trace is non-singleton. See sync/verify\_projected\_positive\_noclashing\_gap.py.
\end{proof}

\subsection{Positive Recursive Teaching Dimension of the Projected Positive No-Clashing Teaching Separation Class}\label{val:positive_recursive_teaching_dimension_projected_positive_noclashing_gap}
\noindent\textbf{Statement.} The Positive Recursive Teaching Dimension of Projected Positive No-Clashing Teaching Separation Class is $1$.
\begin{proof}
Remove $001$, $010$, then $111$, using respectively their unique positive coordinates 2, 1, and 0 among the remaining concepts. Width zero is impossible for a non-singleton class. See sync/verify\_projected\_positive\_noclashing\_gap.py.
\end{proof}

\subsection{Projected Positive Recursive Teaching Dimension of the Projected Positive No-Clashing Teaching Separation Class}\label{val:projected_positive_recursive_teaching_dimension_projected_positive_noclashing_gap}
\noindent\textbf{Statement.} The Projected Positive Recursive Teaching Dimension of Projected Positive No-Clashing Teaching Separation Class is $2$.
\begin{proof}
Projection to coordinates $\{1,2\}$ gives $\{01,10,11\}$, whose positive recursive teaching dimension is two. Exhaustive enumeration verifies that no projection has a larger value. See sync/verify\_projected\_positive\_noclashing\_gap.py.
\end{proof}

\subsection{Projected Positive Recursive Teaching Dimension of Singletons}\label{val:projected_positive_recursive_teaching_dimension_singletons}
\noindent\textbf{Statement.} The Projected Positive Recursive Teaching Dimension of Singletons is $1$.
\begin{proof}
Every nonempty coordinate projection is a class of distinct singleton concepts, possibly together with the empty concept. Remove its singleton concepts using their unique positive coordinates, then remove the empty concept. Thus every projection has positive recursive teaching dimension at most one, and the identity projection is non-singleton.
\end{proof}

\subsection{Maximum Positive Degree of the Projected Positive No-Clashing Teaching Separation Class}\label{val:maximum_positive_degree_projected_positive_noclashing_gap}
\noindent\textbf{Statement.} The Maximum Positive Degree of Projected Positive No-Clashing Teaching Separation Class is $0$.
\begin{proof}
No two distinct concepts in $\{001,010,111\}$ differ on exactly one coordinate, so its one-inclusion graph has no edges and its maximum positive degree is zero. See sync/verify\_projected\_positive\_noclashing\_gap.py.
\end{proof}

\subsection{Projected Maximum Positive Degree of the Projected Positive No-Clashing Teaching Separation Class}\label{val:projected_maximum_positive_degree_projected_positive_noclashing_gap}
\noindent\textbf{Statement.} The Projected Maximum Positive Degree of Projected Positive No-Clashing Teaching Separation Class is $2$.
\begin{proof}
Projection to coordinates $\{1,2\}$ gives $\{01,10,11\}$, where $11$ has two smaller one-coordinate neighbours. Exhaustive enumeration verifies that no projection has larger maximum positive degree. See sync/verify\_projected\_positive\_noclashing\_gap.py.
\end{proof}

\subsection{Littlestone Dimension of Singletons plus empty set}\label{val:littlestone_dimension_singletons_plus_empty_set}
\noindent\textbf{Statement.} The Littlestone Dimension of Singletons plus empty set is $1$.
\begin{proof}
A depth-one mistake tree is realized by the empty concept and any singleton. At any root coordinate, the positive branch contains only that singleton, so no complete depth-two tree is possible.
\end{proof}

\subsection{VC Dimension of the Projected Positive No-Clashing Teaching Separation Class}\label{val:vc_dimension_projected_positive_noclashing_gap}
\noindent\textbf{Statement.} The VC Dimension of Projected Positive No-Clashing Teaching Separation Class is $1$.
\begin{proof}
Every coordinate realizes both labels, so VC dimension is at least one. No pair of coordinates realizes all four binary traces: the three two-coordinate traces have size at most three. Hence the VC dimension is one. See sync/verify\_projected\_positive\_noclashing\_gap.py.
\end{proof}

\subsection{Projected Maximum Positive Degree of Singletons}\label{val:projected_maximum_positive_degree_singletons}
\noindent\textbf{Statement.} The Projected Maximum Positive Degree of Singletons is $1$ for $n\ge2$.
\begin{proof}
Every coordinate projection consists of singleton traces, possibly together with the empty trace. Each singleton has at most one smaller one-coordinate neighbor, namely the empty trace. For $n\ge2$, a one-coordinate projection contains both its singleton and the empty trace, so maximum positive degree one is attained. At $n=1$, every projection has one concept and the projected maximum positive degree is zero.
\end{proof}

\subsection{Projected No-Clashing Teaching Dimension of Singletons}\label{val:projected_noclashing_teaching_dimension_singletons}
\noindent\textbf{Statement.} The Projected No-Clashing Teaching Dimension of Singletons is $1$.
\begin{proof}
Every coordinate projection of the singleton class is a singleton class on the retained coordinates, possibly together with the empty concept after duplicate traces are identified. Map each remaining singleton to its unique positive example and the empty concept to the empty sample. This is no-clashing and has width one; the identity projection gives the matching lower bound.
\end{proof}
\noindent\textbf{Reference.} Direct calculation; compare \cite{KSZ2019} for no-clashing teaching on finite classes..

\subsection{Projected Positive No-Clashing Teaching Dimension of Singletons}\label{val:projected_positive_noclashing_teaching_dimension_singletons}
\noindent\textbf{Statement.} The Projected Positive No-Clashing Teaching Dimension of Singletons is $1$.
\begin{proof}
Every coordinate projection is a singleton class on the retained coordinates, possibly together with the empty concept. Assign each singleton its unique positive example and the empty concept the empty sample. The map is positive and no-clashing with width one, while the identity projection is non-singleton.
\end{proof}
\noindent\textbf{Reference.} Direct calculation; compare \cite{KSZ2019}..

\subsection{Projected Positive No-Clashing Teaching Dimension of Full Cube}\label{val:projected_positive_noclashing_teaching_dimension_full_cube}
\noindent\textbf{Statement.} The Projected Positive No-Clashing Teaching Dimension of Full Cube is $n$.
\begin{proof}
Every coordinate projection of the full n-cube is a full m-cube for some $m\le n$. Theorem 19 gives positive no-clashing teaching dimension $m$ on that trace, and the full projection attains $m=n$.
\end{proof}
\noindent\textbf{Reference.} \cite{KSZ2019}.

\subsection{Projected Positive Recursive Teaching Dimension of Full Cube}\label{val:projected_positive_recursive_teaching_dimension_full_cube}
\noindent\textbf{Statement.} The Projected Positive Recursive Teaching Dimension of Full Cube is $n$.
\begin{proof}
Every coordinate projection of the full n-cube is a full m-cube for some $m\le n$. Positive recursive teaching dimension is $m$ on that trace, and the full projection attains $m=n$.
\end{proof}
\noindent\textbf{Reference.} \cite{fallat2023batch}.

\subsection{Interpolation Degree of Cube with Apex}\label{val:interpolation_degree_cube_with_apex}
\noindent\textbf{Statement.} The Interpolation degree of Cube with Apex is $n$.
\begin{proof}
The full $n$-cube base subfamily forces interpolation degree at least $n$. Conversely, the $2^n$ monomials $x_S$ for $S\subseteq[n]$ interpolate arbitrary functions on the base and are zero at the apex except for the constant monomial. Adding the apex coordinate $x_\star$ separates its value from the base without increasing the maximum degree. Hence degree $n$ suffices and is necessary.
\end{proof}

\subsection{Largest Strongly Shattered Set of Antipodally Punctured Cube}\label{val:largest_strongly_shattered_set_antipodally_punctured_cube}
\noindent\textbf{Statement.} The Largest Strongly Shattered Set of Antipodally Punctured Cube is $n-2$.
\begin{proof}
Fix two coordinates, one to zero and one to one. The resulting $(n-2)$-face avoids both deleted antipodes, so it is a contained cube and strongly shatters an $(n-2)$-set. Every $(n-1)$-face fixes one coordinate and therefore contains one of the two deleted antipodes, so no such face is contained. Thus the value is exactly $n-2$.
\end{proof}

\subsection{Minimum Degree of Antipodally Punctured Cube}\label{val:minimum_degree_antipodally_punctured_cube}
\noindent\textbf{Statement.} The Minimum Degree of Antipodally Punctured Cube is $n-1$.
\begin{proof}
In the full cube every vertex has degree $n$. A remaining vertex adjacent to one deleted antipode loses exactly that neighbor and has degree $n-1$; all other remaining vertices retain degree $n$. Such adjacent vertices exist for $n\geq3$, so the minimum degree is exactly $n-1$.
\end{proof}

\subsection{Minimum Teaching-Set Size of Cube with Apex}\label{val:minimum_teaching_set_size_cube_with_apex}
\noindent\textbf{Statement.} The Minimum Teaching Set Size of Cube with Apex is $1$.
\begin{proof}
The apex is uniquely identified by its positive label at $\star$, because every base-cube concept is zero there. Thus the minimum teaching-set size is at most one. It cannot be zero in a class with more than one concept, so it is exactly one.
\end{proof}

\subsection{Littlestone Dimension of Singletons}\label{val:littlestone_dimension_singletons}
\noindent\textbf{Statement.} The Littlestone Dimension of Singletons is $1$.
\begin{proof}
Any coordinate gives a shattered tree of depth one. Its positive branch contains only one singleton, so no tree of depth two is shattered.
\end{proof}

\subsection{Degeneracy of Cube with Apex}\label{val:degeneracy_cube_with_apex}
\noindent\textbf{Statement.} The Degeneracy of Cube with Apex is $n$.
\begin{proof}
The base concepts induce a full $n$-cube, so degeneracy is at least $n$. Every nonempty base vertex has degree $n$, the empty base vertex has degree $n+1$ because it is also adjacent to the apex, and the apex has degree one. Thus every nontrivial induced subgraph either contains a nonempty base vertex of degree at most $n$ or is supported only on the apex and the empty base vertex, where its minimum degree is at most one. Hence degeneracy is exactly $n$.
\end{proof}

\subsection{Maximum Positive Degree of Addressing}\label{val:maximum_positive_degree_addressing}
\noindent\textbf{Statement.} The Maximum Positive Degree of Addressing is $0$.
\begin{proof}
Two distinct addressing concepts disagree at both of their distinct private address coordinates: the concept indexed by $i$ is one at address point $i$ and zero at address point $j$, while the concept indexed by $j$ has the opposite labels. Thus no two concepts differ on exactly one coordinate, the $1$-inclusion graph is edgeless, and its maximum positive degree is zero.
\end{proof}
\noindent\textbf{Reference.} Directly from the addressing-class definition..

\subsection{Antichain Number of Quadratic-Residue Tournament Classes}\label{val:antichain_number_quadratic_residue_tournament_classes}
\noindent\textbf{Statement.} The Antichain Number of Quadratic-Residue Tournament Classes is $1$.
\begin{proof}
For the tournament class $\mathcal Q_p=\{C_a:a\in\mathbb F_p\}$, assign $C_a$ its consistent singleton sample at coordinate $a$. For distinct $a,b$, the two singleton samples are distinct, hence incomparable, so this is an antichain teacher mapping of width one. Since $p>1$, width zero would assign the same empty sample to two targets and violate the antichain condition. Thus $\mathrm{AN}(\mathcal Q_p)=1$.
\end{proof}
\noindent\textbf{Reference.} Direct calculation from the tournament-class definition in \cite{simon2023tournaments}..

\subsection{Antichain Number of the No-Clashing/Antichain Separation Class}\label{val:antichain_number_noclashing_antichain_gap}
\noindent\textbf{Statement.} The Antichain Number of No-Clashing/Antichain Separation Class is $1$.
\begin{proof}
Assign respectively to $000,001,010,011,100,101$ the singleton samples $(x_0,0),(x_1,0),(x_1,1),(x_2,1),(x_2,0),(x_0,1)$. These six samples are distinct and have equal size, hence form an antichain. Width zero is impossible for a non-singleton class. Thus $\mathrm{AN}=1$; the assignment and lower bound are checked in $\texttt{sync/verify\_noclashing\_antichain\_gap.py}$.
\end{proof}
\noindent\textbf{Reference.} Direct calculation from the class definition..

\subsection{No-Clashing Teaching Dimension of the No-Clashing/Antichain Separation Class}\label{val:noclashing_teaching_dimension_noclashing_antichain_gap}
\noindent\textbf{Statement.} The No-Clashing Teaching Dimension of No-Clashing/Antichain Separation Class is $2$.
\begin{proof}
The samples $\varnothing,(x_2,1),(x_1,1),\{(x_1,1),(x_2,1)\},(x_0,1),\{(x_0,1),(x_2,1)\}$, assigned in the class-definition order, form a no-clashing teacher map of width two. Exhaustive enumeration rules out width one, so $\mathrm{NCTD}=2$; see $\texttt{sync/verify\_noclashing\_antichain\_gap.py}$.
\end{proof}
\noindent\textbf{Reference.} Direct calculation from the class definition..

\subsection{Dual Littlestone Dimension of Half-intervals}\label{val:dual_littlestone_dimension_halfintervals}
\noindent\textbf{Statement.} The Dual Littlestone Dimension of Half-intervals is $\lfloor\log_2 n\rfloor$.
\begin{proof}
The dual consists of the $n$ terminal segments $\{i,\ldots,n\}$, $1\le i\le n$. Querying an index at a median endpoint splits the remaining $m$ consecutive possible endpoints into groups of sizes $\lfloor m/2\rfloor$ and $\lceil m/2\rceil$. Recursing gives a complete tree of depth $\lfloor\log_2 n\rfloor$, including depth zero for $n=1$. Conversely a depth-$d$ shattered tree requires $2^d$ distinct terminal segments, so $d\le\lfloor\log_2 n\rfloor$.
\end{proof}
\noindent\textbf{Reference.} Direct dualization of the standard threshold-tree calculation in \cite{littlestone1988learning}..

\subsection{Dual Littlestone Dimension of Full Cube}\label{val:dual_littlestone_dimension_full_cube}
\noindent\textbf{Statement.} The Dual Littlestone Dimension of Full Cube is $\lfloor\log_2 n\rfloor$.
\begin{proof}
The n coordinate functions in the dual class realize every binary pattern on a selected collection of d concepts provided $2^d\le n$, so dual VC dimension gives $\mathrm L^*\ge\lfloor\log_2n\rfloor$. Conversely a depth-d shattered tree needs $2^d$ distinct coordinate functions, and there are only n. Hence equality holds.
\end{proof}
\noindent\textbf{Reference.} Direct calculation from the full-cube definition and the Littlestone tree definition..

\subsection{Dual Littlestone Dimension of Universal Class}\label{val:dual_littlestone_dimension_universal_class}
\noindent\textbf{Statement.} The Dual Littlestone Dimension of Universal Class is $n$.
\begin{proof}
The dual of the universal class is the n-dimensional full cube. Its Littlestone dimension is n, witnessed by querying the n hypothesis coordinates in a fixed order.
\end{proof}
\noindent\textbf{Reference.} Direct transposition of the universal-class incidence matrix..

\subsection{Dual Littlestone Dimension of Singletons}\label{val:dual_littlestone_dimension_singletons}
\noindent\textbf{Statement.} The Dual Littlestone Dimension of Singletons is $0$ for $n=1$; $1$ for $n\ge2$.
\begin{proof}
The singleton class is self-dual. For $n=1$ it has one concept and dimension zero. For $n\ge2$ a coordinate has both labels, giving depth one. At every possible root coordinate, however, the label-one branch contains exactly one concept and cannot split again. Hence no complete depth-two Littlestone tree is shattered.
\end{proof}
\noindent\textbf{Reference.} Direct calculation from the singleton-class definition..

\subsection{VC Dimension of the Cube–Half-interval Product}\label{val:vc_dimension_cube_halfinterval_product}
\noindent\textbf{Statement.} The VC Dimension of Cube–Half-interval Product is $m+1$.
\begin{proof}
All $m$ cube coordinates together with the second half-interval coordinate are shattered. Conversely, no shattered set contains two half-interval coordinates: their positive sets are nested and miss a label pair. It contains at most $m$ cube coordinates, giving VC exactly $m+1$.
\end{proof}
\noindent\textbf{Reference.} Direct finite product calculation; independent exact regression in sync/verify\_cube\_halfinterval\_trees.py..

\subsection{Littlestone Dimension of Trivial concepts}\label{val:littlestone_dimension_trivial_concepts}
\noindent\textbf{Statement.} The Littlestone Dimension of Trivial concepts is $1$.
\begin{proof}
Either coordinate gives a shattered tree of depth one. After observing either label, only one of the two constant concepts remains, so no depth-two tree is shattered.
\end{proof}

\subsection{Littlestone Dimension of the Cube–Half-interval Product}\label{val:littlestone_dimension_cube_halfinterval_product}
\noindent\textbf{Statement.} The Littlestone Dimension of Cube–Half-interval Product is $m+\lfloor\log_2 n\rfloor$.
\begin{proof}
There are $2^m n$ concepts, so distinct realizers of the leaves of a depth-$d$ complete shattered tree imply $d\le m+\lfloor\log_2n\rfloor$. Query the $m$ cube coordinates in a fixed order, then use a balanced threshold-search tree of depth $\lfloor\log_2n\rfloor$ on the $n$ ordered endpoints. Independence of the factors realizes every branch and attains the bound.
\end{proof}
\noindent\textbf{Reference.} Direct finite product calculation; independent exact regression in sync/verify\_cube\_halfinterval\_trees.py..

\subsection{2-Block Littlestone Dimension of the Cube–Half-interval Product}\label{val:2block_littlestone_dimension_cube_halfinterval_product}
\noindent\textbf{Statement.} The 2-Block Littlestone Dimension of Cube–Half-interval Product is $m+\min\{\lfloor\log_2 n\rfloor,m+1\}$.
\begin{proof}
Specialize the general $k$-block calculation for this same family (value record 1082) to $k=2$: its argument applies to every integer $k\ge2$, even though that catalogue parameter is indexed by $k>2$.
\end{proof}
\noindent\textbf{Reference.} Direct finite product calculation; independent exact regression in sync/verify\_cube\_halfinterval\_trees.py..

\subsection{$k$-Block Littlestone Dimension of the Cube–Half-interval Product}\label{val:kblock_littlestone_dimension_kgt2_cube_halfinterval_product}
\noindent\textbf{Statement.} The $k$-Block Littlestone Dimension ($k>2$) of Cube–Half-interval Product is $m+\min\{\lfloor\log_2 n\rfloor,\lfloor m/(k-1)\rfloor+1\}$.
\begin{proof}
Write $\ell=\lfloor\log_2 n\rfloor$ and $a=\lfloor m/(k-1)\rfloor$. In any nonadaptive block, at most one queried coordinate can come from the half-interval factor: two distinct such coordinates have nested positive sets and cannot realize all four label pairs. Coordinates cannot repeat along a shattered path. A depth $d=qk+r$, $0\le r<k$, therefore consumes at least $q(k-1)+\max\{r-1,0\}$ distinct cube coordinates on each path. Hence $q\le a$ and $d\le m+a+1$ (if $r=0$, even $d\le m+a$). Also the $2^d$ distinct leaf realizers give $d\le\lfloor\log_2(2^m n)\rfloor=m+\ell$. For the reverse bound put $t=\min\{\ell,a+1\}\ge1$. Use $t-1$ full blocks, each querying $k-1$ unused cube coordinates followed by one query of a balanced threshold search. Then query all remaining cube coordinates, and finally the $t$th threshold-search query; the intervening tail may contain cube-only full blocks. There is at most one threshold query in any block, and its choice depends only on earlier threshold answers, which occur in preceding blocks. Cube coordinates are predetermined. A balanced threshold tree of depth $t\le\ell$ exists on $n$ ordered endpoints, and every combination with the $m$ independent cube labels is realized. This is a valid shattered $k$-block tree of depth $m+t$, including its possibly truncated final block.
\end{proof}
\noindent\textbf{Reference.} Direct finite product calculation; independent exact regression in sync/verify\_cube\_halfinterval\_trees.py..

\subsection{Prefix-$k$ Littlestone Dimension of the Cube–Half-interval Product}\label{val:prefix_littlestone_dimension_cube_halfinterval_product}
\noindent\textbf{Statement.} The Prefix-$k$ Littlestone Dimension ($k>1$) of Cube–Half-interval Product is $m+\lfloor\log_2n\rfloor$ if $2\le k\le m+1$; $m+1$ if $k>m+1$.
\begin{proof}
The VC dimension is $m+1$: all cube coordinates and one varying threshold coordinate are shattered, but two threshold coordinates cannot be shattered together. If $k>m+1$, a prefix tree of depth at least $k$ would shatter the $k$ coordinates on its first $k$ constant levels, a contradiction. Any shorter tree is entirely level-constant and has depth at most VC; a shattered set attains $m+1$. If $2\le k\le m+1$, query the $m$ cube coordinates in a fixed order and then perform a balanced threshold search of depth $\lfloor\log_2n\rfloor$. The first $m+1$ levels, including the first threshold query, are horizontally constant. This is an admissible prefix tree and attains the ordinary leaf-count upper bound $m+\lfloor\log_2n\rfloor$.
\end{proof}
\noindent\textbf{Reference.} Direct finite product calculation; independent exact regression in sync/verify\_cube\_halfinterval\_trees.py..

\subsection{Class Size of the Cube–Half-interval Product}\label{val:size_cube_halfinterval_product}
\noindent\textbf{Statement.} The Size of Cube–Half-interval Product is $2^m n$.
\begin{proof}
The $2^m$ choices of cube support and $n$ distinct nonempty initial segments are independent and give distinct concepts.
\end{proof}
\noindent\textbf{Reference.} Direct finite product calculation; independent exact regression in sync/verify\_cube\_halfinterval\_trees.py..

\subsection{Effective Range of the Cube–Half-interval Product}\label{val:effective_range_cube_halfinterval_product}
\noindent\textbf{Statement.} The Effective Range of Cube–Half-interval Product is $m+n-1$.
\begin{proof}
Every cube coordinate varies. In the half-interval factor, the first coordinate is constantly one, while each of the remaining $n-1$ coordinates takes both labels across the initial segments. Thus exactly $m+n-1$ domain coordinates are nonconstant.
\end{proof}
\noindent\textbf{Reference.} Direct finite product calculation; independent exact regression in sync/verify\_cube\_halfinterval\_trees.py..

\subsection{VC Dimension of the Branching Batch-Tree Class}\label{val:vc_dimension_branching_batch_tree}
\noindent\textbf{Statement.} The VC Dimension of Branching Batch-Tree Class is $b$.
\begin{proof}
The $b$ coordinates at the root realize all outgoing bit vectors, so VC is at least $b$. Consider coordinates $(u,i)$ and $(v,j)$ belonging to different tree nodes. If neither node is a prefix of the other, no leaf path visits both; their label pair cannot be $(1,1)$. If $u$ is a proper prefix of $v$, the condition $h_w(v,j)=1$ forces the path to visit $v$, and hence forces the outgoing vector at $u$ to the one fixed by $v$. In particular $h_w(u,i)$ has one fixed value whenever $h_w(v,j)=1$, so one of the two pairs with second label one is missing. The case where $v$ is a prefix of $u$ is the same statement with the coordinates exchanged. Thus any shattered set has all its coordinates at one tree node and has size at most $b$.
\end{proof}
\noindent\textbf{Reference.} Direct finite tree construction, with independent checks in the branching-batch-tree regression..

\subsection{Littlestone Dimension of the Branching Batch-Tree Class}\label{val:littlestone_dimension_branching_batch_tree}
\noindent\textbf{Statement.} The Littlestone Dimension of Branching Batch-Tree Class is $bn$.
\begin{proof}
Query the $b$ coordinates at the root in a fixed order, then those at the child indexed by the observed outgoing vector, and repeat for $n$ visited nodes. Every sequence of $bn$ answers is realized by its defining word, giving a shattered tree of depth $bn$. Distinct words yield distinct concepts: at the first differing outgoing vector they disagree at a coordinate of their common prefix node. There are therefore exactly $2^{bn}$ concepts, and distinct realizers at all leaves of any shattered depth-$d$ tree imply $2^d\le2^{bn}$, hence $d\le bn$.
\end{proof}
\noindent\textbf{Reference.} Direct finite tree construction, with independent checks in the branching-batch-tree regression..

\subsection{$k$-Block Littlestone Dimension of the Branching Batch-Tree Class}\label{val:kblock_littlestone_dimension_kgt2_branching_batch_tree}
\noindent\textbf{Statement.} The $k$-Block Littlestone Dimension ($k>2$) of Branching Batch-Tree Class is $bn$ if $k\mid b$; $b$ if $k>b$; otherwise $\max\{b,kn\lfloor b/k\rfloor\}\le B_k\le bn$.
\begin{proof}
The argument applies to every $k\ge2$. The ordinary dimension is $bn$, from the canonical path tree and the $2^{bn}$ leaf-count upper bound, so $B_k\le bn$. The VC value is $b$ (value record 1086). If $k>b$, a complete first $k$-block would shatter $k$ coordinates, which is impossible. Every shorter constrained tree is entirely level-constant, so $B_k=b$. For $k\le b$, at every visited tree node retain $k\lfloor b/k\rfloor$ of its $b$ coordinates, split them into successive $k$-batches, and fix the remaining outgoing bits to zero when choosing the next node. This describes a subclass and a valid shattered tree of depth $kn\lfloor b/k\rfloor$: each chosen batch has all $2^k$ outcomes, and the next tree-node label depends only on earlier retained outcomes and the fixed zero bits. The VC-shattered root coordinates also give depth $b$, including truncation if necessary. These prove the displayed lower bound. If $k$ divides $b$, it reaches $bn$, so equality holds.
\end{proof}
\noindent\textbf{Reference.} Direct finite tree construction, with independent checks in the branching-batch-tree regression..

\subsection{2-Block Littlestone Dimension of the Branching Batch-Tree Class}\label{val:2block_littlestone_dimension_branching_batch_tree}
\noindent\textbf{Statement.} The 2-Block Littlestone Dimension of Branching Batch-Tree Class is $bn$ if $b$ is even; $1$ if $b=1$; otherwise $\max\{b,(b-1)n\}\le B_2\le bn$.
\begin{proof}
Apply the general block argument for this family (value record 1088), which is proved for all $k\ge2$, at $k=2$. In particular the slice $b=2$ has exact value $2n$. No exact value is asserted here for odd $b>1$, only the displayed bounds.
\end{proof}
\noindent\textbf{Reference.} Direct finite tree construction, with independent checks in the branching-batch-tree regression..

\subsection{Proper Equivalence Queries Complexity of Full Cube}\label{val:proper_equivalence_queries_complexity_full_cube}
\noindent\textbf{Statement.} The Proper Equivalence Queries Complexity of Full Cube is $n$.
\begin{proof}
The full cube is itself the proper query class.  After each failed equivalence query, the returned counterexample fixes one previously unknown target bit; proposing any completion consistent with all returned bits therefore identifies the target in at most $n$ queries.  Conversely, while two or more target bits remain unfixed, an adversary can return one coordinate on which the proposal disagrees and retain both possible target completions.  Thus at least $n$ equivalence queries are required in the worst case.
\end{proof}
\noindent\textbf{Reference.} \cite{littlestone1988learning}.

\subsection{Prefix-$k$ Littlestone Dimension of the Branching Batch-Tree Class}\label{val:prefix_littlestone_dimension_branching_batch_tree}
\noindent\textbf{Statement.} The Prefix-$k$ Littlestone Dimension ($k>1$) of Branching Batch-Tree Class is $bn$ if $2\le k\le b$; $b$ if $k>b$.
\begin{proof}
The canonical depth-$bn$ path tree queries all $b$ root coordinates in a fixed order, so its first $b$ levels are horizontally constant. It therefore attains the ordinary leaf-count upper bound for every prefix size $k\le b$. If $k>b$, depth at least $k$ would shatter the $k$ coordinates on the first $k$ constant levels, contradicting VC $b$. Any shorter admissible tree is entirely level-constant and has depth at most $b$. The $b$ root coordinates attain that bound.
\end{proof}
\noindent\textbf{Reference.} Direct finite tree construction, with independent checks in the branching-batch-tree regression..

\subsection{Class Size of the Branching Batch-Tree Class}\label{val:size_branching_batch_tree}
\noindent\textbf{Statement.} The Size of Branching Batch-Tree Class is $2^{bn}$.
\begin{proof}
There are $|A|^n=2^{bn}$ defining words. At the first position where two words differ, the corresponding concepts disagree on at least one coordinate of their common prefix node, proving injectivity.
\end{proof}
\noindent\textbf{Reference.} Direct finite tree construction, with independent checks in the branching-batch-tree regression..

\subsection{Effective Range of the Branching Batch-Tree Class}\label{val:effective_range_branching_batch_tree}
\noindent\textbf{Statement.} The Effective Range of Branching Batch-Tree Class is $b(2^{bn}-1)/(2^b-1)$.
\begin{proof}
There are $\sum_{t=0}^{n-1}2^{bt}=(2^{bn}-1)/(2^b-1)$ tree nodes, each with $b$ coordinates. Every coordinate varies: choose a word visiting its node with the corresponding next bit zero, and another with that next bit one. Thus the entire domain is effective.
\end{proof}
\noindent\textbf{Reference.} Direct finite tree construction, with independent checks in the branching-batch-tree regression..

\subsection{Characteristic-p Yang Dimension of the Anchored Torsion Class}\label{val:characteristic_p_yang_dimension_anchored_torsion}
\noindent\textbf{Statement.} The Characteristic-p Yang Dimension of Anchored Torsion Class is $4$ if $p=2$; $3$ for every odd prime $p$.
\begin{proof}
Let $K$ be the triangle complex on $V=\{1,\ldots,6\}$ with the ten facets in the class definition. Its one-skeleton is the complete graph. Contract the spanning tree with edges $1i$ for $2\le i\le6$; the remaining oriented edges $e_{ij}$, $2\le i<j\le6$, form an integral cycle basis. The first five triangle boundaries set $e_{23}=e_{24}=e_{35}=e_{46}=e_{56}=0$. The other five impose $e_{36}=e_{26}$, $e_{45}=e_{25}$, $e_{25}=e_{26}$, $e_{34}=-e_{45}$, and $e_{34}=e_{36}$. Thus the first integral homology is generated by $t=e_{26}$ with exactly the relation $2t=0$. There are ten cycle generators and ten triangle boundaries, so the boundary matrix has full rational rank (its determinant in this basis is $-2$), and integral $H_2$ vanishes. There are no higher-dimensional faces. Consequently the reduced homology of $K$ vanishes over characteristic zero and every odd characteristic. Over characteristic two, universal coefficients give one-dimensional reduced homology in degrees one and two, and zero in other degrees.

For a version space $S$, use its nonspanning ambiguity complex $N(S)$: a subfamily is a face if it is constant on a coordinate that varies on $S$. Ideal Betti degree $i$ is reduced homology degree $i-1$; the Taylor/cube-envelope argument is stated in survey Proposition prop:yang-balanced-partition-masking. At the whole class, $N=\Delta_V\cup(h_a*K)$, where $\Delta_V$ is the full simplex on the six old concepts and $h_a*K$ is the cone on $K$. Indeed, the coordinate halves are $V$, $\{h_a\}$, and the pairs $\{h_a\}\cup F_i$, $V\setminus F_i$; the latter non-cone faces are absorbed in $\Delta_V$. The two pieces are contractible and intersect in $K$. Reduced Mayer--Vietoris therefore gives $\widetilde H_j(N;\mathbb F)\cong\widetilde H_{j-1}(K;\mathbb F)$. At the top envelope this contributes ideal degrees three and four in characteristic two, and no positive degree in characteristic zero or odd characteristic.

Every proper version space has ideal contributions of degree at most three. Conditioning only $x_0=1$ gives the six-concept class with all balanced three/three partitions: the ten $F_i$ and their complements exhaust all triples. Its ambiguity complex is the full two-skeleton and contributes degree three over every field. The $x_0=0$ space is a singleton. Any conditioning using some $x_i$, $i\ge1$, leaves at most four concepts, since its coordinate halves have sizes four and three. The Taylor resolution of an ideal with at most four minimal generators has length at most three. Further conditioning cannot increase the number of surviving concepts. The whole-class and all proper-envelope contributions are now exhausted: the maximum is four in characteristic two and exactly three in characteristic zero and every odd characteristic.
\end{proof}
\noindent\textbf{Reference.} Direct integral boundary and Mayer--Vietoris calculation; universal coefficients as in \cite{hatcher2002algebraic}, with the ideal convention of \cite{yang2018homological}. Verified by sync/verify\_yang\_characteristic\_witness.py..

\subsection{Yang Dimension of the Anchored Torsion Class}\label{val:yang_dimension_anchored_torsion}
\noindent\textbf{Statement.} The Yang dimension of Anchored Torsion Class is $4$.
\begin{proof}
The catalogue's default Yang dimension is its characteristic-two specialization. The complete coefficient-field calculation for this class (value record 1093) proves that this value is four.
\end{proof}
\noindent\textbf{Reference.} Direct calculation from the database-owned class definition and its characteristic computation..

\subsection{Characteristic-zero Yang Dimension of the Anchored Torsion Class}\label{val:characteristic_zero_yang_dimension_anchored_torsion}
\noindent\textbf{Statement.} The Characteristic-zero Yang Dimension of Anchored Torsion Class is $3$.
\begin{proof}
The full integral-homology and proper-version-space calculation for this class (value record 1093) also treats characteristic zero and gives exactly three over $\mathbb Q$.
\end{proof}
\noindent\textbf{Reference.} Direct calculation from the database-owned class definition and its characteristic computation..

\subsection{VC Dimension of the Anchored Torsion Class}\label{val:vc_dimension_anchored_torsion}
\noindent\textbf{Statement.} The VC Dimension of Anchored Torsion Class is $2$.
\begin{proof}
Coordinates $x_1,x_2$ are shattered: concepts $h_1,h_3,h_4,h_5$ realize $00,01,10,11$, respectively. No three coordinates can be shattered because the entire class has only seven concepts, fewer than the eight distinct realizers required.
\end{proof}
\noindent\textbf{Reference.} Direct calculation from the database-owned class definition and its characteristic computation..

\subsection{Class Size of the Anchored Torsion Class}\label{val:size_anchored_torsion}
\noindent\textbf{Statement.} The Size of Anchored Torsion Class is $7$.
\begin{proof}
The added zero concept is distinguished from the other six by $x_0$. Distinct old vertices are distinguished by a triple containing one but not the other: the ten triples and their complements give every balanced partition of the six vertices. Thus all seven listed concepts are distinct.
\end{proof}
\noindent\textbf{Reference.} Direct calculation from the database-owned class definition and its characteristic computation..

\subsection{Effective Range of the Anchored Torsion Class}\label{val:effective_range_anchored_torsion}
\noindent\textbf{Statement.} The Effective Range of Anchored Torsion Class is $11$.
\begin{proof}
Coordinate $x_0$ has labels zero and one. Every $x_i$, $1\le i\le10$, has four concepts on its zero side and three on its one side. All eleven coordinates are therefore nonconstant.
\end{proof}
\noindent\textbf{Reference.} Direct calculation from the database-owned class definition and its characteristic computation..

\subsection{Littlestone Dimension of the Anchored Torsion Class}\label{val:littlestone_dimension_anchored_torsion}
\noindent\textbf{Statement.} The Littlestone Dimension of Anchored Torsion Class is $2$.
\begin{proof}
The two VC-shattered coordinates give a depth-two level-constant Littlestone tree. A depth-three complete shattered tree would require eight distinct concepts, but the class has only seven. Hence the value is exactly two.
\end{proof}
\noindent\textbf{Reference.} Direct calculation from the database-owned class definition and its characteristic computation..

\subsection{Proper Equivalence Queries Complexity of Addressing}\label{val:proper_equivalence_queries_complexity_addressing}
\noindent\textbf{Statement.} The Proper Equivalence Queries Complexity of Addressing is $n$.
\begin{proof}
A learner can enumerate the $n$ proper hypotheses, so $n$ queries suffice.  Conversely, when the learner proposes the addressing concept with primary coordinate $j$, every other target has label $0$ at $j$ while the proposal has label $1$.  An adversary can always return $j$ as the counterexample and retain every other target.  Thus each failed proper equivalence query rules out only its proposed hypothesis, and the last remaining target still requires its accepting query.  The worst-case complexity is $n$.
\end{proof}
\noindent\textbf{Reference.} \cite{angluin1988queries,maass1992lower}.

\subsection{Characteristic-p Yang Dimension of the Odd-Torsion Dual Class}\label{val:characteristic_p_yang_dimension_odd_torsion_dual}
\noindent\textbf{Statement.} The Characteristic-p Yang Dimension of Odd-Torsion Dual Class is $10$ if $p=3$; $9$ for every prime $p\ne3$.
\begin{proof}
We also prove that the characteristic-zero value is nine. First consider the initial complex in the class definition. Its first two triangles for each index triangulate the mapping cylinder from the nine-edge cycle on $3,\ldots,11$ to the three-edge cycle on $0,1,2$, sending $3+i$ to $i\bmod3$. The last triangle cones each edge of the nine-cycle to vertex 12. Thus the realization is the mapping cone of a degree-three circle map. The cellular boundary in positive degrees is multiplication by three from degree two to degree one. Its integral homology is $H_1=\mathbb Z/3$, generated by the oriented cycle $0,1,2,0$, and $H_2=0$. Each subsequent edge/triangle addition is an elementary expansion: the newly added edge is a free face of its triangle at the time of addition. Reversing the additions gives collapses back to the initial complex. Hence the completed two-dimensional complex $M$ has the same homology, with the same generator. In particular it cannot contain the full boundary of a tetrahedron, which would be a nonzero integral two-cycle. Its edge set is complete, so its minimal nonfaces are exactly its missing triangles. The initial face counts are $(13,39,27)$; completing the graph to 78 edges adds 39 triangles, giving 66 triangles and $\binom{13}{3}-66=220$ missing triangles. An independent exact check contracts the star spanning tree at vertex zero and obtains a $66\times66$ integral boundary matrix of determinant three.

Let $K=M^*=\{A\subseteq V:V\setminus A\notin M\}$ be the combinatorial Alexander dual on the same thirteen vertices. Its facets are the complements of the missing triangles, all of size ten. Since $M$ has no faces of size four, $K$ contains every subset of size at most nine. The coordinate halves of the Boolean class are precisely the missing triples and their ten-element complements. Every triple is already a face of $K$, so their union as simplices is exactly $K$: this is the whole-class nonspanning ambiguity complex. Combinatorial Alexander duality gives $\widetilde H_j(K;\mathbb F)\cong\widetilde H^{10-j}(M;\mathbb F)$. One can obtain this directly by taking complements of faces: the augmented chain complex of $K$ identifies, with orientation signs, with the relative cochain complex of the full simplex modulo $M$, reversing degrees via $j\mapsto11-j$; the relative cohomology is the reduced cohomology of $M$ one degree lower. Universal coefficients now show that $K$ has reduced homology of rank one in degrees eight and nine over characteristic three, and is reduced-acyclic over characteristic zero and every other prime. The cube-envelope/Taylor formula in survey Proposition prop:yang-balanced-partition-masking shifts ambiguity homology degree $j$ to ideal Betti degree $j+1$. The whole class therefore contributes ideal degree ten in characteristic three and no positive degree in other characteristics. Every proper nonempty version space lies in a coordinate half of size at most ten, so its Taylor resolution has length at most nine. This proves the required upper bounds on the global maximum, not just on one envelope.

For the matching field-independent lower bound, the triangle $\{0,1,2\}$ stays missing: all three edges were present initially, so no expansion can add it. Condition its coordinate to zero, giving the ten-concept space $F=\{3,\ldots,12\}$. For every $v\in F$, at least one of $\{0,1,v\},\{0,2,v\},\{1,2,v\}$ is missing. Otherwise their oriented sum would bound the generator cycle $0,1,2,0$, contradicting its nonzero class in $H_1(M;\mathbb Z)$. The corresponding coordinate restricts to the singleton $\{v\}$ on $F$, and its other half is $F\setminus\{v\}$. Hence the ambiguity complex on $F$ is the full boundary of the nine-dimensional simplex. Its reduced homology in degree eight contributes ideal degree nine over every field. All claimed Yang values follow. These coordinates also distinguish the ten concepts in $F$, and the coordinate $\{0,1,2\}$ distinguishes them from the other three. The triples $\{0,2,4\}$ and $\{0,3,4\}$ are also missing: they are absent initially, and their only initially missing edge $\{0,4\}$ is the first edge added, with common neighbor one, creating $\{0,1,4\}$. After that, all edges of either proposed triangle exist, so the expansion rule cannot add it. Their coordinates distinguish vertices zero, one and two from one another, confirming that the class really has thirteen distinct concepts. The deterministic verifier checks these missing triples, every version-space size, the simplex boundary, and the actual top ambiguity boundary ranks independently of the Alexander-duality prediction.
\end{proof}
\noindent\textbf{Reference.} Direct mapping-cone, elementary-expansion and complement-chain calculations; universal coefficients as in \cite{hatcher2002algebraic}, with the ideal convention of \cite{yang2018homological}. Reproducible checks: sync/verify\_odd\_torsion\_yang.py..

\subsection{Yang Dimension of the Odd-Torsion Dual Class}\label{val:yang_dimension_odd_torsion_dual}
\noindent\textbf{Statement.} The Yang dimension of Odd-Torsion Dual Class is 9.
\begin{proof}
The default is $\mathrm Y=\mathrm Y_2$. The database-owned proof in value record 1100 gives value nine over every characteristic other than three, including characteristic two.
\end{proof}
\noindent\textbf{Reference.} See value record 1100 and its direct integral proof..

\subsection{Characteristic-zero Yang Dimension of the Odd-Torsion Dual Class}\label{val:characteristic_zero_yang_dimension_odd_torsion_dual}
\noindent\textbf{Statement.} The Characteristic-zero Yang Dimension of Odd-Torsion Dual Class is 9.
\begin{proof}
The database-owned integral proof in value record 1100 separately covers characteristic zero: the whole-class ambiguity complex is rationally acyclic, proper spaces contribute at most degree nine, and an explicit ten-concept conditioning contributes degree nine over every field.
\end{proof}
\noindent\textbf{Reference.} See value record 1100..

\subsection{Size of the Odd-Torsion Dual Class}\label{val:size_odd_torsion_dual}
\noindent\textbf{Statement.} The Size of Odd-Torsion Dual Class is 13.
\begin{proof}
There is one concept per vertex. The proof in value record 1100 explicitly distinguishes all thirteen concepts using missing-triple coordinates.
\end{proof}
\noindent\textbf{Reference.} Direct construction and distinguishing-coordinate proof..

\subsection{Effective Range of the Odd-Torsion Dual Class}\label{val:effective_range_odd_torsion_dual}
\noindent\textbf{Statement.} The Effective Range of Odd-Torsion Dual Class is 220.
\begin{proof}
The face count proved in 1100 leaves 220 missing triangles, hence 220 coordinates. Every coordinate has three concepts labeled one and ten labeled zero, so all are nonconstant.
\end{proof}
\noindent\textbf{Reference.} Direct face count and coordinate definition..

\subsection{VC Dimension of the Odd-Torsion Dual Class}\label{val:vc_dimension_odd_torsion_dual}
\noindent\textbf{Statement.} The VC Dimension of Odd-Torsion Dual Class is 2.
\begin{proof}
Every coordinate has exactly three positive concepts. A shattered triple of coordinates would require four distinct concepts positive on its first coordinate, so VC is at most two. The missing-triple coordinates $\{0,1,2\}$ and $\{0,2,4\}$ are shattered: vertices zero, one, four and three realize the pairs $(1,1),(1,0),(0,1),(0,0)$, respectively.
\end{proof}
\noindent\textbf{Reference.} Direct shattering argument; the deterministic construction verifier checks the two coordinates..

\subsection{Littlestone Dimension of the Odd-Torsion Dual Class}\label{val:littlestone_dimension_odd_torsion_dual}
\noindent\textbf{Statement.} The Littlestone Dimension of Odd-Torsion Dual Class is 2.
\begin{proof}
The established VC value two is a lower bound. In any putative depth-three complete mistake tree, the positive branch at the root would need at least four distinct concepts to realize a depth-two subtree. Every coordinate has only three positive concepts. Thus depth three is impossible and Littlestone dimension is two.
\end{proof}
\noindent\textbf{Reference.} Direct complete-tree leaf count and the explicit shattered pair..

\subsection{VC Dimension of Chen–Cheng–Tang Teaching Products}\label{val:vc_dimension_chen_teaching_products}
\noindent\textbf{Statement.} The VC Dimension of Chen–Cheng–Tang Teaching Products is $3n$.
\begin{proof}
The base shatters coordinates 1,2,3. Direct enumeration of all 495 four-coordinate sets finds fewer than sixteen traces on each, so its VC dimension is three. VC is additive on disjoint Cartesian products: a coordinate set is shattered if and only if each block restriction is shattered. Thus the value is $3n$.
\end{proof}
\noindent\textbf{Reference.} \cite{chen2016recursive}, Section 3 and Figure 4; \cite{doliwa2014recursive}, Lemma 16. Definition-driven finite verification: sync/verify\_chen\_teaching\_products.py and sync/chen\_teaching\_products.json..

\subsection{Minimum Teaching Set Size of Chen–Cheng–Tang Teaching Products}\label{val:minimum_teaching_set_size_chen_teaching_products}
\noindent\textbf{Statement.} The Minimum Teaching Set Size of Chen–Cheng–Tang Teaching Products is $5n$.
\begin{proof}
In the base class, none of the 495 four-coordinate restrictions has a singleton fibre. Smaller teaching sets would extend to a four-coordinate teaching set. Enumerating all 792 five-coordinate sets supplies a teaching set for every one of the 100 concepts, so every target has minimum teaching size exactly five. A sample teaches a product target if and only if it teaches each component separately (hold all other components fixed for necessity). Hence every target in the $n$-fold product has teaching size $5n$, in particular the minimum does.
\end{proof}
\noindent\textbf{Reference.} \cite{chen2016recursive}, Section 3 and Figure 4; \cite{doliwa2014recursive}, Lemma 16. Definition-driven finite verification: sync/verify\_chen\_teaching\_products.py and sync/chen\_teaching\_products.json..

\subsection{Teaching Dimension of Chen–Cheng–Tang Teaching Products}\label{val:teaching_dimension_chen_teaching_products}
\noindent\textbf{Statement.} The Teaching Dimension of Chen–Cheng–Tang Teaching Products is $5n$.
\begin{proof}
The database's minimum-teaching value proof for this family establishes that every product target, not merely an easiest target, has minimum teaching size $5n$. Taking the maximum over targets gives teaching dimension $5n$.
\end{proof}
\noindent\textbf{Reference.} \cite{chen2016recursive}, Section 3 and Figure 4; \cite{doliwa2014recursive}, Lemma 16. Definition-driven finite verification: sync/verify\_chen\_teaching\_products.py and sync/chen\_teaching\_products.json..

\subsection{Proper Equivalence Queries Complexity of Tagged Singletons}\label{val:proper_equivalence_queries_complexity_tagged_singletons}
\noindent\textbf{Statement.} The Proper Equivalence Queries Complexity of Tagged Singletons is $2$.
\begin{proof}
Query $\emptyset$ first.  A counterexample at an ordinary coordinate identifies the unique tagged pair containing it.  A counterexample at a tag coordinate leaves only that tag singleton and the tagged pairs carrying that tag; querying the tag singleton then either succeeds or returns the ordinary coordinate identifying the pair.  One query cannot identify every target: for any proper query, an adversary can choose two distinct remaining targets with a common disagreeing coordinate.  Thus the proper equivalence-query complexity is exactly two.
\end{proof}
\noindent\textbf{Reference.} \cite{maass1992lower}.

\subsection{Recursive Teaching Dimension of Chen–Cheng–Tang Teaching Products}\label{val:recursive_teaching_dimension_chen_teaching_products}
\noindent\textbf{Statement.} The Recursive Teaching Dimension of Chen–Cheng–Tang Teaching Products is $5n$.
\begin{proof}
For every finite class, minimum teaching size is at most RTD, which is at most teaching dimension. The minimum-teaching and teaching-dimension values for this product family are both $5n$, giving equality. The construction and ordinary RTD/VC product values are Figure 4 and Section 3 of \cite{chen2016recursive}.
\end{proof}
\noindent\textbf{Reference.} \cite{chen2016recursive}, Section 3 and Figure 4; \cite{doliwa2014recursive}, Lemma 16. Definition-driven finite verification: sync/verify\_chen\_teaching\_products.py and sync/chen\_teaching\_products.json..

\subsection{Projected Recursive Teaching Dimension of Chen–Cheng–Tang Teaching Products}\label{val:monotone_recursive_teaching_dimension_chen_teaching_products}
\noindent\textbf{Statement.} The Projected Recursive Teaching Dimension of Chen–Cheng–Tang Teaching Products is $5n$.
\begin{proof}
The full projection gives the lower bound from RTD $=5n$. For the base, sync/verify\_chen\_teaching\_products.py checks all 4096 projections. On each domain of size greater than five it repeatedly removes a projected concept that is the unique remaining trace on five chosen coordinates; every projection peels to empty. Domains of size at most five have RTD at most five directly. Such a removal sequence is a teaching plan of width at most five, so the base has projected RTD exactly five. Projected RTD is additive on Cartesian products: every projected product splits into projected factors, ordinary RTD is additive, and the maxima over projected domains separate. The full argument is survey Proposition prop:teaching-product-additivity.
\end{proof}
\noindent\textbf{Reference.} \cite{chen2016recursive}, Section 3 and Figure 4; \cite{doliwa2014recursive}, Lemma 16. Definition-driven finite verification: sync/verify\_chen\_teaching\_products.py and sync/chen\_teaching\_products.json..

\subsection{Effective Range of Chen–Cheng–Tang Teaching Products}\label{val:effective_range_chen_teaching_products}
\noindent\textbf{Statement.} The Effective Range of Chen–Cheng–Tang Teaching Products is $12n$.
\begin{proof}
A nonconstant base word and all its cyclic rotations make each of the twelve coordinates attain both labels. Independent copies preserve variation in every one of the $12n$ domain coordinates.
\end{proof}
\noindent\textbf{Reference.} \cite{chen2016recursive}, Section 3 and Figure 4; \cite{doliwa2014recursive}, Lemma 16. Definition-driven finite verification: sync/verify\_chen\_teaching\_products.py and sync/chen\_teaching\_products.json..

\subsection{Size of Chen–Cheng–Tang Teaching Products}\label{val:size_chen_teaching_products}
\noindent\textbf{Statement.} The Size of Chen–Cheng–Tang Teaching Products is $100^n$.
\begin{proof}
The first eight seed words have pairwise disjoint rotation orbits of length twelve; the ninth has a disjoint orbit of length four. The base therefore has $8\cdot12+4=100$ concepts. Independent choices on $n$ disjoint blocks give $100^n$ concepts.
\end{proof}
\noindent\textbf{Reference.} \cite{chen2016recursive}, Section 3 and Figure 4; \cite{doliwa2014recursive}, Lemma 16. Definition-driven finite verification: sync/verify\_chen\_teaching\_products.py and sync/chen\_teaching\_products.json..

\subsection{No-Clashing Teaching Dimension of Chen–Cheng–Tang Teaching Products}\label{val:noclashing_teaching_dimension_chen_teaching_products}
\noindent\textbf{Statement.} The No-Clashing Teaching Dimension of Chen–Cheng–Tang Teaching Products is $\le3n$.
\begin{proof}
Index the nine seed words as in the class definition. Assign them the coordinate sets $$\begin{gathered}S_1=\{2,7,12\},\quad S_2=\{1,2,4\},\quad S_3=\{2,5,7\},\\ S_4=\{1,2,3\},\quad S_5=\{4,5,9\},\quad S_6=\{4,7,8\},\\ S_7=\{2,4\},\quad S_8=\{3,7\},\quad S_9=\{1,5,9\}.\end{gathered}$$ Label each selected coordinate by its target's label and rotate both target and coordinate set by the same cyclic permutation. This is well-defined: the first eight target orbits have length twelve and $S_9$ is invariant under rotation by four. The resulting 100 samples have size at most three. For each of the 4950 distinct target pairs, their disagreement coordinates meet at least one of the two assigned coordinate sets; sync/verify\_cct\_noclash.py checks this directly from the definition and these nine sets. Thus no pair clashes. On a product target, take the union of its block samples. Two distinct product targets differ in some block, and the base teacher separates them in that block; hence width at most $3n$. For the stronger lower bound, encode a target $h$ by the integer $\sum_{j=1}^{12}h(j)2^{j-1}$, and let $\rho$ move coordinate $j$ to $j+1$ cyclically. The following triples $(u,v,q)$ specify pair-orbit weights: $$\begin{gathered}(85,87,8),\quad (85,340,4),\quad (85,349,4),\\ (85,684,2),\quad (85,696,2),\quad (85,1621,4),\\ (85,2133,8),\quad (87,171,1),\quad (87,174,1),\\ (87,343,8),\quad (87,699,1),\quad (87,1111,8),\\ (87,2738,2),\quad (87,2754,1),\quad (87,2798,2),\\ (171,349,2),\quad (171,683,8),\quad (171,1377,1),\\ (171,1399,2),\quad (171,1497,1),\quad (171,2219,8),\\ (343,349,4),\quad (343,1399,4),\quad (343,1879,4),\\ (343,2389,4),\quad (343,2859,2),\quad (343,3003,2),\\ (349,699,2),\quad (349,1369,4),\quad (349,1501,4),\\ (349,2397,8),\quad (683,699,8),\quad (683,1893,2),\\ (683,2987,4),\quad (699,1497,1),\quad (699,1911,1),\\ (699,2747,8),\quad (1399,1911,8),\quad (1911,3003,1)\end{gathered}$$ Give each distinct unordered pair in $\{\{\rho^t u,\rho^t v\}:0\le t<12\}$ weight $q$, counted once, and every other pair weight zero. The 39 orbits are disjoint; the last has four pairs and the others twelve. Write $w_{hg}$ for these weights and $a_h=\max_x\sum_{g:g(x)\ne h(x)}w_{hg}$. Direct integer replay from the class definition verifies that the total pair weight is 1780 and $a_h=8$ for each of the 100 targets. For any no-clashing teacher $T$, each pair is distinguished by a selected coordinate at one endpoint. Multiplying this covering condition by $w_{hg}$ and summing gives $1780\le\sum_h\sum_{x\in T(h)}\sum_{g:g(x)\ne h(x)}w_{hg}\le8\sum_h|T(h)|$. Thus its maximum size is at least $1780/800=89/40>2$. Together with the upper teacher this proves the exact base value three. The base total $\sum_h|T(h)|$ is an integer, so the same weighted count actually forces at least $\lceil1780/8\rceil=223$ selected coordinates across its 100 targets. For the $n$-fold product, fix a coordinate block and all other $n-1$ base targets. Restricting any product no-clashing teacher to that block gives a base no-clashing teacher: targets in this fiber agree on every other coordinate, so they must be distinguished inside the block. Each of these $100^{n-1}$ fibers has integer block total at least 223. Summing over fibers and then over the $n$ blocks gives total product teaching cost at least $223n100^{n-1}$, hence maximum size at least $223n/100$. Rounding this maximum yields $\lceil223n/100\rceil$. This is the fiber-total principle of Proposition \ref{prop:noclash-total-product}; it does not assume additivity of maximum NCTD or a product form for the teacher. It improves the former fractional bound $\lceil89n/40\rceil$, while the exact base value three and upper bound $3n$ are unchanged. No exact minimum base total or product value is claimed. The solver-free verifier sync/verify\_cct\_noclash.py replays both certificates; no unverified numerical optimum is used.
\end{proof}
\noindent\textbf{Reference.} Explicit computer-assisted upper and lower constructions, with independent integer pairwise replay in sync/verify\_cct\_noclash.py. The class is from \cite{chen2016recursive}, Figure 4; the no-clashing value and product bounds are local verified calculations, not claims made by that paper..

\subsection{Double Density or Average Degree of Chen–Cheng–Tang Teaching Products}\label{val:double_density_or_average_degree_chen_teaching_products}
\noindent\textbf{Statement.} The Double Density or Average Degree of Chen–Cheng–Tang Teaching Products is $12n/5$.
\begin{proof}
Direct enumeration of unordered base concept pairs counts exactly 120 pairs at Hamming distance one, on 100 vertices. Two product concepts are adjacent exactly when one block is an adjacent base pair and all other blocks agree. The product therefore has $120n\,100^{n-1}$ edges and $100^n$ vertices. Its average degree is twice their ratio, namely $12n/5$.
\end{proof}
\noindent\textbf{Reference.} Definition-driven edge count in sync/verify\_cct\_noclash.py; the base class is Figure 4 of \cite{chen2016recursive}..

\subsection{VC Dimension of the Chen Neighbor Puncture}\label{val:vc_dimension_chen_neighbor_puncture}
\noindent\textbf{Statement.} The VC Dimension of Chen Neighbor Puncture is $4$.
\begin{proof}
Coordinates $1,2,3,9$ realize all sixteen binary patterns. Direct enumeration of all $\binom{12}{5}$ coordinate sets finds no shattered five-set, so the VC dimension is exactly four. The database verifier reads the authoritative class definitions and independently checks these exact traces.
\end{proof}
\noindent\textbf{Reference.} Definition-driven certificate and independent replay: sync/verify\_chen\_neighbor\_puncture.py; sync/chen\_neighbor\_puncture\_certificate.json.gz. The original Chen base is from \cite{chen2016recursive}; the puncture modification and its separation proof are local results..

\subsection{Recursive Teaching Dimension of the Chen Neighbor Puncture}\label{val:recursive_teaching_dimension_chen_neighbor_puncture}
\noindent\textbf{Statement.} The Recursive Teaching Dimension of Chen Neighbor Puncture is $4$.
\begin{proof}
The full-domain removal plan in sync/chen\_neighbor\_puncture\_certificate.json.gz is a recursive teaching plan of width at most four. Its checker compares the actual remaining words at every removal and verifies that the recorded sample uniquely identifies the removed word. VC dimension four gives the matching lower bound.
\end{proof}
\noindent\textbf{Reference.} Definition-driven certificate and independent replay: sync/verify\_chen\_neighbor\_puncture.py; sync/chen\_neighbor\_puncture\_certificate.json.gz. The original Chen base is from \cite{chen2016recursive}; the puncture modification and its separation proof are local results..

\subsection{Projected Recursive Teaching Dimension of the Chen Neighbor Puncture}\label{val:monotone_recursive_teaching_dimension_chen_neighbor_puncture}
\noindent\textbf{Statement.} The Projected Recursive Teaching Dimension of Chen Neighbor Puncture is $4$.
\begin{proof}
This is a computer-assisted finite proof. The certificate sync/chen\_neighbor\_puncture\_certificate.json.gz contains a width-four removal plan for each of the 3302 coordinate sets of size greater than four. The independent checker in sync/verify\_chen\_neighbor\_puncture.py reconstructs the distinct projected words and, at each of 135515 removals, directly checks that the recorded sample lies in the projected domain, has size at most four, and uniquely identifies the removed word among the remaining words. Every plan ends with no words left. The other 794 projections have at most four coordinates and are taught using their entire domain. Hence every projection has RTD at most four. The shattered four-set gives the matching lower bound. The checker does not use the producer's bitset partitions and rejects an intentionally corrupted teaching sample.
\end{proof}
\noindent\textbf{Reference.} Definition-driven certificate and independent replay: sync/verify\_chen\_neighbor\_puncture.py; sync/chen\_neighbor\_puncture\_certificate.json.gz. The original Chen base is from \cite{chen2016recursive}; the puncture modification and its separation proof are local results..

\subsection{Order Sample Compression of the Chen Neighbor Puncture}\label{val:order_sample_compression_chen_neighbor_puncture}
\noindent\textbf{Statement.} The Order Sample Compression of Chen Neighbor Puncture is $\ge5$.
\begin{proof}
Write $\mathcal N=\mathcal N_{\mathrm{CCT}}$, and use $\mathcal C_0,h$ from its definition. Every five-coordinate sample of $h$ is realized by a neighbor $h^{(i)}\in\mathcal N$ at an omitted coordinate. Therefore $\mathcal C_0\subseteq C_5(\mathcal N)$. If $\mathrm{OSC}(\mathcal N)\le4$, survey Corollary cor:order-compression-trace-completion gives $\mathrm{OSC}(C_5(\mathcal N))\le4$. This contradicts $\mathrm{OSC}\ge\mathrm{RTD}$ and $\mathrm{RTD}(\mathcal C_0)=5$, using subclass monotonicity. Thus OSC is at least five. For the upper bound, every word lies within Hamming distance eight of $c=\mathtt{110000110000}$. Order the full cube by increasing distance to $c$, breaking ties arbitrarily. The first hypothesis consistent with a sample is $c$ modified only at its sampled disagreement coordinates. Keeping exactly those examples gives an order scheme of size at most eight. Hence $5\le\mathrm{OSC}(\mathcal N)\le8$; no exact value is asserted.
\end{proof}
\noindent\textbf{Reference.} Definition-driven certificate and independent replay: sync/verify\_chen\_neighbor\_puncture.py; sync/chen\_neighbor\_puncture\_certificate.json.gz. The original Chen base is from \cite{chen2016recursive}; the puncture modification and its separation proof are local results..

\subsection{Proper Equivalence Queries Complexity of Half-intervals}\label{val:proper_equivalence_queries_complexity_halfintervals}
\noindent\textbf{Statement.} The Proper Equivalence Queries Complexity of Half-intervals is $\log n$.
\begin{proof}
Maintain the interval of feasible right endpoints and query its median initial segment.  A counterexample tells on which side of the queried boundary the target endpoint lies, so the interval is halved at each query.  Conversely, against a query $[k]$, an adversary can answer at $k$ and retain all targets ending before $k$, or answer at $k+1$ and retain all targets ending at or after $k+1$; it chooses the larger side.  Thus logarithmically many proper equivalence queries are both sufficient and necessary.
\end{proof}
\noindent\textbf{Reference.} \cite{angluin1988queries}.

\subsection{Size of the Chen Neighbor Puncture}\label{val:size_chen_neighbor_puncture}
\noindent\textbf{Statement.} The Size of Chen Neighbor Puncture is $109$.
\begin{proof}
The base has 100 words. The removed word is in the base, and exactly two of its twelve Hamming neighbors are already in the remaining base. The union therefore has $99+12-2=109$ distinct words. The definition-driven verifier checks the count.
\end{proof}
\noindent\textbf{Reference.} Definition-driven certificate and independent replay: sync/verify\_chen\_neighbor\_puncture.py; sync/chen\_neighbor\_puncture\_certificate.json.gz. The original Chen base is from \cite{chen2016recursive}; the puncture modification and its separation proof are local results..

\subsection{Effective Range of the Chen Neighbor Puncture}\label{val:effective_range_chen_neighbor_puncture}
\noindent\textbf{Statement.} The Effective Range of Chen Neighbor Puncture is $12$.
\begin{proof}
The domain has twelve coordinates and both labels occur at every coordinate, as directly checked from the defining words. Thus all twelve coordinates are active.
\end{proof}
\noindent\textbf{Reference.} Definition-driven certificate and independent replay: sync/verify\_chen\_neighbor\_puncture.py; sync/chen\_neighbor\_puncture\_certificate.json.gz. The original Chen base is from \cite{chen2016recursive}; the puncture modification and its separation proof are local results..

\subsection{Hamming Radius of the Chen Neighbor Puncture}\label{val:hamming_radius_chen_neighbor_puncture}
\noindent\textbf{Statement.} The Hamming Radius of Chen Neighbor Puncture is $8$.
\begin{proof}
The center $\mathtt{110000110000}$ encloses the class in a Hamming ball of radius eight. Exhaustive evaluation of all 4096 possible centers gives minimum maximum distance eight. The checker computes the distances directly from the authoritative class definition, using exact integer arithmetic.
\end{proof}
\noindent\textbf{Reference.} Definition-driven certificate and independent replay: sync/verify\_chen\_neighbor\_puncture.py; sync/chen\_neighbor\_puncture\_certificate.json.gz. The original Chen base is from \cite{chen2016recursive}; the puncture modification and its separation proof are local results..

\subsection{VC Dimension of Quadratic-Residue Tournament Classes}\label{val:vc_dimension_quadratic_residue_tournament_classes}
\noindent\textbf{Statement.} The VC Dimension of Quadratic-Residue Tournament Classes is $\Theta(\log_2 p)$.
\begin{proof}
Theorem 13 of Simon gives the strong signed extension property when $p>k^2 2^{2k-2}$: for any distinct coordinates $a_1,\ldots,a_k$ and prescribed signs, there is a vertex $x$ outside these coordinates with the prescribed orientations between $x$ and each $a_j$. In the catalogue convention $C_x(a_j)=1$ exactly when $x\to a_j$, so all binary patterns occur. Thus every such coordinate set is shattered. There are at most $p$ concepts, giving the upper bound by counting the $2^d$ traces of a shattered $d$-set. The inequality defining $k_p$ yields $k_p=\tfrac12\log_2p-O(\log_2\log_2p)$: for sufficiently large $p$, the integer $\lfloor\tfrac12\log_2p-\log_2\log_2p\rfloor$ is admissible, and $k_p\le1+\tfrac12\log_2p$. This proves the asserted order.
\end{proof}
\noindent\textbf{Reference.} \cite[Theorem 13]{simon2023tournaments}; the VC consequence follows directly from the signed extension property..

\subsection{Projected Recursive Teaching Dimension of Quadratic-Residue Tournament Classes}\label{val:projected_recursive_teaching_dimension_quadratic_residue_tournament_classes}
\noindent\textbf{Statement.} The Projected Recursive Teaching Dimension of Quadratic-Residue Tournament Classes is $\Theta(\log_2 p)$.
\begin{proof}
The VC value proof for this class shows that every $k_p$-set is shattered. Projecting onto one gives a full $k_p$-cube. Each cube vertex requires all $k_p$ labels for its first teaching set, and domain size bounds every later teaching set by $k_p$, so its RTD is exactly $k_p$. For the upper bound, every projection has at most $p$ concepts. Any finite class of size $M>1$ has a teaching sample of size at most $\lfloor\log_2M\rfloor$: repeatedly choose a nonconstant coordinate and its smaller nonempty label fibre until one concept remains. Each choice at least halves the residual size; the retained labels identify the survivor. Apply this to every remaining subclass during recursive elimination. Thus every projection has RTD at most $\lfloor\log_2p\rfloor$. The bounds on $k_p$ in the VC value proof give the asymptotic order.
\end{proof}
\noindent\textbf{Reference.} Derived from \cite[Theorem 13]{simon2023tournaments}, the full-cube calculation, and the elementary halving bound..

\subsection{Projected No-Clashing Teaching Dimension of Quadratic-Residue Tournament Classes}\label{val:projected_noclashing_teaching_dimension_quadratic_residue_tournament_classes}
\noindent\textbf{Statement.} The Projected No-Clashing Teaching Dimension of Quadratic-Residue Tournament Classes is $\Theta(\log_2 p)$.
\begin{proof}
The VC value proof for this class supplies a projection equal to the full $k_p$-cube. For a no-clashing teacher on a $k$-cube, each of its $k2^{k-1}$ edges must have the coordinate on which its endpoints differ selected by at least one endpoint's teacher. The sum of teacher sizes is at most $2^k$ times the maximum teacher size. Counting these edge incidences gives maximum size at least $\lceil k/2\rceil$. Apply this to the shattered projection. For the upper bound, NCTD is at most RTD on every projection: in a recursive teaching order, the earlier of two concepts receives a sample inconsistent with the later one. The preceding projected RTD value proof therefore bounds projected NCTD by $\lfloor\log_2p\rfloor$. Since $k_p=\tfrac12\log_2p-O(\log_2\log_2p)$, both bounds have logarithmic order.
\end{proof}
\noindent\textbf{Reference.} Derived from \cite[Theorem 13]{simon2023tournaments} and the cube edge-counting argument for no-clashing teaching; see also \cite{KSZ2019}..

\subsection{Labeled Sample Compression of Full Cube}\label{val:labeled_sample_compression_full_cube}
\noindent\textbf{Statement.} The Labeled Sample Compression of Full Cube is $\Theta(n)$.
\begin{proof}
Every binary total reconstruction belongs to a full cube, so ordinary LSC equals proper LSC. The database-owned proper labeled full-cube proof gives the explicit six-coordinate width-two decoder, its fixed-block extension of size $2\lfloor n/6\rfloor+a_{n\bmod6}$ with $(a_0,\ldots,a_5)=(0,1,1,2,2,2)$, and the counting lower bound $b_n$ defined there. These give the stated exact values and $\Theta(n)$ growth. In particular the six-cube has LSC two whereas stable labeled compression and NCTD are three. The ordinary and stable values must not be identified. The full decoder and case proof are maintained only in the proper labeled value record.
\end{proof}
\noindent\textbf{Reference.} Database-owned proper labeled full-cube value proof (value 743); exact replay in sync/verify\_six\_cube\_compression.py..

\subsection{VC Dimension of Singleton–Cosingleton Classes}\label{val:vc_dimension_singleton_cosingleton}
\noindent\textbf{Statement.} The VC Dimension of Singleton–Cosingleton Classes is $3$.
\begin{proof}
Every three-coordinate set is shattered: singletons give its one-element traces, cosingletons its two-element traces, and an index outside the triple supplies the empty and full traces. No four-coordinate set is shattered, since all its traces have sizes zero, one, three or four, never two. For $n\ge4$ this proves VC exactly three.
\end{proof}
\noindent\textbf{Reference.} Local definition-driven proof and exact finite replay in sync/verify\_singleton\_cosingleton.py. The class was already used in the survey's signed intersection-completion argument..

\subsection{Order Sample Compression of Singleton–Cosingleton Classes}\label{val:order_sample_compression_singleton_cosingleton}
\noindent\textbf{Statement.} The Order Sample Compression of Singleton–Cosingleton Classes is $3$.
\begin{proof}
Use the fixed decoder order $\varnothing,[n],\{1\},\ldots,\{n\},[n]\setminus\{1\},\ldots,[n]\setminus\{n\}$, with all other binary hypotheses placed afterward if a full-cube order is desired. Write $P,N$ for the positive and negative coordinates of a realizable sample. If $P=\varnothing$, retain nothing and decode to $\varnothing$. If $P\ne\varnothing$ and $N=\varnothing$, retain the first positive and decode to $[n]$. If $|P|=1$ and $N\ne\varnothing$, retain the positive and the first negative; the first consistent hypothesis is that positive's singleton. Otherwise $|P|\ge2$, and realizability forces $|N|=1$: retain the first two positives and the negative, and decode to the cosingleton missing the negative coordinate. Each retained sample has the same first consistent hypothesis as the original sample. Deleting any retained example introduces an earlier consistent hypothesis, so these are minimal order kernels. This proves OSC at most three. The general bound $\mathrm{OSC}\ge\mathrm{RTD}^p\ge\mathrm{VC}$ and the VC value proof give the matching lower bound. Reconstructions $\varnothing$ and $[n]$ are outside the class, so this is an improper construction, not a proper-order bound.
\end{proof}
\noindent\textbf{Reference.} Local definition-driven proof and exact finite replay in sync/verify\_singleton\_cosingleton.py. The class was already used in the survey's signed intersection-completion argument..

\subsection{Projected Positive Recursive Teaching Dimension of Singleton–Cosingleton Classes}\label{val:projected_positive_recursive_teaching_dimension_singleton_cosingleton}
\noindent\textbf{Statement.} The Projected Positive Recursive Teaching Dimension of Singleton–Cosingleton Classes is $n-1$.
\begin{proof}
The nonempty intersection closure is exactly $I(\mathcal A_n)=2^{[n]}\setminus\{[n]\}$. Indeed every proper $T\subsetneq[n]$ is the nonempty intersection $\bigcap_{i\notin T}([n]\setminus\{i\})$, and no intersection of proper supports can equal $[n]$. This closure has VC $n-1$. Survey Proposition \ref{prop:intersection-closure-positive-teaching} identifies its VC with $\mathrm{RTD}^{+p}(\mathcal A_n)$, proving the stated value.
\end{proof}
\noindent\textbf{Reference.} Local definition-driven proof and exact finite replay in sync/verify\_singleton\_cosingleton.py. The class was already used in the survey's signed intersection-completion argument..

\subsection{Proper Equivalence Queries Complexity of Singletons plus empty set}\label{val:proper_equivalence_queries_complexity_singletons_plus_empty_set}
\noindent\textbf{Statement.} The Proper Equivalence Queries Complexity of Singletons plus empty set is $1$.
\begin{proof}
Query the empty concept. If it is not the target, a counterexample is its unique positive coordinate, identifying the target singleton. One query is necessary for a non-singleton class.
\end{proof}

\subsection{Minimum Teaching Set Size of Singleton–Cosingleton Classes}\label{val:minimum_teaching_set_size_singleton_cosingleton}
\noindent\textbf{Statement.} The Minimum Teaching Set Size of Singleton–Cosingleton Classes is $3$.
\begin{proof}
Every target has teaching size three. To teach singleton $\{i\}$, label $i$ positive and two distinct other coordinates negative. The positive excludes the other singletons, and the two negatives exclude all cosingletons. A sample of size at most two with no positive label is consistent with at least $n-2\ge2$ singletons. A sample with a positive label must label $i$ positive; if it also labels $j$ negative, both $\{i\}$ and $[n]\setminus\{j\}$ survive. With no negative label there are again multiple realizers. Thus no such sample teaches the singleton. Complementing all labels proves the same statement for every cosingleton. Taking the minimum gives three.
\end{proof}
\noindent\textbf{Reference.} Local definition-driven proof and exact finite replay in sync/verify\_singleton\_cosingleton.py. The class was already used in the survey's signed intersection-completion argument..

\subsection{Recursive Teaching Dimension of Singleton–Cosingleton Classes}\label{val:recursive_teaching_dimension_singleton_cosingleton}
\noindent\textbf{Statement.} The Recursive Teaching Dimension of Singleton–Cosingleton Classes is $3$.
\begin{proof}
The minimum-teaching value proof \ref{val:minimum_teaching_set_size_singleton_cosingleton} shows that every target initially has teaching size three. Hence the first recursive teaching step costs at least three. The same size-three teachers remain valid after targets are removed, giving a recursive teaching plan of width at most three. Therefore RTD is exactly three.
\end{proof}
\noindent\textbf{Reference.} Local definition-driven proof and exact finite replay in sync/verify\_singleton\_cosingleton.py. The class was already used in the survey's signed intersection-completion argument..

\subsection{Teaching Dimension of Singleton–Cosingleton Classes}\label{val:teaching_dimension_singleton_cosingleton}
\noindent\textbf{Statement.} The Teaching Dimension of Singleton–Cosingleton Classes is $3$.
\begin{proof}
The minimum-teaching value proof \ref{val:minimum_teaching_set_size_singleton_cosingleton} establishes teaching size three for every target, not just for the easiest target. Taking the maximum also gives three.
\end{proof}
\noindent\textbf{Reference.} Local definition-driven proof and exact finite replay in sync/verify\_singleton\_cosingleton.py. The class was already used in the survey's signed intersection-completion argument..

\subsection{Size of Singleton–Cosingleton Classes}\label{val:size_singleton_cosingleton}
\noindent\textbf{Statement.} The Size of Singleton–Cosingleton Classes is $2n$.
\begin{proof}
There are $n$ distinct singletons and $n$ distinct cosingletons. Their support sizes one and $n-1$ differ for $n\ge4$, so the two groups are disjoint.
\end{proof}
\noindent\textbf{Reference.} Local definition-driven proof and exact finite replay in sync/verify\_singleton\_cosingleton.py. The class was already used in the survey's signed intersection-completion argument..

\subsection{Effective Range of Singleton–Cosingleton Classes}\label{val:effective_range_singleton_cosingleton}
\noindent\textbf{Statement.} The Effective Range of Singleton–Cosingleton Classes is $n$.
\begin{proof}
The domain has $n$ coordinates. Each coordinate takes value one on its singleton and zero on a different singleton, so every coordinate is nonconstant.
\end{proof}
\noindent\textbf{Reference.} Local definition-driven proof and exact finite replay in sync/verify\_singleton\_cosingleton.py. The class was already used in the survey's signed intersection-completion argument..

\subsection{co-VC Dimension of Singleton–Cosingleton Classes}\label{val:covc_dimension_singleton_cosingleton}
\noindent\textbf{Statement.} The co-VC dimension of Singleton–Cosingleton Classes is $n$.
\begin{proof}
The all-zero labeling on $[n]$ is absent, and all its one-coordinate flips are the singleton concepts. It is therefore a hollow $n$-star. No projected hollow star can have more than the $n$ domain coordinates, giving equality.
\end{proof}
\noindent\textbf{Reference.} Local definition-driven proof and exact finite replay in sync/verify\_singleton\_cosingleton.py. The class was already used in the survey's signed intersection-completion argument..

\subsection{Proper Ordered Sample Compression of Singleton–Cosingleton Classes}\label{val:proper_ordered_sample_compression_singleton_cosingleton}
\noindent\textbf{Statement.} The Proper Ordered Sample Compression of Singleton–Cosingleton Classes is $n-1$.
\begin{proof}
The hollow-star lower bound in Proposition \ref{prop:proper-compression-covc}, together with the class's co-VC value $n$, gives at least $n-1$. For the upper bound order all singletons by their coordinate, followed by all cosingletons in the same order. Write $P,N$ for the positive and negative coordinates of a realizable sample. If $P=\varnothing$, let $i$ be the first coordinate not in $N$ and retain all negatives before $i$; the first compatible singleton is $\{i\}$. If $|P|=1$, retain that positive alone, selecting its singleton. If $|P|\ge2$ and $N\ne\varnothing$, realizability forces $N=\{j\}$; retain two positives and the negative at $j$, selecting the cosingleton missing $j$. This uses three examples, at most $n-1$. Finally if $|P|\ge2$ and $N=\varnothing$, let $j$ be the least coordinate not in $P$; such a coordinate exists because the full set is not a class member. Retain every positive before $j$, adding further positives if needed to retain at least two. This uses $\max(j-1,2)\le n-1$ examples, rules out every singleton and every earlier cosingleton, and selects the cosingleton missing $j$. In every case the retained sample preserves the first consistent proper hypothesis. Removing any redundant retained examples gives a minimal order kernel without increasing its size. Thus the bound holds for every $n\ge4$.
\end{proof}
\noindent\textbf{Reference.} Local proof using the survey comparison theorems and database-owned endpoint values..

\subsection{Projected Recursive Teaching Dimension of Singleton–Cosingleton Classes}\label{val:monotone_recursive_teaching_dimension_singleton_cosingleton}
\noindent\textbf{Statement.} The Projected Recursive Teaching Dimension of Singleton–Cosingleton Classes is $3$.
\begin{proof}
The general inequalities $\mathrm{VC}\le\mathrm{RTD}^p\le\mathrm{OSC}$ place projected RTD between the two established values three on this class. See the database-owned VC and OSC value proofs. This uses the projected, not just the ordinary, recursive teaching parameter.
\end{proof}
\noindent\textbf{Reference.} Local proof using the survey comparison theorems and database-owned endpoint values..

\subsection{Labeled Sample Compression of Warmuth's C5 Class}\label{val:labeled_sample_compression_warmuth_c5}
\noindent\textbf{Statement.} The Labeled Sample Compression of Warmuth's $C_5$ class is $2$.
\begin{proof}
Index coordinates modulo five. The ten concepts are precisely the indicators of the five nonadjacent pairs and of the five consecutive triples. For a key of size zero or one, copy its labels and put zero elsewhere. For a two-example key with coordinate set $U$, copy its labels and predict one at an unretained coordinate $x$ exactly when $U\cup\{x\}$ is a consecutive triple in the five-cycle; predict zero otherwise. This is one fixed decoder depending only on the retained labeled key. A consecutive triple cannot have all-zero labels in a realizable sample, because its complementary adjacent pair contains neither a nonadjacent positive pair nor a positive triple. A nonconsecutive triple cannot have all-one labels, because it is contained in neither kind of positive support. These missing patterns are respectively $000$ and $111$, so Proposition~\ref{prop:short-sample-missing-pattern} covers every sample of size at most three. For a full concept supported on $\{i,i+2\}$, retain the adjacent negative pair $\{i+3,i+4\}$: the decoder predicts one exactly on $\{i,i+2\}$. For a full concept supported on $\{i,i+1,i+2\}$, retain its two positive endpoints $\{i,i+2\}$: the only further positive prediction is the middle coordinate. Finally, take any four-example sample and a full concept extending it. If that concept's just-described key is observed, retain it. Otherwise the omitted coordinate lies in that key. Up to a rotation or reflection of the five-cycle, omission from the negative-pair key gives observed labeling $1010$ on coordinates $0,1,2,3$, while omission from the positive-endpoint key gives $1100$ on those coordinates. For $1100$, retain coordinates $0,2$ with labels $1,0$; the decoder returns $11000$. For $1010$, retain coordinates $0,1$ with labels $1,0$; the decoder returns $10101$. Both reconstructions match the four observed labels. The decoder commutes with cycle rotations and reflections, so these cases cover every four-example sample. This proves the upper bound without exhaustive search. The 51 keys and 176 samples are also independently replayed with bit masks and literal sets; a separate check verifies the two four-sample orbits. Outputs such as $11000$ and $10101$ are outside the class, so this is an improper scheme. For the lower bound, consider the realizable samples on triples of coordinates. On a consecutive triple $\{i,i+1,i+2\}$ give a labeling weight equal to its number of ones at the two endpoints $i,i+2$. On a nonconsecutive triple give it weight equal to its number of zeros at its unique adjacent pair. Every positive-weight labeling is realizable. For example, on $\{0,1,2\}$ the seven nonzero patterns $100,010,001,110,101,011,111$ are realized by supports $\{0,3\},\{1,3\},\{2,4\},\{4,0,1\},\{0,2\},\{1,2,3\},\{0,1,2\}$, respectively. Multiplication of coordinates by two modulo five followed by complementation preserves the class and exchanges the two kinds of triples. The omitted pattern has weight zero in each case. The eight binary labelings on any triple have total weight eight, so the total weight is $10\cdot8=80$, carried by 60 samples. For any total binary output $h$, the sum of weights of its ten triple restrictions is exactly ten: its ones are counted over all five nonadjacent pairs and its zeros over all five adjacent pairs, with each coordinate counted twice in each collection. At a singleton key, only the six triples containing its coordinate contribute. In their weight sum that coordinate occurs twice with a positive coefficient and twice with a negative coefficient, and each other coordinate occurs once with each sign; the constant terms sum to six. Thus every output extending that singleton key covers weight exactly six. This statement holds for all total binary outputs, not only class members. The empty key and ten singleton keys can therefore cover total weight at most $10+10\cdot6=70<80$. Some realizable sample must be missed, ruling out every width-one decoder without an injection search. The same count bounds minimum expected coverage under arbitrary per-key laws by $7/8$. The program sync/verify\_c5\_labeled\_compression.py replays the 60 weighted samples and all 192 key/output loads in bitmask and literal-set representations. Its older exhaustive lower check remains an independent regression (39811 states and 5848 terminal injections), not a dependency of this counting proof. Together the direct upper and lower bounds give LSC equal to two.
\end{proof}
\noindent\textbf{Reference.} Direct cyclic construction and weighted-counting obstruction, independently replayed in sync/verify\_c5\_labeled\_compression.py; the general short-sample lemma is proved in the survey. The earlier exhaustive lower test is retained only as a regression. No external theorem or solver UNSAT status is used as the lower-bound proof..

\subsection{VC Dimension of Private-Coordinate Cubes}\label{val:vc_dimension_private_coordinate_cube}
\noindent\textbf{Statement.} The VC Dimension of Private-Coordinate Cubes is $n$.
\begin{proof}
Projection onto the original coordinates is the full $n$-cube, so those coordinates are shattered. The class has $2^n$ distinct concepts, and shattering $d$ coordinates requires at least $2^d$ concepts. Hence VC is exactly $n$.
\end{proof}
\noindent\textbf{Reference.} Direct counting and shattering proof; \cite{KSZ2019}, Proposition 26, records the underlying teaching/VC separation and credits Goldman--Kearns..

\subsection{Proper Equivalence Queries Complexity of Singletons}\label{val:proper_equivalence_queries_complexity_singletons}
\noindent\textbf{Statement.} The Proper Equivalence Queries Complexity of Singletons is $n-1$.
\begin{proof}
Query singleton hypotheses.  For each wrong query, an adversary can return the queried coordinate as a negative counterexample, eliminating only that singleton.  After $n-1$ rejected hypotheses, the remaining singleton is forced, so $n-1$ queries are both necessary and sufficient.
\end{proof}

\subsection{No-Clashing Teaching Dimension of Private-Coordinate Cubes}\label{val:noclashing_teaching_dimension_private_coordinate_cube}
\noindent\textbf{Statement.} The No-Clashing Teaching Dimension of Private-Coordinate Cubes is $1$.
\begin{proof}
Give $h_v$ the singleton teacher $\{(p_v,1)\}$. Every other concept labels $p_v$ zero, so these are ordinary teaching sets and in particular form a no-clashing assignment. Width zero is impossible for a class with at least two concepts: two empty teachers would clash. Thus NCTD is exactly one.
\end{proof}
\noindent\textbf{Reference.} \cite{KSZ2019}, Proposition 28(1), and the explicit private-positive teaching assignment above..

\subsection{Projected No-Clashing Teaching Dimension of Private-Coordinate Cubes}\label{val:projected_noclashing_teaching_dimension_private_coordinate_cube}
\noindent\textbf{Statement.} The Projected No-Clashing Teaching Dimension of Private-Coordinate Cubes is $\lceil n/2\rceil$.
\begin{proof}
Projection onto the original coordinates gives $Q_n$, with NCTD $\lceil n/2\rceil$ by \cite{KSZ2019}, Theorem 19. For the reverse bound consider an arbitrary projection retaining original coordinates $S\subseteq[n]$ and some private coordinates. Each concept with a retained private positive has a unique projected trace; assign it that singleton positive teacher. All other distinct traces have zero on every retained private coordinate and are distinguished by their restrictions to $S$. Use a no-clashing assignment of size at most $\lceil |S|/2\rceil$ for this subfamily of the full cube on $S$ (subclass monotonicity and Theorem 19). Two private-positive traces are separated by their private teachers; a private-positive trace and a zero-private trace are separated by the former's positive teacher; pairs of zero-private traces are separated by the chosen cube assignment. If $S$ is empty there is at most one zero-private trace and its teacher is empty. The width is at most $\max(1,\lceil |S|/2\rceil)\le\lceil n/2\rceil$ for $n\ge1$.
\end{proof}
\noindent\textbf{Reference.} Direct projection calculation above, using the full-cube value in \cite{KSZ2019}, Theorem 19. This refinement is not attributed to Proposition 28. Regression: sync/verify\_private\_coordinate\_cube.py..

\subsection{Labeled Sample Compression of Private-Coordinate Cubes}\label{val:labeled_sample_compression_private_coordinate_cube}
\noindent\textbf{Statement.} The Labeled Sample Compression of Private-Coordinate Cubes is $\Theta(n)$.
\begin{proof}
In fact $\mathrm{LSC}(\mathcal P_n)=\mathrm{LSC}(Q_n)$. A compressor for $\mathcal P_n$ restricts to one for its cube projection: feed it samples using only original coordinates, so its retained key also uses only those coordinates, and restrict the reconstruction. Conversely let $(\kappa,\rho)$ be any width-$k$ compressor for $Q_n$, where $k\ge1$ since $n\ge1$. If a realizable sample on $X_n$ contains $(p_v,1)$, retain that singleton and reconstruct $h_v$; realizability makes it consistent with the whole sample. If there is no private positive, compress the original-coordinate subsample using $\kappa$; reconstruct $\rho$ on the original coordinates and zero on every private coordinate. This agrees with every observed private label. The decoder distinguishes a private-positive singleton from a key on the original domain, so there is no auxiliary message, and the width is at most $k$. The second reconstruction may be improper, which is allowed for LSC. The cube value proof gives $n/5<\mathrm{LSC}(Q_n)\le2\lfloor n/5\rfloor+\lceil(n\bmod5)/2\rceil$, establishing the stated asymptotic order.
\end{proof}
\noindent\textbf{Reference.} \cite{KSZ2019}, Proposition 28(1), for the separation. The exact cube-compression identity is proved directly above, not attributed to that proposition. Cube bounds are database-owned in the proper labeled full-cube value proof. Regression: sync/verify\_private\_coordinate\_cube.py..

\subsection{Size of Private-Coordinate Cubes}\label{val:size_private_coordinate_cube}
\noindent\textbf{Statement.} The Size of Private-Coordinate Cubes is $2^n$.
\begin{proof}
There is one concept $h_v$ for each of the $2^n$ original binary words. They are distinct on the original coordinates, as well as on their private coordinates.
\end{proof}
\noindent\textbf{Reference.} Direct count from the database-owned class definition..

\subsection{Effective Range of Private-Coordinate Cubes}\label{val:effective_range_private_coordinate_cube}
\noindent\textbf{Statement.} The Effective Range of Private-Coordinate Cubes is $n+2^n$.
\begin{proof}
Every original coordinate varies because all binary words occur. Every private coordinate $p_v$ is one on $h_v$ and zero on any different concept, which exists for $n\ge1$. Thus all $n+2^n$ coordinates belong to the effective range.
\end{proof}
\noindent\textbf{Reference.} Direct count of varying coordinates from the database-owned class definition..

\subsection{Projected No-Clashing Teaching Dimension of Antipodal Six-Code Class}\label{val:projected_noclashing_teaching_dimension_antipodal_six_code}
\noindent\textbf{Statement.} The Projected No-Clashing Teaching Dimension of Antipodal Six-Code Class is $1$.
\begin{proof}
On the full class assign singleton teachers as follows, listing word, coordinate and label: $$\begin{array}{c|c|c}h&x&h(x)\\\hline 0000&0&0\\1100&3&0\\1010&1&0\\0101&1&1\\0011&3&1\\1111&0&1\end{array}$$ Every pair of distinct concepts disagrees on at least one of its two selected coordinates, as direct inspection of the fifteen pairs verifies. Deleting any one coordinate leaves six of the eight three-bit words, with the two missing words complementary. After a fixed label flip, this is the six-cycle $000,100,110,111,011,001$. Assign each cycle vertex the coordinate of its outgoing edge in this cyclic order, with that vertex's label. Between any two vertices, one of the two oriented arcs has length at most three; it flips distinct coordinates, including the outgoing coordinate of its first vertex. That teacher distinguishes the endpoints, so the map is no-clashing. Any projection onto at most two coordinates is a subclass of a full two-cube and has NCTD at most one, by restricting the standard full-cube singleton assignment. The empty projection has one concept and value zero. Thus every projection has NCTD at most one. Since the original class has more than one concept, its NCTD, and hence projected NCTD, is at least one.
\end{proof}
\noindent\textbf{Reference.} Direct finite proof and independent solver-free verification in sync/verify\_projected\_nctd\_compression\_gap.py. The stable pair decoder is owned by the proper stable labeled full-cube value proof..

\subsection{No-Clashing Teaching Dimension of Antipodal Six-Code Class}\label{val:noclashing_teaching_dimension_antipodal_six_code}
\noindent\textbf{Statement.} The No-Clashing Teaching Dimension of Antipodal Six-Code Class is $1$.
\begin{proof}
The database-owned projected-NCTD value proof gives a complete singleton no-clashing assignment on this class. Width zero would give identical empty teachers to distinct concepts and cause a clash. Hence NCTD is one.
\end{proof}
\noindent\textbf{Reference.} Direct finite proof and independent solver-free verification in sync/verify\_projected\_nctd\_compression\_gap.py. The stable pair decoder is owned by the proper stable labeled full-cube value proof..

\subsection{Labeled Sample Compression of Antipodal Six-Code Class}\label{val:labeled_sample_compression_antipodal_six_code}
\noindent\textbf{Statement.} The Labeled Sample Compression of Antipodal Six-Code Class is $2$.
\begin{proof}
For the upper bound, restrict the database-owned stable pair scheme for the full four-cube to the realizable samples of this subclass. Its width is two; improper outputs are allowed. For the lower bound there are nine possible empty/singleton labeled keys and sixteen possible binary total outputs per key. An output of a key may without loss be required to extend that key: any used key must do so, and unused keys can be changed arbitrarily. Enumerate the 63 realizable partial samples. A decoder is feasible exactly when every sample has an available key whose output extends it. The solver-free checker sync/verify\_projected\_nctd\_compression\_gap.py exhausts this finite condition in two independent ways. First, keep each key's possible-output set. A sample is satisfied if one available key's entire remaining output set extends it; reject if no available key has any compatible output. If only one key can cover a sample, restrict its output set to the compatible words. Otherwise partition a key's output set into the words that cover a chosen unresolved sample and those that do not, and recurse on both parts. These are exhaustive disjoint branches, each strictly decreasing a finite domain. All 108 terminal branches are rejected (215 states, 107 splits, 824 forced domain restrictions). Thus no width-one decoder exists. Independently inject the six full concepts into six distinct available keys, since different full inputs require different reconstructed words. For every surviving injection enumerate all possible total outputs of the three remaining keys, pruning only when an uncovered sample has no available unassigned key. Literal-set availability checks reject all 656 surviving complete injections and their 17840 residual states. Neither checker relies on a solver UNSAT answer; both accept the full two-cube width-one scheme as a control. This proves the lower bound two, including improper reconstruction, with freely chosen keys and outputs. The same program constructs and directly checks width-one decoders after deletion of any one concept or any one coordinate. Thus the exhibited obstruction is minimal under either deletion operation, without any claim of global uniqueness.
\end{proof}
\noindent\textbf{Reference.} Direct finite proof and independent solver-free verification in sync/verify\_projected\_nctd\_compression\_gap.py. The stable pair decoder is owned by the proper stable labeled full-cube value proof..

\subsection{Stable Labeled Sample Compression of Antipodal Six-Code Class}\label{val:stable_labeled_sample_compression_antipodal_six_code}
\noindent\textbf{Statement.} The Stable Labeled Sample Compression of Antipodal Six-Code Class is $2$.
\begin{proof}
The ordinary labeled compression value is two, so the stable value is at least two. The database-owned stable pair scheme on the full four-cube has width two. Restrict its encoder to samples realizable by this subclass and leave its possibly improper decoder unchanged. Consistency and reconstruction stability remain valid on all such samples and their subsamples. This gives the matching stable upper bound.
\end{proof}
\noindent\textbf{Reference.} Direct finite proof and independent solver-free verification in sync/verify\_projected\_nctd\_compression\_gap.py. The stable pair decoder is owned by the proper stable labeled full-cube value proof..

\subsection{VC Dimension of Antipodal Six-Code Class}\label{val:vc_dimension_antipodal_six_code}
\noindent\textbf{Statement.} The VC Dimension of Antipodal Six-Code Class is $2$.
\begin{proof}
The first two coordinates realize all four labelings, so VC is at least two. Shattering three coordinates would require eight distinct concepts, whereas this class has six. Hence VC is exactly two.
\end{proof}
\noindent\textbf{Reference.} Direct finite proof and independent solver-free verification in sync/verify\_projected\_nctd\_compression\_gap.py. The stable pair decoder is owned by the proper stable labeled full-cube value proof..

\subsection{Proper Equivalence Queries Complexity of Trivial concepts}\label{val:proper_equivalence_queries_complexity_trivial_concepts}
\noindent\textbf{Statement.} The Proper Equivalence Queries Complexity of Trivial concepts is $1$.
\begin{proof}
Query either of the two concepts. A failed equivalence query has a counterexample, which identifies the other constant concept.
\end{proof}

\subsection{Proper Labeled Sample Compression of Chen–Cheng–Tang Teaching Products}\label{val:proper_labeled_sample_compression_chen_teaching_products}
\noindent\textbf{Statement.} The Proper Labeled Sample Compression of Chen–Cheng–Tang Teaching Products is $\le3n$.
\begin{proof}
For the twelve-coordinate base, the database-owned certificate sync/cct\_compression\_certificate.json specifies one reconstruction for each of the 2049 realizable labeled keys of size at most three. Each row $(m,y,h)$ encodes its support, positive labels and reconstructed concept by binary integers, with coordinate 1 as the least significant bit; thus $y\subseteq m$ in support notation. The solver-free checker sync/verify\_cct\_compression.py reconstructs all 100 concepts from the authoritative class definition, verifies $h$ is in that class and $h\cap m=y$ for every key, and independently verifies all 135073 realizable partial samples. Forward bit-mask enumeration checks that every sample contains a key whose output extends it. A second reverse literal-set enumeration takes, for each key and its output, every intermediate partial sample; their union is exactly the independently reconstructed sample family. Define the encoder by the first consistent key in the fixed order (key size, support mask, positive-label mask). This is a proper, message-free labeled compression scheme of size at most three. The certificate is the proof input; its discovery by a bounded SAT search is not relied on for correctness. For $n$ disjoint blocks, compress each block sample by this same base scheme and retain the union of its labeled keys. The decoder partitions the retained coordinates by their known block indices and applies the base decoder in each block, including the empty key in blocks with no retained examples. Every output block belongs to the base class and extends its whole block sample; hence the product output belongs to $\mathcal C^{\mathrm{CCT}}_n$ and is consistent. The total retained size is at most $3n$. This proves the upper bound for every $n$, without enumerating products or transmitting auxiliary block messages. The particular first-key encoder is not reconstruction-stable, as the verifier checks; this does not rule out a different stable scheme.
\end{proof}
\noindent\textbf{Reference.} Explicit finite certificate and independent replay in sync/cct\_compression\_certificate.json and sync/verify\_cct\_compression.py. The product argument is given above. No external compression theorem or novelty claim is used..

\subsection{Labeled Sample Compression of Chen–Cheng–Tang Teaching Products}\label{val:labeled_sample_compression_chen_teaching_products}
\noindent\textbf{Statement.} The Labeled Sample Compression of Chen–Cheng–Tang Teaching Products is $\le3n$.
\begin{proof}
The proper width-$3n$ scheme and its full certificate are in the companion value proof \ref{val:proper_labeled_sample_compression_chen_teaching_products}; allowing improper outputs cannot increase the optimum. For the lower bound there are $100^n$ distinct full concepts on $12n$ coordinates. Two different full concepts must compress to different labeled keys, because one key has only one reconstruction. If at most $n$ examples were retained, there would be at most $\sum_{i=0}^n2^i\binom{12n}{i}$ possible keys. For $t=1/24$, the inequality $t^i\ge t^n$ for $i\le n$ gives $$\sum_{i=0}^n2^i\binom{12n}{i}\le24^n(1+1/12)^{12n}=\bigl[24(13/12)^{12}\bigr]^n<100^n,$$ where the last step is the exact integer inequality $24\cdot13^{12}<100\cdot12^{12}$. This contradicts injectivity, so the integer compression size is at least $n+1$. The same lower bound applies to proper compression. At $n=1$ it rules out widths zero and one, but leaves the exact optimum at two or three. The verifier checks the constant inequality and twenty finite counting instances; the displayed argument proves the bound for all $n$.
\end{proof}
\noindent\textbf{Reference.} Database-owned proper decoder certificate and direct counting argument; see sync/verify\_cct\_compression.py..

\subsection{Order Sample Compression of Antipodal Six-Code Class}\label{val:order_sample_compression_antipodal_six_code}
\noindent\textbf{Statement.} The Order Sample Compression of Antipodal Six-Code Class is $2$.
\begin{proof}
Index binary words by integers $j\in\{0,\ldots,15\}$, with label at $x_i$ equal to bit $i$ of $j$ (least significant bit first, matching the displayed coordinate order). The full-cube decoding order is $\langle 0,8,12,14,10,2,3,4,6,7,15,5,1,13,11,9\rangle$. For any realizable sample $S$, let $q(S)$ be the first word in this order extending $S$. Direct finite enumeration of the class's 63 realizable partial samples shows that each has a contained labeled subset $U$ of size at most two with $q(U)=q(S)$. Compress to a smallest such subset, breaking ties by coordinate order, and decode by the same fixed order. This is an order scheme of width at most two; the subset excludes all hypotheses preceding $q(S)$ and is inclusion-minimal with that property. The verification enumerates every actual subset and checks the minimum in two independent bit-mask and literal-set representations; it includes the empty sample, all 24 three-coordinate samples and all six full concepts. All labels and coordinates are retained unchanged, with no auxiliary message. Some reconstructed words are outside the class, which is allowed for OSC. The natural integer order is rejected as a negative control. Conversely every order scheme is a labeled compression scheme, and the independently proved labeled compression value of this class is two. Therefore OSC is exactly two. The order was found by one bounded solver query, but the proof uses only the saved order and solver-free exhaustive replay, not a solver nonexistence verdict. No novelty claim is made.
\end{proof}
\noindent\textbf{Reference.} Direct finite certificate, independently replayed in sync/verify\_projected\_nctd\_compression\_gap.py. The general finite-order formulation is prop:order-compression-finite-order in the survey; the lower-bound proof is the database-owned Labeled Sample Compression value for this class..

\subsection{Proper Ordered Sample Compression of the Branching Batch-Tree Class}\label{val:proper_ordered_sample_compression_branching_batch_tree}
\noindent\textbf{Statement.} The Proper Ordered Sample Compression of Branching Batch-Tree Class is $b$.
\begin{proof}
Identify a concept with its positive support. The class is closed under binary intersections. Indeed, for two distinct leaf words take their common prefix, replace their first differing outgoing vectors $a,a'$ by the bitwise intersection $a\wedge a'$, and append zero vectors to obtain a leaf word of length $n$. Before the divergence the two concepts agree; at the divergence their pointwise intersection is the new vector; below it their paths are disjoint, so the intersection is zero everywhere. This remains valid when $a\wedge a'$ equals one of the two original vectors. Identical words intersect to themselves. Thus every nonempty finite intersection is again a concept. Apply Proposition~\ref{prop:intersection-closed-compression} of the survey (the standard intersection-closed scheme): retain a minimal positive spanning subset of the sample and reconstruct its closure. That proposition proves that its size is at most VC, that the closure is a proper consistent output even in the presence of negative examples, and that a linear extension of inclusion makes it the first compatible hypothesis both for the original sample and for the retained key. It includes the empty positive set; here the closure of that set is the zero concept. The established VC value is $b$, so this gives proper order compression at most $b$. Conversely proper order compression dominates ordinary order compression, which dominates VC. The value is therefore exactly $b$. This is an application of the standard scheme, not a new general theorem; its general proof is not duplicated here.
\end{proof}
\noindent\textbf{Reference.} \cite{darnstadt2016order}, Theorem 28 and Corollary 29, printed pages 13--14 of the cached August 26, 2014 author preprint. The family-specific closure argument and complete derivation above are database-owned; sync/verify\_branching\_batch\_trees.py independently checks intersections and replays the decoder on all samples of its six small-domain cases..

\subsection{Order Sample Compression of the Branching Batch-Tree Class}\label{val:order_sample_compression_branching_batch_tree}
\noindent\textbf{Statement.} The Order Sample Compression of Branching Batch-Tree Class is $b$.
\begin{proof}
The proper order scheme in the database-owned Proper Ordered Sample Compression value of this class has width $b$ and is also admissible for ordinary order compression. Conversely OSC dominates VC, whose value on this class is $b$ because the root batch is shattered. Thus $b\le\mathrm{OSC}(\mathcal T_{b,n})\le b$.
\end{proof}
\noindent\textbf{Reference.} The database-owned Proper Ordered Sample Compression and VC Dimension values of this class; \cite{darnstadt2016order}, Corollaries 25 and 29 of the cached August 26, 2014 author preprint..

\subsection{Recursive Teaching Dimension of the Branching Batch-Tree Class}\label{val:recursive_teaching_dimension_branching_batch_tree}
\noindent\textbf{Statement.} The Recursive Teaching Dimension of Branching Batch-Tree Class is $b$.
\begin{proof}
For every leaf word $w$ and each of the $b$ coordinates in its final visited batch, flipping that bit of the final outgoing vector gives another concept differing from $h_w$ at exactly that one coordinate. Any teaching set for $h_w$ in the full class must include all $b$ such coordinates. Therefore every possible first recursive teaching removal costs at least $b$. For the upper bound, RTD is at most OSC, and the database-owned Order Sample Compression value of this class is $b$. Thus RTD is exactly $b$. The lower argument concerns the full class, not an unproved general inequality between RTD and VC.
\end{proof}
\noindent\textbf{Reference.} Direct last-batch neighbour lower bound; the database-owned Order Sample Compression value of this class and \cite{darnstadt2016order}, Lemma 23 of the cached author preprint..

\subsection{Projected Recursive Teaching Dimension of the Branching Batch-Tree Class}\label{val:monotone_recursive_teaching_dimension_branching_batch_tree}
\noindent\textbf{Statement.} The Projected Recursive Teaching Dimension of Branching Batch-Tree Class is $b$.
\begin{proof}
The identity projection gives $\mathrm{RTD}^p(\mathcal T_{b,n})\ge\mathrm{RTD}(\mathcal T_{b,n})=b$, by the database-owned recursive teaching value. The universal order-compression bound gives $\mathrm{RTD}^p(\mathcal T_{b,n})\le\mathrm{OSC}(\mathcal T_{b,n})=b$, using the database-owned OSC value. Hence the projected value is exactly $b$. No assumption that an arbitrary projection remains a batch-tree class is needed.
\end{proof}
\noindent\textbf{Reference.} The database-owned Recursive Teaching Dimension and Order Sample Compression values of this class; \cite{darnstadt2016order}, Theorem 24 of the cached August 26, 2014 author preprint..

\subsection{Proper Stable Labeled Sample Compression of Singletons}\label{val:proper_stable_labeled_sample_compression_singletons}
\noindent\textbf{Statement.} The Proper Stable Labeled Sample Compression of Singletons is $\Theta(n)$.
\begin{proof}
The database-owned co-VC value of the singleton class is $n$. The proper stable compression lower bound in survey Proposition \ref{prop:proper-compression-covc} therefore gives $\mathrm{psLSC}\ge\lceil(n-1)/2\rceil$. The order-to-stable conversion in Proposition \ref{prop:order-stable-installation} gives $\mathrm{psLSC}\le\mathrm{pOSC}$, whose database-owned singleton value is $n-1$. These bounds prove linear asymptotic order. At $n=1$ both bounds give zero.
\end{proof}
\noindent\textbf{Reference.} The database-owned singleton co-VC and proper ordered compression values; survey Propositions \ref{prop:proper-compression-covc} and \ref{prop:order-stable-installation}..

\subsection{Balbach Teaching Dimension of Blockwise-Addressing}\label{val:balbach_teaching_dimension_blockwiseaddressing}
\noindent\textbf{Statement.} The Balbach Teaching Dimension of Blockwise-Addressing is $2$.
\begin{proof}
For every member of the Blockwise-Addressing family, the database-owned teaching and recursive teaching values are both two. The universal finite-class inequalities $\mathrm{RTD}\le\mathrm{BTD}\le\mathrm{TD}$ therefore give $2\le\mathrm{BTD}\le2$. This uses the same parameter range as the class definition; it does not require a new teaching construction.
\end{proof}
\noindent\textbf{Reference.} Database-owned endpoint proofs in Sections \ref{val:recursive_teaching_dimension_blockwiseaddressing} and \ref{val:teaching_dimension_blockwiseaddressing}. The universal comparisons are relationships 257 and 255; see \cite{zilles2011models,balbach2008measuring}..

\subsection{Projected Recursive Teaching Dimension of Antipodal Six-Code Class}\label{val:monotone_recursive_teaching_dimension_antipodal_six_code}
\noindent\textbf{Statement.} The Projected Recursive Teaching Dimension of Antipodal Six-Code Class is $2$.
\begin{proof}
The database-owned VC and order-compression values of this six-concept class are both two. For finite binary classes, $\mathrm{VC}\le\mathrm{RTD}^{p}\le\mathrm{OSC}$, where the superscript denotes maximization over coordinate projections. Consequently $2\le\mathrm{RTD}^{p}\le2$. The upper endpoint uses the saved full-cube order and its independent exhaustive replay; no new optimal-decoder search or general VC compression bound is asserted.
\end{proof}
\noindent\textbf{Reference.} Database-owned endpoint proofs in Sections \ref{val:vc_dimension_antipodal_six_code} and \ref{val:order_sample_compression_antipodal_six_code}. Relationships 118 and 117; survey Lemma \ref{lem:projected-rtd-vc} and \cite{darnstadt2016order}..

\subsection{No-Clashing Teaching Dimension of Full Cube}\label{val:noclashing_teaching_dimension_full_cube}
\noindent\textbf{Statement.} The No-Clashing Teaching Dimension of Full Cube is $\lceil n/2\rceil$.
\begin{proof}
Theorem 19 gives $\mathrm{NCTD}(2^{[n]})=\lceil n/2\rceil$.
\end{proof}
\noindent\textbf{Reference.} \cite{KSZ2019}.

\subsection{Projected Minimum Degree of the Chen Neighbor Puncture}\label{val:projected_minimum_degree_chen_neighbor_puncture}
\noindent\textbf{Statement.} The Projected Minimum Degree of Chen Neighbor Puncture is $4$.
\begin{proof}
The database-owned endpoint proofs establish $\mathrm{VC}=\mathrm{RTD}^{p}=4$ on this particular 109-concept class. A shattered four-set projects to the four-cube, giving $\mathrm d_{\min}^{p}\ge4$. Conversely Corollary \ref{cor:projected-rtd-projected-min-degree} gives $\mathrm d_{\min}^{p}\le\mathrm{RTD}^{p}=4$. Thus the projected minimum degree is exactly four. The upper endpoint depends on the existing independently replayed finite teaching-plan certificate, not an unverified general equality with VC.
\end{proof}
\noindent\textbf{Reference.} Database-owned endpoint proofs in Sections \ref{val:vc_dimension_chen_neighbor_puncture} and \ref{val:monotone_recursive_teaching_dimension_chen_neighbor_puncture}. Relationships 78 and 498; survey Corollary \ref{cor:projected-rtd-projected-min-degree}. Certificate replay: sync/verify\_chen\_neighbor\_puncture.py..

\subsection{Antichain Number of Singletons plus Empty Set}\label{val:antichain_number_singletons_plus_empty_set}
\noindent\textbf{Statement.} The Antichain Number of Singletons plus empty set is $1$.
\begin{proof}
Assign the singleton concept $\{i\}$ the labeled sample $\{(i,1)\}$, and assign the empty concept $\{(1,0)\}$. Every assigned sample is consistent with its target. All $n+1$ samples are distinct singletons as sets of labeled examples, including the two samples using coordinate one, since their labels differ. Distinct singleton sets are inclusion-incomparable, so this is an antichain teacher of width one. Width zero would assign the same empty sample to the class's at least two distinct concepts, violating incomparability. Thus the antichain number is exactly one.
\end{proof}
\noindent\textbf{Reference.} Direct proof from the database definitions; finite regression in sync/verify\_noclashing\_antichain\_gap.py..

\subsection{Antichain Number of Full Cube}\label{val:antichain_number_full_cube}
\noindent\textbf{Statement.} The Antichain Number of Full Cube is $\min\{k\in\{0,\ldots,n\}:2^k\binom nk\ge2^n\}$.
\begin{proof}
Put $M_j=2^j\binom nj$, the number of labeled samples of size $j$. For the lower bound, any antichain teacher has $2^n$ distinct assigned samples. Choose a uniformly random full binary word and an independent uniformly random coordinate permutation; successively exposing its labels gives a chain of samples containing a fixed size-$j$ sample with probability $1/M_j$. An antichain meets each chain at most once, so its samples $F$ satisfy $\sum_{S\in F}1/M_{|S|}\le1$. If every assigned sample has size at most $k$ and every $M_j$ for $j\le k$ is smaller than $2^n$, this sum is greater than one, a contradiction. Thus the width is at least the displayed minimum. For the upper bound, choose any $k$ satisfying $M_k\ge2^n$. Form a bipartite graph between the full concepts and their consistent size-$k$ samples. Each concept has degree $\binom nk$ and each sample extends to $2^{n-k}$ concepts. For every concept set $A$, counting its incident edges gives $|N(A)|2^{n-k}\ge |A|\binom nk$, whence $|N(A)|\ge|A|$. Hall's theorem gives a matching assigning every concept a distinct consistent size-$k$ sample. Equal-sized distinct samples are incomparable, so this is an antichain teacher. This proves the exact formula, including the one-concept case n=0. For the linear lower bound, the binomial theorem gives $\binom nj\le(5/4)^n4^j$. If $j\le n/5$, then $M_j\le(5/4)^n8^{n/5}<2^n$, since $(5/4)^5 8<2^5$. Hence the minimum is greater than n/5 for n>0; retaining each full concept gives the upper bound n.
\end{proof}
\noindent\textbf{Reference.} Direct antichain-chain counting and Hall matching proof for Definition 32 of \cite{simon2023map}; independent small-cube matching and chain regression in sync/verify\_noclashing\_antichain\_gap.py. No novelty claim is made..

\subsection{Antichain Number of Private-Coordinate Cubes}\label{val:antichain_number_private_coordinate_cube}
\noindent\textbf{Statement.} The Antichain Number of Private-Coordinate Cubes is $1$.
\begin{proof}
For every n>=1, assign concept $h_v$ its private positive example $\{(p_v,1)\}$. These are distinct consistent singleton samples, hence pairwise inclusion-incomparable. The class has $2^n\ge2$ concepts, so assigning the empty sample everywhere cannot satisfy incomparability. Thus AN is exactly one. Projection onto the original n coordinates gives the full n-cube, whose database-owned antichain value grows linearly in n. Consequently projection can increase antichain number by an unbounded factor.
\end{proof}
\noindent\textbf{Reference.} Direct proof from the database class definition; the projected value is proved in Section \ref{val:antichain_number_full_cube}. Finite regression: sync/verify\_noclashing\_antichain\_gap.py..

\subsection{Antichain Number of Repeated-Coordinate Cubes}\label{val:antichain_number_repeated_coordinate_cube}
\noindent\textbf{Statement.} The Antichain Number of Repeated-Coordinate Cubes is $1$.
\begin{proof}
For target $h_v$, put $w=(v_2,\ldots,v_n)$ and assign the sample $\{(z_w,v_1)\}$. It is consistent since every $z_w$ reads the first bit. Distinct bit strings $v$ give distinct labeled singleton samples: either their suffixes choose different coordinate copies, or their first bits give different labels at the same copy. Hence the samples form an antichain of width one. There are $2^n\ge2$ concepts, so width zero is impossible. This proves AN exactly one, including $n=1$.
\end{proof}
\noindent\textbf{Reference.} Local definition-driven proof; bounded independent replay in sync/verify\_noclashing\_antichain\_gap.py..

\subsection{Subset Teaching Dimension of Repeated-Coordinate Cubes}\label{val:subset_teaching_dimension_repeated_coordinate_cube}
\noindent\textbf{Statement.} The Subset Teaching Dimension of Repeated-Coordinate Cubes is $n$.
\begin{proof}
Partition the domain into the first-bit group $\{z_w:w\in W_n\}$ and the singleton groups $\{x_j\}$ for $2\le j\le n$. Each target has independent alternatives obtained by flipping any one underlying bit. Therefore every ordinary teaching sample must meet all n groups; choosing one coordinate from each group is sufficient. Its ordinary minimum teachers are exactly all such choices with the target's labels. A sample can be contained in one of these teachers only if it selects at most one coordinate from each group. If it omits a group, both assignments of that group's underlying bit extend it to ordinary minimum teachers, so it cannot identify a unique target in the first subset-teaching iteration. If it meets every group, it is already one of the target's ordinary minimum teachers and identifies that target uniquely. Thus the first iteration has value n and exactly the same family of minimum samples as stage zero. Induction leaves every later stage unchanged, proving STD exactly n. This verifies the catalogue's iteration over all minimum teachers, not just a chosen teaching map.
\end{proof}
\noindent\textbf{Reference.} Local definition-driven proof; bounded independent replay in sync/verify\_noclashing\_antichain\_gap.py..

\subsection{No-Clashing Teaching Dimension of Repeated-Coordinate Cubes}\label{val:noclashing_teaching_dimension_repeated_coordinate_cube}
\noindent\textbf{Statement.} The No-Clashing Teaching Dimension of Repeated-Coordinate Cubes is $\lceil n/2\rceil$.
\begin{proof}
Map every $z_w$ to the first cube coordinate and each $x_j$ to coordinate j. Replacing the coordinates in any target-consistent sample by these images, and removing repeated labeled examples, preserves consistency and inconsistency with every concept under $h_v\leftrightarrow v$. Thus a no-clashing teacher map on $\mathcal R_n$ folds to a no-clashing map on the n-cube of no larger width. Conversely choose one fixed copy $z_w$ to lift any cube teacher; again all consistency relations are preserved and its width is unchanged. The two no-clashing dimensions are therefore equal. The database-owned cube value in Section \ref{val:noclashing_teaching_dimension_full_cube} is $\lceil n/2\rceil$, giving the assertion for every n>=1.
\end{proof}
\noindent\textbf{Reference.} Local definition-driven proof; bounded independent replay in sync/verify\_noclashing\_antichain\_gap.py..

\subsection{Teaching Dimension of Repeated-Coordinate Cubes}\label{val:teaching_dimension_repeated_coordinate_cube}
\noindent\textbf{Statement.} The Teaching Dimension of Repeated-Coordinate Cubes is $n$.
\begin{proof}
As established in the database-owned subset-teaching proof for this family, every target has ordinary minimum teaching size n: each of the n independent bit groups must be represented, and one coordinate per group suffices. Taking the maximum over targets gives TD exactly n. See Section \ref{val:subset_teaching_dimension_repeated_coordinate_cube}.
\end{proof}
\noindent\textbf{Reference.} Local definition-driven proof; bounded independent replay in sync/verify\_noclashing\_antichain\_gap.py..

\subsection{Recursive Teaching Dimension of Repeated-Coordinate Cubes}\label{val:recursive_teaching_dimension_repeated_coordinate_cube}
\noindent\textbf{Statement.} The Recursive Teaching Dimension of Repeated-Coordinate Cubes is $n$.
\begin{proof}
Every target of the initial class has ordinary minimum teaching size n, by Section \ref{val:subset_teaching_dimension_repeated_coordinate_cube}. Any recursive teaching order must therefore pay n at its first removal. Conversely every remaining subclass is taught using one fixed coordinate per underlying bit, a sample of size n. Hence RTD is exactly n.
\end{proof}
\noindent\textbf{Reference.} Local definition-driven proof; bounded independent replay in sync/verify\_noclashing\_antichain\_gap.py..

\subsection{VC Dimension of Repeated-Coordinate Cubes}\label{val:vc_dimension_repeated_coordinate_cube}
\noindent\textbf{Statement.} The VC Dimension of Repeated-Coordinate Cubes is $n$.
\begin{proof}
One fixed copy of the first bit together with $x_2,\ldots,x_n$ realizes the full n-cube, so VC is at least n. The class has exactly $2^n$ distinct concepts; shattering d coordinates requires $2^d$ distinct traces, so d<=n. Thus VC equals n.
\end{proof}
\noindent\textbf{Reference.} Local definition-driven proof; bounded independent replay in sync/verify\_noclashing\_antichain\_gap.py..

\subsection{Log Size Ratio of Full Cube}\label{val:log_size_ratio_full_cube}
\noindent\textbf{Statement.} The Log Size Ratio of Full Cube is $\frac{\log(2^n-1)}{\log(n+1)}$.
\begin{proof}
The full cube has $2^n$ concepts on a domain of size $n$, so substitution into the definition gives $\log(2^n-1)/\log(n+1)$.
\end{proof}

\subsection{Size of Repeated-Coordinate Cubes}\label{val:size_repeated_coordinate_cube}
\noindent\textbf{Statement.} The Size of Repeated-Coordinate Cubes is $2^n$.
\begin{proof}
The map $v\mapsto h_v$ is injective: each underlying bit is read by at least one coordinate. Its domain is the full n-cube, of cardinality $2^n$, so the class has exactly $2^n$ concepts.
\end{proof}
\noindent\textbf{Reference.} Local definition-driven proof; bounded independent replay in sync/verify\_noclashing\_antichain\_gap.py..

\subsection{Effective Range of Repeated-Coordinate Cubes}\label{val:effective_range_repeated_coordinate_cube}
\noindent\textbf{Statement.} The Effective Range of Repeated-Coordinate Cubes is $n-1+2^{n-1}$.
\begin{proof}
There are n-1 singleton coordinate groups and $2^{n-1}$ copies of the first bit. Every one of these coordinates takes both labels because its underlying bit varies freely. Thus all domain coordinates are effective, giving effective range $n-1+2^{n-1}$. Equal columns still correspond to distinct domain coordinates.
\end{proof}
\noindent\textbf{Reference.} Local definition-driven proof; bounded independent replay in sync/verify\_noclashing\_antichain\_gap.py..

\subsection{Subset Teaching Dimension of Half-intervals}\label{val:subset_teaching_dimension_halfintervals}
\noindent\textbf{Statement.} The Subset Teaching Dimension of Half-intervals is $\begin{cases}0&n=1,\\1&n\ge2.\end{cases}$.
\begin{proof}
For $n=1$ there is only one target, taught by the empty sample. Let $n\ge2$ and write $h_i=[i]$. The unique ordinary minimum sample of $h_1$ is $\{(2,0)\}$, that of $h_n$ is $\{(n,1)\}$, and for $1<i<n$ it is $\{(i,1),(i+1,0)\}$. Adjacent targets force these examples and they suffice to exclude every other target. The constant first coordinate cannot appear in a minimum sample. Every labeled example in this list occurs in the minimum sample of exactly one target. At the first subset-teaching iteration, each interior target therefore has both of its singleton subsamples as minimizers, while the endpoint samples remain unchanged. These pairwise distinct singleton samples form a fixed point: every proper subset is empty and belongs to all targets. Thus STD is one. This is a direct all-n calculation, not a consequence of subclass monotonicity.
\end{proof}

\subsection{Subset Teaching Dimension of Addressing}\label{val:subset_teaching_dimension_addressing}
\noindent\textbf{Statement.} The Subset Teaching Dimension of Addressing is $\begin{cases}0&n=1,\\1&n\ge2.\end{cases}$.
\begin{proof}
Each target has a private positive address coordinate, which teaches it with one ordinary example. For $n\ge2$, all ordinary minimum samples consequently have size one, and no labeled singleton sample can teach two different targets. At the next subset-teaching step these singleton families remain unchanged: an empty sample is contained in every target's family and cannot identify just one. Hence the process is fixed at value one. For $n=1$, the empty sample teaches the sole target and STD is zero.
\end{proof}

\subsection{Subset Teaching Dimension of Private-Coordinate Cubes}\label{val:subset_teaching_dimension_private_coordinate_cube}
\noindent\textbf{Statement.} The Subset Teaching Dimension of Private-Coordinate Cubes is $1$.
\begin{proof}
For every $n\ge1$, the private positive example $(p_v,1)$ teaches $h_v$ among the $2^n$ concepts. Thus every ordinary minimum teacher has size one, and its singleton sample belongs to only one target's initial minimum family. These families are already a subset-teaching fixed point, since their only proper subsample is empty and occurs for every target. There are at least two targets, so value zero is impossible.
\end{proof}

\subsection{Subset Teaching Dimension of Cube–Half-interval Product}\label{val:subset_teaching_dimension_cube_halfinterval_product}
\noindent\textbf{Statement.} The Subset Teaching Dimension of Cube–Half-interval Product is $m+1$.
\begin{proof}
For $m\ge1,n\ge2$, the two factors have STD $m$ and one, respectively, by the full-cube and half-interval calculations. Proposition \ref{prop:subset-teaching-product} proves that STD is additive for nonempty finite classes on disjoint domains: at every iteration, the full family of minimum samples for a product target consists exactly of unions of factor minimizers, because their containment sets factor as a Cartesian product. Applying that identity gives $m+1$. The nonempty initial-segment convention, including its constant first coordinate, is the one used in the half-interval value proof.
\end{proof}

\subsection{Balbach Teaching Dimension of Private-Coordinate Cubes}\label{val:balbach_teaching_dimension_private_coordinate_cube}
\noindent\textbf{Statement.} The Balbach Teaching Dimension of Private-Coordinate Cubes is $1$.
\begin{proof}
For every $n\ge1$, each target $h_v$ is ordinarily taught by its private positive example $(p_v,1)$. The class has $2^n\ge2$ targets, so every ordinary minimum teaching size is exactly one. At a size-one sample the Balbach filter retains all targets; the private example still isolates its target. The size-zero sample is consistent with and retains all targets, so it cannot identify one. Thus the initial all-one vector is already fixed and BTD is exactly one.
\end{proof}

\subsection{Balbach Teaching Dimension of Half-intervals}\label{val:balbach_teaching_dimension_halfintervals}
\noindent\textbf{Statement.} The Balbach Teaching Dimension of Half-intervals is $\min\{2,n-1\}$.
\begin{proof}
For $n=1$ the class is a singleton, with value zero. For $n=2$ both concepts have ordinary teaching size one, and this vector is fixed. For $n\ge3$, the two endpoint concepts $[1]$ and $[n]$ have ordinary teaching size one, using $(2,0)$ and $(n,1)$ respectively. Each interior $[i]$, $1<i<n$, needs two examples: $(i,1),(i+1,0)$ teach it, while any one consistent positive or negative example leaves a larger or smaller initial segment. All initial values are at least one, so at sample size one the filter deletes no competitor; no interior value can drop to one. The size-zero sample retains all concepts. Each original minimum sample retains its target under this unchanged vector. Thus the initial vector is fixed and the maximum is two.
\end{proof}

\subsection{Balbach Teaching Dimension of Common-Core Cube}\label{val:balbach_teaching_dimension_common_core_cube}
\noindent\textbf{Statement.} The Balbach Teaching Dimension of Common-Core Cube is $n$.
\begin{proof}
The class is a full $n$-cube with an additional constant positive coordinate. Proposition \ref{prop:balbach-monotonicity} proves invariance under adjoining constant coordinates, at every iterate and with the sample-size filter retained. The full cube has BTD $n$, giving the result. Equivalently every target has ordinary teaching size $n$: all variable bits must be specified, and the constant coordinate distinguishes none. This uniform vector is fixed by the Balbach recursion.
\end{proof}

\subsection{Balbach Teaching Dimension of Repeated-Coordinate Cubes}\label{val:balbach_teaching_dimension_repeated_coordinate_cube}
\noindent\textbf{Statement.} The Balbach Teaching Dimension of Repeated-Coordinate Cubes is $n$.
\begin{proof}
For every $n\ge1$, an ordinary teacher must query each of the $n-1$ nonrepeated bits and at least one copy of the repeated bit. Otherwise the concept obtained by flipping an unqueried underlying bit is still consistent. Those $n$ queries suffice, so every target has ordinary minimum teaching size $n$. For any sample of size at most $n$, the initial Balbach filter retains every consistent concept, since all their thresholds equal $n$. No sample of size below $n$ can therefore teach its target, whereas its ordinary size-$n$ sample does. The uniform initial vector is fixed, proving BTD $n$ despite the repeated coordinates.
\end{proof}

\subsection{Log Size Ratio of Singletons}\label{val:log_size_ratio_singletons}
\noindent\textbf{Statement.} The Log Size Ratio of Singletons is $\frac{\log(n-1)}{\log(n+1)}$.
\begin{proof}
The singleton class has $n$ concepts on a domain of size $n$, hence the definition gives $\log(n-1)/\log(n+1)$.
\end{proof}

\subsection{Balbach Teaching Dimension of Cube–Half-interval Product}\label{val:balbach_teaching_dimension_cube_halfinterval_product}
\noindent\textbf{Statement.} The Balbach Teaching Dimension of Cube–Half-interval Product is $m+\min\{2,n-1\}$.
\begin{proof}
For $m\ge1,n\ge2$, Proposition \ref{prop:balbach-cube-product} gives BTD $m+\mathrm{BTD}(\mathcal I_n)$: every admissible sample must query all $m$ independent cube coordinates, after which the size filter is exactly that of the half-interval factor. The half-interval value proof in Section \ref{val:balbach_teaching_dimension_halfintervals} gives $\mathrm{BTD}(\mathcal I_n)=\min\{2,n-1\}$. Thus the value is $m+1$ for $n=2$ and $m+2$ for $n\ge3$. The half-interval factor consists of nonempty initial segments, retaining its constant first coordinate.
\end{proof}

\subsection{Effective Range of Balanced Six-Split Class}\label{val:effective_range_balanced_six_split}
\noindent\textbf{Statement.} The Effective Range of Balanced Six-Split Class is $10$.
\begin{proof}
The domain consists of the triples containing 1, so it has $\binom52=10$ coordinates. Each triple gives label one to exactly three of the six concepts and label zero to the other three; every coordinate is active.
\end{proof}

\subsection{Size of Balanced Six-Split Class}\label{val:size_balanced_six_split}
\noindent\textbf{Statement.} The Size of Balanced Six-Split Class is $6$.
\begin{proof}
There are six indexed concepts. Any two vertices of $[6]$ can be placed on opposite sides of a balanced split, so their concepts differ. Thus all six concepts are distinct.
\end{proof}

\subsection{Proper Equivalence Queries Complexity of Balanced Six-Split Class}\label{val:proper_equivalence_queries_complexity_balanced_six_split}
\noindent\textbf{Statement.} The Proper Equivalence Queries Complexity of Balanced Six-Split Class is $3$.
\begin{proof}
We count queries until the target is identified, with no final confirmation required. After any first rejected proper proposal, the counterexample leaves exactly three candidate vertices, since every coordinate is a balanced split. Querying candidates successively then identifies the target in at most two more queries. Conversely the adversary can reject the first query and keep three candidates $V$. A proposal outside $V$ can be answered by the split $V\mid([6]\setminus V)$, retaining all three. For a proposal $v\in V$, choose a balanced split with $V\setminus\{v\}$ and one vertex outside $V$ on one side; the reply leaves the other two candidates. They cannot yet be identified, so at least one further query is necessary. Thus the minimax number is three, even when all original class members are allowed as proposals.
\end{proof}

\subsection{Proper Unlabeled Sample Compression of Balanced Six-Split Class}\label{val:proper_unlabeled_sample_compression_balanced_six_split}
\noindent\textbf{Statement.} The Proper Unlabeled Sample Compression of Balanced Six-Split Class is $+\infty$.
\begin{proof}
Apply survey Proposition \ref{prop:proper-unlabeled-pair-obstruction}. Distinct coordinate triples $A,B$ both contain 1 and have intersection size one or two. The four cells $A\cap B$, $A\setminus B$, $B\setminus A$ and $[6]\setminus(A\cup B)$ are all nonempty, so each coordinate pair is shattered. For two distinct vertices $u,v$, the number of unordered balanced splits separating them is $\binom42=6$: orient such a split by its side containing $u$, exclude $v$, and choose the other two vertices. Our domain contains one representative of each such split, so every two distinct concepts have Hamming distance six. The proposition would force some distance at most $(10+1)/2=5.5$, a contradiction. Hence pUSC is $+\infty$ under the explicit nonexistence convention.
\end{proof}

\subsection{Proper Stable Sample Compression of Balanced Six-Split Class}\label{val:proper_stable_sample_compression_balanced_six_split}
\noindent\textbf{Statement.} The Proper Stable Sample Compression of Balanced Six-Split Class is $+\infty$.
\begin{proof}
Section \ref{val:proper_unlabeled_sample_compression_balanced_six_split} proves that no proper message-free unlabeled scheme exists, regardless of size. A proper reconstruction-stable scheme would in particular be such a scheme. Thus psUSC is also $+\infty$.
\end{proof}

\subsection{VC Dimension of Balanced Six-Split Class}\label{val:vc_dimension_balanced_six_split}
\noindent\textbf{Statement.} The VC Dimension of Balanced Six-Split Class is $2$.
\begin{proof}
Every pair of coordinates is shattered, as verified in Section \ref{val:proper_unlabeled_sample_compression_balanced_six_split}. No triple can be shattered because the class has only six concepts, fewer than the eight distinct traces required. Thus VC dimension is two.
\end{proof}

\subsection{Littlestone Dimension of Balanced Six-Split Class}\label{val:littlestone_dimension_balanced_six_split}
\noindent\textbf{Statement.} The Littlestone Dimension of Balanced Six-Split Class is $2$.
\begin{proof}
A shattered coordinate pair supplies a depth-two Littlestone tree. A depth-three complete tree requires eight distinct concepts realizing its eight paths, whereas this class has six. Thus Littlestone dimension equals two.
\end{proof}

\subsection{Stable Unlabeled Complexity of Balanced Six-Split Class}\label{val:stable_unlabeled_complexity_balanced_six_split}
\noindent\textbf{Statement.} The Stable Unlabeled Complexity of Balanced Six-Split Class is $2$.
\begin{proof}
A shattered pair forces unlabeled compression size at least two by counting its four possible labelings. Survey Proposition \ref{prop:one-pass-stable-unlabeled} gives a possibly improper reconstruction-stable scheme of size at most the Littlestone dimension, which is two by Section \ref{val:littlestone_dimension_balanced_six_split}. Hence sUSC equals two; proper reconstruction is essential to the nonexistence result.
\end{proof}

\subsection{Proper Ordered Sample Compression of Singletons plus Empty Set}\label{val:proper_ordered_sample_compression_singletons_plus_empty_set}
\noindent\textbf{Statement.} The Proper Ordered Sample Compression of Singletons plus empty set is $1$.
\begin{proof}
Order the empty concept first, followed by the singleton concepts in any order. A realizable sample with no positive label has the empty concept as its first consistent hypothesis and compresses to the empty sample. Otherwise its unique positive example is at some i; retaining just this labeled example makes the singleton {i} the first (and only) consistent hypothesis. The decoder depends only on that retained labeled key. Hence the scheme is proper and ordered with width one. Width zero cannot decode both full samples of the empty concept and a singleton, so the value is exactly one for n>=1.
\end{proof}
\noindent\textbf{Reference.} Direct proper-order construction; replayed by sync/verify\_upper\_branch\_comparisons.py..

\subsection{Log Size Ratio of Trivial concepts}\label{val:log_size_ratio_trivial_concepts}
\noindent\textbf{Statement.} The Log Size Ratio of Trivial concepts is $0$.
\begin{proof}
The class has two concepts, so $\log(|\mathcal H|-1)=\log 1=0$.
\end{proof}

\subsection{Combinatorial Eluder Dimension of the Reverse Eluder--Range Separation}\label{val:combinatorial_eluder_dimension_eluder_range_reverse_separation}
\noindent\textbf{Statement.} The Combinatorial Eluder Dimension of Reverse Eluder--Range Separation Class is $3$.
\begin{proof}
Anchor at 000 and query $x_1$,$x_2$,$x_3$ in that order, with witness concepts 110,011,001 respectively. Each witness differs at its own queried coordinate and agrees with 000 at every earlier coordinate. This is an eluder sequence of length three. No coordinate can repeat in such a sequence, so the three-coordinate domain also gives the matching upper bound.
\end{proof}
\noindent\textbf{Reference.} Direct sequence; independently replayed by sync/verify\_upper\_branch\_comparisons.py..

\subsection{Projected Distinguishing Range of the Reverse Eluder--Range Separation}\label{val:projected_distinguishing_range_eluder_range_reverse_separation}
\noindent\textbf{Statement.} The Projected Distinguishing Range of Reverse Eluder--Range Separation Class is $2$.
\begin{proof}
On coordinates $x_2$,$x_3$ the four concepts have the distinct traces 00,01,11,10. Thus two coordinates distinguish the full class, whereas one binary coordinate cannot distinguish four concepts. Its distinguishing range is two. Every proper projection has at most two coordinates and therefore distinguishing range at most two; the full projection attains two. Maximizing over all projections proves the value.
\end{proof}
\noindent\textbf{Reference.} Direct projection argument; independently replayed by sync/verify\_upper\_branch\_comparisons.py..

\subsection{Projected Distinguishing Range of the Paired-Code Family}\label{val:projected_distinguishing_range_paired_code_separation}
\noindent\textbf{Statement.} The Projected Distinguishing Range of Paired-Code Separation Family is $2^n$.
\begin{proof}
Write N=$2^n$. For every i in [N], the class contains both $w_i$ and $w_i$ XOR $e_i$, which are distinct and differ only at coordinate i. Therefore every coordinate set distinguishing all concepts must contain i. All N coordinates are necessary and suffice, so the full-domain distinguishing range equals N. The full domain is an allowed projection, and no projection has more than N coordinates, proving that the maximum over projections is also N. This argument does not require different pairs to have disjoint sets of concepts.
\end{proof}
\noindent\textbf{Reference.} Direct essential-coordinate argument; admissibility and existence are proved in survey prop:paired-code-eluder..

\subsection{Combinatorial Eluder Dimension of the Paired-Code Family}\label{val:combinatorial_eluder_dimension_paired_code_separation}
\noindent\textbf{Statement.} The Combinatorial Eluder Dimension of Paired-Code Separation Family is $\Theta(n)$.
\begin{proof}
The upper bound is Proposition \ref{prop:paired-code-eluder}: any eluder sequence of length 2m would contain an admissibility-forbidden subsequence with m=4n+2 distinct nonanchor pairs, so Edim is at most 2m-1=8n+3. For the lower bound, Section \ref{val:size_paired_code_separation} gives size at least $2^n$+1. The established general inequality Edim >= $log_2$(size), together with integrality, gives Edim >= n+1.
\end{proof}
\noindent\textbf{Reference.} Survey prop:paired-code-eluder and the established eluder/log-size inequality (relationship 265)..

\subsection{Size of the Paired-Code Family}\label{val:size_paired_code_separation}
\noindent\textbf{Statement.} The Size of Paired-Code Separation Family is $\Theta(2^n)$.
\begin{proof}
There are N=$2^n$ pairs, giving at most 2N distinct words. For the lower bound, regard the words as vectors over $F_2$. For each i, the difference of the two words in pair i is $e_i$. Thus the span of all pairwise differences of class members contains all N standard basis vectors. Fixing one class member $h_0$, this same span is generated by the vectors h XOR $h_0$ over all other class members h, since h XOR g=(h XOR $h_0$)+(g XOR $h_0$). A spanning set of an N-dimensional space has at least N members. Consequently the class has at least N+1 members.
\end{proof}
\noindent\textbf{Reference.} Direct finite-field dimension argument..

\subsection{Effective Range of the Paired-Code Family}\label{val:effective_range_paired_code_separation}
\noindent\textbf{Statement.} The Effective Range of Paired-Code Separation Family is $2^n$.
\begin{proof}
The ambient domain has N=$2^n$ coordinates. Each coordinate i varies on the pair $w_i$,$w_i$ XOR $e_i$ contained in the class. Thus every domain coordinate contributes to the effective range, which equals N.
\end{proof}
\noindent\textbf{Reference.} Direct coordinate-pair argument..

\subsection{Combinatorial Eluder Dimension of the Triangular-Code Family}\label{val:combinatorial_eluder_dimension_triangular_code_separation}
\noindent\textbf{Statement.} The Combinatorial Eluder Dimension of Triangular-Code Separation Family is $2^n$.
\begin{proof}
Anchor at the all-zero word. Query coordinates 1,...,N in that order and use $h_i$ as the witness at coordinate i. Its diagonal bit is one and its earlier bits are all zero, so this is an eluder sequence of length N=$2^n$. Coordinates cannot repeat in an eluder sequence, because a later witness must agree with the anchor at each earlier coordinate. Thus the N-coordinate domain gives the matching upper bound.
\end{proof}
\noindent\textbf{Reference.} Direct diagonal sequence; admissibility and existence are proved in survey prop:triangular-code-range..

\subsection{Projected Distinguishing Range of the Triangular-Code Family}\label{val:projected_distinguishing_range_triangular_code_separation}
\noindent\textbf{Statement.} The Projected Distinguishing Range of Triangular-Code Separation Family is $\Theta(n)$.
\begin{proof}
Proposition \ref{prop:triangular-code-range} proves the upper bound 6n+5 using the absence of an essential projection of size 6n+6 and the general characterization in Lemma \ref{lem:essential-range-projections}. For the lower bound, Section \ref{val:size_triangular_code_separation} gives exactly $2^n$+1 distinct concepts. Any d binary coordinates have at most $2^d$ traces, so distinguishing this full class requires d>=ceil($log_2$($2^n$+1))=n+1. The full projection is included in the maximum defining projected distinguishing range.
\end{proof}
\noindent\textbf{Reference.} Survey prop:triangular-code-range and direct binary trace counting..

\subsection{Size of the Triangular-Code Family}\label{val:size_triangular_code_separation}
\noindent\textbf{Statement.} The Size of Triangular-Code Separation Family is $2^n+1$.
\begin{proof}
There are N=$2^n$ triangular rows and the zero word. For i<j, coordinate i separates $h_i$ from $h_j$, since their values are one and zero. Each $h_i$ also differs from the zero word at its diagonal coordinate i. Hence all N+1 displayed words are distinct.
\end{proof}
\noindent\textbf{Reference.} Direct triangular-matrix argument..

\subsection{Effective Range of the Triangular-Code Family}\label{val:effective_range_triangular_code_separation}
\noindent\textbf{Statement.} The Effective Range of Triangular-Code Separation Family is $2^n$.
\begin{proof}
The domain has N=$2^n$ coordinates. Coordinate i takes value one on $h_i$ and zero on the all-zero word, so every domain coordinate is varying and contributes to the effective range.
\end{proof}
\noindent\textbf{Reference.} Direct diagonal-coordinate argument..

\subsection{Log Size of Addressing}\label{val:log_size_addressing}
\noindent\textbf{Statement.} The Log Size of Addressing is $\log_2 n$.
\begin{proof}
The class has exactly n distinct indexed concepts. Substitution into the definition of Log Size gives $\log_2n$, including zero at n=1.
\end{proof}
\noindent\textbf{Reference.} \cite{maass1992lower}.

\subsection{Yang Dimension of Parity Functions}\label{val:yang_dimension_parity_functions}
\noindent\textbf{Statement.} The Yang dimension of Parity Functions is $2n$.
\begin{proof}
The evaluation map $a\mapsto h_a$ is linear and injective, since the $2n$ coordinate-basis input vectors recover all coefficients. Its image is therefore a binary linear class of dimension $2n$. Proposition \ref{prop:affine-class-yang} gives Yang dimension exactly $2n$ over every coefficient field, including the catalogue's default characteristic two.
\end{proof}
\noindent\textbf{Reference.} Survey prop:affine-class-yang..

\subsection{Monotonic Yang Dimension of Parity Functions}\label{val:monotonic_yang_dimension_parity_functions}
\noindent\textbf{Statement.} The Monotonic Yang Dimension of Parity Functions is $2n$.
\begin{proof}
As in Section \ref{val:yang_dimension_parity_functions}, this is a binary linear class of dimension 2n. Proposition \ref{prop:affine-class-yang} bounds Yang dimension on each arbitrary subclass by 2n, through the affine-span closure of every version-space ambiguity complex. The whole class has Yang dimension 2n, so the maximum over subclasses is exactly 2n.
\end{proof}
\noindent\textbf{Reference.} Survey prop:affine-class-yang..

\subsection{Sign Rank of Parity Functions}\label{val:sign_rank_parity_functions}
\noindent\textbf{Statement.} The Sign Rank of Parity Functions is $\ge2^n$.
\begin{proof}
Let $N=4^n$ and set $W_{a,x}=(-1)^{a\cdot x}$. For $a\ne b$, choose a coordinate on which $a+b$ is nonzero and pair each $x$ with its translate by that basis vector; the two character products cancel. Thus $WW^{\mathsf T}=NI$, and its spectral norm is $\sqrt N$. Forster's bound gives $\mathrm{SR}(W)\ge N/\lVert W\rVert=\sqrt N=2^n$. The catalogue incidence sign matrix differs by at most a global sign, which preserves sign rank. The matrix itself has rank $N$, giving the upper bound $4^n$. No matching exponential rate is claimed.
\end{proof}
\noindent\textbf{Reference.} \cite{forster2002linear}; the spectral consequence is also explicit in \cite[Theorem 13 and Section A.2]{alon2016sign}..

\subsection{Star Number of Parity Functions}\label{val:star_number_parity_functions}
\noindent\textbf{Statement.} The Star number of Parity Functions is $2n$.
\begin{proof}
On the $2n$ coordinate-basis input vectors, $h_0$ and the functions indexed by those basis vectors give a realized $2n$-star. Conversely a star's spoke witnesses, queried in any order, form an eluder sequence anchored at its centre. Section \ref{val:combinatorial_eluder_dimension_parity_functions} proves eluder dimension $2n$, giving the matching upper bound.
\end{proof}
\noindent\textbf{Reference.} Direct basis construction and the database-owned parity eluder proof..

\subsection{Combinatorial Eluder Dimension of Parity Functions}\label{val:combinatorial_eluder_dimension_parity_functions}
\noindent\textbf{Statement.} The Combinatorial Eluder Dimension of Parity Functions is $2n$.
\begin{proof}
Fix an anchor $h_a$. At step $i$, an eluder witness $h_b$ agrees on earlier queried vectors $x_j$ and differs on $x_i$, so $(a+b)\cdot x_j=0$ for $j<i$ and $(a+b)\cdot x_i=1$. Therefore $x_i$ is not in the binary linear span of the earlier vectors. There are at most $2n$ such steps. The $2n$ basis vectors, anchored at $h_0$ with the corresponding basis-indexed witnesses, attain this length.
\end{proof}
\noindent\textbf{Reference.} Direct linear-independence proof..

\subsection{VC Dimension of Parity Functions}\label{val:vc_dimension_parity_functions}
\noindent\textbf{Statement.} The VC Dimension of Parity Functions is $2n$.
\begin{proof}
The $2n$ coordinate-basis input vectors are shattered, because their labels are precisely the arbitrary coefficients $a_i$. A shattered set of $t$ points requires at least $2^t$ distinct functions, while the class has $4^n=2^{2n}$ functions. Hence $t\le2n$.
\end{proof}
\noindent\textbf{Reference.} Direct shattering and counting..

\subsection{Size of Parity Functions}\label{val:size_parity_functions}
\noindent\textbf{Statement.} The Size of Parity Functions is $4^n$.
\begin{proof}
There are $4^n$ coefficient vectors $a$. Distinct vectors differ at some coordinate $i$, so their functions differ at the basis input $e_i$. Thus all $4^n$ functions are distinct.
\end{proof}
\noindent\textbf{Reference.} Direct coefficient enumeration and injectivity..

\subsection{Effective Range of Parity Functions}\label{val:effective_range_parity_functions}
\noindent\textbf{Statement.} The Effective Range of Parity Functions is $4^n-1$.
\begin{proof}
The input zero is identically labeled zero. At every nonzero input $x$, $h_0$ has label zero and a coefficient-basis function supported on a nonzero coordinate of $x$ has label one. Thus exactly the $4^n-1$ nonzero inputs vary.
\end{proof}
\noindent\textbf{Reference.} Direct evaluation..

\subsection{Proper Stable Sample Compression of Singletons}\label{val:proper_stable_sample_compression_singletons}
\noindent\textbf{Statement.} The Proper Stable Sample Compression of Singletons is $\Theta(n)$.
\begin{proof}
Identify each singleton with its index in $[n]$. For $n=1$ use the empty key and the sole concept. For $n\ge2$, fix $\rho(\varnothing)=1$, $\rho(\{i\})=i$ for $i\ge2$, and $\rho([r-1])=r$ for $2\le r\le n$; these prescriptions do not conflict. Send unused keys to 1. If a sample contains its positive point at 1, retain nothing; if its positive point is $i\ge2$, retain $\{i\}$. An all-negative realizable sample on $A\subsetneq[n]$ has a least missing point $r$; retain $[r-1]\subseteq A$ and reconstruct singleton $r$. Each reconstruction is proper and consistent, and the key has size at most $n-1$. In the positive case $i\ge2$, every intermediate sample retaining the key still has positive point $i$. In the positive-at-1 case every subsample either retains that positive or has only negatives away from 1, so still selects the empty key. In the all-negative case an intermediate sample containing $[r-1]$ still omits $r$, so its least missing point remains $r$. Thus the key, and hence reconstruction, is stable on every interval. Forgetting the retained labels in a labeled decoder simulates any unlabeled scheme, so psUSC is at least psLSC; Section \ref{val:proper_stable_labeled_sample_compression_singletons} gives the lower bound $\lceil(n-1)/2\rceil$. The exact query value $n-1$ is recorded in Section \ref{val:proper_equivalence_queries_complexity_singletons}.
\end{proof}
\noindent\textbf{Reference.} Explicit prefix decoder; the database-owned proper stable labeled compression lower bound and proper-query value..

\subsection{Proper Equivalence Queries Complexity of K4 Incidence with Empty Set}\label{val:proper_equivalence_queries_complexity_k4_incidence_with_empty}
\noindent\textbf{Statement.} The Proper Equivalence Queries Complexity of K4 Incidence with Empty Set is $2$.
\begin{proof}
Query $h_0$. Acceptance finishes; a counterexample edge $uv$ leaves exactly $h_u,h_v$. One further proposal identifies the target, either by acceptance or by elimination, without a final confirmation query. One query cannot suffice: if the first proposal is $h_0$, any edge leaves two targets. If it is $h_a$, an incident edge $ab$ can be returned negatively, leaving $h_0,h_c,h_d$. Vertex symmetry covers the other proposals. Thus the optimum is two.
\end{proof}

\subsection{Log Size of Majority}\label{val:log_size_majority}
\noindent\textbf{Statement.} The Log Size of Majority is $n$.
\begin{proof}
The class has exactly $2^n$ concepts, so its binary logarithmic size is $n$.
\end{proof}
\noindent\textbf{Reference.} \cite{maass1992lower}.

\subsection{Proper Unlabeled Sample Compression of K4 Incidence with Empty Set}\label{val:proper_unlabeled_sample_compression_k4_incidence_with_empty}
\noindent\textbf{Statement.} The Proper Unlabeled Sample Compression of K4 Incidence with Empty Set is $3$.
\begin{proof}
For the lower bound suppose a proper width-two decoder $\rho$ exists, and write $d_e=\rho(\{e\})$. Every pair of adjacent edges is shattered: its common endpoint supplies 11, its two other endpoints supply 10 and 01, and $h_0$ supplies 00. Consequently the four keys on this pair must decode to four distinct pair traces. If $\rho(\varnothing)=h_0$, a positive singleton forces $d_{uv}\in\{h_u,h_v\}$. Orient each edge toward this selected endpoint. Two edges cannot both point toward their common endpoint, since their singleton outputs would coincide. Six edges would therefore have to enter four vertices with indegree at most one, a contradiction. Otherwise, by vertex symmetry put $\rho(\varnothing)=h_a$. On the triangle $bc,cd,db$, singleton outputs again select endpoints with indegree at most one, so their orientations form a directed cycle. Relabel $b,c,d$ so that $d_{bc}=h_c,d_{cd}=h_d,d_{db}=h_b$. The singleton sample opposite to $h_a$ on $ab$ permits only $h_0,h_c,h_d$; comparing keys on the shattered pair $ab,bc$ excludes $h_c$. Cyclically, $d_{ab}\in\{h_0,h_d\}$, $d_{ac}\in\{h_0,h_b\}$ and $d_{ad}\in\{h_0,h_c\}$. On $ab,ac$, the empty output has trace 11 and $d_{ab}$ has trace 00; distinctness therefore forces $d_{ac}=h_b$. The pairs $ac,ad$ and $ad,ab$ similarly force $d_{ad}=h_c$ and $d_{ab}=h_d$, respectively. On each pair of the three spokes $ab,ac,ad$, these empty and singleton outputs already include trace 00. Thus its pair key cannot decode to $h_0$. Every key of size at most two within the three spokes therefore decodes to a nonzero vertex star, whereas the all-negative sample on those spokes is realizable only by $h_0$. This contradicts coverage. Hence pUSC is at least three.

For a stable upper bound first remove $h_0$, leaving the four vertex stars. Flip every coordinate by its value under $h_a$, and identify the three opposite-edge pairs $\{ac,bd\}$, $\{ad,bc\}$, $\{ab,cd\}$ as coordinate types 1, 2, 3. The vertex stars become the even-parity words $000,110,011,101$. Apply the exact-key-stable cyclic parity scheme of survey Proposition \ref{prop:parity-stable-unlabeled}, pulled back by choosing the first observed copy of every selected type. Its width is two and it never selects both edges of an opposite pair. Reserve all three opposite-edge pairs and all four three-edge stars as new keys decoding to $h_0$; none is used by that parity encoder. A sample realizable by $h_0$ but by no vertex star is all-negative on an edge set covering all four vertices. Every such edge cover contains either an opposite-edge pair or a three-edge star: without disjoint edges, a pairwise-intersecting edge family covering four vertices must be a star. Each reserved key, with all-negative labels, already excludes every vertex star. The reserved-key extension lemma \ref{lem:stable-reserved-key-extension} now supplies a proper stable scheme of width three. Explicitly, use the old encoder on vertex-realizable samples, and otherwise choose the first contained reserved key in a fixed order. Intermediate samples retaining a new key still exclude every vertex star and keep that same first key. The old branch retains its exact-key stability. This proves pUSC at most three, with a stable upper scheme.
\end{proof}
\noindent\textbf{Reference.} Direct pair-trace lower argument; survey Proposition \ref{prop:parity-stable-unlabeled} and Lemma \ref{lem:stable-reserved-key-extension}..

\subsection{Proper Stable Sample Compression of K4 Incidence with Empty Set}\label{val:proper_stable_sample_compression_k4_incidence_with_empty}
\noindent\textbf{Statement.} The Proper Stable Sample Compression of K4 Incidence with Empty Set is $3$.
\begin{proof}
Section \ref{val:proper_unlabeled_sample_compression_k4_incidence_with_empty} proves that every proper unlabeled scheme has width at least three and constructs a proper width-three scheme preserving the selected key on every intermediate subsample. Its lower argument does not assume stability, while its upper scheme satisfies a stronger condition than reconstruction stability. Therefore psUSC is exactly three.
\end{proof}

\subsection{VC Dimension of K4 Incidence with Empty Set}\label{val:vc_dimension_k4_incidence_with_empty}
\noindent\textbf{Statement.} The VC Dimension of K4 Incidence with Empty Set is $2$.
\begin{proof}
Two adjacent edges $ab,ac$ realize all four patterns: $h_a$ gives 11, $h_b$ gives 10, $h_c$ gives 01, and $h_0$ gives 00. No three coordinates can be shattered because the class has only five concepts, fewer than $2^3=8$.
\end{proof}

\subsection{Littlestone Dimension of K4 Incidence with Empty Set}\label{val:littlestone_dimension_k4_incidence_with_empty}
\noindent\textbf{Statement.} The Littlestone Dimension of K4 Incidence with Empty Set is $2$.
\begin{proof}
The adjacent-edge shattered pair supplies a depth-two Littlestone tree. A shattered binary tree of depth three would require at least eight distinct concepts, one for each leaf: two different paths disagree at their first split. The class has only five concepts.
\end{proof}

\subsection{Stable Unlabeled Complexity of K4 Incidence with Empty Set}\label{val:stable_unlabeled_complexity_k4_incidence_with_empty}
\noindent\textbf{Statement.} The Stable Unlabeled Complexity of K4 Incidence with Empty Set is $2$.
\begin{proof}
The one-pass stable unlabeled construction in survey Proposition \ref{prop:one-pass-stable-unlabeled} bounds sUSC by Littlestone dimension, which is two on this class. Conversely a shattered pair has four labelings but a width-one unlabeled encoder can use only its empty and two singleton keys, so width one cannot cover all four labelings. Improper reconstruction is permitted for this parameter.
\end{proof}

\subsection{Effective Range of K4 Incidence with Empty Set}\label{val:effective_range_k4_incidence_with_empty}
\noindent\textbf{Statement.} The Effective Range of K4 Incidence with Empty Set is $6$.
\begin{proof}
The domain contains the six edges of $K_4$. Every edge is positive in its two endpoint stars and negative in $h_0$, so all six coordinates vary.
\end{proof}

\subsection{Size of K4 Incidence with Empty Set}\label{val:size_k4_incidence_with_empty}
\noindent\textbf{Statement.} The Size of K4 Incidence with Empty Set is $5$.
\begin{proof}
The four vertex stars are distinct: for $u\ne v$, an edge from $u$ to a third vertex belongs to the star of $u$ but not to that of $v$. Each star is nonempty, so adjoining the empty concept gives exactly five concepts.
\end{proof}

\subsection{Log Size of Tagged Singletons}\label{val:log_size_tagged_singletons}
\noindent\textbf{Statement.} The Log Size of Tagged Singletons is $\log(n+\lceil\log n\rceil+1)$.
\begin{proof}
Taking the binary logarithm of the exact class size $n+\lceil\log n\rceil+1$ gives a quantity of order $\log n$.
\end{proof}
\noindent\textbf{Reference.} \cite{maass1992lower}.

\subsection{Log Size of Half-intervals}\label{val:log_size_halfintervals}
\noindent\textbf{Statement.} The Log Size of Half-intervals is $\log n$.
\begin{proof}
This is the binary logarithm of the $n$ initial segments.
\end{proof}

\subsection{Log Size of Singletons plus empty set}\label{val:log_size_singletons_plus_empty_set}
\noindent\textbf{Statement.} The Log Size of Singletons plus empty set is $\log(n+1)$.
\begin{proof}
This is the binary logarithm of the class size $n+1$.
\end{proof}

\subsection{Log Size of Full Cube}\label{val:log_size_full_cube}
\noindent\textbf{Statement.} The Log Size of Full Cube is $n$.
\begin{proof}
The class has $2^n$ concepts, hence binary logarithmic size $n$.
\end{proof}

\subsection{Log Size of Singletons}\label{val:log_size_singletons}
\noindent\textbf{Statement.} The Log Size of Singletons is $\log n$.
\begin{proof}
The class has $n$ concepts, so its binary logarithmic size is $\log_2 n$.
\end{proof}

\subsection{Log Size of Trivial concepts}\label{val:log_size_trivial_concepts}
\noindent\textbf{Statement.} The Log Size of Trivial concepts is $1$.
\begin{proof}
The class has two concepts, so its binary logarithmic size is $\log_2 2=1$.
\end{proof}

\subsection{Threshold Dimension of Addressing}\label{val:threshold_dimension_addressing}
\noindent\textbf{Statement.} The Threshold Dimension of Addressing is $\log_2 n+1\quad(n=2^r,\ r\ge1)$.
\begin{proof}
Restrict here to $n=2^r$ with all r-bit words present, so $\lfloor\log_2n\rfloor=\log_2n=r$. Projecting to the $d=\log_2 n$ address coordinates gives the full $d$-cube.  Its chain $\emptyset,\{a_1\},\ldots,\{a_1,\ldots,a_d\}$ is realized by the address concepts; adding the primary coordinate of the all-one address concept gives one further realized trace.  Hence the threshold dimension is at least $d+1$.  Conversely, the final trace of a threshold chain is contained in one addressing concept and so uses at most its $d$ address coordinates and its single primary coordinate.  Thus no chain has dimension exceeding $d+1$.
\end{proof}
\noindent\textbf{Reference.} \cite{maass1992lower}.

\subsection{Threshold Dimension of Majority}\label{val:threshold_dimension_majority}
\noindent\textbf{Statement.} The Threshold Dimension of Majority is $n$.
\begin{proof}
Projecting to the first $n$ coordinates gives the full cube, since every base subset occurs with its uniquely determined majority bit. Hence an arbitrary length-$n$ threshold chain is present, and no chain can use more than the $n$ projected coordinates.
\end{proof}
\noindent\textbf{Reference.} \cite{maass1992lower}.

\subsection{Threshold Dimension of Half-intervals}\label{val:threshold_dimension_halfintervals}
\noindent\textbf{Statement.} The Threshold Dimension of Half-intervals is $n-1$.
\begin{proof}
The chain of initial segments from $[1]$ through $[n]$ gives a threshold sequence on $n-1$ successive changes. A longer threshold sequence would require more than the $n$ distinct concepts in the class.
\end{proof}

\subsection{Oriented Diameter of Half-intervals}\label{val:oriented_diameter_halfintervals}
\noindent\textbf{Statement.} The Oriented Diameter of Half-intervals is $n-1$.
\begin{proof}
Restrict the two endpoint concepts $[1]$ and $[n]$ to coordinates 2 through $n$; their traces are the empty and full traces, respectively, on $n-1$ coordinates. No oriented pair can differ on more than the diameter $n-1$.
\end{proof}

\subsection{VC Dimension of Blockwise-Addressing}\label{val:vc_dimension_blockwiseaddressing}
\noindent\textbf{Statement.} The VC Dimension of Blockwise-Addressing is $d$.
\begin{proof}
Each block of d coordinates realizes all $2^d$ binary patterns among the $f_i$, so the VC dimension is at least d.  The blockwise-addressing construction is the standard separating family whose VC dimension is d; the matching upper bound is proved in the cited source.
\end{proof}
\noindent\textbf{Reference.} \cite{ben1997online}, p. 54.

\subsection{Recursive Teaching Dimension of Blockwise-Addressing}\label{val:recursive_teaching_dimension_blockwiseaddressing}
\noindent\textbf{Statement.} The Recursive Teaching Dimension of Blockwise-Addressing is $2$.
\begin{proof}
For $f_i$, the labeled pair $(z,0),(x_i,1)$ distinguishes it from every $f_j$ and $g_j$; symmetrically $(z,1),(y_i,1)$ teaches $g_i$.  The same bound holds in every subfamily, so $RTD\le2$.  Conversely, in the full class no one-coordinate sample identifies a concept: at $z$ there are many concepts with each label; at a private coordinate a concept from the opposite family can match; and at a block coordinate every bit occurs for many members of the same family.  Thus $TS_{\min}\ge2$, hence $RTD\ge2$.
\end{proof}
\noindent\textbf{Reference.} \cite{ben1997online}, Example 2; direct teaching-set calculation.

\subsection{Diameter of Singletons plus empty set}\label{val:diameter_singletons_plus_empty_set}
\noindent\textbf{Statement.} The Diameter of Singletons plus empty set is $2$.
\begin{proof}
Two distinct singleton concepts have symmetric difference of size two. No two concepts in this class can differ on more than two coordinates.
\end{proof}

\subsection{Hamming Radius of Singletons plus empty set}\label{val:hamming_radius_singletons_plus_empty_set}
\noindent\textbf{Statement.} The Hamming Radius of Singletons plus empty set is $1$.
\begin{proof}
The Hamming ball of radius one centred at the empty set contains the empty set and every singleton. Radius zero cannot contain this nonconstant class.
\end{proof}

\subsection{Path Dimension of Singletons plus empty set}\label{val:path_dimension_singletons_plus_empty_set}
\noindent\textbf{Statement.} The Path Dimension of Singletons plus empty set is $2$.
\begin{proof}
For distinct $i,j$, start at the singleton at $i$, flip $i$ to obtain the empty set, then flip $j$ to obtain its singleton: this is a path of length two. A third distinct coordinate would produce a doubleton after the next flip, which is absent from the class.
\end{proof}

\subsection{Log Shattering of Singletons plus empty set}\label{val:log_shattering_singletons_plus_empty_set}
\noindent\textbf{Statement.} The Log Shattering of Singletons plus empty set is $\log(n+1)$.
\begin{proof}
The shattered sets are exactly the empty set and the $n$ singleton coordinate sets. No two-point set is shattered because the doubleton trace is missing, giving $n+1$ shattered sets.
\end{proof}

\subsection{Diameter of Half-intervals}\label{val:diameter_halfintervals}
\noindent\textbf{Statement.} The Diameter of Half-intervals is $n-1$.
\begin{proof}
The endpoint concepts $[1]$ and $[n]$ differ exactly on coordinates 2 through $n$. Every pair of nested initial segments differs on at most those $n-1$ coordinates.
\end{proof}

\subsection{Hamming Radius of Half-intervals}\label{val:hamming_radius_halfintervals}
\noindent\textbf{Statement.} The Hamming Radius of Half-intervals is $\lceil(n-1)/2\rceil$.
\begin{proof}
The endpoint concepts are at distance $n-1$, so every enclosing Hamming ball has radius at least $\lceil(n-1)/2\rceil$. An initial segment with a middle-sized endpoint attains this radius for every member of the chain.
\end{proof}

\subsection{Path Dimension of Half-intervals}\label{val:path_dimension_halfintervals}
\noindent\textbf{Statement.} The Path Dimension of Half-intervals is $n-1$.
\begin{proof}
Successive initial segments from $[1]$ to $[n]$ form a path of length $n-1$, flipping one new coordinate at each step. A longer path would require more than the $n$ concepts in the class.
\end{proof}

\subsection{Log Shattering of Half-intervals}\label{val:log_shattering_halfintervals}
\noindent\textbf{Statement.} The Log Shattering of Half-intervals is $\log n$.
\begin{proof}
The empty coordinate set is shattered, as is each singleton coordinate set from 2 through $n$. No two-point set is shattered by a chain, and coordinate 1 is constant. Thus there are exactly $n$ shattered sets.
\end{proof}

\subsection{VC Dimension of Half-intervals}\label{val:vc_dimension_halfintervals}
\noindent\textbf{Statement.} The VC Dimension of Half-intervals is $1$.
\begin{proof}
Any coordinate other than 1 is realized both positively and negatively by the chain. A two-point set cannot be shattered because the traces of a chain are totally ordered and therefore omit one incomparable pattern.
\end{proof}

\subsection{Diameter of Full Cube}\label{val:diameter_full_cube}
\noindent\textbf{Statement.} The Diameter of Full Cube is $n$.
\begin{proof}
The empty and full subsets differ on all $n$ coordinates. No pair of subsets of an $n$-element domain can differ on more coordinates.
\end{proof}

\subsection{Hamming Radius of Full Cube}\label{val:hamming_radius_full_cube}
\noindent\textbf{Statement.} The Hamming Radius of Full Cube is $n$.
\begin{proof}
For every proposed centre, its complement is a cube vertex at Hamming distance $n$. Thus no smaller enclosing ball exists, and radius $n$ trivially contains the cube.
\end{proof}

\subsection{Path Dimension of Full Cube}\label{val:path_dimension_full_cube}
\noindent\textbf{Statement.} The Path Dimension of Full Cube is $n$.
\begin{proof}
Starting with the empty subset and adding the coordinates in any order gives a path of length $n$. A path dimension cannot exceed the number of domain coordinates.
\end{proof}

\subsection{Log Shattering of Full Cube}\label{val:log_shattering_full_cube}
\noindent\textbf{Statement.} The Log Shattering of Full Cube is $n$.
\begin{proof}
Every subset of the $n$-element domain is shattered by the full cube, so there are $2^n$ shattered sets and their binary logarithm is $n$.
\end{proof}

\subsection{Effective VC Radius of Full Cube}\label{val:effective_vc_radius_full_cube}
\noindent\textbf{Statement.} The Effective VC Radius of Full Cube is $n$.
\begin{proof}
The effective range is the full domain and that entire domain is shattered by the cube. Hence the largest permitted radius is $n$.
\end{proof}

\subsection{Diameter of Singletons}\label{val:diameter_singletons}
\noindent\textbf{Statement.} The Diameter of Singletons is $0$ for $n=1$; $2$ for $n\ge2$.
\begin{proof}
For $n=1$ the class has one concept, so its diameter is zero. For $n\ge2$, every pair of distinct singleton concepts differs on exactly its two respective coordinates; such a pair exists, giving diameter two.
\end{proof}

\subsection{Hamming Radius of Singletons}\label{val:hamming_radius_singletons}
\noindent\textbf{Statement.} The Hamming Radius of Singletons is $0$ for $n=1$; $1$ for $n\ge2$.
\begin{proof}
For $n=1$ use the sole concept as the center of a radius-zero ball. For $n\ge2$, the empty center has distance one to every singleton, while no radius-zero ball can contain two distinct concepts. The radius is therefore one.
\end{proof}

\subsection{Path Dimension of Singletons}\label{val:path_dimension_singletons}
\noindent\textbf{Statement.} The Path Dimension of Singletons is $\min\{n-1,2\}$.
\begin{proof}
For $n=1$ there is only one concept and no nonempty path. For $n=2$, a one-coordinate projection has an edge, while the full projection consists of two nonadjacent singleton traces, so the value is one. For $n\ge3$, project to two coordinates $i,j$, omitting another coordinate. The traces $\{i\},\emptyset,\{j\}$ form a path obtained by flipping $i$ and then $j$, giving dimension two. Every proper projection consists of the empty trace and singleton traces; its one-inclusion graph is a star, whose longest simple path has two edges. The full projection is edgeless. Thus no longer coordinate-flip path exists.
\end{proof}

\subsection{VC Dimension of Majority}\label{val:vc_dimension_majority}
\noindent\textbf{Statement.} The VC Dimension of Majority is $n$.
\begin{proof}
Projection to the first $n$ coordinates realizes every binary pattern, so those coordinates are shattered. The class has only $2^n$ concepts and therefore cannot shatter $n+1$ coordinates.
\end{proof}
\noindent\textbf{Reference.} \cite{maass1992lower}.

\subsection{Diameter of Majority}\label{val:diameter_majority}
\noindent\textbf{Statement.} The Diameter of Majority is $n+1$.
\begin{proof}
The concepts associated with the empty and full base subsets differ on every base coordinate and on the majority coordinate. This is the full $n+1$-coordinate diameter.
\end{proof}
\noindent\textbf{Reference.} \cite{maass1992lower}.

\subsection{Extended Teaching Dimension of Trivial concepts}\label{val:extended_teaching_dimension_trivial_concepts}
\noindent\textbf{Statement.} The Extended Teaching Dimension of Trivial concepts is $1$.
\begin{proof}
At any one coordinate, the queried label selects exactly one of the two constant concepts.  Thus one coordinate specifies every target labeling with respect to the class.  No empty specifying set can do so, since both constants would remain consistent.  Hence the extended teaching dimension is $1$.
\end{proof}
\noindent\textbf{Reference.} \cite{hegedHus1995generalized}.

\subsection{Extended Teaching Dimension of Singletons}\label{val:extended_teaching_dimension_singletons}
\noindent\textbf{Statement.} The Extended Teaching Dimension of Singletons is $n-1$.
\begin{proof}
For the all-zero target, querying $n-1$ zero coordinates leaves only the singleton on the remaining coordinate; fewer queries leave at least two singletons.  Any target having a positive coordinate is specified by that coordinate, while a target with at least two positive coordinates is excluded by querying two of them.  Thus, for $n\ge2$, the maximum is $n-1$.
\end{proof}
\noindent\textbf{Reference.} \cite{hegedHus1995generalized}.

\subsection{Extended Teaching Dimension of Full Cube}\label{val:extended_teaching_dimension_full_cube}
\noindent\textbf{Statement.} The Extended Teaching Dimension of Full Cube is $n$.
\begin{proof}
Every labeling is itself a concept of the full cube.  Specifying it uniquely requires querying every one of the $n$ coordinates, and all $n$ coordinates suffice.
\end{proof}
\noindent\textbf{Reference.} \cite{hegedHus1995generalized}.

\subsection{Extended Teaching Dimension of Singletons plus empty set}\label{val:extended_teaching_dimension_singletons_plus_empty_set}
\noindent\textbf{Statement.} The Extended Teaching Dimension of Singletons plus empty set is $n$.
\begin{proof}
For the all-zero target, every coordinate must be queried to eliminate each singleton and leave only $\emptyset$.  Querying all coordinates specifies every target, so the extended teaching dimension is $n$.
\end{proof}
\noindent\textbf{Reference.} \cite{hegedHus1995generalized}.

\subsection{Membership Query Complexity of Full Cube}\label{val:membership_query_complexity_full_cube}
\noindent\textbf{Statement.} The Membership Query Complexity of Full Cube is $n$.
\begin{proof}
Querying all $n$ coordinates identifies the target.  If any coordinate is unqueried, two concepts differing only there remain consistent with every answer, so $n$ queries are necessary.
\end{proof}

\subsection{Membership Query Complexity of Singletons plus empty set}\label{val:membership_query_complexity_singletons_plus_empty_set}
\noindent\textbf{Statement.} The Membership Query Complexity of Singletons plus empty set is $n$.
\begin{proof}
For the empty target, every coordinate must be queried negatively to rule out its singleton; hence $n$ queries are necessary.  Querying all coordinates always identifies the target.
\end{proof}

\subsection{Teaching Dimension of Full Cube}\label{val:teaching_dimension_full_cube}
\noindent\textbf{Statement.} The Teaching Dimension of Full Cube is $n$.
\begin{proof}
For any target subset, every coordinate must be included in a teaching set: changing an omitted coordinate gives a distinct full-cube concept consistent with the sample.  All $n$ coordinates suffice.
\end{proof}

\subsection{Teaching Dimension of Singletons plus empty set}\label{val:teaching_dimension_singletons_plus_empty_set}
\noindent\textbf{Statement.} The Teaching Dimension of Singletons plus empty set is $n$.
\begin{proof}
The empty concept requires a negative example at every coordinate to exclude every singleton, so it has teaching-set size $n$.  Every singleton is taught by its positive coordinate; therefore the maximum teaching-set size is $n$.
\end{proof}

\subsection{Teaching Dimension of Half-intervals}\label{val:teaching_dimension_halfintervals}
\noindent\textbf{Statement.} The Teaching Dimension of Half-intervals is $2$.
\begin{proof}
For an interior initial segment $[i]$, the positive example at $i$ and negative example at $i+1$ identify it.  Neither one alone identifies an interior segment, since its adjacent segment remains consistent.  The two endpoint segments need only one example, so for $n\ge3$ the teaching dimension is $2$.
\end{proof}

\subsection{Minimum Teaching Set Size of Full Cube}\label{val:minimum_teaching_set_size_full_cube}
\noindent\textbf{Statement.} The Minimum Teaching Set Size of Full Cube is $n$.
\begin{proof}
Every concept in the full cube has a neighbor differing on each coordinate.  Hence any teaching set must include all $n$ coordinates, and the full labeled concept is a teaching set.
\end{proof}

\subsection{Minimum Teaching Set Size of Singletons plus empty set}\label{val:minimum_teaching_set_size_singletons_plus_empty_set}
\noindent\textbf{Statement.} The Minimum Teaching Set Size of Singletons plus empty set is $1$.
\begin{proof}
Any singleton is uniquely identified by its positive coordinate, while no concept is identified by the empty sample.  Thus the minimum teaching-set size is $1$.
\end{proof}

\subsection{Minimum Teaching Set Size of Half-intervals}\label{val:minimum_teaching_set_size_halfintervals}
\noindent\textbf{Statement.} The Minimum Teaching Set Size of Half-intervals is $1$.
\begin{proof}
The endpoint segment $[1]$ is uniquely identified by the negative label at coordinate $2$ (and $[n]$ by the positive label at $n$).  No nontrivial class has a zero-size teaching set.
\end{proof}

\subsection{Maximum Teaching Set Size of Trivial concepts}\label{val:maximum_teaching_set_size_trivial_concepts}
\noindent\textbf{Statement.} The Maximum Teaching Set Size of Trivial concepts is $1$.
\begin{proof}
By definition $\mathrm{TS}_{\max}=\mathrm{TD}$, whose value on the two trivial concepts is $1$.
\end{proof}

\subsection{Maximum Teaching Set Size of Singletons}\label{val:maximum_teaching_set_size_singletons}
\noindent\textbf{Statement.} The Maximum Teaching Set Size of Singletons is $1$.
\begin{proof}
By definition $\mathrm{TS}_{\max}=\mathrm{TD}$, and each singleton is taught by its unique positive coordinate.
\end{proof}

\subsection{Maximum Teaching Set Size of Full Cube}\label{val:maximum_teaching_set_size_full_cube}
\noindent\textbf{Statement.} The Maximum Teaching Set Size of Full Cube is $n$.
\begin{proof}
By definition $\mathrm{TS}_{\max}=\mathrm{TD}$, and every full-cube concept needs all $n$ coordinates to be taught.
\end{proof}

\subsection{Maximum Teaching Set Size of Singletons plus empty set}\label{val:maximum_teaching_set_size_singletons_plus_empty_set}
\noindent\textbf{Statement.} The Maximum Teaching Set Size of Singletons plus empty set is $n$.
\begin{proof}
By definition $\mathrm{TS}_{\max}=\mathrm{TD}$; the empty concept needs all $n$ negative coordinates to exclude the singletons.
\end{proof}

\subsection{Maximum Teaching Set Size of Half-intervals}\label{val:maximum_teaching_set_size_halfintervals}
\noindent\textbf{Statement.} The Maximum Teaching Set Size of Half-intervals is $2$.
\begin{proof}
By definition $\mathrm{TS}_{\max}=\mathrm{TD}$; an interior initial segment needs its final positive and next negative coordinates.
\end{proof}

\subsection{Positive No-Clashing Teaching Dimension of Full Cube}\label{val:positive_noclashing_teaching_dimension_full_cube}
\noindent\textbf{Statement.} The Positive No-Clashing Teaching Dimension of Full Cube is $n$.
\begin{proof}
Theorem 19 gives $\mathrm{NCTD}^+(2^{[n]})=n$.
\end{proof}
\noindent\textbf{Reference.} \cite{KSZ2019}.

\subsection{No-Clashing Teaching Dimension of Singletons}\label{val:noclashing_teaching_dimension_singletons}
\noindent\textbf{Statement.} The No-Clashing Teaching Dimension of Singletons is $1$.
\begin{proof}
Map the singleton $\{i\}$ to its positive example $(i,1)$.  This is non-clashing, so the value is at most $1$.  A size-$0$ mapping cannot distinguish two distinct concepts, so the value is exactly $1$.
\end{proof}

\subsection{Positive No-Clashing Teaching Dimension of Singletons}\label{val:positive_noclashing_teaching_dimension_singletons}
\noindent\textbf{Statement.} The Positive No-Clashing Teaching Dimension of Singletons is $1$.
\begin{proof}
The mapping $\{i\}\mapsto\{(i,1)\}$ is positive and non-clashing.  As the class has at least two concepts, an empty teaching set for every concept would clash, giving the matching lower bound.
\end{proof}

\subsection{No-Clashing Teaching Dimension of Singletons plus empty set}\label{val:noclashing_teaching_dimension_singletons_plus_empty_set}
\noindent\textbf{Statement.} The No-Clashing Teaching Dimension of Singletons plus empty set is $1$.
\begin{proof}
Map each singleton $\{i\}$ to $(i,1)$ and map $\emptyset$ to $(1,0)$.  Every pair is non-clashing: the positive example rules out the empty concept, and $(1,0)$ rules out $\{1\}$.  As there are distinct concepts, size $0$ is impossible.
\end{proof}

\subsection{Positive No-Clashing Teaching Dimension of Singletons plus empty set}\label{val:positive_noclashing_teaching_dimension_singletons_plus_empty_set}
\noindent\textbf{Statement.} The Positive No-Clashing Teaching Dimension of Singletons plus empty set is $1$.
\begin{proof}
Map $\emptyset$ to the empty set and each singleton $\{i\}$ to $(i,1)$.  For a singleton and $\emptyset$, the positive example is inconsistent with $\emptyset$; two different singletons are separated by either positive example.  Thus the mapping is positive and non-clashing.  A size-$0$ mapping cannot work for distinct concepts.
\end{proof}

\subsection{No-Clashing Teaching Dimension of Trivial concepts}\label{val:noclashing_teaching_dimension_trivial_concepts}
\noindent\textbf{Statement.} The No-Clashing Teaching Dimension of Trivial concepts is $1$.
\begin{proof}
At one coordinate, teach $\emptyset$ with label $0$ and the full concept with label $1$.  The two singleton teaching sets are inconsistent with the opposite concept, so this is non-clashing.  The two distinct concepts rule out value $0$.
\end{proof}

\subsection{Positive No-Clashing Teaching Dimension of Trivial concepts}\label{val:positive_noclashing_teaching_dimension_trivial_concepts}
\noindent\textbf{Statement.} The Positive No-Clashing Teaching Dimension of Trivial concepts is $1$.
\begin{proof}
Teach $\emptyset$ by the empty set and the full concept by one positive example.  The latter is inconsistent with $\emptyset$, so the mapping is positive and non-clashing.  Since the class has two distinct concepts, value $0$ is impossible.
\end{proof}

\subsection{Log Size Ratio of Singletons plus empty set}\label{val:log_size_ratio_singletons_plus_empty_set}
\noindent\textbf{Statement.} The Log Size Ratio of Singletons plus empty set is $\frac{\log n}{\log(n+1)}$.
\begin{proof}
This class has $n+1$ concepts on a domain of size $n$.  Substitution into $\log(|\mathcal H|-1)/\log(n+1)$ gives $\log n/\log(n+1)$.
\end{proof}

\subsection{Log Size Ratio of Half-intervals}\label{val:log_size_ratio_halfintervals}
\noindent\textbf{Statement.} The Log Size Ratio of Half-intervals is $\frac{\log(n-1)}{\log(n+1)}$.
\begin{proof}
The half-interval class has $n$ initial segments on a domain of size $n$.  Hence the definition gives $\log(n-1)/\log(n+1)$.
\end{proof}

\subsection{Log Size Ratio of Majority}\label{val:log_size_ratio_majority}
\noindent\textbf{Statement.} The Log Size Ratio of Majority is $\frac{\log(2^n-1)}{\log(n+2)}$.
\begin{proof}
The Majority family has $2^n$ concepts on its $n+1$ point domain.  Substituting these quantities into the definition gives the stated value.
\end{proof}

\subsection{Log Size Ratio of Addressing}\label{val:log_size_ratio_addressing}
\noindent\textbf{Statement.} The Log Size Ratio of Addressing is $\frac{\log(n-1)}{\log(n+\lceil\log_2 n\rceil+1)}$.
\begin{proof}
The class has n concepts on n+r coordinates, where $r=\lceil\log_2n\rceil$. Substituting these two counts into the definition of Log Size Ratio gives the displayed expression. It lies between zero and one.
\end{proof}

\subsection{Log Size Ratio of Tagged Singletons}\label{val:log_size_ratio_tagged_singletons}
\noindent\textbf{Statement.} The Log Size Ratio of Tagged Singletons is $\frac{\log(n+\lceil\log n\rceil)}{\log(n+\lceil\log n\rceil+1)}$.
\begin{proof}
The class consists of the empty concept, $\lceil\log n\rceil$ tag singletons, and $n$ tagged concepts, hence has $n+\lceil\log n\rceil+1$ concepts.  Substitution into the definition yields the stated expression.
\end{proof}

\subsection{Minimum Degree of Full Cube}\label{val:minimum_degree_full_cube}
\noindent\textbf{Statement.} The Minimum Degree of Full Cube is $n$.
\begin{proof}
The $1$-inclusion graph of the full cube is the $n$-dimensional hypercube, which is $n$-regular.
\end{proof}

\subsection{Degeneracy of Full Cube}\label{val:degeneracy_full_cube}
\noindent\textbf{Statement.} The Degeneracy of Full Cube is $n$.
\begin{proof}
The full $n$-cube is an induced subgraph with minimum degree $n$, while no induced subgraph can have minimum degree above the ambient maximum degree $n$.
\end{proof}

\subsection{Densest Subgraph (twice) of Full Cube}\label{val:densest_subgraph_twice_full_cube}
\noindent\textbf{Statement.} The Densest Subgraph (twice) of Full Cube is $n$.
\begin{proof}
The $n$-cube itself has average degree $n$, and every subgraph has average degree at most the ambient maximum degree $n$.
\end{proof}

\subsection{Double Density or Average Degree of Full Cube}\label{val:double_density_or_average_degree_full_cube}
\noindent\textbf{Statement.} The Double Density or Average Degree of Full Cube is $n$.
\begin{proof}
The $1$-inclusion graph is the $n$-regular hypercube, so its average degree is $n$.
\end{proof}

\subsection{Maximum Positive Degree of Full Cube}\label{val:maximum_positive_degree_full_cube}
\noindent\textbf{Statement.} The Maximum Positive Degree of Full Cube is $n$.
\begin{proof}
The full set has $n$ smaller $1$-inclusion neighbours, one obtained by deleting each coordinate, and no vertex can have more than $n$ neighbours.
\end{proof}

\subsection{Projected Maximal Degree of Full Cube}\label{val:projected_maximal_degree_full_cube}
\noindent\textbf{Statement.} The Projected Maximal Degree of Full Cube is $n$.
\begin{proof}
The full-domain projection is the $n$-cube and has maximum degree $n$.  Every projection onto at most $n$ coordinates has maximum degree at most $n$.
\end{proof}

\subsection{Projected Double Density or Projected Average Degree of Full Cube}\label{val:projected_double_density_or_projected_average_degree_full_cube}
\noindent\textbf{Statement.} The Projected Double Density or Projected Average Degree of Full Cube is $n$.
\begin{proof}
The full-domain projection has average degree $n$.  A projection onto $k\le n$ coordinates is a $k$-cube, whose average degree is $k$.
\end{proof}

\subsection{Projected Minimum Degree of Full Cube}\label{val:projected_minimum_degree_full_cube}
\noindent\textbf{Statement.} The Projected Minimum Degree of Full Cube is $n$.
\begin{proof}
The full-domain projection has minimum degree $n$.  Every projection onto at most $n$ coordinates has minimum degree at most $n$.
\end{proof}

\subsection{Minimum Degree of Singletons}\label{val:minimum_degree_singletons}
\noindent\textbf{Statement.} The Minimum Degree of Singletons is $0$.
\begin{proof}
Two distinct singleton concepts differ on two coordinates, so the original $1$-inclusion graph has no edges and every vertex has degree zero.
\end{proof}

\subsection{Degeneracy of Singletons}\label{val:degeneracy_singletons}
\noindent\textbf{Statement.} The Degeneracy of Singletons is $0$.
\begin{proof}
The $1$-inclusion graph has no edges, so every induced subgraph has minimum degree zero.
\end{proof}

\subsection{Densest Subgraph (twice) of Singletons}\label{val:densest_subgraph_twice_singletons}
\noindent\textbf{Statement.} The Densest Subgraph (twice) of Singletons is $0$.
\begin{proof}
The original $1$-inclusion graph is edgeless, as are all of its subgraphs, so their average degrees are zero.
\end{proof}

\subsection{Double Density or Average Degree of Singletons}\label{val:double_density_or_average_degree_singletons}
\noindent\textbf{Statement.} The Double Density or Average Degree of Singletons is $0$.
\begin{proof}
Distinct singleton concepts differ on two coordinates, hence the $1$-inclusion graph is edgeless and has average degree zero.
\end{proof}

\subsection{Maximum Positive Degree of Singletons}\label{val:maximum_positive_degree_singletons}
\noindent\textbf{Statement.} The Maximum Positive Degree of Singletons is $0$.
\begin{proof}
There are no $1$-inclusion neighbours in the original singleton graph, so every dominance is zero.
\end{proof}

\subsection{Projected Maximal Degree of Singletons}\label{val:projected_maximal_degree_singletons}
\noindent\textbf{Statement.} The Projected Maximal Degree of Singletons is $n-1$.
\begin{proof}
Projecting onto a proper $k$-coordinate subset gives the empty trace (from any omitted singleton) and the $k$ projected singletons, whose $1$-inclusion graph is a $k$-leaf star.  To retain an empty trace one must have $k<n$, so the largest such degree is $n-1$.  The full-domain projection has no edges.
\end{proof}

\subsection{Projected Double Density or Projected Average Degree of Singletons}\label{val:projected_double_density_or_projected_average_degree_singletons}
\noindent\textbf{Statement.} The Projected Double Density or Projected Average Degree of Singletons is $\frac{2(n-1)}{n}$.
\begin{proof}
A proper projection onto $k<n$ coordinates consists of the empty trace and its $k$ singleton traces. Its one-inclusion graph is a $k$-leaf star with average degree $2k/(k+1)$. This increases with $k$ and is maximized at $k=n-1$, giving $2(n-1)/n$. The full projection has only singleton traces and is edgeless, so cannot improve the maximum. For $n=1$ both possible projections are edgeless and the same formula gives zero.
\end{proof}

\subsection{Projected Minimum Degree of Singletons}\label{val:projected_minimum_degree_singletons}
\noindent\textbf{Statement.} The Projected Minimum Degree of Singletons is $0$ for $n=1$; $1$ for $n\ge2$.
\begin{proof}
Every proper projection onto $k\ge1$ coordinates is a $k$-leaf star, whose minimum degree is one. The empty projection and the full projection are edgeless. A nonempty proper projection exists exactly when $n\ge2$. Thus the maximum of these minimum degrees is zero at $n=1$ and one otherwise.
\end{proof}

\subsection{Minimum Degree of Singletons plus empty set}\label{val:minimum_degree_singletons_plus_empty_set}
\noindent\textbf{Statement.} The Minimum Degree of Singletons plus empty set is $1$.
\begin{proof}
The $1$-inclusion graph is the star centered at $\emptyset$ with one leaf for each singleton.  Every leaf has degree one.
\end{proof}

\subsection{Degeneracy of Singletons plus empty set}\label{val:degeneracy_singletons_plus_empty_set}
\noindent\textbf{Statement.} The Degeneracy of Singletons plus empty set is $1$.
\begin{proof}
Every nonempty subgraph of a star has a leaf of degree at most one, while the star itself has minimum degree one.
\end{proof}

\subsection{Densest Subgraph (twice) of Singletons plus empty set}\label{val:densest_subgraph_twice_singletons_plus_empty_set}
\noindent\textbf{Statement.} The Densest Subgraph (twice) of Singletons plus empty set is $\frac{2n}{n+1}$.
\begin{proof}
The $n$-leaf star has average degree $2n/(n+1)$.  A substar with $k$ leaves has average degree $2k/(k+1)$, which is maximized at $k=n$.
\end{proof}

\subsection{Double Density or Average Degree of Singletons plus empty set}\label{val:double_density_or_average_degree_singletons_plus_empty_set}
\noindent\textbf{Statement.} The Double Density or Average Degree of Singletons plus empty set is $\frac{2n}{n+1}$.
\begin{proof}
The $1$-inclusion graph is an $n$-leaf star with $n$ edges and $n+1$ vertices, hence average degree $2n/(n+1)$.
\end{proof}

\subsection{Maximum Positive Degree of Singletons plus empty set}\label{val:maximum_positive_degree_singletons_plus_empty_set}
\noindent\textbf{Statement.} The Maximum Positive Degree of Singletons plus empty set is $1$.
\begin{proof}
Each singleton has exactly one smaller neighbour, $\emptyset$, and $\emptyset$ has none.
\end{proof}

\subsection{Projected Maximal Degree of Singletons plus empty set}\label{val:projected_maximal_degree_singletons_plus_empty_set}
\noindent\textbf{Statement.} The Projected Maximal Degree of Singletons plus empty set is $n$.
\begin{proof}
The full-domain projection is the $n$-leaf star and has maximum degree $n$.  A projection onto $k$ coordinates is a $k$-leaf star, whose maximum degree is $k\le n$.
\end{proof}

\subsection{Projected Double Density or Projected Average Degree of Singletons plus empty set}\label{val:projected_double_density_or_projected_average_degree_singletons_plus_empty_set}
\noindent\textbf{Statement.} The Projected Double Density or Projected Average Degree of Singletons plus empty set is $\frac{2n}{n+1}$.
\begin{proof}
A projection onto $k$ coordinates gives a $k$-leaf star, with average degree $2k/(k+1)$.  This is maximized at the full projection $k=n$.
\end{proof}

\subsection{Projected Minimum Degree of Singletons plus empty set}\label{val:projected_minimum_degree_singletons_plus_empty_set}
\noindent\textbf{Statement.} The Projected Minimum Degree of Singletons plus empty set is $1$.
\begin{proof}
Every nonempty projection is a star and has minimum degree one at its leaves.  The empty projection has degree zero, so the maximum over projections is one.
\end{proof}

\subsection{Minimum Degree of Half-intervals}\label{val:minimum_degree_halfintervals}
\noindent\textbf{Statement.} The Minimum Degree of Half-intervals is $1$.
\begin{proof}
Consecutive initial segments form a path on $n$ vertices in the $1$-inclusion graph.  Its endpoints have degree one.
\end{proof}

\subsection{Degeneracy of Half-intervals}\label{val:degeneracy_halfintervals}
\noindent\textbf{Statement.} The Degeneracy of Half-intervals is $1$.
\begin{proof}
Every nonempty subgraph of a path has a vertex of degree at most one, while the path itself has minimum degree one.
\end{proof}

\subsection{Densest Subgraph (twice) of Half-intervals}\label{val:densest_subgraph_twice_halfintervals}
\noindent\textbf{Statement.} The Densest Subgraph (twice) of Half-intervals is $\frac{2(n-1)}{n}$.
\begin{proof}
The $1$-inclusion graph is the path on $n$ vertices, with average degree $2(n-1)/n$.  A subpath on $k$ vertices has average degree $2(k-1)/k$, increasing in $k$, so the whole path is densest.
\end{proof}

\subsection{Double Density or Average Degree of Half-intervals}\label{val:double_density_or_average_degree_halfintervals}
\noindent\textbf{Statement.} The Double Density or Average Degree of Half-intervals is $\frac{2(n-1)}{n}$.
\begin{proof}
The path on $n$ vertices has $n-1$ edges, so its average degree is $2(n-1)/n$.
\end{proof}

\subsection{Maximum Positive Degree of Half-intervals}\label{val:maximum_positive_degree_halfintervals}
\noindent\textbf{Statement.} The Maximum Positive Degree of Half-intervals is $1$.
\begin{proof}
Along the inclusion-ordered path, every nonminimal initial segment has exactly one smaller $1$-inclusion neighbour, and the minimal segment has none.
\end{proof}

\subsection{Minimum Degree of Trivial concepts}\label{val:minimum_degree_trivial_concepts}
\noindent\textbf{Statement.} The Minimum Degree of Trivial concepts is $0$.
\begin{proof}
The empty and full concepts differ on all $n>1$ coordinates, so their $1$-inclusion graph has no edge.
\end{proof}

\subsection{Degeneracy of Trivial concepts}\label{val:degeneracy_trivial_concepts}
\noindent\textbf{Statement.} The Degeneracy of Trivial concepts is $0$.
\begin{proof}
The $1$-inclusion graph is edgeless, and so are all its induced subgraphs.
\end{proof}

\subsection{Densest Subgraph (twice) of Trivial concepts}\label{val:densest_subgraph_twice_trivial_concepts}
\noindent\textbf{Statement.} The Densest Subgraph (twice) of Trivial concepts is $0$.
\begin{proof}
The original $1$-inclusion graph and every subgraph are edgeless, hence have average degree zero.
\end{proof}

\subsection{Double Density or Average Degree of Trivial concepts}\label{val:double_density_or_average_degree_trivial_concepts}
\noindent\textbf{Statement.} The Double Density or Average Degree of Trivial concepts is $0$.
\begin{proof}
The two concepts are not $1$-inclusion neighbours when $n>1$, so the graph has average degree zero.
\end{proof}

\subsection{Maximum Positive Degree of Trivial concepts}\label{val:maximum_positive_degree_trivial_concepts}
\noindent\textbf{Statement.} The Maximum Positive Degree of Trivial concepts is $0$.
\begin{proof}
There is no $1$-inclusion edge, hence no smaller neighbour of either concept.
\end{proof}

\subsection{Projected Maximal Degree of Trivial concepts}\label{val:projected_maximal_degree_trivial_concepts}
\noindent\textbf{Statement.} The Projected Maximal Degree of Trivial concepts is $1$.
\begin{proof}
Projection onto any one coordinate gives the two concepts $0$ and $1$, whose $1$-inclusion graph is a single edge.  Every projection has maximum degree at most one.
\end{proof}

\subsection{Projected Double Density or Projected Average Degree of Trivial concepts}\label{val:projected_double_density_or_projected_average_degree_trivial_concepts}
\noindent\textbf{Statement.} The Projected Double Density or Projected Average Degree of Trivial concepts is $1$.
\begin{proof}
A one-coordinate projection is a single edge and has average degree one.  No projection of two concepts can have average degree above one.
\end{proof}

\subsection{Projected Minimum Degree of Trivial concepts}\label{val:projected_minimum_degree_trivial_concepts}
\noindent\textbf{Statement.} The Projected Minimum Degree of Trivial concepts is $1$.
\begin{proof}
A one-coordinate projection is a single edge, with minimum degree one.  Any projection has at most two vertices and cannot have minimum degree above one.
\end{proof}

\subsection{VC Dimension of Warmuth's $C_5$ class}\label{val:vc_dimension_warmuth_c5}
\noindent\textbf{Statement.} The VC Dimension of Warmuth's $C_5$ class is $2$.
\begin{proof}
Every pair of positions is shattered, whereas every three-position trace misses one labeling: $000$ for three consecutive positions and $111$ for three non-consecutive positions.
\end{proof}
\noindent\textbf{Reference.} \cite{palvolgyi2021unlabeled}.

\subsection{Unlabeled Sample Compression of Warmuth's $C_5$ class}\label{val:unlabeled_sample_compression_warmuth_c5}
\noindent\textbf{Statement.} The Unlabeled sample compression of Warmuth's $C_5$ class is $3$.
\begin{proof}
Keeping all positively labeled positions gives a size-three scheme.  Pálvölgyi and Tardos prove that no size-two unlabeled scheme exists.
\end{proof}
\noindent\textbf{Reference.} \cite{palvolgyi2021unlabeled}.

\subsection{Size of Warmuth's $C_5$ class}\label{val:size_warmuth_c5}
\noindent\textbf{Statement.} The Size of Warmuth's $C_5$ class is $10$.
\begin{proof}
There are five cyclic choices of the four consecutive prescribed positions, and the remaining position may take either value, giving $5\cdot 2=10$ functions.
\end{proof}
\noindent\textbf{Reference.} \cite{palvolgyi2021unlabeled}.

\subsection{Effective Range of Warmuth's $C_5$ class}\label{val:effective_range_warmuth_c5}
\noindent\textbf{Statement.} The Effective Range of Warmuth's $C_5$ class is $5$.
\begin{proof}
Each of the five cyclic positions is the unconstrained position for one of the defining four-position patterns, so two concepts differ there.  Hence every domain coordinate belongs to the effective range.
\end{proof}
\noindent\textbf{Reference.} \cite{palvolgyi2021unlabeled}.

\subsection{Log Size of Warmuth's $C_5$ class}\label{val:log_size_warmuth_c5}
\noindent\textbf{Statement.} The Log Size of Warmuth's $C_5$ class is $\log_2 10$.
\begin{proof}
The class has exactly ten concepts, so its binary logarithmic size is $\log_2 10$.
\end{proof}
\noindent\textbf{Reference.} \cite{palvolgyi2021unlabeled}.

\subsection{Largest Shattered Set of Tagged Singletons}\label{val:largest_shattered_set_tagged_singletons}
\noindent\textbf{Statement.} The Largest Shattered Set of Tagged Singletons is $1$.
\begin{proof}
Largest Shattered Set is, by definition, the VC dimension. The established VC-dimension value for tagged singletons is one.
\end{proof}

\subsection{Largest Shattered Set of Singletons plus Empty Set}\label{val:largest_shattered_set_singletons_plus_empty_set}
\noindent\textbf{Statement.} The Largest Shattered Set of Singletons plus empty set is $1$.
\begin{proof}
Largest Shattered Set is, by definition, the VC dimension. The established VC-dimension value for singletons plus the empty set is one.
\end{proof}

\subsection{Largest Shattered Set of the Full Cube}\label{val:largest_shattered_set_full_cube}
\noindent\textbf{Statement.} The Largest Shattered Set of Full Cube is $n$.
\begin{proof}
Largest Shattered Set is, by definition, the VC dimension. The established VC-dimension value for the full cube is $n$.
\end{proof}

\subsection{Largest Shattered Set of Singletons}\label{val:largest_shattered_set_singletons}
\noindent\textbf{Statement.} The Largest Shattered Set of Singletons is $0$ for $n=1$; $1$ for $n\ge2$.
\begin{proof}
Largest Shattered Set is the VC dimension by definition. For $n=1$ no coordinate is shattered. For $n\ge2$ every one-coordinate set is shattered, and no pair is shattered because the pattern $11$ is absent.
\end{proof}

\subsection{Largest Shattered Set of Trivial Concepts}\label{val:largest_shattered_set_trivial_concepts}
\noindent\textbf{Statement.} The Largest Shattered Set of Trivial concepts is $1$.
\begin{proof}
Largest Shattered Set is, by definition, the VC dimension. The established VC-dimension value for the two trivial concepts is one.
\end{proof}

\subsection{Largest Shattered Set of Addressing}\label{val:largest_shattered_set_addressing}
\noindent\textbf{Statement.} The Largest Shattered Set of Addressing is $\lfloor\log n\rfloor\quad(n=2^r,\ r\ge1)$.
\begin{proof}
Restrict here to $n=2^r$ with all r-bit words present, so $\lfloor\log_2n\rfloor=\log_2n=r$. Largest Shattered Set is, by definition, the VC dimension. The established VC-dimension value for the addressing class is $\lfloor\log n\rfloor$.
\end{proof}

\subsection{Largest Shattered Set of Blockwise-Addressing}\label{val:largest_shattered_set_blockwiseaddressing}
\noindent\textbf{Statement.} The Largest Shattered Set of Blockwise-Addressing is $d$.
\begin{proof}
Largest Shattered Set is, by definition, the VC dimension. The established VC-dimension value for Blockwise-Addressing is $d$.
\end{proof}

\subsection{Largest Shattered Set of Half-intervals}\label{val:largest_shattered_set_halfintervals}
\noindent\textbf{Statement.} The Largest Shattered Set of Half-intervals is $1$.
\begin{proof}
Largest Shattered Set is, by definition, the VC dimension. The established VC-dimension value for half-intervals is one.
\end{proof}

\subsection{Largest Shattered Set of Majority}\label{val:largest_shattered_set_majority}
\noindent\textbf{Statement.} The Largest Shattered Set of Majority is $n$.
\begin{proof}
Largest Shattered Set is, by definition, the VC dimension. The established VC-dimension value for the majority class is $n$.
\end{proof}

\subsection{Largest Shattered Set of Warmuth's $C_5$ Class}\label{val:largest_shattered_set_warmuth_c5}
\noindent\textbf{Statement.} The Largest Shattered Set of Warmuth's $C_5$ class is $2$.
\begin{proof}
Largest Shattered Set is, by definition, the VC dimension. The established VC-dimension value for Warmuth's $C_5$ class is two.
\end{proof}

\subsection{Relative Hitting Size of Trivial Concepts}\label{val:relative_hitting_size_trivial_concepts}
\noindent\textbf{Statement.} The Relative Hitting Size of Trivial concepts is $1$.
\begin{proof}
Relative hitting size is, by definition, the maximum teaching-set size. The established teaching dimension of the two trivial concepts is one.
\end{proof}

\subsection{Relative Hitting Size of Singletons}\label{val:relative_hitting_size_singletons}
\noindent\textbf{Statement.} The Relative Hitting Size of Singletons is $1$.
\begin{proof}
Relative hitting size is, by definition, the maximum teaching-set size. The established teaching dimension of singletons is one.
\end{proof}

\subsection{Relative Hitting Size of the Full Cube}\label{val:relative_hitting_size_full_cube}
\noindent\textbf{Statement.} The Relative Hitting Size of Full Cube is $n$.
\begin{proof}
Relative hitting size is, by definition, the maximum teaching-set size. The established teaching dimension of the full cube is $n$.
\end{proof}

\subsection{Relative Hitting Size of Singletons plus Empty Set}\label{val:relative_hitting_size_singletons_plus_empty_set}
\noindent\textbf{Statement.} The Relative Hitting Size of Singletons plus empty set is $n$.
\begin{proof}
Relative hitting size is, by definition, the maximum teaching-set size. The established teaching dimension of singletons plus the empty set is $n$.
\end{proof}

\subsection{Relative Hitting Size of Half-intervals}\label{val:relative_hitting_size_halfintervals}
\noindent\textbf{Statement.} The Relative Hitting Size of Half-intervals is $2$.
\begin{proof}
Relative hitting size is, by definition, the maximum teaching-set size. The established teaching dimension of half-intervals is two.
\end{proof}

\subsection{Maximum Largest Order Shattered Set of Tagged Singletons}\label{val:maximum_largest_order_shattered_set_tagged_singletons}
\noindent\textbf{Statement.} The Maximum Largest Order Shattered Set of Tagged Singletons is $1$.
\begin{proof}
Maximum Largest Order Shattered Set is, by definition, the VC dimension. Tagged singletons have established VC dimension one.
\end{proof}

\subsection{Maximum Largest Order Shattered Set of Singletons plus Empty Set}\label{val:maximum_largest_order_shattered_set_singletons_plus_empty_set}
\noindent\textbf{Statement.} The Maximum Largest Order Shattered Set of Singletons plus empty set is $1$.
\begin{proof}
Maximum Largest Order Shattered Set is, by definition, the VC dimension. Singletons plus the empty set have established VC dimension one.
\end{proof}

\subsection{Maximum Largest Order Shattered Set of the Full Cube}\label{val:maximum_largest_order_shattered_set_full_cube}
\noindent\textbf{Statement.} The Maximum Largest Order Shattered Set of Full Cube is $n$.
\begin{proof}
Maximum Largest Order Shattered Set is, by definition, the VC dimension. The full cube has established VC dimension $n$.
\end{proof}

\subsection{Maximum Largest Order Shattered Set of Singletons}\label{val:maximum_largest_order_shattered_set_singletons}
\noindent\textbf{Statement.} The Maximum Largest Order Shattered Set of Singletons is $0$ for $n=1$; $1$ for $n\ge2$.
\begin{proof}
Apply downshifts in any coordinate order. The first shift replaces its coordinate singleton by the empty concept; every later shift leaves the other $n-1$ singletons unchanged because the empty concept is already present. The resulting downset is the empty concept together with $n-1$ singletons. Its largest shattered set has size zero for $n=1$ and one for $n\ge2$, independently of the order.
\end{proof}

\subsection{Maximum Largest Order Shattered Set of Trivial Concepts}\label{val:maximum_largest_order_shattered_set_trivial_concepts}
\noindent\textbf{Statement.} The Maximum Largest Order Shattered Set of Trivial concepts is $1$.
\begin{proof}
Maximum Largest Order Shattered Set is, by definition, the VC dimension. The two trivial concepts have established VC dimension one.
\end{proof}

\subsection{Maximum Largest Order Shattered Set of Addressing}\label{val:maximum_largest_order_shattered_set_addressing}
\noindent\textbf{Statement.} The Maximum Largest Order Shattered Set of Addressing is $\lfloor\log n\rfloor\quad(n=2^r,\ r\ge1)$.
\begin{proof}
Restrict here to $n=2^r$ with all r-bit words present, so $\lfloor\log_2n\rfloor=\log_2n=r$. Maximum Largest Order Shattered Set is, by definition, the VC dimension. The addressing class has established VC dimension $\lfloor\log n\rfloor$.
\end{proof}

\subsection{Maximum Largest Order Shattered Set of Blockwise-Addressing}\label{val:maximum_largest_order_shattered_set_blockwiseaddressing}
\noindent\textbf{Statement.} The Maximum Largest Order Shattered Set of Blockwise-Addressing is $d$.
\begin{proof}
Maximum Largest Order Shattered Set is, by definition, the VC dimension. Blockwise-Addressing has established VC dimension $d$.
\end{proof}

\subsection{Maximum Largest Order Shattered Set of Half-intervals}\label{val:maximum_largest_order_shattered_set_halfintervals}
\noindent\textbf{Statement.} The Maximum Largest Order Shattered Set of Half-intervals is $1$.
\begin{proof}
Maximum Largest Order Shattered Set is, by definition, the VC dimension. Half-intervals have established VC dimension one.
\end{proof}

\subsection{Maximum Largest Order Shattered Set of Majority}\label{val:maximum_largest_order_shattered_set_majority}
\noindent\textbf{Statement.} The Maximum Largest Order Shattered Set of Majority is $n$.
\begin{proof}
Maximum Largest Order Shattered Set is, by definition, the VC dimension. The majority class has established VC dimension $n$.
\end{proof}

\subsection{Maximum Largest Order Shattered Set of Warmuth's $C_5$ Class}\label{val:maximum_largest_order_shattered_set_warmuth_c5}
\noindent\textbf{Statement.} The Maximum Largest Order Shattered Set of Warmuth's $C_5$ class is $2$.
\begin{proof}
Maximum Largest Order Shattered Set is, by definition, the VC dimension. Warmuth's $C_5$ class has established VC dimension two.
\end{proof}

\subsection{Minimum Largest Order Shattered Set of the Full Cube}\label{val:minimum_largest_order_shattered_set_full_cube}
\noindent\textbf{Statement.} The Minimum Largest Order Shattered Set of Full Cube is $n$.
\begin{proof}
For every order of the domain, downshifting the full cube leaves the full cube. Thus the entire $n$-point domain is order-shattered for every order, and no larger set exists.
\end{proof}

\subsection{Minimum Largest Order Shattered Set of Trivial Concepts}\label{val:minimum_largest_order_shattered_set_trivial_concepts}
\noindent\textbf{Statement.} The Minimum Largest Order Shattered Set of Trivial concepts is $1$.
\begin{proof}
For every domain order, the two constant concepts remain the empty and full traces after downshifting. Every one-point set is order-shattered, while a two-point cube would require four traces and the class has only two.
\end{proof}

\subsection{Monotonic Minimum Teaching-Set Size of Blockwise-Addressing}\label{val:monotonic_minimum_teaching_set_size_blockwiseaddressing}
\noindent\textbf{Statement.} The Monotonic Minimum Teaching Set Size of Blockwise-Addressing is $2$.
\begin{proof}
The established identity $\mathrm{TS}_{\min}^*=RTD$ applies to every finite class.  The directly verified Blockwise-Addressing value $RTD=2$ therefore gives $\mathrm{TS}_{\min}^*=2$.
\end{proof}
\noindent\textbf{Reference.} \cite{doliwa2014recursive,ben1997online}; derived using the established equality with RTD.

\subsection{Preference-Based Teaching Dimension of Addressing}\label{val:preferencebased_teaching_dimension_addressing}
\noindent\textbf{Statement.} The Preference-Based Teaching Dimension of Addressing is $1\quad(n\ge2)$.
\begin{proof}
Preference-based teaching dimension equals recursive teaching dimension. The addressing class has established recursive teaching dimension one.
\end{proof}
\noindent\textbf{Reference.} \cite{gao2017preference}.

\subsection{Preference-Based Teaching Dimension of Tagged Singletons}\label{val:preferencebased_teaching_dimension_tagged_singletons}
\noindent\textbf{Statement.} The Preference-Based Teaching Dimension of Tagged Singletons is $1$.
\begin{proof}
Preference-based teaching dimension equals recursive teaching dimension. Tagged singletons have established recursive teaching dimension one.
\end{proof}
\noindent\textbf{Reference.} \cite{gao2017preference}.

\subsection{Preference-Based Teaching Dimension of Half-intervals}\label{val:preferencebased_teaching_dimension_halfintervals}
\noindent\textbf{Statement.} The Preference-Based Teaching Dimension of Half-intervals is $1$.
\begin{proof}
Preference-based teaching dimension equals recursive teaching dimension. Half-intervals have established recursive teaching dimension one.
\end{proof}
\noindent\textbf{Reference.} \cite{gao2017preference}.

\subsection{Preference-Based Teaching Dimension of the Full Cube}\label{val:preferencebased_teaching_dimension_full_cube}
\noindent\textbf{Statement.} The Preference-Based Teaching Dimension of Full Cube is $n$.
\begin{proof}
Preference-based teaching dimension equals recursive teaching dimension. The full cube has established recursive teaching dimension $n$.
\end{proof}
\noindent\textbf{Reference.} \cite{gao2017preference}.

\subsection{Preference-Based Teaching Dimension of Singletons plus Empty Set}\label{val:preferencebased_teaching_dimension_singletons_plus_empty_set}
\noindent\textbf{Statement.} The Preference-Based Teaching Dimension of Singletons plus empty set is $1$.
\begin{proof}
Preference-based teaching dimension equals recursive teaching dimension. Singletons plus the empty set have established recursive teaching dimension one.
\end{proof}
\noindent\textbf{Reference.} \cite{gao2017preference}.

\subsection{Preference-Based Teaching Dimension of Singletons}\label{val:preferencebased_teaching_dimension_singletons}
\noindent\textbf{Statement.} The Preference-Based Teaching Dimension of Singletons is $1$.
\begin{proof}
Preference-based teaching dimension equals recursive teaching dimension. Singletons have established recursive teaching dimension one.
\end{proof}
\noindent\textbf{Reference.} \cite{gao2017preference}.

\subsection{Preference-Based Teaching Dimension of Trivial Concepts}\label{val:preferencebased_teaching_dimension_trivial_concepts}
\noindent\textbf{Statement.} The Preference-Based Teaching Dimension of Trivial concepts is $1$.
\begin{proof}
Preference-based teaching dimension equals recursive teaching dimension. The two trivial concepts have established recursive teaching dimension one.
\end{proof}
\noindent\textbf{Reference.} \cite{gao2017preference}.

\subsection{Preference-Based Teaching Dimension of Blockwise-Addressing}\label{val:preferencebased_teaching_dimension_blockwiseaddressing}
\noindent\textbf{Statement.} The Preference-Based Teaching Dimension of Blockwise-Addressing is $2$.
\begin{proof}
For finite classes, preference-based teaching dimension equals recursive teaching dimension.  The directly verified Blockwise-Addressing value $RTD=2$ therefore gives $\mathrm{PBTD}=2$.
\end{proof}
\noindent\textbf{Reference.} \cite{gao2017preference,ben1997online}; derived using the established equality with RTD.

\subsection{Minimum Projected Teaching-Set Size of the Full Cube}\label{val:minimum_projected_teaching_set_size_full_cube}
\noindent\textbf{Statement.} The Minimum Projected Teaching Set Size of Full Cube is $n$.
\begin{proof}
For every target of the full cube, projection to the full $n$-coordinate domain is again a full cube, and all $n$ coordinates are required to distinguish the target. Thus every target has worst projected teaching-set size at least $n$; no projection has more coordinates, so the minimum of these worst-case sizes is $n$.
\end{proof}

\subsection{Minimum Projected Teaching-Set Size of Trivial Concepts}\label{val:minimum_projected_teaching_set_size_trivial_concepts}
\noindent\textbf{Statement.} The Minimum Projected Teaching Set Size of Trivial concepts is $1$.
\begin{proof}
Every nonempty projection of the two constant concepts is again two complementary constants, and one labeled coordinate teaches either target. The empty projection needs no examples. Hence each target has worst projected teaching-set size one.
\end{proof}

\subsection{Minimum Projected Teaching-Set Size of Singletons}\label{val:minimum_projected_teaching_set_size_singletons}
\noindent\textbf{Statement.} The Minimum Projected Teaching Set Size of Singletons is $n-1$.
\begin{proof}
For a singleton target $\{i\}$, project to the other $n-1$ coordinates. Its trace is empty, while the projected class contains the empty trace and all $n-1$ singleton traces, so all $n-1$ negative coordinates are needed to teach it. Any projection has at most $n-1$ relevant competing singleton traces after omitting the target, and projections containing $i$ teach the target with its positive coordinate. Thus every target has worst projected teaching-set size exactly $n-1$.
\end{proof}

\subsection{Minimum Projected Teaching-Set Size of Singletons plus Empty Set}\label{val:minimum_projected_teaching_set_size_singletons_plus_empty_set}
\noindent\textbf{Statement.} The Minimum Projected Teaching Set Size of Singletons plus empty set is $n-1$.
\begin{proof}
For a singleton target $\{i\}$, projection to the other $n-1$ coordinates turns it into the empty trace; together with the original empty concept and the other singleton traces, this requires all $n-1$ negative coordinates to teach. No projection can give that target more than $n-1$ competing singleton traces, and a projection containing $i$ is taught by its positive coordinate. The empty target has worst projected teaching-set size $n$, but the parameter takes the minimum over targets, hence its value is $n-1$.
\end{proof}

\subsection{Monotonic Projected Minimum Teaching-Set Size of the Full Cube}\label{val:monotonic_projected_minimum_teaching_set_size_full_cube}
\noindent\textbf{Statement.} The Monotonic Projected Minimum Teaching Set Size of Full Cube is $n$.
\begin{proof}
Taking the full cube itself as the subfamily gives minimum projected teaching-set size $n$. Conversely, every subfamily and every projection live on at most $n$ coordinates, so no teaching set can have size greater than $n$.
\end{proof}

\subsection{Monotonic Projected Minimum Teaching-Set Size of Trivial Concepts}\label{val:monotonic_projected_minimum_teaching_set_size_trivial_concepts}
\noindent\textbf{Statement.} The Monotonic Projected Minimum Teaching Set Size of Trivial concepts is $1$.
\begin{proof}
The full two-concept class has minimum projected teaching-set size one. Every subclass has at most two distinct projected traces, so its minimum projected teaching-set size is at most one; hence the subfamily maximum is one.
\end{proof}

\subsection{Equivalence Queries Complexity of Addressing}\label{val:equivalence_queries_complexity_addressing}
\noindent\textbf{Statement.} The Equivalence Queries Complexity of Addressing is $\lfloor\log n\rfloor\quad(n=2^r,\ r\ge1)$.
\begin{proof}
Restrict here to $n=2^r$ with all r-bit words present, so $\lfloor\log_2n\rfloor=\log_2n=r$. The established equality $\mathrm{Q}_{\mathrm{eq}}=\mathrm{Ldim}$ reduces this to the Littlestone-dimension value for Addressing.
\end{proof}
\noindent\textbf{Reference.} \cite{littlestone1988learning,maass1992lower}.

\subsection{Equivalence Queries Complexity of Tagged Singletons}\label{val:equivalence_queries_complexity_tagged_singletons}
\noindent\textbf{Statement.} The Equivalence Queries Complexity of Tagged Singletons is $2$.
\begin{proof}
The established equality $\mathrm{Q}_{\mathrm{eq}}=\mathrm{Ldim}$ reduces this to the Littlestone-dimension value for Tagged Singletons.
\end{proof}
\noindent\textbf{Reference.} \cite{littlestone1988learning,maass1992lower}.

\subsection{Equivalence Queries Complexity of Half-intervals}\label{val:equivalence_queries_complexity_halfintervals}
\noindent\textbf{Statement.} The Equivalence Queries Complexity of Half-intervals is $\log n$.
\begin{proof}
The established equality $\mathrm{Q}_{\mathrm{eq}}=\mathrm{Ldim}$ reduces this to the Littlestone-dimension value for Half-intervals.
\end{proof}
\noindent\textbf{Reference.} \cite{littlestone1988learning}.

\subsection{Equivalence Queries Complexity of Full Cube}\label{val:equivalence_queries_complexity_full_cube}
\noindent\textbf{Statement.} The Equivalence Queries Complexity of Full Cube is $n$.
\begin{proof}
The established equality $\mathrm{Q}_{\mathrm{eq}}=\mathrm{Ldim}$ reduces this to the Littlestone-dimension value for the Full Cube.
\end{proof}
\noindent\textbf{Reference.} \cite{littlestone1988learning}.

\subsection{Equivalence Queries Complexity of Singletons plus Empty Set}\label{val:equivalence_queries_complexity_singletons_plus_empty_set}
\noindent\textbf{Statement.} The Equivalence Queries Complexity of Singletons plus empty set is $1$.
\begin{proof}
The established equality $\mathrm{Q}_{\mathrm{eq}}=\mathrm{Ldim}$ reduces this to the Littlestone-dimension value for Singletons plus Empty Set.
\end{proof}
\noindent\textbf{Reference.} \cite{littlestone1988learning}.

\subsection{Equivalence Queries Complexity of Singletons}\label{val:equivalence_queries_complexity_singletons}
\noindent\textbf{Statement.} The Equivalence Queries Complexity of Singletons is $1$.
\begin{proof}
The established equality $\mathrm{Q}_{\mathrm{eq}}=\mathrm{Ldim}$ reduces this to the Littlestone-dimension value for Singletons.
\end{proof}
\noindent\textbf{Reference.} \cite{littlestone1988learning}.

\subsection{Equivalence Queries Complexity of Trivial Concepts}\label{val:equivalence_queries_complexity_trivial_concepts}
\noindent\textbf{Statement.} The Equivalence Queries Complexity of Trivial concepts is $1$.
\begin{proof}
The established equality $\mathrm{Q}_{\mathrm{eq}}=\mathrm{Ldim}$ reduces this to the Littlestone-dimension value for Trivial Concepts.
\end{proof}
\noindent\textbf{Reference.} \cite{littlestone1988learning}.

\subsection{Hitting Size of Trivial Concepts}\label{val:hitting_size_trivial_concepts}
\noindent\textbf{Statement.} The Hitting Size of Trivial concepts is $1$.
\begin{proof}
The full concept is the only nonempty member, and any one coordinate hits it. The empty set cannot hit it.
\end{proof}

\subsection{Hitting Size of Singletons}\label{val:hitting_size_singletons}
\noindent\textbf{Statement.} The Hitting Size of Singletons is $n$.
\begin{proof}
Every singleton is a nonempty concept, so a hitting set must contain every coordinate. The full domain meets every singleton.
\end{proof}

\subsection{Hitting Size of Full Cube}\label{val:hitting_size_full_cube}
\noindent\textbf{Statement.} The Hitting Size of Full Cube is $n$.
\begin{proof}
The Full Cube contains every singleton, forcing every coordinate into a hitting set. The full domain is a hitting set.
\end{proof}

\subsection{Hitting Size of Singletons plus Empty Set}\label{val:hitting_size_singletons_plus_empty_set}
\noindent\textbf{Statement.} The Hitting Size of Singletons plus empty set is $n$.
\begin{proof}
Every singleton is a nonempty concept, so a hitting set must contain every coordinate. The empty concept imposes no condition.
\end{proof}

\subsection{Hitting Size of Half-intervals}\label{val:hitting_size_halfintervals}
\noindent\textbf{Statement.} The Hitting Size of Half-intervals is $1$.
\begin{proof}
Every nonempty initial segment contains its first coordinate, so that one coordinate is a hitting set. A hitting set cannot be empty.
\end{proof}

\subsection{Hitting Size of Tagged Singletons}\label{val:hitting_size_tagged_singletons}
\noindent\textbf{Statement.} The Hitting Size of Tagged Singletons is $\lceil\log_2 n\rceil$.
\begin{proof}
Each tag coordinate occurs as a singleton concept, so every hitting set contains all $\lceil\log_2 n\rceil$ tag coordinates. Those coordinates meet every tagged pair as well, hence form a hitting set.
\end{proof}

\subsection{Hitting Size of Majority}\label{val:hitting_size_majority}
\noindent\textbf{Statement.} The Hitting Size of Majority is $n$.
\begin{proof}
For every base coordinate $i$, the concept induced by $\{i\}$ is a singleton, so all $n$ base coordinates are forced. They meet every nonempty Majority concept; the extra majority coordinate is never needed.
\end{proof}

\subsection{Hitting Size of Addressing}\label{val:hitting_size_addressing}
\noindent\textbf{Statement.} The Hitting Size of Addressing is $\log_2n+1\quad(n=2^r)$.
\begin{proof}
Let $B\subseteq\{0,1\}^r$ be the chosen set of n codewords. Fix the selected bit coordinates T of a hitting set. Every target whose code is zero on T must then be hit by its own private coordinate, and those are the only remaining coordinates that can hit it. Conversely, selecting those private coordinates suffices. Hence $\mathrm{HT}=\min_{T\subseteq[r]}(|T|+|\{b\in B:b|_T=0\}|)$. For complete codes and t selected bits, exactly $2^{r-t}$ targets remain unhit, giving minimum $\min_{0\le t\le r}(t+2^{r-t})=r+1$: write k=r-t and use $2^k\ge k+1$, with equality at k=0 or k=1. Thus the old claimed value n is false for r>=2. At r=0 the sole nonempty concept has hitting size one.
\end{proof}

\subsection{Hitting Size of Warmuth's $C_5$ Class}\label{val:hitting_size_warmuth_c5}
\noindent\textbf{Statement.} The Hitting Size of Warmuth's $C_5$ class is $3$.
\begin{proof}
The five concepts with their unrestricted coordinate set to zero have supports $\{0,3\},\{1,4\},\{0,2\},\{1,3\},\{2,4\}$. These are the edges of a $5$-cycle, whose minimum vertex cover has size three, so every hitting set has size at least three. Any three-coordinate vertex cover of this cycle meets those five supports and hence also their one-coordinate extensions, giving a hitting set.
\end{proof}

\subsection{Membership Query Complexity of Half-intervals}\label{val:membership_query_complexity_halfintervals}
\noindent\textbf{Statement.} The Membership Query Complexity of Half-intervals is $\lceil\log_2 n\rceil$.
\begin{proof}
A membership query at a midpoint says on which side of that point the unknown initial-segment endpoint lies, so binary search identifies one of the $n$ concepts in $\lceil\log_2 n\rceil$ queries. Any deterministic membership-query decision tree distinguishing $n$ targets has depth at least $\lceil\log_2 n\rceil$.
\end{proof}
\noindent\textbf{Reference.} \cite{angluin1988queries}.

\subsection{Membership Query Complexity of Addressing}\label{val:membership_query_complexity_addressing}
\noindent\textbf{Statement.} The Membership Query Complexity of Addressing is $\lceil\log_2 n\rceil$.
\begin{proof}
The r bit coordinates, with $r=\lceil\log_2n\rceil$, identify the target because its encoding is injective. Conversely, a binary membership-query decision tree of worst-case depth q has at most $2^q$ leaves, and must distinguish all n targets, so $q\ge\lceil\log_2n\rceil$. This includes zero queries at n=1.
\end{proof}
\noindent\textbf{Reference.} \cite{angluin1988queries}.

\subsection{Membership Query Complexity of Majority}\label{val:membership_query_complexity_majority}
\noindent\textbf{Statement.} The Membership Query Complexity of Majority is $n$.
\begin{proof}
The labels on the $n$ base coordinates are arbitrary, and the final majority coordinate is determined by them. Querying all base coordinates identifies the target. Since all $2^n$ base labelings occur, a binary decision tree requires at least $n$ queries.
\end{proof}
\noindent\textbf{Reference.} \cite{angluin1988queries}.

\subsection{2-Block Littlestone Dimension of Trivial Concepts}\label{val:2block_littlestone_dimension_trivial_concepts}
\noindent\textbf{Statement.} The 2-Block Littlestone Dimension of Trivial concepts is $1$.
\begin{proof}
The two constant hypotheses shatter a tree of depth one. A shattered binary tree of depth two has four leaves, and each leaf needs a distinct realizing hypothesis, while this class has only two. Thus $\mathrm B_2=1$.
\end{proof}

\subsection{Best Mistakes of Trivial Concepts}\label{val:best_mistakes_trivial_concepts}
\noindent\textbf{Statement.} The Best Mistakes of Trivial concepts is $1$.
\begin{proof}
At any coordinate the two constant concepts have opposite labels, so no learner can guarantee zero mistakes against both targets. After the first labeled presentation, the target is known, giving an upper bound of one.
\end{proof}
\noindent\textbf{Reference.} \cite{ben1997online}.

\subsection{Worst Mistakes of Trivial Concepts}\label{val:worst_mistakes_trivial_concepts}
\noindent\textbf{Statement.} The Worst Mistakes of Trivial concepts is $1$.
\begin{proof}
No strategy can guarantee zero mistakes against the two opposite constant targets. Regardless of the adversarial presentation, the first observed label identifies the target, so at most one mistake is possible.
\end{proof}
\noindent\textbf{Reference.} \cite{ben1997online}.

\subsection{Self-Directed Queries Complexity of Trivial Concepts}\label{val:selfdirected_queries_complexity_trivial_concepts}
\noindent\textbf{Statement.} The Self-Directed Queries Complexity of Trivial concepts is $1$.
\begin{proof}
The learner chooses any coordinate. Its first prediction may be wrong, but the returned label identifies which constant is the target and prevents all later mistakes. Zero is impossible for two opposite labelings.
\end{proof}
\noindent\textbf{Reference.} \cite{goldman1993learning,goldman1994power}.

\subsection{Occasional One-Sided Getback of Trivial Concepts}\label{val:occasional_onesided_getback_trivial_concepts}
\noindent\textbf{Statement.} The Occasional One-Sided Getback of Trivial concepts is $1$.
\begin{proof}
No strategy can guarantee zero mistakes against the two opposite constants. If a mistake-only signal is ever returned on the first presented coordinate, it identifies the other constant target; if no signal is returned, no mistake is counted. Thus at most one mistake occurs.
\end{proof}

\subsection{Distinguishing Range of Trivial Concepts}\label{val:distinguishing_range_trivial_concepts}
\noindent\textbf{Statement.} The Distinguishing Range of Trivial concepts is $1$.
\begin{proof}
On any one coordinate the empty and full concepts have opposite traces. The empty projection cannot distinguish two concepts.
\end{proof}

\subsection{Distinguishing Range of Singletons}\label{val:distinguishing_range_singletons}
\noindent\textbf{Statement.} The Distinguishing Range of Singletons is $n-1$.
\begin{proof}
If two coordinates are omitted, their singleton concepts have the same empty trace, so at least $n-1$ coordinates are necessary. Retaining any $n-1$ coordinates gives distinct traces for every singleton: the omitted singleton has the empty trace and all retained singletons have their unique positive coordinate.
\end{proof}

\subsection{Distinguishing Range of Full Cube}\label{val:distinguishing_range_full_cube}
\noindent\textbf{Statement.} The Distinguishing Range of Full Cube is $n$.
\begin{proof}
If a coordinate is omitted, two cube concepts differing only on it have identical traces. Retaining all $n$ coordinates is injective.
\end{proof}

\subsection{Distinguishing Range of Singletons plus Empty Set}\label{val:distinguishing_range_singletons_plus_empty_set}
\noindent\textbf{Statement.} The Distinguishing Range of Singletons plus empty set is $n$.
\begin{proof}
If coordinate $i$ is omitted, the empty concept and singleton $\{i\}$ have identical empty traces. Thus every coordinate is necessary, and the full domain distinguishes all concepts.
\end{proof}

\subsection{Distinguishing Range of Half-intervals}\label{val:distinguishing_range_halfintervals}
\noindent\textbf{Statement.} The Distinguishing Range of Half-intervals is $n-1$.
\begin{proof}
For any retained coordinates, the traces of initial segments change only when an endpoint crosses a retained coordinate, producing at most one more trace than the number retained. Distinguishing all $n$ initial segments therefore needs at least $n-1$ coordinates. Coordinates $2,\ldots,n$ attain this bound.
\end{proof}

\subsection{Distinguishing Range of Majority}\label{val:distinguishing_range_majority}
\noindent\textbf{Statement.} The Distinguishing Range of Majority is $n$.
\begin{proof}
All $2^n$ labelings of the $n$ base coordinates occur, and the final majority coordinate is determined by them. Omitting a base coordinate identifies two concepts that differ only there, while the $n$ base coordinates distinguish every concept.
\end{proof}

\subsection{Distinguishing Range of Addressing}\label{val:distinguishing_range_addressing}
\noindent\textbf{Statement.} The Distinguishing Range of Addressing is $\lceil\log_2 n\rceil$.
\begin{proof}
The r bit coordinates, where $r=\lceil\log_2n\rceil$, distinguish every target. A set of q binary coordinates has at most $2^q$ distinct traces; to distinguish n concepts requires $q\ge\lceil\log_2n\rceil$. Thus the minimum distinguishing range is exactly r, including zero at n=1.
\end{proof}

\subsection{Size of Blockwise-Addressing}\label{val:size_blockwiseaddressing}
\noindent\textbf{Statement.} The Size of Blockwise-Addressing is $2^{d+1}$.
\begin{proof}
The construction contains the $2^d$ functions $f_i$ and the $2^d$ functions $g_i$, for a total of $2^{d+1}$ concepts.
\end{proof}
\noindent\textbf{Reference.} \cite{ben1997online}.

\subsection{Effective Range of Blockwise-Addressing}\label{val:effective_range_blockwiseaddressing}
\noindent\textbf{Statement.} The Effective Range of Blockwise-Addressing is $2^{d+1}+1$.
\begin{proof}
The domain consists of $z$, the $2^d$ coordinates $x_i$, and the $2^d$ coordinates $y_i$. Each coordinate changes value among the displayed $f_i,g_i$, so every domain coordinate is effective.
\end{proof}
\noindent\textbf{Reference.} \cite{ben1997online}.

\subsection{Log Size of Blockwise-Addressing}\label{val:log_size_blockwiseaddressing}
\noindent\textbf{Statement.} The Log Size of Blockwise-Addressing is $d+1$.
\begin{proof}
The class has size $2^{d+1}$, so its binary log-size is $d+1$.
\end{proof}
\noindent\textbf{Reference.} \cite{ben1997online}.

\subsection{Sign Rank of Trivial Concepts}\label{val:sign_rank_trivial_concepts}
\noindent\textbf{Statement.} The Sign Rank of Trivial concepts is $1$.
\begin{proof}
The two sign rows are constant opposites. They are realized by the rank-one matrix with one all-positive row and one all-negative row; rank zero cannot realize a nonzero sign pattern.
\end{proof}
\noindent\textbf{Reference.} \cite{alon2016sign}.

\subsection{Dual Sign Rank of Trivial Concepts}\label{val:dual_sign_rank_trivial_concepts}
\noindent\textbf{Statement.} The Dual Sign Rank of Trivial concepts is $1$.
\begin{proof}
Every realizing matrix has nonzero columns, so one fixed column is independent. Two columns are not forced independent, since the rank-one realization with identical positive column magnitudes realizes the same sign matrix.
\end{proof}
\noindent\textbf{Reference.} \cite{alon2016sign}.

\subsection{Sign Rank of Full Cube}\label{val:sign_rank_full_cube}
\noindent\textbf{Statement.} The Sign Rank of Full Cube is $n$.
\begin{proof}
The $2^n$ cube sign patterns form a realizing matrix with $n$ linearly independent columns, so sign rank is at most $n$. Dual sign rank is at least VC dimension, which is $n$, and dual sign rank never exceeds sign rank. Hence sign rank is exactly $n$.
\end{proof}
\noindent\textbf{Reference.} \cite{alon2016sign}.

\subsection{Dual Sign Rank of Full Cube}\label{val:dual_sign_rank_full_cube}
\noindent\textbf{Statement.} The Dual Sign Rank of Full Cube is $n$.
\begin{proof}
Dual sign rank is at least VC dimension, which is $n$ for the Full Cube. It is at most sign rank, and the cube sign matrix itself has rank $n$; therefore the value is exactly $n$.
\end{proof}
\noindent\textbf{Reference.} \cite{alon2016sign}.

\subsection{Interpolation Degree of Trivial Concepts}\label{val:interpolation_degree_trivial_concepts}
\noindent\textbf{Statement.} The Interpolation degree of Trivial concepts is $1$.
\begin{proof}
On either coordinate, the two concepts have values zero and one. Constant and linear polynomials therefore interpolate every function on the class, while constants alone cannot distinguish its two members.
\end{proof}

\subsection{Interpolation Degree of Singletons}\label{val:interpolation_degree_singletons}
\noindent\textbf{Statement.} The Interpolation degree of Singletons is $1$.
\begin{proof}
For prescribed values $a_i$ on singleton $\{i\}$, the linear polynomial $\sum_i a_i x_i$ has exactly those values. A degree-zero polynomial cannot interpolate distinct values on two singleton concepts.
\end{proof}

\subsection{Interpolation Degree of Full Cube}\label{val:interpolation_degree_full_cube}
\noindent\textbf{Statement.} The Interpolation degree of Full Cube is $n$.
\begin{proof}
Every function on the Boolean cube has a multilinear polynomial representation of degree at most $n$. The parity function has nonzero coefficient on the degree-$n$ monomial in its unique multilinear representation, so degree below $n$ cannot suffice.
\end{proof}

\subsection{Interpolation Degree of Singletons plus Empty Set}\label{val:interpolation_degree_singletons_plus_empty_set}
\noindent\textbf{Statement.} The Interpolation degree of Singletons plus empty set is $1$.
\begin{proof}
For prescribed values $a_0$ on the empty concept and $a_i$ on $\{i\}$, the affine polynomial $a_0+\sum_i(a_i-a_0)x_i$ interpolates them. Constants alone do not distinguish the concepts.
\end{proof}

\subsection{Interpolation Degree of Half-intervals}\label{val:interpolation_degree_halfintervals}
\noindent\textbf{Statement.} The Interpolation degree of Half-intervals is $1$.
\begin{proof}
For target values $a_i$ on initial segments $[i]$, the linear polynomial $a_1x_1+\sum_{i=2}^n(a_i-a_{i-1})x_i$ evaluates to $a_i$ on $[i]$. Constants do not interpolate all functions on this non-singleton class.
\end{proof}

\subsection{Interpolation Degree of Tagged Singletons}\label{val:interpolation_degree_tagged_singletons}
\noindent\textbf{Statement.} The Interpolation degree of Tagged Singletons is $1$.
\begin{proof}
The empty trace, the tag singletons, and the tagged pairs are affinely independent: subtracting each tag singleton from a pair with that tag yields the corresponding ordinary singleton direction. Thus affine-linear polynomials interpolate every function on the class, and degree zero is insufficient.
\end{proof}

\subsection{Interpolation Degree of Addressing}\label{val:interpolation_degree_addressing}
\noindent\textbf{Statement.} The Interpolation degree of Addressing is $1\quad(n\ge2)$.
\begin{proof}
Each concept has a unique addressed primary coordinate. Assigning a coefficient to each such coordinate gives a linear polynomial with arbitrary prescribed values on the corresponding concepts. A constant polynomial cannot distinguish the nontrivial class.
\end{proof}

\subsection{Minimum Largest Order Shattered Set of Singletons}\label{val:minimum_largest_order_shattered_set_singletons}
\noindent\textbf{Statement.} The Minimum Largest Order Shattered Set of Singletons is $1$.
\begin{proof}
For every coordinate order, full downshifting produces a nontrivial downset, hence it contains the empty trace and a singleton trace, a one-cube. Downshifting preserves VC dimension one, so no two-cube can occur.
\end{proof}

\subsection{Minimum Largest Order Shattered Set of Singletons plus Empty Set}\label{val:minimum_largest_order_shattered_set_singletons_plus_empty_set}
\noindent\textbf{Statement.} The Minimum Largest Order Shattered Set of Singletons plus empty set is $1$.
\begin{proof}
The class is already a nontrivial downset, so it contains a one-cube. Its VC dimension is one, and order shattering cannot contain a two-cube after any downshift.
\end{proof}

\subsection{Minimum Largest Order Shattered Set of Half-intervals}\label{val:minimum_largest_order_shattered_set_halfintervals}
\noindent\textbf{Statement.} The Minimum Largest Order Shattered Set of Half-intervals is $1$.
\begin{proof}
For every coordinate order, the downshifted family is a nontrivial downset and hence contains a one-cube. The Half-interval class has VC dimension one, which forbids a two-cube after downshifting.
\end{proof}

\subsection{Effective VC Radius of Singletons plus Empty Set}\label{val:effective_vc_radius_singletons_plus_empty_set}
\noindent\textbf{Statement.} The Effective VC Radius of Singletons plus empty set is $1$.
\begin{proof}
Every effective coordinate is shattered by the empty concept and its singleton. A two-coordinate trace misses the pattern with both coordinates positive, so radius two fails.
\end{proof}

\subsection{Effective VC Radius of Half-intervals}\label{val:effective_vc_radius_halfintervals}
\noindent\textbf{Statement.} The Effective VC Radius of Half-intervals is $1$.
\begin{proof}
Each coordinate is shattered by an initial segment ending before it and one ending at or after it. For two ordered coordinates $i<j$, initial segments realize only $00,10,11$, never $01$, so radius two fails.
\end{proof}

\subsection{Effective VC Radius of Tagged Singletons}\label{val:effective_vc_radius_tagged_singletons}
\noindent\textbf{Statement.} The Effective VC Radius of Tagged Singletons is $1$.
\begin{proof}
Every ordinary or tag coordinate varies and is therefore shattered. Two distinct ordinary coordinates never occur together in a concept, so their trace misses $11$; hence radius two fails.
\end{proof}

\subsection{Effective VC Radius of Majority}\label{val:effective_vc_radius_majority}
\noindent\textbf{Statement.} The Effective VC Radius of Majority is $\lceil n/2\rceil$.
\begin{proof}
Any subset of at most $\lceil n/2\rceil$ coordinates is shattered. If it omits the majority coordinate this is immediate from the free base bits. Otherwise, after fixing its $k-1$ selected base bits, the remaining $n-k+1$ bits can make the majority bit either value exactly when $k\le\lceil n/2\rceil$. For $k=\lceil n/2\rceil+1$, choose the majority coordinate and $k-1$ base coordinates all labeled zero: even setting every remaining base bit to one does not reach a strict majority, so this set is not shattered.
\end{proof}

\subsection{Effective VC Radius of Addressing}\label{val:effective_vc_radius_addressing}
\noindent\textbf{Statement.} The Effective VC Radius of Addressing is $1\quad(n\ge2)$.
\begin{proof}
Every effective coordinate varies and so is shattered. Any two distinct primary coordinates never occur together in one addressing concept, so their trace misses $11$ and radius two fails.
\end{proof}

\subsection{Minimum Teaching-Set Size of Tagged Singletons}\label{val:minimum_teaching_set_size_tagged_singletons}
\noindent\textbf{Statement.} The Minimum Teaching Set Size of Tagged Singletons is $1$.
\begin{proof}
A tagged pair contains an ordinary coordinate $i$ that occurs in no other concept. Its one positive example $(i,1)$ therefore teaches that pair. The nontrivial class cannot have minimum teaching-set size zero.
\end{proof}

\subsection{Minimum Teaching-Set Size of Addressing}\label{val:minimum_teaching_set_size_addressing}
\noindent\textbf{Statement.} The Minimum Teaching Set Size of Addressing is $1\quad(n\ge2)$.
\begin{proof}
Every addressing concept has its own addressed primary coordinate, which is positive for that concept and negative for every other one. Hence one positive example teaches it, while size zero cannot distinguish a nontrivial class.
\end{proof}

\subsection{Positive Teaching Dimension of Trivial Concepts}\label{val:positive_teaching_dimension_trivial_concepts}
\noindent\textbf{Statement.} The Positive Teaching Dimension of Trivial concepts is $\infty$.
\begin{proof}
The empty concept has no positive examples. Its only positive teaching set is therefore empty, which is also consistent with the full concept, so it has no finite positive teaching set.
\end{proof}
\noindent\textbf{Reference.} \cite{balbach2008measuring}.

\subsection{Positive Teaching Dimension of Singletons}\label{val:positive_teaching_dimension_singletons}
\noindent\textbf{Statement.} The Positive Teaching Dimension of Singletons is $1$.
\begin{proof}
The unique positive coordinate of a singleton distinguishes it from every other singleton. A nontrivial class cannot have positive teaching dimension zero.
\end{proof}
\noindent\textbf{Reference.} \cite{balbach2008measuring}.

\subsection{Positive Teaching Dimension of Full Cube}\label{val:positive_teaching_dimension_full_cube}
\noindent\textbf{Statement.} The Positive Teaching Dimension of Full Cube is $\infty$.
\begin{proof}
The empty cube concept has no positive examples, and its empty positive sample is consistent with every cube concept. Thus it has no finite positive teaching set.
\end{proof}
\noindent\textbf{Reference.} \cite{balbach2008measuring}.

\subsection{Positive Teaching Dimension of Singletons plus Empty Set}\label{val:positive_teaching_dimension_singletons_plus_empty_set}
\noindent\textbf{Statement.} The Positive Teaching Dimension of Singletons plus empty set is $\infty$.
\begin{proof}
The empty concept has no positive examples, so its only positive sample is empty and remains consistent with every singleton. Hence it has no finite positive teaching set.
\end{proof}
\noindent\textbf{Reference.} \cite{balbach2008measuring}.

\subsection{Positive Teaching Dimension of Half-intervals}\label{val:positive_teaching_dimension_halfintervals}
\noindent\textbf{Statement.} The Positive Teaching Dimension of Half-intervals is $\infty$.
\begin{proof}
For any proper initial segment, every positive example also belongs to each larger initial segment. In particular, no positive sample can distinguish $[1]$ from its proper supersets, so the positive teaching dimension is infinite.
\end{proof}
\noindent\textbf{Reference.} \cite{balbach2008measuring}.

\subsection{Positive Teaching Dimension of Tagged Singletons}\label{val:positive_teaching_dimension_tagged_singletons}
\noindent\textbf{Statement.} The Positive Teaching Dimension of Tagged Singletons is $\infty$.
\begin{proof}
The class contains the empty concept and other concepts. Its empty target has no positive examples, so its only positive sample is consistent with every other concept and cannot teach it.
\end{proof}
\noindent\textbf{Reference.} \cite{balbach2008measuring}.

\subsection{Positive Teaching Dimension of Majority}\label{val:positive_teaching_dimension_majority}
\noindent\textbf{Statement.} The Positive Teaching Dimension of Majority is $\infty$.
\begin{proof}
The base all-zero labeling is the empty concept in the class. It has no positive examples, and its empty positive sample is consistent with the other Majority concepts, so it cannot be positively taught.
\end{proof}
\noindent\textbf{Reference.} \cite{balbach2008measuring}.

\subsection{Positive Teaching Dimension of Addressing}\label{val:positive_teaching_dimension_addressing}
\noindent\textbf{Statement.} The Positive Teaching Dimension of Addressing is $1\quad(n\ge2)$.
\begin{proof}
Each addressing concept has a unique positive primary coordinate, so that one positive example teaches it. The class is nontrivial, so zero examples cannot teach every target.
\end{proof}
\noindent\textbf{Reference.} \cite{balbach2008measuring}.

\subsection{Positive Recursive Teaching Dimension of Trivial Concepts}\label{val:positive_recursive_teaching_dimension_trivial_concepts}
\noindent\textbf{Statement.} The Positive Recursive Teaching Dimension of Trivial concepts is $1$.
\begin{proof}
Remove the full concept using one positive coordinate, then the empty concept remains alone and needs no sample. A nontrivial class cannot have recursive value zero.
\end{proof}
\noindent\textbf{Reference.} \cite{fallat2023batch}.

\subsection{Positive Recursive Teaching Dimension of Singletons}\label{val:positive_recursive_teaching_dimension_singletons}
\noindent\textbf{Statement.} The Positive Recursive Teaching Dimension of Singletons is $1$.
\begin{proof}
At every nontrivial stage, any singleton is taught by its unique positive coordinate. Thus all concepts can be removed with positive samples of size one.
\end{proof}
\noindent\textbf{Reference.} \cite{fallat2023batch}.

\subsection{Positive Recursive Teaching Dimension of Full Cube}\label{val:positive_recursive_teaching_dimension_full_cube}
\noindent\textbf{Statement.} The Positive Recursive Teaching Dimension of Full Cube is $n$.
\begin{proof}
Remove cube concepts in decreasing inclusion order. A concept is positively taught among the remaining subsets by all of its positive coordinates, giving samples of size at most $n$. Initially the full concept needs all $n$ positive coordinates to exclude each neighbour missing one coordinate, so the bound is exact.
\end{proof}
\noindent\textbf{Reference.} \cite{fallat2023batch}.

\subsection{Positive Recursive Teaching Dimension of Singletons plus Empty Set}\label{val:positive_recursive_teaching_dimension_singletons_plus_empty_set}
\noindent\textbf{Statement.} The Positive Recursive Teaching Dimension of Singletons plus empty set is $1$.
\begin{proof}
Remove every singleton using its positive coordinate; the empty concept is then alone and needs no sample. Since the class is nontrivial, zero is impossible.
\end{proof}
\noindent\textbf{Reference.} \cite{fallat2023batch}.

\subsection{Positive Recursive Teaching Dimension of Half-intervals}\label{val:positive_recursive_teaching_dimension_halfintervals}
\noindent\textbf{Statement.} The Positive Recursive Teaching Dimension of Half-intervals is $1$.
\begin{proof}
Remove initial segments from largest to smallest. At each stage, the current largest endpoint is a positive coordinate appearing in no remaining smaller segment, so one positive example teaches it.
\end{proof}
\noindent\textbf{Reference.} \cite{fallat2023batch}.

\subsection{Positive Recursive Teaching Dimension of Tagged Singletons}\label{val:positive_recursive_teaching_dimension_tagged_singletons}
\noindent\textbf{Statement.} The Positive Recursive Teaching Dimension of Tagged Singletons is $1$.
\begin{proof}
First remove tagged pairs using their unique positive ordinary coordinate. Then remove tag singletons using their tag coordinate, and finally remove the empty concept. Every nonempty stage uses a positive sample of size at most one.
\end{proof}
\noindent\textbf{Reference.} \cite{fallat2023batch}.

\subsection{Positive Recursive Teaching Dimension of Addressing}\label{val:positive_recursive_teaching_dimension_addressing}
\noindent\textbf{Statement.} The Positive Recursive Teaching Dimension of Addressing is $1\quad(n\ge2)$.
\begin{proof}
Every addressing concept has a unique positive primary coordinate, so the concepts may be removed in any order with one positive teaching example each.
\end{proof}
\noindent\textbf{Reference.} \cite{fallat2023batch}.

\subsection{Positive No-Clashing Teaching Dimension of Half-intervals}\label{val:positive_noclashing_teaching_dimension_halfintervals}
\noindent\textbf{Statement.} The Positive No-Clashing Teaching Dimension of Half-intervals is $1$.
\begin{proof}
Assign initial segment $[i]$ its positive endpoint example $(i,1)$. For $i<j$, $[i]$ is inconsistent with $(j,1)$, so no two distinct concepts clash. A nontrivial class cannot have width zero.
\end{proof}
\noindent\textbf{Reference.} \cite{KSZ2019}.

\subsection{Positive No-Clashing Teaching Dimension of Tagged Singletons}\label{val:positive_noclashing_teaching_dimension_tagged_singletons}
\noindent\textbf{Statement.} The Positive No-Clashing Teaching Dimension of Tagged Singletons is $1$.
\begin{proof}
Assign a tagged pair its unique positive ordinary coordinate, a tag singleton its tag coordinate, and the empty concept the empty sample. For every distinct pair, one assigned positive example is absent from the other concept, so the mapping is positive and non-clashing. Width zero is impossible for a nontrivial class.
\end{proof}
\noindent\textbf{Reference.} \cite{KSZ2019}.

\subsection{Positive No-Clashing Teaching Dimension of Addressing}\label{val:positive_noclashing_teaching_dimension_addressing}
\noindent\textbf{Statement.} The Positive No-Clashing Teaching Dimension of Addressing is $1\quad(n\ge2)$.
\begin{proof}
Assign every concept its unique positive primary coordinate. Any other addressing concept is negative there, so the resulting width-one mapping is positive and non-clashing. The class is nontrivial, ruling out zero.
\end{proof}
\noindent\textbf{Reference.} \cite{KSZ2019}.

\subsection{No-Clashing Teaching Dimension of Half-intervals}\label{val:noclashing_teaching_dimension_halfintervals}
\noindent\textbf{Statement.} The No-Clashing Teaching Dimension of Half-intervals is $1$.
\begin{proof}
The endpoint assignment $[i]\mapsto(i,1)$ is already a positive non-clashing teacher mapping of width one. Width zero cannot distinguish a nontrivial class.
\end{proof}
\noindent\textbf{Reference.} \cite{KSZ2019}.

\subsection{No-Clashing Teaching Dimension of Tagged Singletons}\label{val:noclashing_teaching_dimension_tagged_singletons}
\noindent\textbf{Statement.} The No-Clashing Teaching Dimension of Tagged Singletons is $1$.
\begin{proof}
The positive width-one mapping assigning ordinary coordinates to tagged pairs, tag coordinates to tag singletons, and the empty sample to the empty concept is non-clashing. Width zero is impossible.
\end{proof}
\noindent\textbf{Reference.} \cite{KSZ2019}.

\subsection{No-Clashing Teaching Dimension of Addressing}\label{val:noclashing_teaching_dimension_addressing}
\noindent\textbf{Statement.} The No-Clashing Teaching Dimension of Addressing is $1\quad(n\ge2)$.
\begin{proof}
Assigning every concept its unique positive primary coordinate is a width-one non-clashing mapping. The nontrivial class rules out width zero.
\end{proof}
\noindent\textbf{Reference.} \cite{KSZ2019}.

\subsection{Labeled Sample Compression of Trivial concepts}\label{val:labeled_sample_compression_trivial_concepts}
\noindent\textbf{Statement.} The Labeled Sample Compression of Trivial concepts is $1$.
\begin{proof}
Every nonempty realizable sample has a uniform label; retain one labeled point and reconstruct the corresponding constant concept.  The empty sample has a fixed reconstruction.  Size zero is impossible because the two constant concepts have incompatible one-point samples.
\end{proof}
\noindent\textbf{Reference.} \cite{FW1995}.

\subsection{Labeled Sample Compression of Singletons}\label{val:labeled_sample_compression_singletons}
\noindent\textbf{Statement.} The Labeled Sample Compression of Singletons is $0$ for $n=1$; $1$ for $n\ge2$.
\begin{proof}
For $n=1$, always reconstruct the sole concept from the empty key. For $n\ge2$, fix a cyclic successor map. If a sample has a positive point, retain its labeled positive point and reconstruct that singleton. For a nonempty all-negative sample on a proper subset $A$ of the domain, retain a negative point of $A$ whose successor lies outside $A$, and reconstruct the successor singleton. Such a point exists because the cycle is not contained in $A$. Reconstruct any fixed singleton at the empty key, which handles the empty sample. All reconstructions are consistent. Two distinct full concepts rule out width zero.
\end{proof}
\noindent\textbf{Reference.} \cite{FW1995}.

\subsection{Labeled Sample Compression of Singletons plus empty set}\label{val:labeled_sample_compression_singletons_plus_empty_set}
\noindent\textbf{Statement.} The Labeled Sample Compression of Singletons plus empty set is $1$.
\begin{proof}
Use the singleton construction when a positive point is present.  For an all-negative sample, use the empty reconstruction when no point is retained; if a singleton must be reconstructed, the same cyclic-boundary choice as for singletons supplies one retained negative point.  The scheme is consistent, and size zero cannot handle incompatible positive and negative one-point samples.
\end{proof}
\noindent\textbf{Reference.} \cite{FW1995}.

\subsection{Labeled Sample Compression of Half-intervals}\label{val:labeled_sample_compression_halfintervals}
\noindent\textbf{Statement.} The Labeled Sample Compression of Half-intervals is $1$.
\begin{proof}
For a sample with a positive point, retain its greatest positive coordinate and reconstruct the corresponding initial segment.  For an all-negative nonempty sample, retain its least negative coordinate and reconstruct the preceding initial segment; use a fixed interval for the empty sample.  These reconstructions are consistent.  Zero is impossible for a nontrivial class.
\end{proof}
\noindent\textbf{Reference.} \cite{FW1995}.

\subsection{Proper Labeled Sample Compression of Trivial concepts}\label{val:proper_labeled_sample_compression_trivial_concepts}
\noindent\textbf{Statement.} The Proper Labeled Sample Compression of Trivial concepts is $1$.
\begin{proof}
The size-one labeled scheme reconstructs one of the two constant concepts, hence is proper.  The same incompatible one-point samples rule out size zero.
\end{proof}
\noindent\textbf{Reference.} \cite{FW1995}.

\subsection{Proper Labeled Sample Compression of Singletons}\label{val:proper_labeled_sample_compression_singletons}
\noindent\textbf{Statement.} The Proper Labeled Sample Compression of Singletons is $0$ for $n=1$; $1$ for $n\ge2$.
\begin{proof}
The cyclic-boundary labeled scheme reconstructs only singleton concepts, so it is proper and has width one for $n\ge2$. Two distinct full concepts rule out width zero. For $n=1$, the empty key always reconstructs the sole concept, giving width zero.
\end{proof}
\noindent\textbf{Reference.} \cite{FW1995}.

\subsection{Proper Labeled Sample Compression of Singletons plus empty set}\label{val:proper_labeled_sample_compression_singletons_plus_empty_set}
\noindent\textbf{Statement.} The Proper Labeled Sample Compression of Singletons plus empty set is $1$.
\begin{proof}
The size-one construction reconstructs either a singleton or the empty concept, so every reconstruction lies in the class.  Incompatible one-point samples rule out size zero.
\end{proof}
\noindent\textbf{Reference.} \cite{FW1995}.

\subsection{Proper Labeled Sample Compression of Half-intervals}\label{val:proper_labeled_sample_compression_halfintervals}
\noindent\textbf{Statement.} The Proper Labeled Sample Compression of Half-intervals is $1$.
\begin{proof}
The endpoint construction reconstructs an initial segment in every case, hence is a proper size-one scheme.  Since the class is nontrivial, zero is impossible.
\end{proof}
\noindent\textbf{Reference.} \cite{FW1995}.

\subsection{Stable Labeled Sample Compression of Trivial concepts}\label{val:stable_labeled_sample_compression_trivial_concepts}
\noindent\textbf{Statement.} The Stable Labeled Sample Compression of Trivial concepts is $1$.
\begin{proof}
Fix an order of the domain and retain the least sampled coordinate.  Deleting any unselected sample point leaves this retained point and its reconstruction unchanged, so the size-one constant-concept scheme is stable.  Zero is impossible.
\end{proof}
\noindent\textbf{Reference.} \cite{FW1995}.

\subsection{Stable Labeled Sample Compression of Half-intervals}\label{val:stable_labeled_sample_compression_halfintervals}
\noindent\textbf{Statement.} The Stable Labeled Sample Compression of Half-intervals is $1$.
\begin{proof}
Retain the greatest positive sample coordinate when one exists, otherwise the least negative coordinate.  Removing an unselected point preserves that selected extremum and the endpoint reconstruction, making the size-one scheme stable.  Zero is impossible.
\end{proof}
\noindent\textbf{Reference.} \cite{FW1995}.

\subsection{Proper Stable Labeled Sample Compression of Half-intervals}\label{val:proper_stable_labeled_sample_compression_halfintervals}
\noindent\textbf{Statement.} The Proper Stable Labeled Sample Compression of Half-intervals is $1$.
\begin{proof}
The stable endpoint scheme always reconstructs an initial segment, hence lies within the class.  It is therefore proper and stable of size one; zero is impossible.
\end{proof}
\noindent\textbf{Reference.} \cite{FW1995}.

\subsection{Majority Complexity of Trivial concepts}\label{val:majority_complexity_trivial_concepts}
\noindent\textbf{Statement.} The Majority Complexity of Trivial concepts is $1$.
\begin{proof}
The pointwise-majority equivalence query is one of the two constant concepts.  If it is wrong, any counterexample identifies the other constant concept.  A nontrivial two-concept class requires one query in the worst case.
\end{proof}
\noindent\textbf{Reference.} \cite{maass1992lower}.

\subsection{Majority Complexity of Singletons}\label{val:majority_complexity_singletons}
\noindent\textbf{Statement.} The Majority Complexity of Singletons is $1$.
\begin{proof}
With the zero tie-breaking convention, the pointwise majority of the singleton concepts is the empty hypothesis.  Its equivalence query either succeeds only in the degenerate one-concept case or returns the target's unique positive coordinate, which identifies the singleton.  For a nontrivial singleton class, zero queries cannot identify the target.
\end{proof}
\noindent\textbf{Reference.} \cite{maass1992lower}.

\subsection{Majority Complexity of Singletons plus empty set}\label{val:majority_complexity_singletons_plus_empty_set}
\noindent\textbf{Statement.} The Majority Complexity of Singletons plus empty set is $1$.
\begin{proof}
The pointwise majority is the empty hypothesis.  If it is incorrect, the returned positive counterexample is the unique point of the target singleton.  Thus one query suffices, while zero cannot distinguish the empty concept from a singleton.
\end{proof}
\noindent\textbf{Reference.} \cite{maass1992lower}.

\subsection{Majority Complexity of Full Cube}\label{val:majority_complexity_full_cube}
\noindent\textbf{Statement.} The Majority Complexity of Full Cube is $n$.
\begin{proof}
The halving algorithm makes at most $\log_2|2^{[n]}|=n$ equivalence queries.  Conversely the full cube has Littlestone dimension $n$, and the halving learner is a particular equivalence-query learner, so it needs at least $n$ in the worst case.
\end{proof}
\noindent\textbf{Reference.} \cite{littlestone1988learning,maass1992lower}.

\subsection{Random Equivalence Queries Complexity of Trivial concepts}\label{val:random_equivalence_queries_complexity_trivial_concepts}
\noindent\textbf{Statement.} The Random Equivalence Queries Complexity of Trivial concepts is $1$.
\begin{proof}
Query either constant concept.  If it is wrong, any counterexample establishes that the other constant is the target, independently of the random counterexample rule.  Zero queries cannot distinguish the two targets.
\end{proof}
\noindent\textbf{Reference.} \cite{maass1992lower}.

\subsection{Random Proper Equivalence Queries Complexity of Trivial concepts}\label{val:random_proper_equivalence_queries_complexity_trivial_concepts}
\noindent\textbf{Statement.} The Random Proper Equivalence Queries Complexity of Trivial concepts is $1$.
\begin{proof}
The one-query constant strategy is proper because both proposed constants belong to the class.  A failed query identifies the other concept, and zero queries are impossible.
\end{proof}
\noindent\textbf{Reference.} \cite{maass1992lower}.

\subsection{Random Equivalence Queries Complexity of Singletons plus empty set}\label{val:random_equivalence_queries_complexity_singletons_plus_empty_set}
\noindent\textbf{Statement.} The Random Equivalence Queries Complexity of Singletons plus empty set is $1$.
\begin{proof}
Query the empty concept.  It succeeds for the empty target, and otherwise its unique positive counterexample identifies the target singleton.  Thus the random counterexample rule does not change the one-query bound, while zero queries are impossible.
\end{proof}
\noindent\textbf{Reference.} \cite{maass1992lower}.

\subsection{Random Proper Equivalence Queries Complexity of Singletons plus empty set}\label{val:random_proper_equivalence_queries_complexity_singletons_plus_empty_set}
\noindent\textbf{Statement.} The Random Proper Equivalence Queries Complexity of Singletons plus empty set is $1$.
\begin{proof}
The empty hypothesis is itself in the class, so the preceding one-query strategy is proper.  A positive counterexample identifies the target singleton; zero queries do not distinguish all concepts.
\end{proof}
\noindent\textbf{Reference.} \cite{maass1992lower}.

\subsection{Projected Recursive Teaching Dimension of Full Cube}\label{val:monotone_recursive_teaching_dimension_full_cube}
\noindent\textbf{Statement.} The Projected Recursive Teaching Dimension of Full Cube is $n$.
\begin{proof}
Every projection of the full cube is the full cube on its retained coordinates, whose recursive teaching dimension is its number of coordinates.  The full-domain projection has value $n$, and no projection has more than $n$ coordinates.
\end{proof}
\noindent\textbf{Reference.} \cite{doliwa2014recursive}.

\subsection{Projected Recursive Teaching Dimension of Singletons plus empty set}\label{val:monotone_recursive_teaching_dimension_singletons_plus_empty_set}
\noindent\textbf{Statement.} The Projected Recursive Teaching Dimension of Singletons plus empty set is $1$.
\begin{proof}
A coordinate projection is again a singleton-plus-empty class on a smaller domain.  Remove all singleton concepts using their unique positive coordinate, then the empty concept; each removal step has teaching-set size at most one.  A nontrivial projection rules out zero.
\end{proof}
\noindent\textbf{Reference.} \cite{doliwa2014recursive}.

\subsection{Projected Recursive Teaching Dimension of Half-intervals}\label{val:monotone_recursive_teaching_dimension_halfintervals}
\noindent\textbf{Statement.} The Projected Recursive Teaching Dimension of Half-intervals is $1$.
\begin{proof}
Restricting initial segments to any subset of the ordered domain again gives a chain of initial segments, possibly with duplicates removed.  Remove the chain from either endpoint: each current endpoint is taught by one adjacent coordinate when another concept remains.  Thus every projection has recursive teaching dimension at most one, and a nontrivial projection gives equality.
\end{proof}
\noindent\textbf{Reference.} \cite{doliwa2014recursive}.

\subsection{Projected Membership Queries Complexity of Full Cube}\label{val:projected_membership_queries_complexity_full_cube}
\noindent\textbf{Statement.} The Projected Membership Queries Complexity of Full Cube is $n$.
\begin{proof}
Every projection is a full cube on at most $n$ coordinates, and identifying a target in a $k$-cube requires and is achieved by querying all $k$ coordinates.  The full-domain projection supplies the matching value $n$.
\end{proof}
\noindent\textbf{Reference.} \cite{angluin1988queries}.

\subsection{Projected Membership Queries Complexity of Singletons plus empty set}\label{val:projected_membership_queries_complexity_singletons_plus_empty_set}
\noindent\textbf{Statement.} The Projected Membership Queries Complexity of Singletons plus empty set is $n$.
\begin{proof}
A projection to $k$ coordinates is again the empty concept together with the $k$ singleton concepts.  To identify the empty target requires $k$ negative answers, and querying all coordinates suffices.  The full-domain projection has $k=n$.
\end{proof}
\noindent\textbf{Reference.} \cite{angluin1988queries}.

\subsection{Extended Teaching Dimension of Half-intervals}\label{val:extended_teaching_dimension_halfintervals}
\noindent\textbf{Statement.} The Extended Teaching Dimension of Half-intervals is $2$.
\begin{proof}
For an arbitrary target labeling, if its first label is zero then that one example excludes every initial segment; if its last label is one, that example isolates the full interval.  Otherwise a consecutive $1$-then-$0$ transition exists, and its two labeled coordinates isolate the corresponding initial segment.  Conversely, an interior initial segment has ordinary teaching dimension two, so XTD is two (for $n\ge3$).
\end{proof}
\noindent\textbf{Reference.} \cite{hegedHus1995generalized}.

\subsection{Membership and Equivalence Queries Complexity of Singletons plus empty set}\label{val:membership_and_equivalence_queries_complexity_singletons_plus_empty_set}
\noindent\textbf{Statement.} The Membership and Equivalence Queries Complexity of Singletons plus empty set is $1$.
\begin{proof}
Ask the equivalence query for the empty concept.  It succeeds for the empty target; otherwise the sole positive counterexample is the unique point of the target singleton.  Thus one query suffices, while zero cannot distinguish the nontrivial class.
\end{proof}
\noindent\textbf{Reference.} \cite{angluin1988queries}.

\subsection{Membership and Proper Equivalence Queries Complexity of Singletons plus empty set}\label{val:membership_and_proper_equivalence_queries_complexity_singletons_plus_empty_set}
\noindent\textbf{Statement.} The Membership and Proper Equivalence Queries Complexity of Singletons plus empty set is $1$.
\begin{proof}
The empty concept is a member of the class, so the preceding one-query equivalence strategy is proper.  Its positive counterexample identifies the singleton target, and zero queries are impossible for a nontrivial class.
\end{proof}
\noindent\textbf{Reference.} \cite{angluin1988queries}.

\subsection{Partial Equivalence Queries Complexity of Singletons plus empty set}\label{val:partial_equivalence_queries_complexity_singletons_plus_empty_set}
\noindent\textbf{Statement.} The Partial Equivalence Queries Complexity of Singletons plus empty set is $1$.
\begin{proof}
Ask the total partial-equivalence query labeling every coordinate zero.  A consistent reply identifies the empty target; otherwise the supplied disagreement is the unique positive coordinate of the singleton target.  A nontrivial class cannot be identified with zero queries.
\end{proof}
\noindent\textbf{Reference.} \cite{maass1992lower}.

\subsection{Projected Minimum Degree of Half-intervals}\label{val:projected_minimum_degree_halfintervals}
\noindent\textbf{Statement.} The Projected Minimum Degree of Half-intervals is $1$.
\begin{proof}
Every nonempty coordinate projection of the ordered initial segments is again a chain of initial segments.  Its $1$-inclusion graph is a path, so its endpoint vertices have degree one and its minimum degree is one.  The empty projection has degree zero, hence the maximum over projections is one.
\end{proof}

\subsection{Projected Double Density or Projected Average Degree of Half-intervals}\label{val:projected_double_density_or_projected_average_degree_halfintervals}
\noindent\textbf{Statement.} The Projected Double Density or Projected Average Degree of Half-intervals is $\frac{2(n-1)}{n}$.
\begin{proof}
Every coordinate projection of the initial-segment class is a chain, so its $1$-inclusion graph is a path.  It has at most $n$ distinct vertices, whose average degree is at most $2(n-1)/n$.  The full projection is already a path on $n$ intervals (and so is a projection omitting the first coordinate), attaining this value.
\end{proof}

\subsection{Order Sample Compression of Half-intervals}\label{val:order_sample_compression_halfintervals}
\noindent\textbf{Statement.} The Order Sample Compression of Half-intervals is $1$.
\begin{proof}
Half-intervals have VC dimension one.  The order-compression construction for VC-dimension-one classes gives size at most one.  Size zero is impossible for this nontrivial class: an empty compression set has one fixed reconstruction and cannot be consistent with samples distinguishing two concepts.
\end{proof}
\noindent\textbf{Reference.} \cite{darnstadt2016order}.

\subsection{Proper Ordered Sample Compression of Half-intervals}\label{val:proper_ordered_sample_compression_halfintervals}
\noindent\textbf{Statement.} The Proper Ordered Sample Compression of Half-intervals is $1$.
\begin{proof}
The half-interval class is intersection-closed.  The proper order-compression construction for intersection-closed classes has size at most their VC dimension, which is one here.  A nontrivial class rules out size zero.
\end{proof}
\noindent\textbf{Reference.} \cite{darnstadt2016order}.

\subsection{Unlabeled Sample Compression of Half-intervals}\label{val:unlabeled_sample_compression_halfintervals}
\noindent\textbf{Statement.} The Unlabeled sample compression of Half-intervals is $1$.
\begin{proof}
The corrected unlabeled-compression construction for every maximum class has size equal to its VC dimension.  Applied to the effective-coordinate realization of this class, it gives size at most one.  The VC-dimension lower bound for unlabeled compression, together with $\mathrm{VC}=1$, rules out size zero.
\end{proof}
\noindent\textbf{Reference.} \cite{chalopin2022unlabeled,palvolgyi2021unlabeled}.

\subsection{Proper Unlabeled Sample Compression of Half-intervals}\label{val:proper_unlabeled_sample_compression_halfintervals}
\noindent\textbf{Statement.} The Proper Unlabeled Sample Compression of Half-intervals is $1$.
\begin{proof}
For a maximum class, the representation-map construction selects a subset of the sample domain and reconstructs the unique represented concept consistent with the full sample; it is therefore proper.  On the effective-coordinate half-interval class this gives a proper unlabeled scheme of size one.  Proper unlabeled compression is at least ordinary unlabeled compression, whose value here is one.
\end{proof}
\noindent\textbf{Reference.} \cite{chalopin2022unlabeled,palvolgyi2021unlabeled}.

\subsection{Unlabeled Sample Compression of Full Cube}\label{val:unlabeled_sample_compression_full_cube}
\noindent\textbf{Statement.} The Unlabeled sample compression of Full Cube is $n$.
\begin{proof}
The corrected maximum-class construction gives an unlabeled compression scheme of size equal to the VC dimension, hence at most $n$.  Conversely, the full domain is shattered, so the general VC-dimension lower bound for unlabeled compression gives at least $n$.
\end{proof}
\noindent\textbf{Reference.} \cite{chalopin2022unlabeled,palvolgyi2021unlabeled}.

\subsection{Proper Unlabeled Sample Compression of Full Cube}\label{val:proper_unlabeled_sample_compression_full_cube}
\noindent\textbf{Statement.} The Proper Unlabeled Sample Compression of Full Cube is $n$.
\begin{proof}
The maximum-class representation-map construction reconstructs a consistent concept.  Since the full cube contains every concept on its domain, this is a proper unlabeled scheme of size at most $n$.  Proper unlabeled compression is at least ordinary unlabeled compression, whose exact value here is $n$.
\end{proof}
\noindent\textbf{Reference.} \cite{chalopin2022unlabeled,palvolgyi2021unlabeled}.

\subsection{Minimum Degree of Tagged Singletons}\label{val:minimum_degree_tagged_singletons}
\noindent\textbf{Statement.} The Minimum Degree of Tagged Singletons is $1$.
\begin{proof}
Each tagged pair $\{i,n+f(i)\}$ differs in one coordinate only from its tag singleton, and in at least two coordinates from every other class member, so it has degree one.  Every tag singleton is adjacent to $\emptyset$, and $\emptyset$ is adjacent to every tag singleton; therefore there are no isolated vertices.
\end{proof}
\noindent\textbf{Reference.} \cite{maass1992lower}.

\subsection{Double Density or Average Degree of Tagged Singletons}\label{val:double_density_or_average_degree_tagged_singletons}
\noindent\textbf{Statement.} The Double Density or Average Degree of Tagged Singletons is $\frac{2(n+\lceil\log n\rceil)}{n+\lceil\log n\rceil+1}$.
\begin{proof}
Writing $m=\lceil\log n\rceil$, the $1$-inclusion graph has the empty concept, $m$ tag singletons, and $n$ tagged pairs: hence $n+m+1$ vertices.  Its only edges are the $m$ edges from $\emptyset$ to tag singletons and the $n$ edges from a tagged pair to its tag singleton.  It therefore has $n+m$ edges, and its average degree is $2(n+m)/(n+m+1)$.
\end{proof}
\noindent\textbf{Reference.} \cite{maass1992lower}.

\subsection{Densest Subgraph (twice) of Tagged Singletons}\label{val:densest_subgraph_twice_tagged_singletons}
\noindent\textbf{Statement.} The Densest Subgraph (twice) of Tagged Singletons is $\frac{2(n+\lceil\log n\rceil)}{n+\lceil\log n\rceil+1}$.
\begin{proof}
The $1$-inclusion graph is a tree with $n+\lceil\log n\rceil+1$ vertices.  Every subgraph is a forest, so a subgraph on $v$ vertices has average degree at most $2(v-1)/v$, which increases with $v$.  The full tree attains this bound, so it is a densest subgraph and has the stated double density.
\end{proof}
\noindent\textbf{Reference.} \cite{maass1992lower}.

\subsection{Minimum Degree of Majority}\label{val:minimum_degree_majority}
\noindent\textbf{Statement.} The Minimum Degree of Majority is $\left\lfloor\frac{n-1}{2}\right\rfloor$.
\begin{proof}
The $1$-inclusion graph is the $n$-cube on the base coordinates, except that every edge from layer $t=\lfloor n/2\rfloor$ to layer $t+1$ is removed because it changes the majority bit.  A vertex in layer $t$ has degree $t$, while one in layer $t+1$ has degree $n-t-1$; all other vertices retain degree $n$.  The minimum is therefore $\min(t,n-t-1)=\lfloor(n-1)/2\rfloor$.
\end{proof}
\noindent\textbf{Reference.} \cite{maass1992lower}.

\subsection{Double Density or Average Degree of Majority}\label{val:double_density_or_average_degree_majority}
\noindent\textbf{Statement.} The Double Density or Average Degree of Majority is $n-\frac{(n-\lfloor n/2\rfloor)\binom{n}{\lfloor n/2\rfloor}}{2^{n-1}}$.
\begin{proof}
There is one concept for each base subset, so the graph has $2^n$ vertices.  Starting from the $n2^{n-1}$ edges of the $n$-cube, exactly $(n-t)\binom nt$ edges between layers $t=\lfloor n/2\rfloor$ and $t+1$ are deleted because their endpoints have different majority bits.  Dividing twice the remaining edge count by $2^n$ gives the displayed average degree.
\end{proof}
\noindent\textbf{Reference.} \cite{maass1992lower}.

\subsection{Minimum Degree of Addressing}\label{val:minimum_degree_addressing}
\noindent\textbf{Statement.} The Minimum Degree of Addressing is $0$.
\begin{proof}
Distinct addressing concepts have different addressed primary coordinates and different address words, so their symmetric difference has size at least two.  Thus the $1$-inclusion graph is edgeless and every vertex has degree zero.
\end{proof}
\noindent\textbf{Reference.} \cite{maass1992lower}.

\subsection{Double Density or Average Degree of Addressing}\label{val:double_density_or_average_degree_addressing}
\noindent\textbf{Statement.} The Double Density or Average Degree of Addressing is $0$.
\begin{proof}
The unprojected $1$-inclusion graph is edgeless: two distinct addressing concepts differ both in their primary coordinate and in at least one address coordinate.  Therefore its average degree is zero.
\end{proof}
\noindent\textbf{Reference.} \cite{maass1992lower}.

\subsection{Densest Subgraph (twice) of Addressing}\label{val:densest_subgraph_twice_addressing}
\noindent\textbf{Statement.} The Densest Subgraph (twice) of Addressing is $0$.
\begin{proof}
Every subgraph of the edgeless unprojected $1$-inclusion graph is edgeless.  Hence its maximum subgraph average degree, equivalently densest-subgraph double density, is zero.
\end{proof}
\noindent\textbf{Reference.} \cite{maass1992lower}.

\subsection{Maximum Degree of Warmuth's $C_5$ Class}\label{val:maximum_degree_warmuth_c5}
\noindent\textbf{Statement.} The Maximum Degree of Warmuth's $C_5$ class is $1$.
\begin{proof}
For each cyclic placement of the prescribed pattern $1001$, the two choices for the one unrestricted bit differ in exactly that bit, giving one edge.  Direct comparison of the five cyclic placements shows that concepts from different placements differ in at least two bits.  Thus the $1$-inclusion graph is a matching of five edges and has maximum degree one.
\end{proof}
\noindent\textbf{Reference.} \cite{palvolgyi2021unlabeled}.

\subsection{Minimum Degree of Warmuth's $C_5$ Class}\label{val:minimum_degree_warmuth_c5}
\noindent\textbf{Statement.} The Minimum Degree of Warmuth's $C_5$ class is $1$.
\begin{proof}
The $1$-inclusion graph pairs the two concepts obtained from each fixed cyclic placement of $1001$ by its unrestricted bit.  These five edges are disjoint and exhaust the graph, so every vertex has degree one.
\end{proof}
\noindent\textbf{Reference.} \cite{palvolgyi2021unlabeled}.

\subsection{Double Density or Average Degree of Warmuth's $C_5$ Class}\label{val:double_density_or_average_degree_warmuth_c5}
\noindent\textbf{Statement.} The Double Density or Average Degree of Warmuth's $C_5$ class is $1$.
\begin{proof}
Its $1$-inclusion graph is a matching of five edges on ten vertices.  Hence its average degree is $2\cdot5/10=1$.
\end{proof}
\noindent\textbf{Reference.} \cite{palvolgyi2021unlabeled}.

\subsection{Densest Subgraph (twice) of Warmuth's $C_5$ Class}\label{val:densest_subgraph_twice_warmuth_c5}
\noindent\textbf{Statement.} The Densest Subgraph (twice) of Warmuth's $C_5$ class is $1$.
\begin{proof}
Every subgraph of a matching has maximum degree at most one and therefore average degree at most one.  Any one of the five edges attains average degree one, so the densest-subgraph double density is exactly one.
\end{proof}
\noindent\textbf{Reference.} \cite{palvolgyi2021unlabeled}.

\subsection{Hamming Radius of Warmuth's $C_5$ Class}\label{val:hamming_radius_warmuth_c5}
\noindent\textbf{Statement.} The Hamming Radius of Warmuth's $C_5$ class is $3$.
\begin{proof}
The all-zero center is within distance three of every concept, since the ten $C_5$ strings have Hamming weight two or three.  Checking the 32 binary centers against the five rotations of $10010$ and the five rotations of $10011$ shows that each center is at distance at least three from one of these ten strings.  Thus radius two is impossible and the radius is exactly three.
\end{proof}
\noindent\textbf{Reference.} \cite{palvolgyi2021unlabeled}.

\subsection{Oriented Diameter of Warmuth's $C_5$ Class}\label{val:oriented_diameter_warmuth_c5}
\noindent\textbf{Statement.} The Oriented Diameter of Warmuth's $C_5$ class is $2$.
\begin{proof}
The two concepts $10010$ and $01001$ have an oriented disagreement on two coordinates, giving a witness of size two.  Direct comparison of the ten cyclic $C_5$ strings shows that no ordered pair has three coordinates on which the first is $0$ and the second is $1$; hence no larger oriented pair occurs in a projection.
\end{proof}
\noindent\textbf{Reference.} \cite{palvolgyi2021unlabeled}.

\subsection{Threshold Dimension of Warmuth's $C_5$ Class}\label{val:threshold_dimension_warmuth_c5}
\noindent\textbf{Statement.} The Threshold Dimension of Warmuth's $C_5$ class is $2$.
\begin{proof}
On two adjacent cyclic coordinates, the ten concepts realize the traces $00,10,11$, yielding a threshold chain of length two.  Enumerating the traces on each of the ten three-coordinate projections shows that none contains all four initial-segment traces of a length-three threshold chain.  Thus the threshold dimension is exactly two.
\end{proof}
\noindent\textbf{Reference.} \cite{palvolgyi2021unlabeled}.

\subsection{Path Dimension of Warmuth's $C_5$ Class}\label{val:path_dimension_warmuth_c5}
\noindent\textbf{Statement.} The Path Dimension of Warmuth's $C_5$ class is $4$.
\begin{proof}
On the coordinate set $\{0,1,2,4\}$, the ordering $0,1,2,4$ and starting trace $0110$ give a full path through five realized traces.  A length-five path would have to use the full five-coordinate projection; enumerating its $5!$ coordinate orderings and starting traces finds no full path of length five.  Hence the path dimension is four.
\end{proof}
\noindent\textbf{Reference.} \cite{palvolgyi2021unlabeled}.

\subsection{Degeneracy of Warmuth's $C_5$ Class}\label{val:degeneracy_warmuth_c5}
\noindent\textbf{Statement.} The Degeneracy of Warmuth's $C_5$ class is $1$.
\begin{proof}
The $1$-inclusion graph is a matching of five edges.  Every nonempty subgraph therefore has a vertex of degree at most one, and an edge subgraph has minimum degree one.  Thus its degeneracy is exactly one.
\end{proof}
\noindent\textbf{Reference.} \cite{palvolgyi2021unlabeled}.

\subsection{Largest Strongly Shattered Set of Warmuth's $C_5$ Class}\label{val:largest_strongly_shattered_set_warmuth_c5}
\noindent\textbf{Statement.} The Largest Strongly Shattered Set of Warmuth's $C_5$ class is $1$.
\begin{proof}
Each $1$-inclusion edge is a one-dimensional subcube, so a singleton coordinate set is strongly shattered.  A strongly shattered two-set would produce a two-dimensional cube and hence a $4$-cycle in the $1$-inclusion graph, impossible because that graph is a matching.  Therefore the largest strongly shattered set has size one.
\end{proof}
\noindent\textbf{Reference.} \cite{palvolgyi2021unlabeled}.

\subsection{Maximum Positive Degree of Warmuth's $C_5$ Class}\label{val:maximum_positive_degree_warmuth_c5}
\noindent\textbf{Statement.} The Maximum Positive Degree of Warmuth's $C_5$ class is $1$.
\begin{proof}
The raw $1$-inclusion graph is a matching.  In each edge, the concept obtained by setting the unrestricted bit to one properly contains the one with that bit set to zero, so it has one smaller neighbor.  No vertex can have more, because the graph has maximum degree one.
\end{proof}
\noindent\textbf{Reference.} \cite{palvolgyi2021unlabeled}.

\subsection{Projected Maximal Degree of Warmuth's $C_5$ Class}\label{val:projected_maximal_degree_warmuth_c5}
\noindent\textbf{Statement.} The Projected Maximal Degree of Warmuth's $C_5$ class is $3$.
\begin{proof}
Projecting to three consecutive cyclic coordinates yields the seven traces $001,010,011,100,101,110,111$; four of these have degree three in the projected $1$-inclusion graph.  Checking all $31$ coordinate projections of the ten $C_5$ strings gives no vertex of degree larger than three.  Thus the projected maximal degree is three.
\end{proof}
\noindent\textbf{Reference.} \cite{palvolgyi2021unlabeled}.

\subsection{Projected Minimum Degree of Warmuth's $C_5$ Class}\label{val:projected_minimum_degree_warmuth_c5}
\noindent\textbf{Statement.} The Projected Minimum Degree of Warmuth's $C_5$ class is $2$.
\begin{proof}
Projecting to two adjacent cyclic coordinates gives the full two-cube, whose minimum degree is two.  Checking the remaining coordinate projections shows that each has a vertex of degree at most two, so no projection has larger minimum degree.
\end{proof}
\noindent\textbf{Reference.} \cite{palvolgyi2021unlabeled}.

\subsection{Projected Double Density or Projected Average Degree of Warmuth's $C_5$ Class}\label{val:projected_double_density_or_projected_average_degree_warmuth_c5}
\noindent\textbf{Statement.} The Projected Double Density or Projected Average Degree of Warmuth's $C_5$ class is $\frac{18}{7}$.
\begin{proof}
On three consecutive cyclic coordinates, the projection has seven vertices and nine $1$-inclusion edges, hence average degree $18/7$.  Enumerating all coordinate projections of the ten $C_5$ strings shows that none has higher average degree.
\end{proof}
\noindent\textbf{Reference.} \cite{palvolgyi2021unlabeled}.

\subsection{Effective VC Radius of Warmuth's $C_5$ Class}\label{val:effective_vc_radius_warmuth_c5}
\noindent\textbf{Statement.} The Effective VC Radius of Warmuth's $C_5$ class is $2$.
\begin{proof}
The effective range is all five coordinates, and every pair of cyclic positions is shattered by $C_5$.  No three-set is shattered, since the VC dimension is two.  Therefore the largest $k$ for which every effective $k$-set is shattered is two.
\end{proof}
\noindent\textbf{Reference.} \cite{palvolgyi2021unlabeled}.

\subsection{Minimum Largest Order Shattered Set of Warmuth's $C_5$ Class}\label{val:minimum_largest_order_shattered_set_warmuth_c5}
\noindent\textbf{Statement.} The Minimum Largest Order Shattered Set of Warmuth's $C_5$ class is $2$.
\begin{proof}
The effective VC radius is two, and every coordinate order therefore has an order-shattered set of size two.  Conversely every order-shattered set is shattered, so its size is at most the VC dimension, also two.  The minimum over orders is consequently exactly two.
\end{proof}
\noindent\textbf{Reference.} \cite{bollobas1995defect,palvolgyi2021unlabeled}.

\subsection{Minimum Teaching Set Size of Warmuth's $C_5$ Class}\label{val:minimum_teaching_set_size_warmuth_c5}
\noindent\textbf{Statement.} The Minimum Teaching Set Size of Warmuth's $C_5$ class is $3$.
\begin{proof}
A direct check of the ten cyclic $C_5$ strings shows that no restriction to one or two coordinates distinguishes any one string from all the others.  Each string is distinguished by a restriction to three coordinates (for example, $10010$ is taught on coordinates $0,2,4$).  Thus every concept has minimum teaching-set size three, and so does their minimum.
\end{proof}
\noindent\textbf{Reference.} \cite{palvolgyi2021unlabeled}.

\subsection{Maximum Teaching Set Size of Warmuth's $C_5$ Class}\label{val:maximum_teaching_set_size_warmuth_c5}
\noindent\textbf{Statement.} The Maximum Teaching Set Size of Warmuth's $C_5$ class is $3$.
\begin{proof}
Every one of the ten concepts has minimum teaching-set size three: no two-coordinate restriction is unique, and a three-coordinate restriction can be chosen for each cyclic placement.  Taking the maximum over concepts therefore gives three.
\end{proof}
\noindent\textbf{Reference.} \cite{palvolgyi2021unlabeled}.

\subsection{Teaching Dimension of Warmuth's $C_5$ Class}\label{val:teaching_dimension_warmuth_c5}
\noindent\textbf{Statement.} The Teaching Dimension of Warmuth's $C_5$ class is $3$.
\begin{proof}
Teaching dimension is the maximum minimum teaching-set size.  The direct computation for the ten $C_5$ concepts gives size three for each of them, hence the teaching dimension is three.
\end{proof}
\noindent\textbf{Reference.} \cite{palvolgyi2021unlabeled}.

\subsection{Recursive Teaching Dimension of Warmuth's $C_5$ Class}\label{val:recursive_teaching_dimension_warmuth_c5}
\noindent\textbf{Statement.} The Recursive Teaching Dimension of Warmuth's $C_5$ class is $3$.
\begin{proof}
Every concept has teaching-set size three in the original class, so they may all be removed in one recursive-teaching layer of width three.  No width-two first layer is possible because no concept has a teaching set of size at most two initially.  Therefore the recursive teaching dimension is three.
\end{proof}
\noindent\textbf{Reference.} \cite{palvolgyi2021unlabeled}.

\subsection{Projected Recursive Teaching Dimension of Warmuth's $C_5$ Class}\label{val:monotone_recursive_teaching_dimension_warmuth_c5}
\noindent\textbf{Statement.} The Projected Recursive Teaching Dimension of Warmuth's $C_5$ class is $3$.
\begin{proof}
The full projection has recursive teaching dimension three.  Enumerating the other $30$ coordinate projections of the ten $C_5$ strings and their recursive teaching layers gives value at most three in every case.  Thus the maximum over projections is exactly three.
\end{proof}
\noindent\textbf{Reference.} \cite{doliwa2014recursive,palvolgyi2021unlabeled}.

\subsection{Monotonic Minimum Teaching-Set Size of Warmuth's $C_5$ Class}\label{val:monotonic_minimum_teaching_set_size_warmuth_c5}
\noindent\textbf{Statement.} The Monotonic Minimum Teaching Set Size of Warmuth's $C_5$ class is $3$.
\begin{proof}
Monotonic minimum teaching-set size equals recursive teaching dimension.  The latter is three for $C_5$, so this value is three as well.
\end{proof}
\noindent\textbf{Reference.} \cite{doliwa2014recursive,palvolgyi2021unlabeled}.

\subsection{Preference-Based Teaching Dimension of Warmuth's $C_5$ Class}\label{val:preferencebased_teaching_dimension_warmuth_c5}
\noindent\textbf{Statement.} The Preference-Based Teaching Dimension of Warmuth's $C_5$ class is $3$.
\begin{proof}
For finite classes, preference-based teaching dimension equals recursive teaching dimension.  Since the recursive teaching dimension of $C_5$ is three, its preference-based teaching dimension is three.
\end{proof}
\noindent\textbf{Reference.} \cite{gao2017preference,palvolgyi2021unlabeled}.

\subsection{Star Number of Warmuth's $C_5$ Class}\label{val:star_number_warmuth_c5}
\noindent\textbf{Statement.} The Star number of Warmuth's $C_5$ class is $3$.
\begin{proof}
Star number equals projected maximal degree.  The three-consecutive-coordinate projection of $C_5$ has a vertex of degree three, and no coordinate projection has larger degree.  Hence the star number is exactly three.
\end{proof}
\noindent\textbf{Reference.} \cite{hanneke2015minimax,palvolgyi2021unlabeled}.

\subsection{Maximum Projected Teaching-Set Size of Warmuth's $C_5$ Class}\label{val:maximum_projected_teaching_set_size_warmuth_c5}
\noindent\textbf{Statement.} The Maximum Projected Teaching Set Size of Warmuth's $C_5$ class is $3$.
\begin{proof}
Maximum projected teaching-set size is equivalent to star number.  The latter has value three for $C_5$, so this parameter is also three.
\end{proof}
\noindent\textbf{Reference.} \cite{hanneke2015minimax,palvolgyi2021unlabeled}.

\subsection{Log Shattering of Warmuth's $C_5$ Class}\label{val:log_shattering_warmuth_c5}
\noindent\textbf{Statement.} The Log Shattering of Warmuth's $C_5$ class is $4$.
\begin{proof}
The empty set, every one-coordinate set, and every two-coordinate set are shattered, giving $1+5+\binom52=16$ shattered sets.  No three-coordinate set is shattered because the VC dimension is two.  The binary logarithm of the number of shattered sets is therefore $\log_2 16=4$.
\end{proof}
\noindent\textbf{Reference.} \cite{palvolgyi2021unlabeled}.

\subsection{Log Size Ratio of Warmuth's $C_5$ Class}\label{val:log_size_ratio_warmuth_c5}
\noindent\textbf{Statement.} The Log Size Ratio of Warmuth's $C_5$ class is $\frac{\log 9}{\log 6}$.
\begin{proof}
The class has ten concepts on a five-coordinate domain.  Substituting $|\mathcal H|=10$ and $n=5$ into $\log(|\mathcal H|-1)/\log(n+1)$ gives $\log 9/\log 6$.
\end{proof}
\noindent\textbf{Reference.} \cite{palvolgyi2021unlabeled}.

\subsection{Membership Query Complexity of Warmuth's $C_5$ Class}\label{val:membership_query_complexity_warmuth_c5}
\noindent\textbf{Statement.} The Membership Query Complexity of Warmuth's $C_5$ class is $4$.
\begin{proof}
At least four binary answers are necessary to identify one of ten concepts.  Four suffice: query coordinate $0$, then coordinate $1$; on branch $00$, query $3$; on branch $01$, query $2$ and, only after answer $0$, query $3$; on branch $10$, query $2$ and, only after answer $0$, query $4$; on branch $11$, query $2$.  The resulting leaves are the ten distinct $C_5$ strings, so the exact complexity is four.
\end{proof}
\noindent\textbf{Reference.} \cite{palvolgyi2021unlabeled}.

\subsection{Littlestone Dimension of Warmuth's $C_5$ Class}\label{val:littlestone_dimension_warmuth_c5}
\noindent\textbf{Statement.} The Littlestone Dimension of Warmuth's $C_5$ class is $3$.
\begin{proof}
A depth-three shattered mistake tree is obtained by querying coordinate $0$, then coordinate $1$, and on the four resulting branches querying respectively $3,2,2,2$.  A shattered tree of depth four would require $2^4=16$ distinct concepts, while $C_5$ has only ten.  Hence its Littlestone dimension is exactly three.
\end{proof}
\noindent\textbf{Reference.} \cite{littlestone1988learning,palvolgyi2021unlabeled}.

\subsection{Equivalence Queries Complexity of Warmuth's $C_5$ Class}\label{val:equivalence_queries_complexity_warmuth_c5}
\noindent\textbf{Statement.} The Equivalence Queries Complexity of Warmuth's $C_5$ class is $3$.
\begin{proof}
Equivalence-query complexity equals Littlestone dimension.  The latter has exact value three for $C_5$, so its equivalence-query complexity is three.
\end{proof}
\noindent\textbf{Reference.} \cite{littlestone1988learning,palvolgyi2021unlabeled}.

\subsection{Interpolation Degree of Warmuth's $C_5$ Class}\label{val:interpolation_degree_warmuth_c5}
\noindent\textbf{Statement.} The Interpolation degree of Warmuth's $C_5$ class is $2$.
\begin{proof}
On the ten $C_5$ strings, the constant and five coordinate monomials have evaluation rank six, so degree one cannot interpolate every function on the class.  The constant, linear, and quadratic monomials have evaluation rank ten, the full number of concepts, and therefore span all functions on $C_5$.  Hence the interpolation degree is exactly two.
\end{proof}
\noindent\textbf{Reference.} \cite{palvolgyi2021unlabeled}.

\subsection{Diameter of Warmuth's $C_5$ class}\label{val:diameter_warmuth_c5}
\noindent\textbf{Statement.} The Diameter of Warmuth's $C_5$ class is $4$.
\begin{proof}
Enumerating the ten traces satisfying the cyclic 1001 condition, every pair differs on at most four coordinates. The pair 00101 and 01010 differs on four coordinates, so the upper bound is attained.
\end{proof}

\subsection{Oriented Diameter of Majority}\label{val:oriented_diameter_majority}
\noindent\textbf{Statement.} The Oriented Diameter of Majority is $n+1$.
\begin{proof}
The empty ordinary subset with majority bit 0 and the full ordinary subset with majority bit 1 are both class members and disagree on all n+1 coordinates. This witnesses oriented diameter n+1, which is also the largest possible value on this domain.
\end{proof}

\subsection{Hamming Radius of Majority}\label{val:hamming_radius_majority}
\noindent\textbf{Statement.} The Hamming Radius of Majority is $n$.
\begin{proof}
For any proposed center, complementing its ordinary-coordinate trace gives a class member that differs on all n ordinary coordinates, so radius at least n is necessary. Conversely, take the center with all ordinary coordinates equal to 1 and majority coordinate 0. A class member with ordinary subset A has distance n-|A| plus at most one for its majority bit; this is at most n. Hence the radius is exactly n.
\end{proof}

\subsection{Littlestone Dimension of Majority}\label{val:littlestone_dimension_majority}
\noindent\textbf{Statement.} The Littlestone Dimension of Majority is $n$.
\begin{proof}
Query the n ordinary coordinates in a fixed order. Every binary sequence of answers is realized by the concept whose ordinary-coordinate subset is that sequence, so this is a shattered tree of depth n. Conversely, a class of size $2^n$ cannot shatter a complete binary tree of depth greater than n: distinct leaves require distinct realizing concepts. Hence the Littlestone dimension is exactly n.
\end{proof}
\noindent\textbf{Reference.} \cite{maass1992lower,littlestone1988learning}.

\subsection{Equivalence Queries Complexity of Majority}\label{val:equivalence_queries_complexity_majority}
\noindent\textbf{Statement.} The Equivalence Queries Complexity of Majority is $n$.
\begin{proof}
Unrestricted equivalence-query complexity equals Littlestone dimension. The majority class has Littlestone dimension n, witnessed by its full-cube ordinary-coordinate projection and bounded by its $2^n$ concepts. Therefore its equivalence-query complexity is n.
\end{proof}
\noindent\textbf{Reference.} \cite{littlestone1988learning,maass1992lower}.

\subsection{Log Size Ratio of Blockwise-Addressing}\label{val:log_size_ratio_blockwiseaddressing}
\noindent\textbf{Statement.} The Log Size Ratio of Blockwise-Addressing is $\frac{\log(2^{d+1}-1)}{\log(2^{d+1}+2)}$.
\begin{proof}
There are $2^d$ functions $f_i$ and $2^d$ functions $g_i$. The domain consists of z and two blocks of $2^d$ coordinates, so it has $2^{d+1}+1$ elements. Substitution into $\log(|\mathcal H|-1)/\log(|\mathcal X|+1)$ gives the stated expression.
\end{proof}
\noindent\textbf{Reference.} \cite{ben1997online}.

\subsection{Log Shattering of Tagged Singletons}\label{val:log_shattering_tagged_singletons}
\noindent\textbf{Statement.} The Log Shattering of Tagged Singletons is $\log(n+\lceil\log n\rceil+1)$.
\begin{proof}
The empty set and every singleton coordinate set are shattered: each ordinary and tag coordinate varies among the displayed concepts. No two-coordinate set is shattered, since no concept contains two ordinary coordinates, no concept contains two tags, and an ordinary coordinate appears only with its assigned tag. Thus there are exactly $1+n+\lceil\log n\rceil$ shattered sets, whose binary logarithm is the stated value.
\end{proof}
\noindent\textbf{Reference.} \cite{maass1992lower}.

\subsection{Largest Strongly Shattered Set of Tagged Singletons}\label{val:largest_strongly_shattered_set_tagged_singletons}
\noindent\textbf{Statement.} The Largest Strongly Shattered Set of Tagged Singletons is $1$.
\begin{proof}
The empty concept and any tag singleton form a contained one-dimensional cube, so the value is at least one. A contained two-cube would require four concepts realizing all patterns on two varying coordinates while agreeing outside them. The construction has no ordinary singleton, no two-tag concept, and each ordinary coordinate occurs only with its assigned tag; consequently no two-coordinate projection contains such a cube. Hence the largest strongly shattered set has size one.
\end{proof}
\noindent\textbf{Reference.} \cite{maass1992lower}.

\subsection{Largest Strongly Shattered Set of Majority}\label{val:largest_strongly_shattered_set_majority}
\noindent\textbf{Statement.} The Largest Strongly Shattered Set of Majority is $\lfloor n/2\rfloor$.
\begin{proof}
Fix any $\lfloor n/2\rfloor$ ordinary coordinates to vary and set every other ordinary coordinate and the majority coordinate to 0. All resulting traces have at most $\lfloor n/2\rfloor$ ordinary ones, so the whole cube is contained in the class. Conversely, a contained cube cannot vary the majority coordinate, since that bit is determined by the ordinary trace. If it varies $k$ ordinary coordinates, the majority bit must be constant throughout the cube. With $r$ fixed ordinary ones this requires either $r+k\le n/2$ or $r>n/2$; the available fixed coordinates then give respectively $k\le\lfloor n/2\rfloor$ or $k\le\lceil n/2\rceil-1$. Hence $k\le\lfloor n/2\rfloor$.
\end{proof}
\noindent\textbf{Reference.} \cite{maass1992lower}.

\subsection{Log Shattering of Majority}\label{val:log_shattering_majority}
\noindent\textbf{Statement.} The Log Shattering of Majority is $\log\left(2^n+\sum_{k=0}^{\lfloor(n-1)/2\rfloor}\binom{n}{k}\right)$.
\begin{proof}
Every subset of the n ordinary coordinates is shattered because every ordinary trace occurs. For a set consisting of the majority coordinate and $k$ ordinary coordinates, all patterns are realizable exactly when $k\le\lfloor(n-1)/2\rfloor$: after fixing any pattern on those $k$ coordinates, set the remaining ordinary coordinates all to 0 or all to 1 to realize either majority label. Conversely, if $k$ is larger, the all-one ordinary pattern cannot have majority label 0 (or, equivalently, the all-zero pattern cannot have majority label 1). Thus there are $2^n$ shattered sets without the majority coordinate and $\sum_{k=0}^{\lfloor(n-1)/2\rfloor}\binom nk$ with it.
\end{proof}
\noindent\textbf{Reference.} \cite{maass1992lower}.

\subsection{Unlabeled Sample Compression of Trivial concepts}\label{val:unlabeled_sample_compression_trivial_concepts}
\noindent\textbf{Statement.} The Unlabeled sample compression of Trivial concepts is $1$.
\begin{proof}
Fix an order on the domain.  For an all-positive sample, retain its least instance and reconstruct the full concept; for an all-negative sample, retain nothing and reconstruct the empty concept.  This is a size-one unlabeled compression scheme.  Size zero cannot distinguish the incompatible positive and negative one-point samples.
\end{proof}
\noindent\textbf{Reference.} \cite{palvolgyi2021unlabeled}.

\subsection{Proper Unlabeled Sample Compression of Trivial concepts}\label{val:proper_unlabeled_sample_compression_trivial_concepts}
\noindent\textbf{Statement.} The Proper Unlabeled Sample Compression of Trivial concepts is $1$.
\begin{proof}
The size-one unlabeled scheme for the two constant concepts reconstructs either the empty or the full concept, both of which belong to the class.  Thus it is proper; the ordinary unlabeled lower bound of one still applies.
\end{proof}
\noindent\textbf{Reference.} \cite{palvolgyi2021unlabeled}.

\subsection{Stable Unlabeled Complexity of Trivial concepts}\label{val:stable_unlabeled_complexity_trivial_concepts}
\noindent\textbf{Statement.} The Stable Unlabeled Complexity of Trivial concepts is $1$.
\begin{proof}
Use the size-one unlabeled scheme that retains the least instance of an all-positive sample and retains nothing from an all-negative sample.  Any intermediate subsample containing the retained instance has the same least instance; intermediate all-negative subsamples still retain nothing.  Hence the compression map is stable.  Size zero is impossible.
\end{proof}
\noindent\textbf{Reference.} \cite{hanneke2021stable}.

\subsection{Proper Stable Sample Compression of Trivial concepts}\label{val:proper_stable_sample_compression_trivial_concepts}
\noindent\textbf{Statement.} The Proper Stable Sample Compression of Trivial concepts is $1$.
\begin{proof}
The stable size-one unlabeled scheme reconstructs only the two constant concepts in the class, so it is proper.  The ordinary unlabeled lower bound rules out size zero.
\end{proof}
\noindent\textbf{Reference.} \cite{hanneke2021stable}.

\subsection{Unlabeled Sample Compression of Singletons}\label{val:unlabeled_sample_compression_singletons}
\noindent\textbf{Statement.} The Unlabeled sample compression of Singletons is $0$ for $n=1$; $1$ for $n\ge2$.
\begin{proof}
Decode the empty key as the all-zero word, and a singleton key $\{i\}$ as the singleton concept $\{i\}$. Retain the unique positive coordinate if the sample has one, and retain nothing otherwise. The reconstruction is consistent for every realizable sample, including the empty sample, and never reads an omitted label. The empty reconstruction is not a member of the singleton class, so this scheme is improper. Size zero cannot reconstruct both of two distinct full concepts; hence the value is one for $n\ge2$. For $n=1$, the sole concept can always be reconstructed without a key. The former cyclic-successor argument incorrectly let an unlabeled decoder distinguish the sign of its retained point; that argument proves a labeled, not an unlabeled, upper bound. Exact replay: $\texttt{sync/verify\_singleton\_unlabeled\_compression.py}$.
\end{proof}

\subsection{Proper Unlabeled Sample Compression of Singletons}\label{val:proper_unlabeled_sample_compression_singletons}
\noindent\textbf{Statement.} The Proper Unlabeled Sample Compression of Singletons is $0$ for $n=1$; $1$ for $n=2$; $2$ for $n\ge3$.
\begin{proof}
Identify a reconstructed singleton with its index in $[n]$. For $n\ge2$, define $\rho(\varnothing)=1$, $\rho(\{1\})=2$, $\rho(\{i\})=i$ for $i\ge2$, and $\rho(\{1,i\})=i+1$ for $2\le i<n$; send any other unused key to 1. A sample with its positive point at 1 retains nothing; one with positive point $i\ge2$ retains $\{i\}$. An all-negative sample has a proper domain $A\subsetneq[n]$; let $r=\min([n]\setminus A)$. Retain nothing if $r=1$, retain $\{1\}$ if $r=2$, and retain $\{1,r-1\}$ if $r\ge3$. These retained coordinates belong to $A$ by minimality of $r$, and the decoder returns the missing singleton $r$. This proves the proper upper bound two, without labels or messages. For the lower bound, suppose a proper size-one decoder exists and put $a=\rho(\varnothing)$. The one-point positive sample at every $i\ne a$ forces $\rho(\{i\})=i$. Write $b=\rho(\{a\})$, defining unused outputs arbitrarily in the class. If $n\ge3$, choose $v\notin\{a,b\}$. The all-negative sample on $[n]\setminus\{v\}$ is realizable only by singleton $v$, but neither its empty key nor any of its singleton keys decodes to $v$, a contradiction. For $n=2$, the displayed upper scheme uses only keys of size at most one, and two distinct concepts rule out size zero. For $n=1$, always reconstruct the sole concept. The exact verifier $\texttt{sync/verify\_singleton\_unlabeled\_compression.py}$ replays the upper schemes and independently rejects all 81 proper size-one decoder maps on the three-singleton class.
\end{proof}

\subsection{Unlabeled Sample Compression of Singletons plus Empty Set}\label{val:unlabeled_sample_compression_singletons_plus_empty_set}
\noindent\textbf{Statement.} The Unlabeled sample compression of Singletons plus empty set is $1$.
\begin{proof}
Retain the unique positive instance when one is present and reconstruct its singleton; retain nothing from an all-negative sample and reconstruct the empty concept.  This is consistent.  Size zero cannot handle incompatible positive and negative one-point samples.
\end{proof}
\noindent\textbf{Reference.} \cite{palvolgyi2021unlabeled}.

\subsection{Proper Unlabeled Sample Compression of Singletons plus Empty Set}\label{val:proper_unlabeled_sample_compression_singletons_plus_empty_set}
\noindent\textbf{Statement.} The Proper Unlabeled Sample Compression of Singletons plus empty set is $1$.
\begin{proof}
The size-one scheme reconstructs either the empty concept or a singleton, so its reconstruction is always in the class.  The ordinary unlabeled lower bound of one applies.
\end{proof}
\noindent\textbf{Reference.} \cite{palvolgyi2021unlabeled}.

\subsection{Stable Unlabeled Complexity of Singletons plus Empty Set}\label{val:stable_unlabeled_complexity_singletons_plus_empty_set}
\noindent\textbf{Statement.} The Stable Unlabeled Complexity of Singletons plus empty set is $1$.
\begin{proof}
Retain the positive instance exactly when the sample has one; otherwise retain nothing.  Every intermediate subsample containing the retained positive instance still retains it, and every all-negative subsample retains nothing.  Thus the size-one scheme is stable.  Size zero is impossible.
\end{proof}
\noindent\textbf{Reference.} \cite{hanneke2021stable}.

\subsection{Proper Stable Sample Compression of Singletons plus Empty Set}\label{val:proper_stable_sample_compression_singletons_plus_empty_set}
\noindent\textbf{Statement.} The Proper Stable Sample Compression of Singletons plus empty set is $1$.
\begin{proof}
The stable size-one scheme reconstructs only the empty concept and singletons, hence is proper.  The ordinary unlabeled lower bound excludes size zero.
\end{proof}
\noindent\textbf{Reference.} \cite{hanneke2021stable}.

\subsection{Stable Unlabeled Complexity of Full Cube}\label{val:stable_unlabeled_complexity_full_cube}
\noindent\textbf{Statement.} The Stable Unlabeled Complexity of Full Cube is $n$.
\begin{proof}
The Littlestone dimension of the full cube is $n$, and the stable unlabeled-compression construction from a conservative online learner gives a stable scheme of size at most $n$.  Stable unlabeled compression is at least ordinary unlabeled compression, whose exact value for the full cube is $n$.
\end{proof}
\noindent\textbf{Reference.} \cite{hanneke2021stable,chalopin2022unlabeled,palvolgyi2021unlabeled}.

\subsection{Proper Stable Sample Compression of Full Cube}\label{val:proper_stable_sample_compression_full_cube}
\noindent\textbf{Statement.} The Proper Stable Sample Compression of Full Cube is $n$.
\begin{proof}
The stable unlabeled size-$n$ construction reconstructs Boolean concepts on the domain; every such concept belongs to the full cube, so it is proper.  Proper stable compression is at least ordinary unlabeled compression, whose value is $n$.
\end{proof}
\noindent\textbf{Reference.} \cite{hanneke2021stable,chalopin2022unlabeled,palvolgyi2021unlabeled}.

\subsection{Stable Unlabeled Complexity of Singletons}\label{val:stable_unlabeled_complexity_singletons}
\noindent\textbf{Statement.} The Stable Unlabeled Complexity of Singletons is $1$.
\begin{proof}
The stable-unlabeled construction from a conservative online learner has size at most the Littlestone dimension, which is one for singletons.  Stable unlabeled compression is at least ordinary unlabeled compression, whose value is one for this class.
\end{proof}
\noindent\textbf{Reference.} \cite{hanneke2021stable,palvolgyi2021unlabeled}.

\subsection{Relative Hitting Size of Warmuth's $C_5$ class}\label{val:relative_hitting_size_warmuth_c5}
\noindent\textbf{Statement.} The Relative Hitting Size of Warmuth's $C_5$ class is $3$.
\begin{proof}
For every target, a relative hitting set is precisely a teaching set for that target.  The established teaching dimension of Warmuth's $C_5$ class is three, hence its relative hitting size is three as well.
\end{proof}
\noindent\textbf{Reference.} \cite{goldman1995complexity,shinohara1991teachability,palvolgyi2021unlabeled}.

\subsection{Stable Labeled Sample Compression of Singletons}\label{val:stable_labeled_sample_compression_singletons}
\noindent\textbf{Statement.} The Stable Labeled Sample Compression of Singletons is $1$.
\begin{proof}
The established stable unlabeled scheme of size one is also a stable labeled scheme when the retained label is allowed.  Stable labeled compression is at least ordinary labeled compression, whose value for singletons is one.
\end{proof}
\noindent\textbf{Reference.} \cite{FW1995,hanneke2021stable}.

\subsection{Stable Labeled Sample Compression of Singletons plus Empty Set}\label{val:stable_labeled_sample_compression_singletons_plus_empty_set}
\noindent\textbf{Statement.} The Stable Labeled Sample Compression of Singletons plus empty set is $1$.
\begin{proof}
The established stable unlabeled size-one scheme is also a stable labeled scheme.  Ordinary labeled compression has value one, giving the matching lower bound.
\end{proof}
\noindent\textbf{Reference.} \cite{FW1995,hanneke2021stable}.

\subsection{Proper Stable Labeled Sample Compression of Singletons plus Empty Set}\label{val:proper_stable_labeled_sample_compression_singletons_plus_empty_set}
\noindent\textbf{Statement.} The Proper Stable Labeled Sample Compression of Singletons plus empty set is $1$.
\begin{proof}
The stable unlabeled size-one scheme reconstructs only the empty concept and singletons, so it is proper and remains stable when labels may be retained.  Proper labeled compression has value one, giving the lower bound.
\end{proof}
\noindent\textbf{Reference.} \cite{FW1995,hanneke2021stable}.

\subsection{Teaching Dimension of Addressing}\label{val:teaching_dimension_addressing}
\noindent\textbf{Statement.} The Teaching Dimension of Addressing is $1\quad(n\ge2)$.
\begin{proof}
For each concept, its addressed primary coordinate has label one only for that concept, so its single positive example teaches it. The nontrivial class has teaching dimension at least one.
\end{proof}
\noindent\textbf{Reference.} \cite{maass1992lower}.

\subsection{Maximum Teaching-Set Size of Addressing}\label{val:maximum_teaching_set_size_addressing}
\noindent\textbf{Statement.} The Maximum Teaching Set Size of Addressing is $1\quad(n\ge2)$.
\begin{proof}
Maximum teaching-set size is, by definition, teaching dimension. The addressing class has teaching dimension one because every concept has a private positive primary coordinate.
\end{proof}
\noindent\textbf{Reference.} \cite{maass1992lower}.

\subsection{Relative Hitting Size of Addressing}\label{val:relative_hitting_size_addressing}
\noindent\textbf{Statement.} The Relative Hitting Size of Addressing is $1\quad(n\ge2)$.
\begin{proof}
Relative hitting size is the maximum teaching-set size. The addressing class has maximum teaching-set size one, witnessed by each concept's private primary coordinate.
\end{proof}
\noindent\textbf{Reference.} \cite{maass1992lower}.

\subsection{Teaching Dimension of Tagged Singletons}\label{val:teaching_dimension_tagged_singletons}
\noindent\textbf{Statement.} The Teaching Dimension of Tagged Singletons is $n/2$.
\begin{proof}
A tag singleton with the first tag must be distinguished from each of the $n/2$ tagged pairs carrying that tag, requiring a negative example at every corresponding ordinary coordinate. Those $n/2$ negative examples suffice. Every other concept has a teaching set no larger than this, so the maximum teaching-set size is $n/2$.
\end{proof}
\noindent\textbf{Reference.} \cite{maass1992lower}.

\subsection{Maximum Teaching-Set Size of Tagged Singletons}\label{val:maximum_teaching_set_size_tagged_singletons}
\noindent\textbf{Statement.} The Maximum Teaching Set Size of Tagged Singletons is $n/2$.
\begin{proof}
The first tag singleton requires negative examples on all $n/2$ ordinary coordinates carrying that tag, and these examples suffice. All other concepts have no larger teaching sets. Thus the maximum teaching-set size is $n/2$.
\end{proof}
\noindent\textbf{Reference.} \cite{maass1992lower}.

\subsection{Relative Hitting Size of Tagged Singletons}\label{val:relative_hitting_size_tagged_singletons}
\noindent\textbf{Statement.} The Relative Hitting Size of Tagged Singletons is $n/2$.
\begin{proof}
Relative hitting size equals maximum teaching-set size. For tagged singletons, the first tag singleton needs and has a teaching set of size $n/2$, while no target requires more.
\end{proof}
\noindent\textbf{Reference.} \cite{maass1992lower}.

\subsection{$k$-Block Littlestone Dimension of Trivial Concepts}\label{val:kblock_littlestone_dimension_trivial_concepts}
\noindent\textbf{Statement.} The $k$-Block Littlestone Dimension ($k>2$) of Trivial concepts is $1$.
\begin{proof}
The two constant hypotheses shatter a tree of depth one. A shattered binary tree of depth two has four leaves, and each leaf needs a distinct realizing hypothesis, while this class has only two. Hence $\mathrm B_k=1$ for every fixed $k>2$.
\end{proof}

\subsection{Dual VC Dimension of Full Cube}\label{val:dual_vc_dimension_full_cube}
\noindent\textbf{Statement.} The Dual VC Dimension of Full Cube is $\lfloor\log_2 n\rfloor$.
\begin{proof}
The dual of the $n$-dimensional cube is the universal class with $n$ concepts. It shatters $\lfloor\log_2 n\rfloor$ concepts, and no larger set can be shattered because $n$ concepts realize at most $n$ patterns.
\end{proof}
\noindent\textbf{Reference.} \cite{chornomaz2025spherical}.

\subsection{Spherical Dimension of Full Cube}\label{val:spherical_dimension_full_cube}
\noindent\textbf{Statement.} The Spherical Dimension of Full Cube is $n-1$.
\begin{proof}
For the $n$-cube, the realizable-distribution complex is the boundary of the $n$-dimensional cross-polytope, an antipodal $(n-1)$-sphere; the cited paper proves this is exact.
\end{proof}
\noindent\textbf{Reference.} \cite{chornomaz2025spherical}.

\subsection{Dual VC Dimension of Trivial concepts}\label{val:dual_vc_dimension_trivial_concepts}
\noindent\textbf{Statement.} The Dual VC Dimension of Trivial concepts is $0$.
\begin{proof}
Every domain point has the same dual trace, namely the singleton containing the full concept. Thus the dual class has one hypothesis and shatters no nonempty set.
\end{proof}

\subsection{Dual VC Dimension of Singletons}\label{val:dual_vc_dimension_singletons}
\noindent\textbf{Statement.} The Dual VC Dimension of Singletons is $0$ for $n=1$; $1$ for $n\ge2$.
\begin{proof}
The incidence matrix of the singleton class is the identity matrix, so transposition gives the same class on its $n$ hypotheses. For $n=1$ its only concept is fixed and the dual VC dimension is zero. For $n\ge2$, every one-coordinate set takes both labels, while no pair takes the pattern $11$. Hence the dual VC dimension is one.
\end{proof}

\subsection{Dual VC Dimension of Singletons plus empty set}\label{val:dual_vc_dimension_singletons_plus_empty_set}
\noindent\textbf{Statement.} The Dual VC Dimension of Singletons plus empty set is $1$.
\begin{proof}
Each domain point corresponds in the dual to the singleton containing its own singleton hypothesis. These traces shatter one hypothesis but no pair, since no dual concept contains two singleton hypotheses.
\end{proof}

\subsection{Dual VC Dimension of Half-intervals}\label{val:dual_vc_dimension_halfintervals}
\noindent\textbf{Statement.} The Dual VC Dimension of Half-intervals is $1$.
\begin{proof}
After exchanging points and initial segments, the dual class is a chain of terminal segments. It shatters one hypothesis but no pair because every trace in a chain is comparable.
\end{proof}

\subsection{Spherical Dimension of Trivial concepts}\label{val:spherical_dimension_trivial_concepts}
\noindent\textbf{Statement.} The Spherical Dimension of Trivial concepts is $0$.
\begin{proof}
The two constant concepts form a two-element threshold class. The low-VC classification of spherical dimension gives spherical dimension zero for every non-singleton subclass of thresholds.
\end{proof}
\noindent\textbf{Reference.} \cite{chornomaz2025spherical}.

\subsection{Spherical Dimension of Singletons}\label{val:spherical_dimension_singletons}
\noindent\textbf{Statement.} The Spherical Dimension of Singletons is $1$.
\begin{proof}
The antipodal complex is the complete bipartite graph on positive and negative literals with the perfect matching removed. For $n\geq3$ it contains an antipodal six-cycle, giving a 1-sphere, and it is one-dimensional so no larger sphere occurs.
\end{proof}
\noindent\textbf{Reference.} \cite{chornomaz2025spherical}.

\subsection{Spherical Dimension of Half-intervals}\label{val:spherical_dimension_halfintervals}
\noindent\textbf{Statement.} The Spherical Dimension of Half-intervals is $0$.
\begin{proof}
Half-intervals are thresholds. The low-VC classification of spherical dimension gives value zero for every non-singleton subclass of thresholds.
\end{proof}
\noindent\textbf{Reference.} \cite{chornomaz2025spherical}.

\subsection{VC Dimension of Universal Class}\label{val:vc_dimension_universal_class}
\noindent\textbf{Statement.} The VC Dimension of Universal Class is $\lfloor\log_2 n\rfloor$.
\begin{proof}
The $n$ hypotheses can realize at most $n$ label patterns on a chosen set, so no set larger than $\lfloor\log_2 n\rfloor$ is shattered. Choosing one domain point for each binary pattern on that many hypotheses gives equality.
\end{proof}
\noindent\textbf{Reference.} \cite{chornomaz2025spherical}.

\subsection{Dual VC Dimension of Universal Class}\label{val:dual_vc_dimension_universal_class}
\noindent\textbf{Statement.} The Dual VC Dimension of Universal Class is $n$.
\begin{proof}
The dual of the universal class is the full $n$-cube, which shatters all $n$ of its coordinates.
\end{proof}
\noindent\textbf{Reference.} \cite{chornomaz2025spherical}.

\subsection{Spherical Dimension of Universal Class}\label{val:spherical_dimension_universal_class}
\noindent\textbf{Statement.} The Spherical Dimension of Universal Class is $\begin{cases}n-2&n\text{ odd},\\n-2\leq\mathrm{sd}(\mathcal U_n)\leq n-1&n\text{ even}.\end{cases}$.
\begin{proof}
Lemma 2 of the cited paper proves $n-2\leq\mathrm{sd}(\mathcal U_n)\leq n-1$ and, more sharply, $\mathrm{sd}(\mathcal U_n)=n-2$ for odd $n$.
\end{proof}
\noindent\textbf{Reference.} \cite{chornomaz2025spherical}.

\subsection{Dual VC Dimension of Warmuth's $C_5$ class}\label{val:dual_vc_dimension_warmuth_c5}
\noindent\textbf{Statement.} The Dual VC Dimension of Warmuth's $C_5$ class is $2$.
\begin{proof}
The dual class has five concepts, one for each cyclic coordinate, so it cannot shatter three hypotheses. For the two $C_5$ hypotheses $00101$ and $01001$ (in cyclic coordinate order), the five coordinate evaluations include all four patterns $00,01,10,11$; hence those two hypotheses are dually shattered.
\end{proof}

\subsection{Spherical Dimension of Singletons plus empty set}\label{val:spherical_dimension_singletons_plus_empty_set}
\noindent\textbf{Statement.} The Spherical Dimension of Singletons plus empty set is $1$.
\begin{proof}
Its antipodal complex has one positive and one negative literal for each coordinate, with edges precisely between opposite-sign literals from different coordinates: it is the complete bipartite graph minus a perfect matching. It contains an antipodal six-cycle when $n\geq3$, giving a 1-sphere. Since the complex is one-dimensional, no larger sphere occurs.
\end{proof}
\noindent\textbf{Reference.} \cite{chornomaz2025spherical}.

\subsection{Spherical Dimension of Tagged Singletons}\label{val:spherical_dimension_tagged_singletons}
\noindent\textbf{Statement.} The Spherical Dimension of Tagged Singletons is $1$.
\begin{proof}
Projection to any three tag coordinates is the class of three singletons together with the empty concept, whose spherical dimension is one. Projection monotonicity therefore gives $\mathrm{sd}\geq1$. The recorded VC dimension is one, and the low-VC classification in the cited paper gives $\mathrm{sd}\leq1$.
\end{proof}
\noindent\textbf{Reference.} \cite{chornomaz2025spherical}.

\subsection{Dual Antipodal VC Dimension of Universal Class}\label{val:dual_antipodal_vc_dimension_universal_class}
\noindent\textbf{Statement.} The Dual Antipodal VC Dimension of Universal Class is $n$.
\begin{proof}
The dual of the universal class is the full $n$-cube. Its antipodal VC dimension is at least its VC dimension $n$ and at most its domain size $n$, hence equals $n$.
\end{proof}
\noindent\textbf{Reference.} \cite{chornomaz2025spherical}.

\subsection{Extremal-Extension VC Dimension of Full Cube}\label{val:extremal_extension_vc_dimension_full_cube}
\noindent\textbf{Statement.} The Extremal-Extension VC Dimension of Full Cube is $n$.
\begin{proof}
The full cube is extremal, so it is its own extension and gives the upper bound $n$. The lower bound is its VC dimension $n$.
\end{proof}
\noindent\textbf{Reference.} \cite{chornomaz2025spherical}.

\subsection{Extremal-Extension VC Dimension of Trivial concepts}\label{val:extremal_extension_vc_dimension_trivial_concepts}
\noindent\textbf{Statement.} The Extremal-Extension VC Dimension of Trivial concepts is $1$.
\begin{proof}
VC dimension one is necessary because the class shatters a singleton. The full chain of initial segments is an extremal VC-one class containing both constant concepts, giving the matching upper bound.
\end{proof}

\subsection{Extremal-Extension VC Dimension of Singletons}\label{val:extremal_extension_vc_dimension_singletons}
\noindent\textbf{Statement.} The Extremal-Extension VC Dimension of Singletons is $1$.
\begin{proof}
VC dimension one is necessary. Adding the empty concept gives a class of $n+1$ concepts with exactly the empty set and the $n$ singletons shattered, hence an extremal VC-one extension.
\end{proof}

\subsection{Extremal-Extension VC Dimension of Singletons plus empty set}\label{val:extremal_extension_vc_dimension_singletons_plus_empty_set}
\noindent\textbf{Statement.} The Extremal-Extension VC Dimension of Singletons plus empty set is $1$.
\begin{proof}
The class has $n+1$ concepts and exactly $n+1$ shattered sets: the empty set and the $n$ singletons. It is therefore itself extremal of VC dimension one.
\end{proof}

\subsection{Extremal-Extension VC Dimension of Half-intervals}\label{val:extremal_extension_vc_dimension_halfintervals}
\noindent\textbf{Statement.} The Extremal-Extension VC Dimension of Half-intervals is $1$.
\begin{proof}
The chain of half-intervals has VC dimension one and has as many shattered sets as concepts (the empty set and its active singleton coordinates), so it is extremal and is its own extension.
\end{proof}

\subsection{List Replicability Number of Universal Class}\label{val:list_replicability_universal_class}
\noindent\textbf{Statement.} The List Replicability Number of Universal Class is $\lceil(n+1)/2\rceil$.
\begin{proof}
The cited paper combines its list-replicable learner for the universal class with its lower bound from dual antipodal VC dimension to obtain $\mathrm{LR}(\mathcal U_n)=\lceil(n+1)/2\rceil$.
\end{proof}
\noindent\textbf{Reference.} \cite{chornomaz2025spherical}.

\subsection{SCO Dimension of Full Cube}\label{val:sco_dimension_full_cube}
\noindent\textbf{Statement.} The SCO Dimension of Full Cube is $n$.
\begin{proof}
The spherical-dimension lower bound gives $\mathrm{SCO}\geq\mathrm{sd}+1=n$ for the full cube, while the sign-rank upper bound gives $\mathrm{SCO}\leq\mathrm{SR}=n$. Hence both bounds are tight.
\end{proof}
\noindent\textbf{Reference.} \cite{chornomaz2025spherical}.

\subsection{Combinatorial Eluder Dimension of Trivial Concepts}\label{val:combinatorial_eluder_dimension_trivial_concepts}
\noindent\textbf{Statement.} The Combinatorial Eluder Dimension of Trivial concepts is $1$.
\begin{proof}
Either coordinate and the other constant hypothesis give a one-step witness. The general bound $\mathrm{Edim}(\mathcal H)\le |\mathcal H|-1$ gives the matching upper bound.
\end{proof}
\noindent\textbf{Reference.} \cite{li2022understanding}.

\subsection{Combinatorial Eluder Dimension of Singletons}\label{val:combinatorial_eluder_dimension_singletons}
\noindent\textbf{Statement.} The Combinatorial Eluder Dimension of Singletons is $n-1$.
\begin{proof}
Take one singleton as the base and list the other $n-1$ coordinates; the singleton at the current coordinate agrees with the base on all earlier coordinates and disagrees at the current one. The class has n hypotheses, so the size bound gives the matching upper bound.
\end{proof}
\noindent\textbf{Reference.} \cite{li2022understanding}.

\subsection{Combinatorial Eluder Dimension of Singletons plus Empty Set}\label{val:combinatorial_eluder_dimension_singletons_plus_empty_set}
\noindent\textbf{Statement.} The Combinatorial Eluder Dimension of Singletons plus empty set is $n$.
\begin{proof}
Use the empty hypothesis as base and list all n coordinates, with the corresponding singleton as each witness. The effective-range bound gives the matching upper bound.
\end{proof}
\noindent\textbf{Reference.} \cite{li2022understanding}.

\subsection{Combinatorial Eluder Dimension of Half-intervals}\label{val:combinatorial_eluder_dimension_halfintervals}
\noindent\textbf{Statement.} The Combinatorial Eluder Dimension of Half-intervals is $n-1$.
\begin{proof}
With $[1]$ as base, list coordinates $2,\ldots,n$ and use $[i]$ at step i. Each witness agrees with the base on earlier coordinates and differs at i. The class has n hypotheses, so the size bound gives the matching upper bound.
\end{proof}
\noindent\textbf{Reference.} \cite{li2022understanding}.

\subsection{Combinatorial Eluder Dimension of the Full Cube}\label{val:combinatorial_eluder_dimension_full_cube}
\noindent\textbf{Statement.} The Combinatorial Eluder Dimension of Full Cube is $n$.
\begin{proof}
Take the empty set as base and at step i use the singleton at the ith coordinate. This gives n steps; the effective range has size n, giving the upper bound.
\end{proof}
\noindent\textbf{Reference.} \cite{li2022understanding}.

\subsection{Combinatorial Eluder Dimension of Warmuth's $C_5$ Class}\label{val:combinatorial_eluder_dimension_warmuth_c5}
\noindent\textbf{Statement.} The Combinatorial Eluder Dimension of Warmuth's $C_5$ class is $4$.
\begin{proof}
Enumerate the ten binary concepts of $C_5$ and the ten choices of base hypothesis in the sequential definition.  For example, with base $00101$ on cyclic coordinates $0,\ldots,4$, the points $0,2,1,3$ have respective witnesses $10100,01001,01110,00111$.  Exhaustive search over the five coordinates and all bases finds no witness sequence of length five, so the value is exactly four.
\end{proof}
\noindent\textbf{Reference.} \cite{li2022understanding}.

\subsection{Combinatorial Eluder Dimension of Tagged Singletons}\label{val:combinatorial_eluder_dimension_tagged_singletons}
\noindent\textbf{Statement.} The Combinatorial Eluder Dimension of Tagged Singletons is $n+\lceil\log_2 n\rceil$.
\begin{proof}
Use the empty hypothesis as base.  First list all n ordinary coordinates, with the tagged singleton containing the current ordinary coordinate as witness; then list the $\lceil\log_2 n\rceil$ tag coordinates, with their tag singletons as witnesses.  Each witness agrees with the empty base on earlier listed coordinates.  All ordinary and tag coordinates are active, so the effective-range upper bound gives equality.
\end{proof}
\noindent\textbf{Reference.} \cite{li2022understanding,maass1992lower}.

\subsection{Combinatorial Eluder Dimension of the Universal Class}\label{val:combinatorial_eluder_dimension_universal_class}
\noindent\textbf{Statement.} The Combinatorial Eluder Dimension of Universal Class is $n-1$.
\begin{proof}
Take $u_n$ as base.  For each $j<n$, use the domain point $\{j\}$ and witness $u_j$: it differs from $u_n$ at $\{j\}$ and agrees with it at all earlier singleton points.  Thus there is a witness of length $n-1$.  Since the class has n hypotheses, the general size bound gives the matching upper bound.
\end{proof}
\noindent\textbf{Reference.} \cite{li2022understanding,chornomaz2025spherical}.

\subsection{Combinatorial Eluder Dimension of Majority}\label{val:combinatorial_eluder_dimension_majority}
\noindent\textbf{Statement.} The Combinatorial Eluder Dimension of Majority is $n+1$.
\begin{proof}
Use the all-zero ordinary trace (whose majority bit is 0) as base.  Query the majority coordinate first and witness it with the all-one trace.  Then query the n ordinary coordinates, using the corresponding singleton traces as witnesses; for $n\ge2$ these have majority bit 0 and agree with the base on all earlier queries.  This produces n+1 steps, and the effective range has exactly these n+1 coordinates.
\end{proof}
\noindent\textbf{Reference.} \cite{li2022understanding,maass1992lower}.

\subsection{Combinatorial Eluder Dimension of Addressing}\label{val:combinatorial_eluder_dimension_addressing}
\noindent\textbf{Statement.} The Combinatorial Eluder Dimension of Addressing is $n-1$.
\begin{proof}
Choose any target as the base and list the private coordinate of each of the other n-1 targets. At its own listed coordinate that target differs from the base, and at every previously listed private coordinate it agrees with the base. This is an eluder sequence of length n-1. Conversely, in an eluder sequence relative to a base target, each successive witness agrees with the base at all earlier coordinates and disagrees at its own coordinate. These witnesses are distinct and different from the base, so there are at most n-1 steps. For n=1 the empty sequence has value zero.
\end{proof}
\noindent\textbf{Reference.} \cite{li2022understanding,maass1992lower}.

\subsection{Combinatorial Eluder Dimension of Blockwise-Addressing}\label{val:combinatorial_eluder_dimension_blockwiseaddressing}
\noindent\textbf{Statement.} The Combinatorial Eluder Dimension of Blockwise-Addressing is $2^{d+1}-1$.
\begin{proof}
Take $f_i$ as base.  List the primary coordinates $x_j$ for $j\ne i$, witnessed by $f_j$, followed by all $y_j$, witnessed by $g_j$.  Each witness is zero at earlier listed coordinates and differs from $f_i$ at its own coordinate.  This gives $(2^d-1)+2^d=2^{d+1}-1$ steps.  The class has $2^{d+1}$ hypotheses, so the size bound is tight.
\end{proof}
\noindent\textbf{Reference.} \cite{li2022understanding,ben1997online}.

\subsection{Characteristic-p Yang Dimension of Trivial concepts}\label{val:characteristic_p_yang_dimension_trivial_concepts}
\noindent\textbf{Statement.} The Characteristic-p Yang Dimension of Trivial concepts is $1$.
\begin{proof}
The Yang complex of the two constant functions is the disjoint union of two simplices. Its only nonzero reduced homology is in degree 0 over every field, so its Yang dimension is 1 in every positive characteristic.
\end{proof}
\noindent\textbf{Reference.} \cite{yang2018homological}.

\subsection{Characteristic-p Yang Dimension of Singletons}\label{val:characteristic_p_yang_dimension_singletons}
\noindent\textbf{Statement.} The Characteristic-p Yang Dimension of Singletons is $n-1$.
\begin{proof}
For the delta-function class, distinct nonempty subfamilies have distinct least containing cubes, so the Taylor resolution of Yang's canonical ideal is minimal over every field. Its top syzygy has homological degree $n-1$.
\end{proof}
\noindent\textbf{Reference.} \cite{yang2018homological}; see the full arXiv version, Section "Delta Functions"..

\subsection{Characteristic-zero Yang Dimension of Trivial concepts}\label{val:characteristic_zero_yang_dimension_trivial_concepts}
\noindent\textbf{Statement.} The Characteristic-zero Yang Dimension of Trivial concepts is $1$.
\begin{proof}
The Yang complex of the two constant functions is the disjoint union of two simplices. Its only nonzero reduced homology is in degree 0 over $\mathbb Q$, so $\mathrm Y_0=1$.
\end{proof}
\noindent\textbf{Reference.} \cite{yang2018homological}.

\subsection{Characteristic-zero Yang Dimension of Singletons}\label{val:characteristic_zero_yang_dimension_singletons}
\noindent\textbf{Statement.} The Characteristic-zero Yang Dimension of Singletons is $n-1$.
\begin{proof}
For the delta-function class, distinct nonempty subfamilies have distinct least containing cubes, so the Taylor resolution of Yang's canonical ideal is minimal over $\mathbb Q$. Its top syzygy has homological degree $n-1$.
\end{proof}
\noindent\textbf{Reference.} \cite{yang2018homological}; see the full arXiv version, Section "Delta Functions"..

\subsection{Yang dimension of Full cube}\label{val:yang_dimension_full_cube}
\noindent\textbf{Statement.} The Yang dimension of Full Cube is $n$.
\begin{proof}
The canonical ideal of the full binary cube is $\prod_{x=1}^n(v_{x,0},v_{x,1})$. Tensoring the length-one resolutions of its n factors gives a minimal free resolution of length n, so $\mathrm Y=n$.
\end{proof}
\noindent\textbf{Reference.} \cite{yang2018homological}.

\subsection{Characteristic-p Yang Dimension of Full cube}\label{val:characteristic_p_yang_dimension_full_cube}
\noindent\textbf{Statement.} The Characteristic-p Yang Dimension of Full Cube is $n$.
\begin{proof}
The canonical ideal is $\prod_{x=1}^n(v_{x,0},v_{x,1})$. The tensor product of the n length-one factor resolutions is minimal over every field and has length n, so $\mathrm Y_p=n$.
\end{proof}
\noindent\textbf{Reference.} \cite{yang2018homological}.

\subsection{Characteristic-zero Yang Dimension of Full cube}\label{val:characteristic_zero_yang_dimension_full_cube}
\noindent\textbf{Statement.} The Characteristic-zero Yang Dimension of Full Cube is $n$.
\begin{proof}
The canonical ideal is $\prod_{x=1}^n(v_{x,0},v_{x,1})$. The tensor product of the n length-one factor resolutions is minimal over $\mathbb Q$ and has length n, so $\mathrm Y_0=n$.
\end{proof}
\noindent\textbf{Reference.} \cite{yang2018homological}.

\subsection{Yang dimension of Singletons plus empty set}\label{val:yang_dimension_singletons_plus_empty_set}
\noindent\textbf{Statement.} The Yang dimension of Singletons plus empty set is $1$.
\begin{proof}
This is Yang's delta-function class together with the all-zero function. Yang proves that its canonical ideal has a minimal one-dimensional cellular resolution (equivalently, its canonical-suboplex ring is Cohen--Macaulay of regularity 1), so $\mathrm Y=1$.
\end{proof}
\noindent\textbf{Reference.} \cite{yang2018homological}; see the full arXiv version, the delta-functions-plus-zero example..

\subsection{Characteristic-p Yang Dimension of Singletons plus empty set}\label{val:characteristic_p_yang_dimension_singletons_plus_empty_set}
\noindent\textbf{Statement.} The Characteristic-p Yang Dimension of Singletons plus empty set is $1$.
\begin{proof}
The delta-function class together with the all-zero function has Yang's explicit minimal one-dimensional cellular resolution over every field. Hence $\mathrm Y_p=1$.
\end{proof}
\noindent\textbf{Reference.} \cite{yang2018homological}; see the full arXiv version, the delta-functions-plus-zero example..

\subsection{Characteristic-zero Yang Dimension of Singletons plus empty set}\label{val:characteristic_zero_yang_dimension_singletons_plus_empty_set}
\noindent\textbf{Statement.} The Characteristic-zero Yang Dimension of Singletons plus empty set is $1$.
\begin{proof}
The delta-function class together with the all-zero function has Yang's explicit minimal one-dimensional cellular resolution over $\mathbb Q$. Hence $\mathrm Y_0=1$.
\end{proof}
\noindent\textbf{Reference.} \cite{yang2018homological}; see the full arXiv version, the delta-functions-plus-zero example..

\subsection{Effective Range of Universal Class}\label{val:effective_range_universal_class}
\noindent\textbf{Statement.} The Effective Range of Universal Class is $2^n-2$.
\begin{proof}
At the domain points $\varnothing$ and $[n]$, every universal hypothesis has the common value 0 or 1, respectively, so neither is effective. For every proper nonempty subset $A\subset[n]$, choose $i\in A$ and $j\notin A$; then $u_i(A)=1$ and $u_j(A)=0$. Thus all and only the other $2^n-2$ points are effective (also for $n=1$, when this number is zero).
\end{proof}
\noindent\textbf{Reference.} \cite{chornomaz2025spherical}.

\subsection{Size of Universal Class}\label{val:size_universal_class}
\noindent\textbf{Statement.} The Size of Universal Class is $n$.
\begin{proof}
By definition $\mathcal U_n=\{u_i:i\in[n]\}$ consists of exactly n distinct hypotheses.
\end{proof}
\noindent\textbf{Reference.} \cite{chornomaz2025spherical}.

\subsection{Log Size of Universal Class}\label{val:log_size_universal_class}
\noindent\textbf{Statement.} The Log Size of Universal Class is $\log_2 n$.
\begin{proof}
The universal class has exactly n hypotheses, and log size is by definition the binary logarithm of the class size.
\end{proof}
\noindent\textbf{Reference.} \cite{chornomaz2025spherical}.

\subsection{Threshold Dimension of Universal Class}\label{val:threshold_dimension_universal_class}
\noindent\textbf{Statement.} The Threshold Dimension of Universal Class is $n-1$.
\begin{proof}
Use the $n-1$ domain points $A_k=\{1,\ldots,k\}$, ordered as $A_{n-1},A_{n-2},\ldots,A_1$. The traces of $u_n,u_{n-1},\ldots,u_1$ are respectively the empty trace and the successive initial segments, giving a threshold witness of dimension n-1. Conversely, a threshold witness of dimension d contains d+1 distinct traces and therefore needs d+1 distinct hypotheses; since $|\mathcal U_n|=n$, d\leq n-1.
\end{proof}
\noindent\textbf{Reference.} \cite{chornomaz2025spherical}.

\subsection{Path Dimension of Universal Class}\label{val:path_dimension_universal_class}
\noindent\textbf{Statement.} The Path Dimension of Universal Class is $n-1$.
\begin{proof}
On the points $A_{n-1},A_{n-2},\ldots,A_1$, the traces of $u_n,u_{n-1},\ldots,u_1$ start at the empty trace and add exactly one coordinate at every step. This is a path of length n-1. A path of length d contains d+1 distinct traces and therefore requires d+1 distinct hypotheses, so d\leq n-1.
\end{proof}
\noindent\textbf{Reference.} \cite{chornomaz2025spherical}.

\subsection{Star Number of Universal Class}\label{val:star_number_universal_class}
\noindent\textbf{Statement.} The Star number of Universal Class is $n-1$.
\begin{proof}
Fix i and project to the n-1 domain points $\{j\}$ for $j\ne i$. Then $u_i$ has the all-zero trace and each $u_j$ has the corresponding unit-vector trace, giving an (n-1)-star. No star can have more leaves than the n-1 hypotheses other than its center.
\end{proof}
\noindent\textbf{Reference.} \cite{hanneke2015minimax,chornomaz2025spherical}.

\subsection{Projected Maximal Degree of Universal Class}\label{val:projected_maximal_degree_universal_class}
\noindent\textbf{Statement.} The Projected Maximal Degree of Universal Class is $n-1$.
\begin{proof}
The singleton-subset projection in the star-number proof has a vertex of degree n-1. Conversely, every projected class has at most n vertices, so its maximum degree is at most n-1.
\end{proof}
\noindent\textbf{Reference.} \cite{hanneke2015minimax,chornomaz2025spherical}.

\subsection{Maximum Projected Teaching Set Size of Universal Class}\label{val:maximum_projected_teaching_set_size_universal_class}
\noindent\textbf{Statement.} The Maximum Projected Teaching Set Size of Universal Class is $n-1$.
\begin{proof}
Maximum projected teaching-set size equals the star number. The singleton-subset projection gives an (n-1)-star, and no projected class with n hypotheses can have a teaching set larger than n-1.
\end{proof}
\noindent\textbf{Reference.} \cite{hanneke2015minimax,chornomaz2025spherical}.

\subsection{Maximum Degree of Universal Class}\label{val:maximum_degree_universal_class}
\noindent\textbf{Statement.} The Maximum Degree of Universal Class is $0$.
\begin{proof}
For distinct i,j, the hypotheses $u_i$ and $u_j$ disagree precisely on the $2^{n-1}$ subsets containing exactly one of i,j. Thus no two distinct hypotheses differ on one coordinate, so the 1-inclusion graph has no edges.
\end{proof}
\noindent\textbf{Reference.} \cite{haussler1994predicting,chornomaz2025spherical}.

\subsection{Minimum Degree of Universal Class}\label{val:minimum_degree_universal_class}
\noindent\textbf{Statement.} The Minimum Degree of Universal Class is $0$.
\begin{proof}
The 1-inclusion graph is edgeless because every distinct pair $u_i,u_j$ differs on $2^{n-1}$ coordinates. Therefore every vertex has degree zero.
\end{proof}
\noindent\textbf{Reference.} \cite{haussler1994predicting,chornomaz2025spherical}.

\subsection{Double Density or Average Degree of Universal Class}\label{val:double_density_or_average_degree_universal_class}
\noindent\textbf{Statement.} The Double Density or Average Degree of Universal Class is $0$.
\begin{proof}
The 1-inclusion graph is edgeless: distinct universal hypotheses differ on $2^{n-1}$ coordinates, never exactly one. Hence its average degree is zero.
\end{proof}
\noindent\textbf{Reference.} \cite{haussler1994predicting,chornomaz2025spherical}.

\subsection{Densest Subgraph (twice) of Universal Class}\label{val:densest_subgraph_twice_universal_class}
\noindent\textbf{Statement.} The Densest Subgraph (twice) of Universal Class is $0$.
\begin{proof}
The 1-inclusion graph is edgeless, and every one of its subgraphs is edgeless. Thus the maximum subgraph average degree is zero.
\end{proof}
\noindent\textbf{Reference.} \cite{haussler1994predicting,chornomaz2025spherical}.

\subsection{Degeneracy of Universal Class}\label{val:degeneracy_universal_class}
\noindent\textbf{Statement.} The Degeneracy of Universal Class is $0$.
\begin{proof}
The 1-inclusion graph is edgeless, so every induced subgraph has minimum degree zero and its degeneracy is zero.
\end{proof}
\noindent\textbf{Reference.} \cite{haussler1988predicting,chornomaz2025spherical}.

\subsection{Balbach Teaching Dimension of Singletons plus Empty Set}\label{val:balbach_teaching_dimension_singletons_plus_empty_set}
\noindent\textbf{Statement.} The Balbach Teaching Dimension of Singletons plus empty set is $2$.
\begin{proof}
Every singleton has a one-point positive teaching set. For the empty concept, one negative example still leaves singleton hypotheses whose ordinary teaching-set size is one, so the size filter does not isolate it; any two negative examples do. Thus the stabilized Balbach maximum is two (and the source notes the separate degenerate $n=1$ case has value one).
\end{proof}
\noindent\textbf{Reference.} \cite{zilles2011models,balbach2008measuring}.

\subsection{Subset Teaching Dimension of Singletons plus Empty Set}\label{val:subset_teaching_dimension_singletons_plus_empty_set}
\noindent\textbf{Statement.} The Subset Teaching Dimension of Singletons plus empty set is $1$.
\begin{proof}
Each singleton has its one-point positive subset teaching set. At the first subset-teaching iteration, every one-point negative sample is a subset teaching set for the empty concept: it is contained in an ordinary teaching set for the empty concept and in none for a singleton. Hence all concepts have subset teaching dimension one.
\end{proof}
\noindent\textbf{Reference.} \cite{zilles2011models}.

\subsection{Monotonic Yang Dimension of Full Cube}\label{val:monotonic_yang_dimension_full_cube}
\noindent\textbf{Statement.} The Monotonic Yang Dimension of Full Cube is $n$.
\begin{proof}
The full cube itself is an allowed subfamily and has Yang dimension $n$, so $\mathrm Y^*\geq n$. Conversely every subfamily has effective range at most $n$, and the established effective-range Yang bound gives Yang dimension at most $n$. Thus the maximum is exactly $n$.
\end{proof}
\noindent\textbf{Reference.} \cite{yang2018homological}; derived using the established effective-range bound.

\subsection{Monotonic Yang Dimension of Singletons}\label{val:monotonic_yang_dimension_singletons}
\noindent\textbf{Statement.} The Monotonic Yang Dimension of Singletons is $n-1$.
\begin{proof}
The whole delta-function class is an allowed subfamily and has Yang dimension $n-1$. Conversely every subfamily has at most $n$ concepts, and the established generator-count bound gives Yang dimension at most $|\mathcal G|-1\leq n-1$. Therefore $\mathrm Y^*=n-1$.
\end{proof}
\noindent\textbf{Reference.} \cite{yang2018homological}; derived using the established generator-count bound.

\subsection{Prefix-$k$ Littlestone Dimension of Full Cube}\label{val:prefix_littlestone_dimension_full_cube}
\noindent\textbf{Statement.} The Prefix-$k$ Littlestone Dimension ($k>1$) of Full Cube is $n$.
\begin{proof}
Query the $n$ coordinates in the same fixed order at every node. This is a prefix-$k$ tree of depth $n$, so $\mathrm{PL}_k\geq n$. No Littlestone tree on an $n$-coordinate class has depth greater than $n$, and $\mathrm{PL}_k\leq\mathrm L$, giving equality.
\end{proof}
\noindent\textbf{Reference.} Database-defined variant; the comparison $\mathrm{PL}_k\leq\mathrm L$ is Proposition \ref{prop:prefix-littlestone-order} in the survey.

\subsection{Prefix-$k$ Littlestone Dimension of Singletons}\label{val:prefix_littlestone_dimension_singletons}
\noindent\textbf{Statement.} The Prefix-$k$ Littlestone Dimension ($k>1$) of Singletons is $1$.
\begin{proof}
Any coordinate gives a prefix-$k$ shattered tree of depth one. The singleton class has Littlestone dimension one, and $\mathrm{PL}_k\leq\mathrm L$, so no deeper prefix-$k$ tree exists.
\end{proof}
\noindent\textbf{Reference.} Database-defined variant; the comparison $\mathrm{PL}_k\leq\mathrm L$ is Proposition \ref{prop:prefix-littlestone-order} in the survey.

\subsection{Prefix-$k$ Littlestone Dimension of Trivial Concepts}\label{val:prefix_littlestone_dimension_trivial_concepts}
\noindent\textbf{Statement.} The Prefix-$k$ Littlestone Dimension ($k>1$) of Trivial concepts is $1$.
\begin{proof}
Either coordinate gives a prefix-$k$ shattered tree of depth one. The two constant concepts have Littlestone dimension one, and $\mathrm{PL}_k\leq\mathrm L$, so the value is exactly one.
\end{proof}
\noindent\textbf{Reference.} Database-defined variant; the comparison $\mathrm{PL}_k\leq\mathrm L$ is Proposition \ref{prop:prefix-littlestone-order} in the survey.

\subsection{Prefix-$k$ Littlestone Dimension of Addressing}\label{val:prefix_littlestone_dimension_addressing}
\noindent\textbf{Statement.} The Prefix-$k$ Littlestone Dimension ($k>1$) of Addressing is $\lfloor\log n\rfloor\quad(n=2^r,\ r\ge1)$.
\begin{proof}
Restrict here to $n=2^r$ with all r-bit words present, so $\lfloor\log_2n\rfloor=\log_2n=r$. On the addressing class, the established VC and Littlestone dimensions both equal $\lfloor\log n\rfloor$. The surveyed sandwich theorem places Prefix-k Littlestone Dimension between them, determining the same value for every fixed k.
\end{proof}
\noindent\textbf{Reference.} Derived from Proposition \ref{prop:prefix-littlestone-order} in the survey and the established addressing benchmark values.

\subsection{Prefix-$k$ Littlestone Dimension of Majority}\label{val:prefix_littlestone_dimension_majority}
\noindent\textbf{Statement.} The Prefix-$k$ Littlestone Dimension ($k>1$) of Majority is $n$.
\begin{proof}
On the majority class, both the established VC and Littlestone dimensions equal $n$. The surveyed sandwich theorem places Prefix-k Littlestone Dimension between them, so its value is $n$ for every fixed k.
\end{proof}
\noindent\textbf{Reference.} Derived from Proposition \ref{prop:prefix-littlestone-order} in the survey and the established majority benchmark values.

\subsection{Prefix-$k$ Littlestone Dimension of Singletons plus Empty Set}\label{val:prefix_littlestone_dimension_singletons_plus_empty_set}
\noindent\textbf{Statement.} The Prefix-$k$ Littlestone Dimension ($k>1$) of Singletons plus empty set is $1$.
\begin{proof}
Singletons plus the empty set has established VC and Littlestone dimension one. The surveyed sandwich theorem therefore forces Prefix-k Littlestone Dimension to equal one for every fixed k.
\end{proof}
\noindent\textbf{Reference.} Derived from Proposition \ref{prop:prefix-littlestone-order} in the survey and the established benchmark values.

\subsection{Largest Shattered Set of Universal Class}\label{val:largest_shattered_set_universal_class}
\noindent\textbf{Statement.} The Largest Shattered Set of Universal Class is $\lfloor\log_2 n\rfloor$.
\begin{proof}
Largest Shattered Set is exactly VC dimension, whose established value on the universal class is $\lfloor\log_2 n\rfloor$.
\end{proof}
\noindent\textbf{Reference.} \cite{chornomaz2025spherical}.

\subsection{Projected VC Radius of Universal Class}\label{val:projected_vc_radius_universal_class}
\noindent\textbf{Statement.} The Projected VC Radius of Universal Class is $\lfloor\log_2 n\rfloor$.
\begin{proof}
Projected VC Radius is exactly VC dimension, whose established value on the universal class is $\lfloor\log_2 n\rfloor$.
\end{proof}
\noindent\textbf{Reference.} \cite{chornomaz2025spherical}.

\subsection{Projected Strongly Shattered Dimension of Universal Class}\label{val:projected_strongly_shattered_dimension_universal_class}
\noindent\textbf{Statement.} The Projected Strongly Shattered Dimension of Universal Class is $\lfloor\log_2 n\rfloor$.
\begin{proof}
Projected Strongly Shattered Dimension is exactly VC dimension, whose established value on the universal class is $\lfloor\log_2 n\rfloor$.
\end{proof}
\noindent\textbf{Reference.} \cite{chornomaz2025spherical}.

\subsection{Maximum Largest Order Shattered Set of Universal Class}\label{val:maximum_largest_order_shattered_set_universal_class}
\noindent\textbf{Statement.} The Maximum Largest Order Shattered Set of Universal Class is $\lfloor\log_2 n\rfloor$.
\begin{proof}
Maximum Largest Order Shattered Set is exactly VC dimension, whose established value on the universal class is $\lfloor\log_2 n\rfloor$.
\end{proof}
\noindent\textbf{Reference.} \cite{chornomaz2025spherical}.

\subsection{2-Block Littlestone Dimension of Full Cube}\label{val:2block_littlestone_dimension_full_cube}
\noindent\textbf{Statement.} The 2-Block Littlestone Dimension of Full Cube is $n$.
\begin{proof}
The established VC and Littlestone dimensions of the full cube are both n. Proposition \ref{prop:block-littlestone-sandwich} places $\mathrm{VC}\leq\mathrm B_2\leq\mathrm L$, so $\mathrm B_2=n$.
\end{proof}

\subsection{$k$-Block Littlestone Dimension of Full Cube}\label{val:kblock_littlestone_dimension_full_cube}
\noindent\textbf{Statement.} The $k$-Block Littlestone Dimension ($k>2$) of Full Cube is $n$.
\begin{proof}
For every fixed $k>2$, the established VC and Littlestone dimensions of the full cube are both n. Proposition \ref{prop:block-littlestone-sandwich} gives $\mathrm{VC}\leq\mathrm B_k\leq\mathrm L$, hence $\mathrm B_k=n$.
\end{proof}

\subsection{2-Block Littlestone Dimension of Singletons}\label{val:2block_littlestone_dimension_singletons}
\noindent\textbf{Statement.} The 2-Block Littlestone Dimension of Singletons is $1$.
\begin{proof}
The established VC and Littlestone dimensions of singletons are both one. Proposition \ref{prop:block-littlestone-sandwich} gives $1\leq\mathrm B_2\leq1$.
\end{proof}

\subsection{$k$-Block Littlestone Dimension of Singletons}\label{val:kblock_littlestone_dimension_singletons}
\noindent\textbf{Statement.} The $k$-Block Littlestone Dimension ($k>2$) of Singletons is $1$.
\begin{proof}
For every fixed $k>2$, the established VC and Littlestone dimensions of singletons are both one. Proposition \ref{prop:block-littlestone-sandwich} gives $1\leq\mathrm B_k\leq1$.
\end{proof}

\subsection{2-Block Littlestone Dimension of Singletons plus Empty Set}\label{val:2block_littlestone_dimension_singletons_plus_empty_set}
\noindent\textbf{Statement.} The 2-Block Littlestone Dimension of Singletons plus empty set is $1$.
\begin{proof}
The established VC and Littlestone dimensions of singletons plus the empty set are both one. Proposition \ref{prop:block-littlestone-sandwich} gives $1\leq\mathrm B_2\leq1$.
\end{proof}

\subsection{$k$-Block Littlestone Dimension of Singletons plus Empty Set}\label{val:kblock_littlestone_dimension_singletons_plus_empty_set}
\noindent\textbf{Statement.} The $k$-Block Littlestone Dimension ($k>2$) of Singletons plus empty set is $1$.
\begin{proof}
For every fixed $k>2$, the established VC and Littlestone dimensions of singletons plus the empty set are both one. Proposition \ref{prop:block-littlestone-sandwich} gives $1\leq\mathrm B_k\leq1$.
\end{proof}

\subsection{2-Block Littlestone Dimension of Majority}\label{val:2block_littlestone_dimension_majority}
\noindent\textbf{Statement.} The 2-Block Littlestone Dimension of Majority is $n$.
\begin{proof}
The established VC and Littlestone dimensions of the majority class are both n. Proposition \ref{prop:block-littlestone-sandwich} gives $n\leq\mathrm B_2\leq n$.
\end{proof}

\subsection{$k$-Block Littlestone Dimension of Majority}\label{val:kblock_littlestone_dimension_majority}
\noindent\textbf{Statement.} The $k$-Block Littlestone Dimension ($k>2$) of Majority is $n$.
\begin{proof}
For every fixed $k>2$, the established VC and Littlestone dimensions of the majority class are both n. Proposition \ref{prop:block-littlestone-sandwich} gives $n\leq\mathrm B_k\leq n$.
\end{proof}

\subsection{2-Block Littlestone Dimension of Addressing}\label{val:2block_littlestone_dimension_addressing}
\noindent\textbf{Statement.} The 2-Block Littlestone Dimension of Addressing is $\lfloor\log n\rfloor\quad(n=2^r,\ r\ge1)$.
\begin{proof}
Restrict here to $n=2^r$ with all r-bit words present, so $\lfloor\log_2n\rfloor=\log_2n=r$. The established VC and Littlestone dimensions of the addressing class are both $\lfloor\log n\rfloor$. Proposition \ref{prop:block-littlestone-sandwich} forces the same value for $\mathrm B_2$.
\end{proof}

\subsection{$k$-Block Littlestone Dimension of Addressing}\label{val:kblock_littlestone_dimension_addressing}
\noindent\textbf{Statement.} The $k$-Block Littlestone Dimension ($k>2$) of Addressing is $\lfloor\log n\rfloor\quad(n=2^r,\ r\ge1)$.
\begin{proof}
Restrict here to $n=2^r$ with all r-bit words present, so $\lfloor\log_2n\rfloor=\log_2n=r$. For every fixed $k>2$, the established VC and Littlestone dimensions of the addressing class are both $\lfloor\log n\rfloor$. Proposition \ref{prop:block-littlestone-sandwich} forces $\mathrm B_k=\lfloor\log n\rfloor$.
\end{proof}

\subsection{Monotonic Yang Dimension of Trivial Concepts}\label{val:monotonic_yang_dimension_trivial_concepts}
\noindent\textbf{Statement.} The Monotonic Yang Dimension of Trivial concepts is $1$.
\begin{proof}
The whole two-constant class is an allowed subfamily and has Yang dimension 1. Conversely every subfamily has at most two concepts, and Yang dimension is at most the number of generators minus one, hence at most 1. Therefore $\mathrm Y^*=1$.
\end{proof}
\noindent\textbf{Reference.} \cite{yang2018homological}; derived using the established generator-count bound.

\subsection{Monotonic Yang Dimension of Singletons plus Empty Set}\label{val:monotonic_yang_dimension_singletons_plus_empty_set}
\noindent\textbf{Statement.} The Monotonic Yang Dimension of Singletons plus empty set is $\max\{1,n-1\}$.
\begin{proof}
The singleton subfamily has Yang dimension $n-1$. The whole class has Yang dimension one for every $n\ge1$, giving the other lower bound. A subfamily containing zero is either a singleton (Yang dimension zero), or a delta-functions-plus-zero class on its active coordinates with additional constant coordinates, and has Yang dimension one. Without zero it has at most $n$ concepts, so the generator-count bound gives at most $n-1$. Thus $\mathrm Y^*=\max\{1,n-1\}$; in particular the value is one, not zero, at $n=1$.
\end{proof}
\noindent\textbf{Reference.} \cite{yang2018homological}; derived using the established delta-functions-plus-zero and generator-count bounds.

\subsection{Best Mistakes of Blockwise-Addressing}\label{val:best_mistakes_blockwiseaddressing}
\noindent\textbf{Statement.} The Best Mistakes of Blockwise-Addressing is $\Theta(d)=\Theta(\log |\mathcal X|)$.
\begin{proof}
Example 2 of Ben-David--Kushilevitz--Mansour proves that every fixed presentation order admits an adversarial target forcing $\Omega(d)$ mistakes.  Its general upper bound, together with the class size $2^{d+1}$, gives $O(d)$; hence the best fixed-order mistake bound is $\Theta(d)=\Theta(\log|\mathcal X|)$.
\end{proof}
\noindent\textbf{Reference.} \cite{ben1997online}, Example 2.

\subsection{Proper Stable Labeled Sample Compression of Trivial Concepts}\label{val:proper_stable_labeled_sample_compression_trivial_concepts}
\noindent\textbf{Statement.} The Proper Stable Labeled Sample Compression of Trivial concepts is $1$.
\begin{proof}
The established proper stable unlabeled scheme of size one is also a proper stable labeled scheme after retaining its label, so the value is at most one. Proper labeled compression has value one for the two constant hypotheses, and proper stable labeled compression is a restriction of that model, giving the matching lower bound.
\end{proof}
\noindent\textbf{Reference.} \cite{FW1995,hanneke2021stable,bousquet2020proper}; derived from the established proper stable unlabeled and proper labeled benchmark values.

\subsection{Balbach Teaching Dimension of Trivial Concepts}\label{val:balbach_teaching_dimension_trivial_concepts}
\noindent\textbf{Statement.} The Balbach Teaching Dimension of Trivial concepts is $1$.
\begin{proof}
For this class, ordinary teaching dimension and recursive teaching dimension are both one. The established sandwich $\mathrm{TD}\ge\mathrm{BTD}\ge\mathrm{RTD}$ therefore forces $\mathrm{BTD}=1$.
\end{proof}
\noindent\textbf{Reference.} \cite{zilles2011models,balbach2008measuring}; derived using established TD and RTD values.

\subsection{Balbach Teaching Dimension of Singletons}\label{val:balbach_teaching_dimension_singletons}
\noindent\textbf{Statement.} The Balbach Teaching Dimension of Singletons is $1$.
\begin{proof}
For the singleton class, ordinary teaching dimension and recursive teaching dimension are both one. The established sandwich $\mathrm{TD}\ge\mathrm{BTD}\ge\mathrm{RTD}$ therefore forces $\mathrm{BTD}=1$.
\end{proof}
\noindent\textbf{Reference.} \cite{zilles2011models,balbach2008measuring}; derived using established TD and RTD values.

\subsection{Balbach Teaching Dimension of Full Cube}\label{val:balbach_teaching_dimension_full_cube}
\noindent\textbf{Statement.} The Balbach Teaching Dimension of Full Cube is $n$.
\begin{proof}
For the full cube, ordinary teaching dimension and recursive teaching dimension are both $n$. The established sandwich $\mathrm{TD}\ge\mathrm{BTD}\ge\mathrm{RTD}$ therefore forces $\mathrm{BTD}=n$.
\end{proof}
\noindent\textbf{Reference.} \cite{zilles2011models,balbach2008measuring}; derived using established TD and RTD values.

\subsection{Balbach Teaching Dimension of Addressing}\label{val:balbach_teaching_dimension_addressing}
\noindent\textbf{Statement.} The Balbach Teaching Dimension of Addressing is $1\quad(n\ge2)$.
\begin{proof}
For Addressing, ordinary teaching dimension and recursive teaching dimension are both one. The established sandwich $\mathrm{TD}\ge\mathrm{BTD}\ge\mathrm{RTD}$ therefore forces $\mathrm{BTD}=1$.
\end{proof}
\noindent\textbf{Reference.} \cite{zilles2011models,balbach2008measuring,maass1992lower}; derived using established TD and RTD values.

\subsection{Balbach Teaching Dimension of Warmuth's $C_5$ Class}\label{val:balbach_teaching_dimension_warmuth_c5}
\noindent\textbf{Statement.} The Balbach Teaching Dimension of Warmuth's $C_5$ class is $3$.
\begin{proof}
For Warmuth's $C_5$ class, ordinary teaching dimension and recursive teaching dimension are both three. The established sandwich $\mathrm{TD}\ge\mathrm{BTD}\ge\mathrm{RTD}$ therefore forces $\mathrm{BTD}=3$.
\end{proof}
\noindent\textbf{Reference.} \cite{zilles2011models,balbach2008measuring,palvolgyi2021unlabeled}; derived using established TD and RTD values.

\subsection{Self-Directed Queries Complexity of Blockwise-Addressing}\label{val:selfdirected_queries_complexity_blockwiseaddressing}
\noindent\textbf{Statement.} The Self-Directed Queries Complexity of Blockwise-Addressing is $\leq 2$.
\begin{proof}
Example 2 first queries $z$, which determines whether the target is an $f_i$ or a $g_i$. It then queries the corresponding private coordinates while predicting 0; the unique positive coordinate identifies the target. At most one mistake occurs at $z$ and at most one afterwards, so $\mathrm{SDC}\le2$.
\end{proof}
\noindent\textbf{Reference.} \cite{ben1997online}, Example 2.

\subsection{Information Complexity of Trivial concepts}\label{val:information_complexity_trivial_concepts}
\noindent\textbf{Statement.} The Information Complexity of Trivial concepts is $0$.
\begin{proof}
On every realizable sample for either constant target, the first observed label determines that target.  Output that constant function.  Conditional on any realizable distribution, the output is the fixed target function and is therefore independent of the random sample, so $I(S;\mathcal A(S))=0$.  Nonnegativity of mutual information gives equality.
\end{proof}
\noindent\textbf{Reference.} \cite{pradeep2022information}.

\subsection{Worst Mistakes of Full Cube}\label{val:worst_mistakes_full_cube}
\noindent\textbf{Statement.} The Worst Mistakes of Full Cube is $n$.
\begin{proof}
Even when the entire instance sequence is known in advance, the learner can err at most once at each coordinate: predict arbitrarily on its first occurrence and reuse the observed label afterwards.  Conversely, present all $n$ coordinates once.  Against any deterministic learner, choose each target label opposite to its first prediction.  These $n$ labels define a full-cube target, so the adversary forces $n$ mistakes.
\end{proof}
\noindent\textbf{Reference.} \cite{ben1997online}.

\subsection{Worst Mistakes of Singletons}\label{val:worst_mistakes_singletons}
\noindent\textbf{Statement.} The Worst Mistakes of Singletons is $1$.
\begin{proof}
Predict $0$ until a positive label appears.  It can appear only at the target singleton's unique coordinate, so this strategy makes at most one mistake and then identifies the target.  Zero mistakes are impossible because at any coordinate some singleton target has label $1$ and another has label $0$.
\end{proof}
\noindent\textbf{Reference.} \cite{ben1997online}.

\subsection{Worst Mistakes of Singletons plus empty set}\label{val:worst_mistakes_singletons_plus_empty_set}
\noindent\textbf{Statement.} The Worst Mistakes of Singletons plus empty set is $1$.
\begin{proof}
Always predict $0$.  The empty target causes no mistake, while a singleton target causes exactly one, at its unique positive coordinate.  Zero mistakes are impossible because the empty target and any singleton disagree at that singleton's coordinate.
\end{proof}
\noindent\textbf{Reference.} \cite{ben1997online}.

\subsection{List Replicability Number of Full Cube}\label{val:list_replicability_full_cube}
\noindent\textbf{Statement.} The List Replicability Number of Full Cube is $n$.
\begin{proof}
The full cube is extremal of VC dimension $n$. The cited theorem gives simplicial covering dimension $n-1$ in this exceptional full-cube case, and its equivalence $\mathrm{LR}=\mathrm{SCdim}+1$ yields $\mathrm{LR}(\{0,1\}^{[n]})=n$.
\end{proof}
\noindent\textbf{Reference.} \cite{blondal2026simplicialcovering}.

\subsection{List Replicability Number of Half-intervals}\label{val:list_replicability_halfintervals}
\noindent\textbf{Statement.} The List Replicability Number of Half-intervals is $2$.
\begin{proof}
The half-interval class is a proper extremal class of VC dimension one. The cited extremal-class theorem gives $\mathrm{SCdim}=1$, and $\mathrm{LR}=\mathrm{SCdim}+1$ gives $\mathrm{LR}=2$.
\end{proof}
\noindent\textbf{Reference.} \cite{blondal2026simplicialcovering}.

\subsection{List Replicability Number of Singletons plus Empty Set}\label{val:list_replicability_singletons_plus_empty_set}
\noindent\textbf{Statement.} The List Replicability Number of Singletons plus empty set is $2$.
\begin{proof}
This is a proper extremal class of VC dimension one: its $n+1$ concepts and its shattered sets are respectively the empty set plus the $n$ singletons. The cited theorem gives $\mathrm{SCdim}=1$, and $\mathrm{LR}=\mathrm{SCdim}+1$ gives $\mathrm{LR}=2$.
\end{proof}
\noindent\textbf{Reference.} \cite{blondal2026simplicialcovering}.

\subsection{Subset Teaching Dimension of Full Cube}\label{val:subset_teaching_dimension_full_cube}
\noindent\textbf{Statement.} The Subset Teaching Dimension of Full Cube is $n$.
\begin{proof}
For a cube concept, its only ordinary minimum teaching sample is its complete labeled graph. A sample is contained in no other concept's such graph exactly when it specifies every coordinate, so the first subset-teaching iteration still has value n and the same unique sample. Hence the process is already stable and $\mathrm{STD}=n$.
\end{proof}
\noindent\textbf{Reference.} \cite{zilles2011models}.

\subsection{Subset Teaching Dimension of Singletons}\label{val:subset_teaching_dimension_singletons}
\noindent\textbf{Statement.} The Subset Teaching Dimension of Singletons is $1$.
\begin{proof}
The sole ordinary minimum teaching sample for singleton $\{i\}$ is $(i,1)$. This sample is contained in that target's teaching sample and in no other singleton's, so it remains a subset teaching sample at every iteration. The empty sample is contained in every target's sample, hence cannot identify one. Thus $\mathrm{STD}=1$.
\end{proof}
\noindent\textbf{Reference.} \cite{zilles2011models}.

\subsection{Subset Teaching Dimension of Trivial Concepts}\label{val:subset_teaching_dimension_trivial_concepts}
\noindent\textbf{Statement.} The Subset Teaching Dimension of Trivial concepts is $1$.
\begin{proof}
Every one-point sample labeled 0 is a minimum teaching sample only for the empty concept, and every one labeled 1 only for the full concept. These samples therefore remain subset teaching samples at every iteration; the empty sample identifies neither of the two targets. Thus $\mathrm{STD}=1$.
\end{proof}
\noindent\textbf{Reference.} \cite{zilles2011models}.

\subsection{Dual Antipodal VC Dimension of Full Cube}\label{val:dual_antipodal_vc_dimension_full_cube}
\noindent\textbf{Statement.} The Dual Antipodal VC Dimension of Full Cube is $\lfloor\log_2 n\rfloor+1$.
\begin{proof}
For a dually antipodally shattered set of $d$ cube functions, the $n$ coordinates must realize all $2^{d-1}$ complementary pairs of $d$-bit patterns, so $2^{d-1}\leq n$. Conversely, for $d=\lfloor\log_2 n\rfloor+1$, choose $d$ functions whose values on $2^{d-1}$ coordinates are one representative from every complementary pair; assign their remaining values arbitrarily. These functions are dually antipodally shattered.
\end{proof}
\noindent\textbf{Reference.} \cite{chornomaz2025spherical}.

\subsection{Dual Antipodal VC Dimension of Singletons}\label{val:dual_antipodal_vc_dimension_singletons}
\noindent\textbf{Statement.} The Dual Antipodal VC Dimension of Singletons is $\begin{cases}1&n\leq2,\\2&n=3,\\3&n\geq4.\end{cases}$.
\begin{proof}
On any selected singleton hypotheses, every coordinate trace is a unit vector, together with the zero vector when an unselected coordinate remains. Thus at most $d+1$ patterns occur on $d$ selected hypotheses, whereas antipodal shattering requires $2^{d-1}$ complementary pairs. This excludes $d\geq4$. The zero vector and three unit vectors witness $d=3$ for $n\geq4$; the same count gives the two smaller boundary values.
\end{proof}
\noindent\textbf{Reference.} \cite{chornomaz2025spherical}.

\subsection{Dual Antipodal VC Dimension of Singletons plus Empty Set}\label{val:dual_antipodal_vc_dimension_singletons_plus_empty_set}
\noindent\textbf{Statement.} The Dual Antipodal VC Dimension of Singletons plus empty set is $\begin{cases}1&n=1,\\2&2\leq n\leq3,\\3&n\geq4.\end{cases}$.
\begin{proof}
For a selected set of singleton hypotheses, coordinate traces are unit vectors and, when available, the zero vector. Including the empty hypothesis only forces one selected coordinate of every trace to be zero, so it cannot create the missing complementary pairs. Hence $d\geq4$ is impossible. For $n\geq4$, three singleton hypotheses and one outside coordinate realize the zero vector and all three unit vectors, giving $d=3$; direct inspection of $n=1,2,3$ gives the stated boundary values.
\end{proof}
\noindent\textbf{Reference.} \cite{chornomaz2025spherical}.

\subsection{Dual Antipodal VC Dimension of Half-intervals}\label{val:dual_antipodal_vc_dimension_halfintervals}
\noindent\textbf{Statement.} The Dual Antipodal VC Dimension of Half-intervals is $\begin{cases}1&n=1,\\2&n\geq2.\end{cases}$.
\begin{proof}
The traces of any chosen initial segments form a chain. For two distinct initial segments, the two endpoint coordinates induce $11$ and one of $01,10$, so both complementary pairs on two hypotheses occur when $n\geq2$. After ordering three chosen segments by reverse inclusion, every coordinate trace is one of $000$, $100$, $110$, and $111$; these cover only three of the four complementary pairs, so three hypotheses cannot be dually antipodally shattered. Thus the value is two for $n\geq2$; the one-point case has value one.
\end{proof}
\noindent\textbf{Reference.} \cite{chornomaz2025spherical}.

\subsection{Dual Antipodal VC Dimension of Trivial concepts}\label{val:dual_antipodal_vc_dimension_trivial_concepts}
\noindent\textbf{Statement.} The Dual Antipodal VC Dimension of Trivial concepts is $1$.
\begin{proof}
Either constant hypothesis alone is dually antipodally shattered, since its one-bit trace and its complement cover both labels. For the two constants together, every point has the same trace $01$ (up to ordering), which covers only one of the two complementary pairs required for two hypotheses. Thus the value is one.
\end{proof}
\noindent\textbf{Reference.} \cite{chornomaz2025spherical}.

\subsection{SCO Dimension of Trivial Concepts}\label{val:sco_dimension_trivial_concepts}
\noindent\textbf{Statement.} The SCO Dimension of Trivial concepts is $1$.
\begin{proof}
For the two constant concepts, the established spherical dimension is $0$, so the SCO lower bound gives $\mathrm{SCO}\geq\mathrm{sd}+1=1$. Their established sign rank is $1$, and the sign-rank SCO reduction gives $\mathrm{SCO}\leq1$. Thus the value is exactly one.
\end{proof}
\noindent\textbf{Reference.} \cite{chornomaz2025spherical,alon2016sign}.

\subsection{Best Mistakes of Full Cube}\label{val:best_mistakes_full_cube}
\noindent\textbf{Statement.} The Best Mistakes of Full Cube is $n$.
\begin{proof}
For any fixed favorable presentation order, predict arbitrarily on a coordinate's first occurrence and reuse its observed label thereafter, so at most one mistake is made at each of the $n$ coordinates. Conversely, the adversary chooses the target's label opposite to the learner's first prediction at every coordinate when it is first presented. These labels define a full-cube target and force $n$ mistakes. Thus the best-mistakes value is $n$.
\end{proof}
\noindent\textbf{Reference.} \cite{ben1997online}.

\subsection{Best Mistakes of Singletons}\label{val:best_mistakes_singletons}
\noindent\textbf{Statement.} The Best Mistakes of Singletons is $\begin{cases}0&n=1,\\1&n\geq2.\end{cases}$.
\begin{proof}
For $n\geq2$, present any $n-1$ coordinates and predict $0$ at each. A positive label identifies its singleton target and causes one mistake; if all are negative, the unpresented coordinate is the target, with no mistake. Hence one mistake suffices. On the first presented coordinate, any learner prediction can be falsified: choose its singleton if it predicts $0$, and a different singleton if it predicts $1$. Thus one mistake is necessary. For $n=1$ the target is already known, so the value is zero.
\end{proof}
\noindent\textbf{Reference.} \cite{ben1997online}.

\subsection{Best Mistakes of Singletons plus Empty Set}\label{val:best_mistakes_singletons_plus_empty_set}
\noindent\textbf{Statement.} The Best Mistakes of Singletons plus empty set is $1$.
\begin{proof}
Present every coordinate and predict $0$. The empty target causes no mistake and each singleton target causes one positive label, which identifies it; hence one mistake suffices. At the first presented coordinate, an initial prediction of $0$ is falsified by its singleton target, while a prediction of $1$ is falsified by the empty target. Therefore one mistake is necessary.
\end{proof}
\noindent\textbf{Reference.} \cite{ben1997online}.

\subsection{Extended Teaching Dimension of Majority}\label{val:extended_teaching_dimension_majority}
\noindent\textbf{Statement.} The Extended Teaching Dimension of Majority is $n$.
\begin{proof}
Extended teaching dimension equals star number for finite binary classes. The established star number of Majority is $n$, so $\mathrm{XTD}=n$.
\end{proof}
\noindent\textbf{Reference.} \cite{hanneke2015minimax,maass1992lower}.

\subsection{Extended Teaching Dimension of Addressing}\label{val:extended_teaching_dimension_addressing}
\noindent\textbf{Statement.} The Extended Teaching Dimension of Addressing is $n-1$.
\begin{proof}
Extended teaching dimension equals star number for finite binary classes. The established star number of Addressing is $n-1$, so $\mathrm{XTD}=n-1$.
\end{proof}
\noindent\textbf{Reference.} \cite{hanneke2015minimax,maass1992lower}.

\subsection{Extended Teaching Dimension of Tagged Singletons}\label{val:extended_teaching_dimension_tagged_singletons}
\noindent\textbf{Statement.} The Extended Teaching Dimension of Tagged Singletons is $n$.
\begin{proof}
Extended teaching dimension equals star number for finite binary classes. The established star number of Tagged Singletons is $n$, so $\mathrm{XTD}=n$.
\end{proof}
\noindent\textbf{Reference.} \cite{hanneke2015minimax,maass1992lower}.

\subsection{Extended Teaching Dimension of Warmuth's $C_5$ Class}\label{val:extended_teaching_dimension_warmuth_c5}
\noindent\textbf{Statement.} The Extended Teaching Dimension of Warmuth's $C_5$ class is $3$.
\begin{proof}
Extended teaching dimension equals star number for finite binary classes. The established star number of Warmuth's $C_5$ class is $3$, so $\mathrm{XTD}=3$.
\end{proof}
\noindent\textbf{Reference.} \cite{hanneke2015minimax,palvolgyi2021unlabeled}.

\subsection{Extended Teaching Dimension of Universal Class}\label{val:extended_teaching_dimension_universal_class}
\noindent\textbf{Statement.} The Extended Teaching Dimension of Universal Class is $n-1$.
\begin{proof}
Extended teaching dimension equals star number for finite binary classes. The established star number of Universal Class is $n-1$, so $\mathrm{XTD}=n-1$.
\end{proof}
\noindent\textbf{Reference.} \cite{hanneke2015minimax,chornomaz2025spherical}.

\subsection{Minimum Star Number of Trivial Concepts}\label{val:minimum_star_number_trivial_concepts}
\noindent\textbf{Statement.} The Minimum Star Number of Trivial concepts is $1$.
\begin{proof}
The general bounds $\mathrm{VC}\leq\mathrm S_{\mathrm{min}}\leq\mathrm S$ are both one for the two constant concepts, so the minimum star number is one.
\end{proof}
\noindent\textbf{Reference.} \cite{hanneke2024star_eluder}.

\subsection{Minimum Star Number of Full Cube}\label{val:minimum_star_number_full_cube}
\noindent\textbf{Statement.} The Minimum Star Number of Full Cube is $n$.
\begin{proof}
The general bounds $\mathrm{VC}\leq\mathrm S_{\mathrm{min}}\leq\mathrm S$ are both n for the full cube, so the minimum star number is n.
\end{proof}
\noindent\textbf{Reference.} \cite{hanneke2024star_eluder}.

\subsection{Minimum Star Number of Singletons plus Empty Set}\label{val:minimum_star_number_singletons_plus_empty_set}
\noindent\textbf{Statement.} The Minimum Star Number of Singletons plus empty set is $1$.
\begin{proof}
This class is intersection-closed and has VC dimension one. Example 2 of the cited paper gives $\mathrm S_{\mathrm{min}}=\mathrm{VC}$ for intersection-closed classes, so the value is one.
\end{proof}
\noindent\textbf{Reference.} \cite{hanneke2024star_eluder}.

\subsection{Minimum Star Number of Half-intervals}\label{val:minimum_star_number_halfintervals}
\noindent\textbf{Statement.} The Minimum Star Number of Half-intervals is $\begin{cases}0&n=1,\\1&n\geq2.\end{cases}$.
\begin{proof}
For $n=1$ there is one initial segment, so no nonempty star exists. For $n\ge2$, the class of initial segments is intersection-closed and has VC dimension one. Example 2 of the cited paper gives $\mathrm S_{\mathrm{min}}=\mathrm{VC}$ for intersection-closed classes, so the value is one.
\end{proof}
\noindent\textbf{Reference.} \cite{hanneke2024star_eluder}.

\subsection{Minimum Star Number of Majority}\label{val:minimum_star_number_majority}
\noindent\textbf{Statement.} The Minimum Star Number of Majority is $n$.
\begin{proof}
The general bounds $\mathrm{VC}\leq\mathrm S_{\mathrm{min}}\leq\mathrm S$ are both n for Majority, so the minimum star number is n.
\end{proof}
\noindent\textbf{Reference.} \cite{hanneke2024star_eluder,maass1992lower}.

\subsection{Minimum Star Number of Singletons}\label{val:minimum_star_number_singletons}
\noindent\textbf{Statement.} The Minimum Star Number of Singletons is $\begin{cases}0&n=1,\\1&n\geq2.\end{cases}$.
\begin{proof}
For $n=1$ there is one hypothesis, so no nonempty star exists. For $n\ge2$, VC dimension one gives $\mathrm S_{\mathrm{min}}\geq1$. For the all-one centre labeling, a centre hypothesis in the singleton class can agree with the centre on at most one selected coordinate, so every centred star has size at most one. Hence $\mathrm S_{\mathrm{min}}=1$.
\end{proof}
\noindent\textbf{Reference.} \cite{hanneke2024star_eluder}.

\subsection{Minimum Star Number of Tagged Singletons}\label{val:minimum_star_number_tagged_singletons}
\noindent\textbf{Statement.} The Minimum Star Number of Tagged Singletons is $1$.
\begin{proof}
Tagged Singletons is intersection-closed: the intersection of two distinct tagged pairs is either empty or their common tag singleton, and all other intersections are also class members. Its VC dimension is one, so Example 2 of the cited paper gives $\mathrm S_{\mathrm{min}}=\mathrm{VC}=1$.
\end{proof}
\noindent\textbf{Reference.} \cite{hanneke2024star_eluder,maass1992lower}.

\subsection{Minimum Star Number of Warmuth's $C_5$ Class}\label{val:minimum_star_number_warmuth_c5}
\noindent\textbf{Statement.} The Minimum Star Number of Warmuth's $C_5$ class is $3$.
\begin{proof}
The checked exact enumerator constructs the ten cyclic $1001$-pattern concepts, enumerates all $2^5$ centre labelings and all coordinate subsets, and implements the centred-star definition directly. Its minimum is $3$.
\end{proof}
\noindent\textbf{Reference.} \cite{hanneke2024star_eluder,palvolgyi2021unlabeled}.

\subsection{Minimum Star Number of Universal Class}\label{val:minimum_star_number_universal_class}
\noindent\textbf{Statement.} The Minimum Star Number of Universal Class is $\tfrac n2-1\leq\mathrm S_{\mathrm{min}}(\mathcal U_n)\leq n-1$.
\begin{proof}
The universal class has dual VC dimension n, so Corollary 25 of the cited paper gives $\mathrm S_{\mathrm{min}}\geq n/2-1$. Its established ordinary star number is $n-1$, and $\mathrm S_{\mathrm{min}}\leq\mathrm S$ gives the upper bound.
\end{proof}
\noindent\textbf{Reference.} \cite{hanneke2024star_eluder,chornomaz2025spherical}.

\subsection{Radon Number of Singletons}\label{val:radon_number_singletons}
\noindent\textbf{Statement.} The Radon Number of Singletons is $n$.
\begin{proof}
Example 1 of the cited paper proves that the n singleton hypotheses are Radon independent, hence the Radon number is n.
\end{proof}
\noindent\textbf{Reference.} \cite{chase2024dual}.

\subsection{Radon Number of Full Cube}\label{val:radon_number_full_cube}
\noindent\textbf{Statement.} The Radon Number of Full Cube is $\lfloor\log_2(2n+2)\rfloor$.
\begin{proof}
Proposition 22 of the cited paper proves that the n-dimensional full cube has Radon number $\lfloor\log_2(2n+2)\rfloor$.
\end{proof}
\noindent\textbf{Reference.} \cite{chase2024dual}.

\subsection{Radon Number of Trivial Concepts}\label{val:radon_number_trivial_concepts}
\noindent\textbf{Statement.} The Radon Number of Trivial concepts is 2.
\begin{proof}
The class consists of the two constant hypotheses.  They form a Radon-independent pair: for their only nontrivial bipartition, any coordinate separates the two singleton convex hulls.  No Radon-independent set can have more than the two hypotheses in the class.
\end{proof}
\noindent\textbf{Reference.} \cite{chase2024dual}.

\subsection{Radon Number of Halfintervals}\label{val:radon_number_halfintervals}
\noindent\textbf{Statement.} The Radon Number of Half-intervals is 2.
\begin{proof}
The halfintervals are linearly ordered by inclusion, and their version-space convexity is the ordinary interval convexity on this chain.  Every three distinct members have a middle one; the partition consisting of that middle member and the other two has intersecting convex hulls.  Any two distinct halfintervals are Radon independent, so the value is 2.
\end{proof}
\noindent\textbf{Reference.} \cite{chase2024dual}.

\subsection{Radon Number of Singletons plus Empty Set}\label{val:radon_number_singletons_plus_empty_set}
\noindent\textbf{Statement.} The Radon Number of Singletons plus empty set is $\begin{cases}2&n\leq2,\\3&n\geq3.\end{cases}$.
\begin{proof}
For $n\geq3$, any three singleton hypotheses are Radon independent: in every bipartition, the singleton side is disjoint from the hull of the other two.  No four hypotheses are Radon independent.  If four singletons are chosen, split them into two pairs; both pair hulls contain the empty hypothesis.  If the empty hypothesis and three singletons are chosen, split them into a pair containing the empty hypothesis and a pair of singletons; again both hulls contain the empty hypothesis.  The classes for $n=1,2$ are checked directly, giving 2.
\end{proof}
\noindent\textbf{Reference.} \cite{chase2024dual}.

\subsection{Sign Rank of Singletons}\label{val:sign_rank_singletons}
\noindent\textbf{Statement.} The Sign Rank of Singletons is $\begin{cases}1&n\leq2,\\2&n=3,\\3&n\geq4.\end{cases}$.
\begin{proof}
The incidence sign matrix of the singleton class is the $n\times n$ signed identity matrix: it has positive diagonal and negative off-diagonal entries.  Its sign rank is 3 for $n\geq4$ (Claim 1 in the cited paper).  For $n=3$, three vectors at angles $0,120,240$ degrees give a rank-two realization; for $n\leq2$, the displayed sign matrix has rank one.
\end{proof}
\noindent\textbf{Reference.} \cite{alon2016sign}.

\subsection{Dual Sign Rank of Singletons}\label{val:dual_sign_rank_singletons}
\noindent\textbf{Statement.} The Dual Sign Rank of Singletons is $\begin{cases}1&n\leq2,\\2&n=3,\\3&n\geq4.\end{cases}$.
\begin{proof}
The singleton incidence sign matrix is the signed identity matrix and is unchanged by transposition.  Therefore its dual sign rank equals its sign rank, whose exact value is the displayed piecewise formula.
\end{proof}
\noindent\textbf{Reference.} \cite{alon2016sign}.

\subsection{Sign Rank of Halfintervals}\label{val:sign_rank_halfintervals}
\noindent\textbf{Statement.} The Sign Rank of Half-intervals is $\begin{cases}1&n=1,\\2&n\geq2.\end{cases}$.
\begin{proof}
Index the initial segments by $i$ and the domain points by $x$, both from 1 to $n$.  The rank-two real matrix $M_{i,x}=i-x+\tfrac12$ is positive exactly when $x\leq i$, so it realizes the incidence signs.  For $n\geq2$, the first two rows are neither equal nor opposite, whereas every nonzero rank-one sign matrix has all rows equal or opposite; hence rank one is impossible.  The $n=1$ matrix has rank one.
\end{proof}
\noindent\textbf{Reference.} \cite{alon2016sign}.

\subsection{Dual Sign Rank of Halfintervals}\label{val:dual_sign_rank_halfintervals}
\noindent\textbf{Statement.} The Dual Sign Rank of Half-intervals is $\begin{cases}1&n=1,\\2&n\geq2.\end{cases}$.
\begin{proof}
Transposing the signed threshold-incidence matrix leaves its rank-two realization argument unchanged.  Thus dual sign rank has the same value as sign rank for halfintervals.
\end{proof}
\noindent\textbf{Reference.} \cite{alon2016sign}.

\subsection{Radon Number of Universal Class}\label{val:radon_number_universal_class}
\noindent\textbf{Statement.} The Radon Number of Universal Class is $n$.
\begin{proof}
Write the hypotheses as $u_i(A)=1$ exactly when $i\in A$, for $A\subseteq[n]$.  For every subset $T\subseteq[n]$, the one-point positive sample at $A=T$ has version space exactly $\{u_i:i\in T\}$.  Hence every subset of hypotheses is convex, every convex hull is the set itself, and all $n$ hypotheses form a Radon-independent set.  The cardinality bound gives equality.
\end{proof}
\noindent\textbf{Reference.} \cite{chase2024dual}.

\subsection{SCO Dimension of Singletons}\label{val:sco_dimension_singletons}
\noindent\textbf{Statement.} The SCO Dimension of Singletons is $\begin{cases}2&n=3,\\2\leq\mathrm{SCO}\leq3&n\geq4.\end{cases}$.
\begin{proof}
For $n\geq3$, the spherical dimension of singletons is 1, so Theorem 12 gives $\mathrm{SCO}\geq2$.  Their sign rank is 2 at n=3 and 3 for $n\geq4$; the sign-rank construction in the same theorem gives the corresponding upper bounds.  Thus n=3 is exact, while the displayed interval is established for larger n.
\end{proof}
\noindent\textbf{Reference.} \cite{chornomaz2025spherical,alon2016sign}.

\subsection{SCO Dimension of Halfintervals}\label{val:sco_dimension_halfintervals}
\noindent\textbf{Statement.} The SCO Dimension of Half-intervals is $1\leq\mathrm{SCO}\leq2$.
\begin{proof}
The spherical dimension of every non-singleton threshold class is 0, so Theorem 12 gives $\mathrm{SCO}\geq1$.  The halfinterval incidence matrix has sign rank 2, and the sign-rank SCO reduction gives $\mathrm{SCO}\leq2$.
\end{proof}
\noindent\textbf{Reference.} \cite{chornomaz2025spherical,alon2016sign}.

\subsection{SCO Dimension of Universal Class}\label{val:sco_dimension_universal_class}
\noindent\textbf{Statement.} The SCO Dimension of Universal Class is $n-1\leq\mathrm{SCO}\leq n$.
\begin{proof}
The established spherical-dimension bound $\mathrm{sd}(\mathcal U_n)\geq n-2$ and Theorem 12 give $\mathrm{SCO}(\mathcal U_n)\geq n-1$.  The universal class is the dual of the full cube, so its sign rank equals the established dual sign rank n of the full cube.  The sign-rank SCO reduction gives $\mathrm{SCO}(\mathcal U_n)\leq n$.
\end{proof}
\noindent\textbf{Reference.} \cite{chornomaz2025spherical,alon2016sign}.

\subsection{SCO Dimension of Singletons plus Empty Set}\label{val:sco_dimension_singletons_plus_empty_set}
\noindent\textbf{Statement.} The SCO Dimension of Singletons plus empty set is $2\leq\mathrm{SCO}\leq3$.
\begin{proof}
For $n\geq3$, the established spherical dimension is 1, so Theorem 12 gives $\mathrm{SCO}\geq2$.  This class has VC dimension one, and the cited sign-rank theorem gives sign rank at most 3 for every VC-one sign matrix.  The sign-rank SCO reduction therefore gives $\mathrm{SCO}\leq3$.
\end{proof}
\noindent\textbf{Reference.} \cite{chornomaz2025spherical,alon2016sign}.

\subsection{SCO Dimension of Tagged Singletons}\label{val:sco_dimension_tagged_singletons}
\noindent\textbf{Statement.} The SCO Dimension of Tagged Singletons is $2\leq\mathrm{SCO}\leq3$.
\begin{proof}
For $n\geq5$, the established spherical dimension of tagged singletons is 1, so Theorem 12 gives $\mathrm{SCO}\geq2$.  Their established VC dimension is one, hence their sign rank is at most 3 by the VC-one sign-rank theorem.  The sign-rank SCO reduction gives $\mathrm{SCO}\leq3$.
\end{proof}
\noindent\textbf{Reference.} \cite{chornomaz2025spherical,alon2016sign}.

\subsection{List Replicability Number of Trivial Concepts}\label{val:list_replicability_trivial_concepts}
\noindent\textbf{Statement.} The List Replicability Number of Trivial concepts is $1$.
\begin{proof}
On a realizable distribution with either constant target, every labeled example has that target's constant label.  From one example, output the corresponding constant hypothesis.  The singleton list containing that hypothesis has error zero and contains the learner output with probability one.  A nonempty list is necessary, so $\mathrm{LR}=1$.
\end{proof}
\noindent\textbf{Reference.} \cite{chornomaz2025spherical}.

\subsection{List Replicability Number of Singletons}\label{val:list_replicability_singletons}
\noindent\textbf{Statement.} The List Replicability Number of Singletons is $\begin{cases}2&n=3,\\2\leq\mathrm{LR}\leq3&n\geq4.\end{cases}$.
\begin{proof}
For $n\geq3$, the established spherical dimension is 1, and Theorem 10 gives $\mathrm{LR}\geq2$.  The sign rank is 2 at n=3 and 3 for $n\geq4$; the sign-rank upper bound on list replicability gives the corresponding upper bounds. Thus n=3 is exact and the displayed interval holds for larger n.
\end{proof}
\noindent\textbf{Reference.} \cite{chornomaz2025spherical,alon2016sign,blondal2026simplicialcovering}.

\subsection{List Replicability Number of Tagged Singletons}\label{val:list_replicability_tagged_singletons}
\noindent\textbf{Statement.} The List Replicability Number of Tagged Singletons is $2\leq\mathrm{LR}\leq3$.
\begin{proof}
For $n\geq5$, the established spherical dimension is 1, so Theorem 10 gives $\mathrm{LR}\geq2$.  Tagged singletons have VC dimension one, hence sign rank at most 3 by the VC-one sign-rank theorem. The sign-rank upper bound on list replicability gives $\mathrm{LR}\leq3$.
\end{proof}
\noindent\textbf{Reference.} \cite{chornomaz2025spherical,alon2016sign,blondal2026simplicialcovering}.

\subsection{Sign Rank of Universal Class}\label{val:sign_rank_universal_class}
\noindent\textbf{Statement.} The Sign Rank of Universal Class is $n$.
\begin{proof}
The incidence sign matrix of $\mathcal U_n$ is the transpose of that of the n-dimensional full cube. Sign rank is invariant under transposition, and the full cube has sign rank n. Hence $\mathrm{SR}(\mathcal U_n)=n$.
\end{proof}
\noindent\textbf{Reference.} \cite{alon2016sign}.

\subsection{Dual Sign Rank of Universal Class}\label{val:dual_sign_rank_universal_class}
\noindent\textbf{Statement.} The Dual Sign Rank of Universal Class is $\lfloor\log_2 n\rfloor\leq\mathrm{DSR}(\mathcal U_n)\leq2\lfloor\log_2 n\rfloor+1$.
\begin{proof}
Dual sign rank is not sign rank of the transposed matrix. Proposition 1 of Alon--Moran--Yehudayoff gives \(\mathrm{VC}(\mathcal H)\leq\mathrm{DSR}(\mathcal H)\leq2\mathrm{VC}(\mathcal H)+1\). Since \(\mathrm{VC}(\mathcal U_n)=\lfloor\log_2n\rfloor\), substitution gives the displayed bounds.
\end{proof}
\noindent\textbf{Reference.} \cite{alon2016sign}.

\subsection{Sign Rank of Singletons plus Empty Set}\label{val:sign_rank_singletons_plus_empty_set}
\noindent\textbf{Statement.} The Sign Rank of Singletons plus empty set is $3$.
\begin{proof}
The rows indexed by the singleton concepts form the n-by-n signed identity submatrix, whose sign rank is 3 for $n\geq4$. Thus the full incidence matrix has sign rank at least 3. This class has VC dimension one, and every VC-one sign matrix has sign rank at most 3. Hence the value is exactly 3.
\end{proof}
\noindent\textbf{Reference.} \cite{alon2016sign}.

\subsection{Dual Sign Rank of Singletons plus Empty Set}\label{val:dual_sign_rank_singletons_plus_empty_set}
\noindent\textbf{Statement.} The Dual Sign Rank of Singletons plus empty set is $3$.
\begin{proof}
After transposition, the columns indexed by singleton concepts still form an n-by-n signed identity submatrix, so the dual sign rank is at least 3 for $n\geq4$. The dual class has VC dimension one, and the VC-one sign-rank theorem gives the upper bound 3. Therefore the dual sign rank is exactly 3.
\end{proof}
\noindent\textbf{Reference.} \cite{alon2016sign}.

\subsection{Maximum Positive Degree of Universal Class}\label{val:maximum_positive_degree_universal_class}
\noindent\textbf{Statement.} The Maximum Positive Degree of Universal Class is $0$.
\begin{proof}
For distinct i,j, the universal hypotheses $u_i$ and $u_j$ disagree on exactly the $2^{n-1}$ subsets containing precisely one of i,j. Hence the 1-inclusion graph has no edges (for n at least 2), so every positive degree is zero.
\end{proof}
\noindent\textbf{Reference.} \cite{haussler1994predicting,chornomaz2025spherical}.

\subsection{Distinguishing Range of Universal Class}\label{val:distinguishing_range_universal_class}
\noindent\textbf{Statement.} The Distinguishing Range of Universal Class is $\lceil\log_2 n\rceil$.
\begin{proof}
Any set of r domain points yields only $2^r$ distinct restrictions, so r must be at least $\lceil\log_2 n\rceil$ to distinguish the n hypotheses. Conversely, encode i in $\lceil\log_2 n\rceil$ bits and, for each bit position, choose the domain point A consisting of precisely the indices whose bit is one. The restrictions of the $u_i$ to these points are their distinct codes.
\end{proof}
\noindent\textbf{Reference.} \cite{chornomaz2025spherical}.

\subsection{Teaching Dimension of Universal Class}\label{val:teaching_dimension_universal_class}
\noindent\textbf{Statement.} The Teaching Dimension of Universal Class is $1$.
\begin{proof}
For each target $u_i$, the single labeled example $(\{i\},1)$ excludes every $u_j$ with j different from i. Since the class has at least two hypotheses, no empty sample identifies a target. Thus every target has minimum teaching-set size one.
\end{proof}
\noindent\textbf{Reference.} \cite{shinohara1991teachability,goldman1995complexity}.

\subsection{Minimum Teaching Set Size of Universal Class}\label{val:minimum_teaching_set_size_universal_class}
\noindent\textbf{Statement.} The Minimum Teaching Set Size of Universal Class is $1$.
\begin{proof}
The one-point positive sample $(\{i\},1)$ teaches each $u_i$, because no other $u_j$ is positive at the domain point $\{i\}$. No target can be taught by the empty sample in a non-singleton class. Hence the minimum attained teaching-set size is one.
\end{proof}
\noindent\textbf{Reference.} \cite{shinohara1991teachability,goldman1995complexity}.

\subsection{Maximum Teaching Set Size of Universal Class}\label{val:maximum_teaching_set_size_universal_class}
\noindent\textbf{Statement.} The Maximum Teaching Set Size of Universal Class is $1$.
\begin{proof}
For every i, $(\{i\},1)$ is a teaching sample for $u_i$, while no empty sample distinguishes it from the other universal hypotheses. Thus every target has teaching-set size exactly one, and so does their maximum.
\end{proof}
\noindent\textbf{Reference.} \cite{shinohara1991teachability,goldman1995complexity}.

\subsection{Recursive Teaching Dimension of Universal Class}\label{val:recursive_teaching_dimension_universal_class}
\noindent\textbf{Statement.} The Recursive Teaching Dimension of Universal Class is $1$.
\begin{proof}
All hypotheses can be removed in the first recursive batch: each $u_i$ has the one-point teaching set $(\{i\},1)$ in the full class. A recursive teaching dimension of zero is impossible for a class with more than one hypothesis. Hence $\mathrm{RTD}=1$.
\end{proof}
\noindent\textbf{Reference.} \cite{zilles2011models}.

\subsection{Monotonic Minimum Teaching Set Size of Universal Class}\label{val:monotonic_minimum_teaching_set_size_universal_class}
\noindent\textbf{Statement.} The Monotonic Minimum Teaching Set Size of Universal Class is $1$.
\begin{proof}
For finite classes, monotonic minimum teaching-set size equals recursive teaching dimension. The universal class has recursive teaching dimension one: each $u_i$ is taught by the sample $(\{i\},1)$.
\end{proof}
\noindent\textbf{Reference.} \cite{doliwa2014recursive,zilles2011models}.

\subsection{Preference-Based Teaching Dimension of Universal Class}\label{val:preferencebased_teaching_dimension_universal_class}
\noindent\textbf{Statement.} The Preference-Based Teaching Dimension of Universal Class is $1$.
\begin{proof}
Preference-based teaching dimension equals recursive teaching dimension for finite classes. The universal class has recursive teaching dimension one because $(\{i\},1)$ identifies every target $u_i$.
\end{proof}
\noindent\textbf{Reference.} \cite{gao2017preference,zilles2011models}.

\subsection{Membership Query Complexity of Universal Class}\label{val:membership_query_complexity_universal_class}
\noindent\textbf{Statement.} The Membership Query Complexity of Universal Class is $\lceil\log_2 n\rceil$.
\begin{proof}
A binary membership-query decision tree of depth r has at most $2^r$ leaves, giving r at least $\lceil\log_2 n\rceil$. Conversely, at every stage choose a subset A of the surviving indices with sizes as balanced as possible and query A. The answer from target $u_i$ says exactly whether i lies in A, so this binary search identifies i in $\lceil\log_2 n\rceil$ queries.
\end{proof}
\noindent\textbf{Reference.} \cite{angluin1988queries}.

\subsection{Maximum Positive Degree of Tagged Singletons}\label{val:maximum_positive_degree_tagged_singletons}
\noindent\textbf{Statement.} The Maximum Positive Degree of Tagged Singletons is $1$.
\begin{proof}
The oriented 1-inclusion graph has only the edges from the empty concept to tag singletons and from a tag singleton to each tagged pair. Every nonempty vertex has exactly one smaller neighbour: the empty concept for a tag singleton, and its tag singleton for a tagged pair. Thus the largest positive degree is one.
\end{proof}
\noindent\textbf{Reference.} \cite{maass1992lower}.

\subsection{Degeneracy of Tagged Singletons}\label{val:degeneracy_tagged_singletons}
\noindent\textbf{Statement.} The Degeneracy of Tagged Singletons is $1$.
\begin{proof}
Its 1-inclusion graph is the tree consisting of the empty concept, the tag singletons, and the tagged pairs. Every nonempty tree has a leaf, so every nonempty induced subgraph has a vertex of degree at most one; an edge shows that degeneracy is not zero. Hence it is exactly one.
\end{proof}
\noindent\textbf{Reference.} \cite{maass1992lower}.

\subsection{Distinguishing Range of Tagged Singletons}\label{val:distinguishing_range_tagged_singletons}
\noindent\textbf{Statement.} The Distinguishing Range of Tagged Singletons is $n+\lceil\log_2 n\rceil$.
\begin{proof}
Every ordinary coordinate i is necessary: the tagged pair $\{i,n+f(i)\}$ and its tag singleton differ only at i. Every tag coordinate $n+\ell$ is necessary as well, because the tag singleton $\{n+\ell\}$ and the empty concept differ only there. Thus all $n+\lceil\log_2 n\rceil$ coordinates are required, and querying all of them plainly distinguishes the class.
\end{proof}
\noindent\textbf{Reference.} \cite{maass1992lower}.

\subsection{Maximum Projected Teaching Set Size of Tagged Singletons}\label{val:maximum_projected_teaching_set_size_tagged_singletons}
\noindent\textbf{Statement.} The Maximum Projected Teaching Set Size of Tagged Singletons is $n$.
\begin{proof}
Maximum projected teaching-set size equals the star number. Projecting tagged singletons to their n ordinary coordinates gives the empty concept and the n singleton leaves, an n-star. The established star-number computation gives the matching upper bound n.
\end{proof}
\noindent\textbf{Reference.} \cite{hanneke2015minimax,maass1992lower}.

\subsection{Balbach Teaching Dimension of Universal Class}\label{val:balbach_teaching_dimension_universal_class}
\noindent\textbf{Statement.} The Balbach Teaching Dimension of Universal Class is $1$.
\begin{proof}
The universal class has both ordinary teaching dimension and recursive teaching dimension equal to one. The established finite-class sandwich $\mathrm{TD}\geq\mathrm{BTD}\geq\mathrm{RTD}$ therefore forces $\mathrm{BTD}=1$.
\end{proof}
\noindent\textbf{Reference.} \cite{zilles2011models,balbach2008measuring}.

\subsection{Subset Teaching Dimension of Universal Class}\label{val:subset_teaching_dimension_universal_class}
\noindent\textbf{Statement.} The Subset Teaching Dimension of Universal Class is $1$.
\begin{proof}
For each target $u_i$, the one-point positive sample $(\{i\},1)$ is an ordinary minimum teaching sample and it is contained in no minimum teaching sample of another $u_j$: no other universal hypothesis is positive at $\{i\}$. It therefore remains a singleton subset-teaching sample at every iteration. An empty sample identifies no target, so $\mathrm{STD}=1$.
\end{proof}
\noindent\textbf{Reference.} \cite{zilles2011models}.

\subsection{Maximum Projected Teaching Set Size of Full Cube}\label{val:maximum_projected_teaching_set_size_full_cube}
\noindent\textbf{Statement.} The Maximum Projected Teaching Set Size of Full Cube is $n$.
\begin{proof}
Maximum projected teaching-set size equals star number. The full cube has an n-star around each vertex, so the value is n.
\end{proof}
\noindent\textbf{Reference.} \cite{hanneke2015minimax}.

\subsection{Maximum Projected Teaching Set Size of Singletons}\label{val:maximum_projected_teaching_set_size_singletons}
\noindent\textbf{Statement.} The Maximum Projected Teaching Set Size of Singletons is $n-1$.
\begin{proof}
Maximum projected teaching-set size equals star number. The established star-number value for singletons is n-1, obtained after projecting away one singleton coordinate.
\end{proof}
\noindent\textbf{Reference.} \cite{hanneke2015minimax}.

\subsection{Maximum Projected Teaching Set Size of Singletons plus Empty Set}\label{val:maximum_projected_teaching_set_size_singletons_plus_empty_set}
\noindent\textbf{Statement.} The Maximum Projected Teaching Set Size of Singletons plus empty set is $n$.
\begin{proof}
Maximum projected teaching-set size equals star number. The empty concept and the n singleton concepts form an n-star already in the original class.
\end{proof}
\noindent\textbf{Reference.} \cite{hanneke2015minimax}.

\subsection{Maximum Projected Teaching Set Size of Half-intervals}\label{val:maximum_projected_teaching_set_size_halfintervals}
\noindent\textbf{Statement.} The Maximum Projected Teaching Set Size of Half-intervals is $2$.
\begin{proof}
Maximum projected teaching-set size equals star number. The 1-inclusion graph of the half-interval chain is a path, whose internal vertices give a two-star and whose maximum degree is two.
\end{proof}
\noindent\textbf{Reference.} \cite{hanneke2015minimax}.

\subsection{Maximum Projected Teaching Set Size of Trivial Concepts}\label{val:maximum_projected_teaching_set_size_trivial_concepts}
\noindent\textbf{Statement.} The Maximum Projected Teaching Set Size of Trivial concepts is $1$.
\begin{proof}
Maximum projected teaching-set size equals star number. The two constant concepts form a one-star on any active coordinate, and no larger star is possible with only two hypotheses.
\end{proof}
\noindent\textbf{Reference.} \cite{hanneke2015minimax}.

\subsection{Maximum Projected Teaching Set Size of Majority}\label{val:maximum_projected_teaching_set_size_majority}
\noindent\textbf{Statement.} The Maximum Projected Teaching Set Size of Majority is $n$.
\begin{proof}
Maximum projected teaching-set size equals star number. The established majority star-number value is n, so the equality gives the same value here.
\end{proof}
\noindent\textbf{Reference.} \cite{hanneke2015minimax,maass1992lower}.

\subsection{Maximum Projected Teaching Set Size of Addressing}\label{val:maximum_projected_teaching_set_size_addressing}
\noindent\textbf{Statement.} The Maximum Projected Teaching Set Size of Addressing is $n-1$.
\begin{proof}
Maximum projected teaching-set size equals star number. The established addressing star-number computation is n-1, using the projection to all but one primary coordinate.
\end{proof}
\noindent\textbf{Reference.} \cite{hanneke2015minimax,maass1992lower}.

\subsection{Projected Maximal Degree of Majority}\label{val:projected_maximal_degree_majority}
\noindent\textbf{Statement.} The Projected Maximal Degree of Majority is $n$.
\begin{proof}
Projected maximal degree equals star number. The established majority star-number value is n, so the projected maximal degree is n as well.
\end{proof}
\noindent\textbf{Reference.} \cite{hanneke2015minimax,maass1992lower}.

\subsection{Projected Maximal Degree of Addressing}\label{val:projected_maximal_degree_addressing}
\noindent\textbf{Statement.} The Projected Maximal Degree of Addressing is $n-1$.
\begin{proof}
Projected maximal degree equals star number. The established addressing star-number computation gives n-1, using the projection to all but one primary coordinate.
\end{proof}
\noindent\textbf{Reference.} \cite{hanneke2015minimax,maass1992lower}.

\subsection{Projected Maximal Degree of Tagged Singletons}\label{val:projected_maximal_degree_tagged_singletons}
\noindent\textbf{Statement.} The Projected Maximal Degree of Tagged Singletons is $n$.
\begin{proof}
Projected maximal degree equals star number. Projecting to the n ordinary coordinates gives the empty trace with all n singleton traces as neighbours, so the established star-number value n is attained.
\end{proof}
\noindent\textbf{Reference.} \cite{hanneke2015minimax,maass1992lower}.

\subsection{Projected Maximal Degree of Half-intervals}\label{val:projected_maximal_degree_halfintervals}
\noindent\textbf{Statement.} The Projected Maximal Degree of Half-intervals is $2$.
\begin{proof}
Projected maximal degree equals star number. The original 1-inclusion graph is a path and an internal vertex has degree two; the established star-number upper bound is two.
\end{proof}
\noindent\textbf{Reference.} \cite{hanneke2015minimax}.

\subsection{Projected VC Radius of Tagged Singletons}\label{val:projected_vc_radius_tagged_singletons}
\noindent\textbf{Statement.} The Projected VC Radius of Tagged Singletons is $1$.
\begin{proof}
Projected VC radius equals VC dimension. Tagged singletons have established VC dimension one, hence projected VC radius one.
\end{proof}
\noindent\textbf{Reference.} \cite{maass1992lower}.

\subsection{Projected VC Radius of Full Cube}\label{val:projected_vc_radius_full_cube}
\noindent\textbf{Statement.} The Projected VC Radius of Full Cube is $n$.
\begin{proof}
Projected VC radius equals VC dimension. The full cube shatters its n-coordinate domain, so the value is n.
\end{proof}

\subsection{Projected VC Radius of Singletons plus Empty Set}\label{val:projected_vc_radius_singletons_plus_empty_set}
\noindent\textbf{Statement.} The Projected VC Radius of Singletons plus empty set is $1$.
\begin{proof}
Projected VC radius equals VC dimension. The established VC dimension of singletons plus the empty concept is one.
\end{proof}

\subsection{Projected VC Radius of Singletons}\label{val:projected_vc_radius_singletons}
\noindent\textbf{Statement.} The Projected VC Radius of Singletons is $1$.
\begin{proof}
Projected VC radius equals VC dimension. The singleton class has established VC dimension one.
\end{proof}

\subsection{Projected VC Radius of Addressing}\label{val:projected_vc_radius_addressing}
\noindent\textbf{Statement.} The Projected VC Radius of Addressing is $\lfloor\log n\rfloor\quad(n=2^r,\ r\ge1)$.
\begin{proof}
Restrict here to $n=2^r$ with all r-bit words present, so $\lfloor\log_2n\rfloor=\log_2n=r$. Projected VC radius equals VC dimension. The address coordinates give the established value $\lfloor\log n\rfloor$.
\end{proof}
\noindent\textbf{Reference.} \cite{maass1992lower}.

\subsection{Projected VC Radius of Blockwise-Addressing}\label{val:projected_vc_radius_blockwiseaddressing}
\noindent\textbf{Statement.} The Projected VC Radius of Blockwise-Addressing is $d$.
\begin{proof}
Projected VC radius equals VC dimension. The established blockwise-addressing VC dimension is d.
\end{proof}
\noindent\textbf{Reference.} \cite{ben1997online}.

\subsection{Projected VC Radius of Half-intervals}\label{val:projected_vc_radius_halfintervals}
\noindent\textbf{Statement.} The Projected VC Radius of Half-intervals is $1$.
\begin{proof}
Projected VC radius equals VC dimension. Every nontrivial half-interval chain has VC dimension one.
\end{proof}

\subsection{Projected VC Radius of Majority}\label{val:projected_vc_radius_majority}
\noindent\textbf{Statement.} The Projected VC Radius of Majority is $n$.
\begin{proof}
Projected VC radius equals VC dimension. The established majority VC dimension is n.
\end{proof}
\noindent\textbf{Reference.} \cite{maass1992lower}.

\subsection{Projected VC Radius of Warmuth's $C_5$ Class}\label{val:projected_vc_radius_warmuth_c5}
\noindent\textbf{Statement.} The Projected VC Radius of Warmuth's $C_5$ class is $2$.
\begin{proof}
Projected VC radius equals VC dimension. Warmuth's $C_5$ class has established VC dimension two.
\end{proof}
\noindent\textbf{Reference.} \cite{palvolgyi2021unlabeled}.

\subsection{Projected Strongly Shattered Dimension of Tagged Singletons}\label{val:projected_strongly_shattered_dimension_tagged_singletons}
\noindent\textbf{Statement.} The Projected Strongly Shattered Dimension of Tagged Singletons is $1$.
\begin{proof}
Projected strongly shattered dimension equals VC dimension. Tagged singletons have established VC dimension one.
\end{proof}
\noindent\textbf{Reference.} \cite{bollobas1995defect,maass1992lower}.

\subsection{Projected Strongly Shattered Dimension of Full Cube}\label{val:projected_strongly_shattered_dimension_full_cube}
\noindent\textbf{Statement.} The Projected Strongly Shattered Dimension of Full Cube is $n$.
\begin{proof}
Projected strongly shattered dimension equals VC dimension. The full cube has VC dimension n.
\end{proof}
\noindent\textbf{Reference.} \cite{bollobas1995defect}.

\subsection{Projected Strongly Shattered Dimension of Singletons plus Empty Set}\label{val:projected_strongly_shattered_dimension_singletons_plus_empty_set}
\noindent\textbf{Statement.} The Projected Strongly Shattered Dimension of Singletons plus empty set is $1$.
\begin{proof}
Projected strongly shattered dimension equals VC dimension. Singletons plus the empty concept have established VC dimension one.
\end{proof}
\noindent\textbf{Reference.} \cite{bollobas1995defect}.

\subsection{Projected Strongly Shattered Dimension of Trivial Concepts}\label{val:projected_strongly_shattered_dimension_trivial_concepts}
\noindent\textbf{Statement.} The Projected Strongly Shattered Dimension of Trivial concepts is $1$.
\begin{proof}
Projected strongly shattered dimension equals VC dimension. The two trivial concepts have established VC dimension one.
\end{proof}
\noindent\textbf{Reference.} \cite{bollobas1995defect}.

\subsection{Projected Strongly Shattered Dimension of Singletons}\label{val:projected_strongly_shattered_dimension_singletons}
\noindent\textbf{Statement.} The Projected Strongly Shattered Dimension of Singletons is $1$.
\begin{proof}
Projected strongly shattered dimension equals VC dimension. The singleton class has established VC dimension one.
\end{proof}
\noindent\textbf{Reference.} \cite{bollobas1995defect}.

\subsection{Projected Strongly Shattered Dimension of Addressing}\label{val:projected_strongly_shattered_dimension_addressing}
\noindent\textbf{Statement.} The Projected Strongly Shattered Dimension of Addressing is $\lfloor\log n\rfloor\quad(n=2^r,\ r\ge1)$.
\begin{proof}
Restrict here to $n=2^r$ with all r-bit words present, so $\lfloor\log_2n\rfloor=\log_2n=r$. Projected strongly shattered dimension equals VC dimension. Addressing has established VC dimension $\lfloor\log n\rfloor$.
\end{proof}
\noindent\textbf{Reference.} \cite{bollobas1995defect,maass1992lower}.

\subsection{Projected Strongly Shattered Dimension of Blockwise-Addressing}\label{val:projected_strongly_shattered_dimension_blockwiseaddressing}
\noindent\textbf{Statement.} The Projected Strongly Shattered Dimension of Blockwise-Addressing is $d$.
\begin{proof}
Projected strongly shattered dimension equals VC dimension. The established blockwise-addressing VC dimension is d.
\end{proof}
\noindent\textbf{Reference.} \cite{bollobas1995defect,ben1997online}.

\subsection{Projected Strongly Shattered Dimension of Half-intervals}\label{val:projected_strongly_shattered_dimension_halfintervals}
\noindent\textbf{Statement.} The Projected Strongly Shattered Dimension of Half-intervals is $1$.
\begin{proof}
Projected strongly shattered dimension equals VC dimension. Every nontrivial half-interval chain has VC dimension one.
\end{proof}
\noindent\textbf{Reference.} \cite{bollobas1995defect}.

\subsection{Projected Strongly Shattered Dimension of Majority}\label{val:projected_strongly_shattered_dimension_majority}
\noindent\textbf{Statement.} The Projected Strongly Shattered Dimension of Majority is $n$.
\begin{proof}
Projected strongly shattered dimension equals VC dimension. The established majority VC dimension is n.
\end{proof}
\noindent\textbf{Reference.} \cite{bollobas1995defect,maass1992lower}.

\subsection{Projected Strongly Shattered Dimension of Warmuth's $C_5$ Class}\label{val:projected_strongly_shattered_dimension_warmuth_c5}
\noindent\textbf{Statement.} The Projected Strongly Shattered Dimension of Warmuth's $C_5$ class is $2$.
\begin{proof}
Projected strongly shattered dimension equals VC dimension. Warmuth's $C_5$ class has established VC dimension two.
\end{proof}
\noindent\textbf{Reference.} \cite{bollobas1995defect,palvolgyi2021unlabeled}.

\subsection{Monotonic Minimum Teaching Set Size of Addressing}\label{val:monotonic_minimum_teaching_set_size_addressing}
\noindent\textbf{Statement.} The Monotonic Minimum Teaching Set Size of Addressing is $1\quad(n\ge2)$.
\begin{proof}
For finite classes, monotonic minimum teaching-set size equals recursive teaching dimension. Addressing has established recursive teaching dimension one.
\end{proof}
\noindent\textbf{Reference.} \cite{doliwa2014recursive,maass1992lower}.

\subsection{Monotonic Minimum Teaching Set Size of Half-intervals}\label{val:monotonic_minimum_teaching_set_size_halfintervals}
\noindent\textbf{Statement.} The Monotonic Minimum Teaching Set Size of Half-intervals is $1$.
\begin{proof}
For finite classes, monotonic minimum teaching-set size equals recursive teaching dimension. The half-interval chain has established recursive teaching dimension one.
\end{proof}
\noindent\textbf{Reference.} \cite{doliwa2014recursive}.

\subsection{Monotonic Minimum Teaching Set Size of Singletons plus Empty Set}\label{val:monotonic_minimum_teaching_set_size_singletons_plus_empty_set}
\noindent\textbf{Statement.} The Monotonic Minimum Teaching Set Size of Singletons plus empty set is $1$.
\begin{proof}
For finite classes, monotonic minimum teaching-set size equals recursive teaching dimension. Singletons plus the empty concept have established recursive teaching dimension one.
\end{proof}
\noindent\textbf{Reference.} \cite{doliwa2014recursive}.

\subsection{Monotonic Minimum Teaching Set Size of Trivial Concepts}\label{val:monotonic_minimum_teaching_set_size_trivial_concepts}
\noindent\textbf{Statement.} The Monotonic Minimum Teaching Set Size of Trivial concepts is $1$.
\begin{proof}
For finite classes, monotonic minimum teaching-set size equals recursive teaching dimension. The two trivial concepts have established recursive teaching dimension one.
\end{proof}
\noindent\textbf{Reference.} \cite{doliwa2014recursive}.

\subsection{Monotonic Minimum Teaching Set Size of Tagged Singletons}\label{val:monotonic_minimum_teaching_set_size_tagged_singletons}
\noindent\textbf{Statement.} The Monotonic Minimum Teaching Set Size of Tagged Singletons is $1$.
\begin{proof}
For finite classes, monotonic minimum teaching-set size equals recursive teaching dimension. Tagged singletons have established recursive teaching dimension one.
\end{proof}
\noindent\textbf{Reference.} \cite{doliwa2014recursive,maass1992lower}.

\subsection{Monotonic Minimum Teaching Set Size of Full Cube}\label{val:monotonic_minimum_teaching_set_size_full_cube}
\noindent\textbf{Statement.} The Monotonic Minimum Teaching Set Size of Full Cube is $n$.
\begin{proof}
For finite classes, monotonic minimum teaching-set size equals recursive teaching dimension. The full cube has established recursive teaching dimension n.
\end{proof}
\noindent\textbf{Reference.} \cite{doliwa2014recursive}.

\subsection{Monotonic Minimum Teaching Set Size of Singletons}\label{val:monotonic_minimum_teaching_set_size_singletons}
\noindent\textbf{Statement.} The Monotonic Minimum Teaching Set Size of Singletons is $1$.
\begin{proof}
For finite classes, monotonic minimum teaching-set size equals recursive teaching dimension. The singleton class has established recursive teaching dimension one.
\end{proof}
\noindent\textbf{Reference.} \cite{doliwa2014recursive}.

\subsection{Littlestone Dimension of Universal Class}\label{val:littlestone_dimension_universal_class}
\noindent\textbf{Statement.} The Littlestone Dimension of Universal Class is $\lfloor\log_2 n\rfloor$.
\begin{proof}
Let d be $\lfloor\log_2 n\rfloor$ and select $2^d$ indices. Assign them distinct d-bit strings; at level t query the domain point consisting of the selected indices whose tth bit is one. Every root-to-leaf label sequence is realized by its coded $u_i$, so $\mathrm L\geq d$. A depth-d mistake tree has $2^d$ distinct realized leaves, whereas the class has n hypotheses, proving the matching upper bound.
\end{proof}
\noindent\textbf{Reference.} \cite{littlestone1988learning}.

\subsection{Equivalence Queries Complexity of Universal Class}\label{val:equivalence_queries_complexity_universal_class}
\noindent\textbf{Statement.} The Equivalence Queries Complexity of Universal Class is $\lfloor\log_2 n\rfloor$.
\begin{proof}
For finite classes in the unrestricted equivalence-query model, equivalence-query complexity equals Littlestone dimension. The universal class has the established value $\lfloor\log_2 n\rfloor$ for the latter.
\end{proof}
\noindent\textbf{Reference.} \cite{littlestone1988learning}.

\subsection{Relative Hitting Size of Universal Class}\label{val:relative_hitting_size_universal_class}
\noindent\textbf{Statement.} The Relative Hitting Size of Universal Class is $1$.
\begin{proof}
For finite concept classes, relative hitting size equals teaching dimension. The universal class has teaching dimension one: $(\{i\},1)$ uniquely identifies each target $u_i$.
\end{proof}
\noindent\textbf{Reference.} \cite{goldman1995complexity,shinohara1991teachability}.

\subsection{Littlestone Dimension of Blockwise-Addressing}\label{val:littlestone_dimension_blockwiseaddressing}
\noindent\textbf{Statement.} The Littlestone Dimension of Blockwise-Addressing is $d+1$.
\begin{proof}
Query z at the root. In the 0-subtree, the $f_i$ realize every d-bit pattern on any y-block, and in the 1-subtree the $g_i$ realize every d-bit pattern on the corresponding x-block. Thus both subtrees support depth d and the class shatters a tree of depth d+1. It has exactly $2^{d+1}$ concepts, so a deeper complete tree would require more distinct leaf realizers.
\end{proof}
\noindent\textbf{Reference.} \cite{ben1997online,littlestone1988learning}.

\subsection{Equivalence Queries Complexity of Blockwise-Addressing}\label{val:equivalence_queries_complexity_blockwiseaddressing}
\noindent\textbf{Statement.} The Equivalence Queries Complexity of Blockwise-Addressing is $d+1$.
\begin{proof}
For finite classes in the unrestricted equivalence-query model, equivalence-query complexity equals Littlestone dimension. Blockwise-addressing has the established Littlestone value d+1.
\end{proof}
\noindent\textbf{Reference.} \cite{ben1997online,littlestone1988learning}.

\subsection{co-VC dimension of Universal Class}\label{val:covc_dimension_universal_class}
\noindent\textbf{Statement.} The co-VC dimension of Universal Class is $n$.
\begin{proof}
Let the $n$ concepts be $u_1$,...,$u_n$. On $n$ domain points, take $A_j$ to contain every i except j. The induced traces are precisely the $n$ one-bit flips of the missing all-one labeling, giving a hollow n-star. Conversely a hollow d-star requires its d distinct one-bit neighbours to be realized by distinct hypotheses, and the universal class has only n hypotheses. Hence coVC($U_n$)=n.
\end{proof}
\noindent\textbf{Reference.} Direct calculation from the definition; see also the hollow-star formulation in \cite{bousquet2020proper}..

\subsection{co-VC dimension of Majority}\label{val:covc_dimension_majority}
\noindent\textbf{Statement.} The co-VC dimension of Majority is $\lfloor n/2\rfloor+2$.
\begin{proof}
Write a Majority concept as its $n$ base bits together with the majority bit.  On any $t=\lfloor n/2\rfloor+1$ base coordinates and the majority coordinate, the all-one base trace with majority bit 0 is absent, whereas flipping the majority bit or any base bit gives a realized trace; hence there is a hollow $(t+1)$-star.  Conversely, a projection omitting the majority coordinate is a full cube.  In a projection using it, suppose first that the missing majority bit is 0. Every selected base bit must be 1 (otherwise its flip cannot create a 0-majority extension), and the two majority inequalities before and after a flip force $t=\lfloor n/2\rfloor+1$.  If the missing bit is 1, all selected base bits must instead be 0; the same inequalities give $t=\lceil n/2\rceil$, which is no larger.  Thus no hollow star has dimension exceeding $\lfloor n/2\rfloor+2$.
\end{proof}
\noindent\textbf{Reference.} Direct calculation from the definition; the Majority class is from \cite{maass1992lower}..

\subsection{co-VC dimension of Warmuth's $C_5$ Class}\label{val:covc_dimension_warmuth_c5}
\noindent\textbf{Statement.} The co-VC dimension of Warmuth's $C_5$ class is $4$.
\begin{proof}
On four consecutive cyclic coordinates, the missing trace $0110$ has all four one-bit neighbours among the ten $C_5$ concepts, so it is a hollow 4-star.  The checked finite enumerator constructs the class directly from its cyclic $1001$ definition and checks every coordinate subset and missing center; it finds no hollow 5-star.  Hence $\mathrm{coVC}(C_5)=4$.
\end{proof}
\noindent\textbf{Reference.} Direct finite calculation from the Warmuth class; see also \cite{palvolgyi2021unlabeled,bousquet2020proper}..

\subsection{co-VC dimension of Blockwise-Addressing}\label{val:covc_dimension_blockwiseaddressing}
\noindent\textbf{Statement.} The co-VC dimension of Blockwise-Addressing is $2^d+1\leq\mathrm{coVC}\leq2^{d+1}$.
\begin{proof}
Put $m=2^d$.  Project to the selector z and all m x-coordinates.  The all-zero trace is absent.  Flipping $x_i$ is realized by $f_i$, while flipping z is realized by the $g_i$ whose d-bit x-block word is all zero.  This is a hollow $(m+1)$-star.  Conversely the class has exactly 2m concepts, and every hollow r-star needs r distinct hypotheses for its realized neighbours; hence $\mathrm{coVC}\leq2m$.
\end{proof}
\noindent\textbf{Reference.} Direct calculation from the Blockwise-Addressing construction in \cite{ben1997online}; the hollow-star formulation is in \cite{bousquet2020proper}..

\subsection{2-Block Littlestone Dimension of Tagged Singletons}\label{val:2block_littlestone_dimension_tagged_singletons}
\noindent\textbf{Statement.} The 2-Block Littlestone Dimension of Tagged Singletons is $1$.
\begin{proof}
VC dimension is one, which gives a depth-one level-constant tree. A 2-block tree of depth two has one common coordinate at its second level, so its four leaves realize all labelings on two fixed coordinates. This would give VC dimension at least two, which is impossible.
\end{proof}
\noindent\textbf{Reference.} Directly from the block-tree definition and the VC calculation for this class in \cite{maass1992lower}..

\subsection{$k$-Block Littlestone Dimension of Tagged Singletons}\label{val:kblock_littlestone_dimension_tagged_singletons}
\noindent\textbf{Statement.} The $k$-Block Littlestone Dimension ($k>2$) of Tagged Singletons is $1$.
\begin{proof}
The depth-one VC witness gives the lower bound. For every fixed $k>2$, the first two levels of a depth-two k-block tree use fixed coordinates and would shatter a two-point set, contradicting VC dimension one.
\end{proof}
\noindent\textbf{Reference.} Directly from the block-tree definition and the VC calculation for this class in \cite{maass1992lower}..

\subsection{Prefix-$k$ Littlestone Dimension of Tagged Singletons}\label{val:prefix_littlestone_dimension_tagged_singletons}
\noindent\textbf{Statement.} The Prefix-$k$ Littlestone Dimension ($k>1$) of Tagged Singletons is $1$.
\begin{proof}
For every fixed $k>1$, a depth-two prefix-k tree is level-constant on its first two levels and hence would shatter two fixed coordinates. Since VC dimension is one, its depth is at most one; the VC witness realizes that depth.
\end{proof}
\noindent\textbf{Reference.} Directly from the prefix-tree definition and the VC/Littlestone calculations for this class in \cite{maass1992lower}..

\subsection{2-Block Littlestone Dimension of Universal Class}\label{val:2block_littlestone_dimension_universal_class}
\noindent\textbf{Statement.} The 2-Block Littlestone Dimension of Universal Class is $\lfloor\log_2 n\rfloor$.
\begin{proof}
The standard coded Littlestone witness labels every level by one domain point, so it is in particular a 2-block tree of depth $\lfloor\log_2 n\rfloor$. The ordinary Littlestone dimension gives the matching upper bound.
\end{proof}
\noindent\textbf{Reference.} Direct calculation from the coded universal-class witness and \cite{littlestone1988learning}..

\subsection{$k$-Block Littlestone Dimension of Universal Class}\label{val:kblock_littlestone_dimension_universal_class}
\noindent\textbf{Statement.} The $k$-Block Littlestone Dimension ($k>2$) of Universal Class is $\lfloor\log_2 n\rfloor$.
\begin{proof}
The standard coded Littlestone witness is horizontally constant on every level, hence is a k-block tree for every k. The ordinary Littlestone dimension gives the matching upper bound.
\end{proof}
\noindent\textbf{Reference.} Direct calculation from the coded universal-class witness and \cite{littlestone1988learning}..

\subsection{Prefix-$k$ Littlestone Dimension of Universal Class}\label{val:prefix_littlestone_dimension_universal_class}
\noindent\textbf{Statement.} The Prefix-$k$ Littlestone Dimension ($k>1$) of Universal Class is $\lfloor\log_2 n\rfloor$.
\begin{proof}
The standard coded Littlestone witness is horizontally constant on every level, so it satisfies every prefix constraint. The ordinary Littlestone dimension gives the matching upper bound.
\end{proof}
\noindent\textbf{Reference.} Direct calculation from the coded universal-class witness and \cite{littlestone1988learning}..

\subsection{2-Block Littlestone Dimension of Warmuth's $C_5$ Class}\label{val:2block_littlestone_dimension_warmuth_c5}
\noindent\textbf{Statement.} The 2-Block Littlestone Dimension of Warmuth's $C_5$ class is $3$.
\begin{proof}
The depth-three Littlestone witness queries coordinates 0 and 1 on its first two levels, then coordinates 3,2,2,2 on the four branches. Its first block therefore satisfies the 2-block constraint, proving the lower bound. The ordinary Littlestone dimension is three, proving equality.
\end{proof}
\noindent\textbf{Reference.} Directly from the $C_5$ Littlestone witness; see \cite{littlestone1988learning,palvolgyi2021unlabeled}..

\subsection{$k$-Block Littlestone Dimension of Warmuth's $C_5$ Class}\label{val:kblock_littlestone_dimension_warmuth_c5}
\noindent\textbf{Statement.} The $k$-Block Littlestone Dimension ($k>2$) of Warmuth's $C_5$ class is $2$.
\begin{proof}
For every fixed $k>2$, a depth-three k-block tree is level-constant on its first three levels and would shatter three coordinates, while VC dimension is two. The VC witness supplies depth two.
\end{proof}
\noindent\textbf{Reference.} Directly from the block-tree definition and the $C_5$ calculations in \cite{palvolgyi2021unlabeled}..

\subsection{Prefix-$k$ Littlestone Dimension of Warmuth's $C_5$ Class}\label{val:prefix_littlestone_dimension_warmuth_c5}
\noindent\textbf{Statement.} The Prefix-$k$ Littlestone Dimension ($k>1$) of Warmuth's $C_5$ class is $\begin{cases}3&k=2,\\2&k\ge3.\end{cases}$.
\begin{proof}
For k=2, the recorded depth-three Littlestone witness has a common coordinate at its second level and is a prefix-2 tree. For k at least three, depth three would be level-constant on its first three levels and would shatter three coordinates, contradicting VC dimension two. The VC witness gives depth two.
\end{proof}
\noindent\textbf{Reference.} Directly from the prefix-tree definition and the $C_5$ calculations in \cite{littlestone1988learning,palvolgyi2021unlabeled}..

\subsection{2-Block Littlestone Dimension of Half-intervals}\label{val:2block_littlestone_dimension_halfintervals}
\noindent\textbf{Statement.} The 2-Block Littlestone Dimension of Half-intervals is $1$.
\begin{proof}
VC dimension one gives a depth-one level-constant tree. A depth-two 2-block tree has a fixed coordinate at its second level and would realize all four labelings on two fixed coordinates, contradicting VC dimension one.
\end{proof}
\noindent\textbf{Reference.} Directly from the block-tree definition and the threshold/VC calculation in \cite{littlestone1988learning}..

\subsection{$k$-Block Littlestone Dimension of Half-intervals}\label{val:kblock_littlestone_dimension_halfintervals}
\noindent\textbf{Statement.} The $k$-Block Littlestone Dimension ($k>2$) of Half-intervals is $1$.
\begin{proof}
The VC witness gives depth one. For every fixed $k>2$, the first two levels of a depth-two k-block tree would be a fixed two-coordinate cube, contradicting VC dimension one.
\end{proof}
\noindent\textbf{Reference.} Directly from the block-tree definition and the threshold/VC calculation in \cite{littlestone1988learning}..

\subsection{Prefix-$k$ Littlestone Dimension of Half-intervals}\label{val:prefix_littlestone_dimension_halfintervals}
\noindent\textbf{Statement.} The Prefix-$k$ Littlestone Dimension ($k>1$) of Half-intervals is $1$.
\begin{proof}
For every fixed $k>1$, depth two would be level-constant on its first two levels and would shatter two coordinates. Since VC dimension is one, the value is at most one; a one-point VC witness gives equality.
\end{proof}
\noindent\textbf{Reference.} Directly from the prefix-tree definition and the threshold calculations in \cite{littlestone1988learning}..

\subsection{Randomized Littlestone Dimension of Full Cube}\label{val:randomized_littlestone_dimension_full_cube}
\noindent\textbf{Statement.} The Randomized Littlestone Dimension of Full Cube is $\frac n2$.
\begin{proof}
The cube has a complete shattered tree of depth n, so its randomized value is at least n/2. Conversely, any shattered tree has at most $2^n$ leaves, and a random branch reaches a leaf with probability $2^{-d}$ at each leaf of depth d; its expected length is the entropy of this dyadic leaf distribution and is at most $\log_2(2^n)=n$. Thus the value is exactly n/2.
\end{proof}
\noindent\textbf{Reference.} \cite{littlestone1988learning,filmus2023randomized}.

\subsection{Randomized Littlestone Dimension of Trivial Concepts}\label{val:randomized_littlestone_dimension_trivial_concepts}
\noindent\textbf{Statement.} The Randomized Littlestone Dimension of Trivial concepts is $1/2$.
\begin{proof}
One query separates the two constant hypotheses, producing a shattered tree whose two branches both have length one and hence randomized value one half. No branch can be extended after that separation, so every shattered tree has average branch length at most one.
\end{proof}
\noindent\textbf{Reference.} \cite{filmus2023randomized}.

\subsection{Randomized Littlestone Dimension of Singletons}\label{val:randomized_littlestone_dimension_singletons}
\noindent\textbf{Statement.} The Randomized Littlestone Dimension of Singletons is $1-2^{1-n}$.
\begin{proof}
At each nontrivial node, querying a coordinate separates exactly one remaining singleton from all the others. Querying n-1 distinct coordinates consecutively gives branch-length expectation $\sum_{j=0}^{n-2}2^{-j}=2-2^{2-n}$. Conversely every shattered tree has this comb form with at most n-1 internal nodes on its continuing branch. Dividing by two gives $1-2^{1-n}$.
\end{proof}
\noindent\textbf{Reference.} \cite{filmus2023randomized}.

\subsection{Randomized Littlestone Dimension of Singletons plus Empty Set}\label{val:randomized_littlestone_dimension_singletons_plus_empty_set}
\noindent\textbf{Statement.} The Randomized Littlestone Dimension of Singletons plus empty set is $1-2^{-n}$.
\begin{proof}
Each nontrivial query separates one singleton on its positive branch, leaving the empty concept and the other singletons on the negative branch. Querying all n coordinates therefore gives a comb with expected branch length $\sum_{j=0}^{n-1}2^{-j}=2-2^{1-n}$. Every shattered tree has at most these n successive splits, giving the matching upper bound. Divide by two.
\end{proof}
\noindent\textbf{Reference.} \cite{filmus2023randomized}.

\subsection{Randomized Littlestone Dimension of Half-intervals}\label{val:randomized_littlestone_dimension_halfintervals}
\noindent\textbf{Statement.} The Randomized Littlestone Dimension of Half-intervals is $\frac12\bigl(k+\frac{n-2^k}{2^k}\bigr)$.
\begin{proof}
A threshold query partitions the feasible initial segments at any chosen cut, so shattered trees for this class are precisely binary comparison trees with n leaves. If $F(n)$ is the largest expected random-branch depth among such trees, then $F(1)=0$ and $F(n)=1+\max_{1\leq a<n}(F(a)+F(n-a))/2$. Induction gives $F(n)=k+(n-2^k)/2^k$ for $k=\lfloor\log_2 n\rfloor$. Dividing by two gives the stated randomized Littlestone value.
\end{proof}
\noindent\textbf{Reference.} \cite{filmus2023randomized,littlestone1988learning}.

\subsection{Randomized Littlestone Dimension of Universal Class}\label{val:randomized_littlestone_dimension_universal_class}
\noindent\textbf{Statement.} The Randomized Littlestone Dimension of Universal Class is $\frac12\bigl(k+\frac{n-2^k}{2^k}\bigr)$.
\begin{proof}
Every subset of the n hypotheses can be selected by a domain point of the universal class. Hence every binary tree with n leaves is shattered, and no shattered tree can have more than n leaves. The optimal expected random-branch depth is therefore F(n)=k+(n-$2^k$)/$2^k$, by the binary-tree recurrence; divide by two.
\end{proof}
\noindent\textbf{Reference.} \cite{filmus2023randomized,chornomaz2025spherical}.

\subsection{Randomized Littlestone Dimension of Addressing}\label{val:randomized_littlestone_dimension_addressing}
\noindent\textbf{Statement.} The Randomized Littlestone Dimension of Addressing is $\frac12\lfloor\log_2 n\rfloor\quad(n=2^r,\ r\ge1)$.
\begin{proof}
Restrict here to $n=2^r$ with all r-bit words present, so $\lfloor\log_2n\rfloor=\log_2n=r$. Write d for $\lfloor\log_2 n\rfloor$. The class has exactly $2^d$ concepts, so every shattered tree has expected random-branch depth at most d by the entropy bound. Querying the d address coordinates gives a complete shattered tree of depth d, hence randomized value d/2.
\end{proof}
\noindent\textbf{Reference.} \cite{filmus2023randomized,maass1992lower}.

\subsection{Randomized Littlestone Dimension of Majority}\label{val:randomized_littlestone_dimension_majority}
\noindent\textbf{Statement.} The Randomized Littlestone Dimension of Majority is $\frac n2$.
\begin{proof}
The class has $2^n$ concepts, so the entropy bound gives expected random-branch depth at most n. Its n freely varying ordinary coordinates give a complete shattered tree of depth n. Thus randomized Littlestone dimension is n/2.
\end{proof}
\noindent\textbf{Reference.} \cite{filmus2023randomized,maass1992lower}.

\subsection{Randomized Littlestone Dimension of Blockwise-Addressing}\label{val:randomized_littlestone_dimension_blockwiseaddressing}
\noindent\textbf{Statement.} The Randomized Littlestone Dimension of Blockwise-Addressing is $\frac{d+1}{2}$.
\begin{proof}
The class has exactly $2^{d+1}$ concepts, so no shattered tree has expected random-branch depth exceeding d+1. The established complete Littlestone tree of depth d+1 has every branch that long, giving randomized value (d+1)/2.
\end{proof}
\noindent\textbf{Reference.} \cite{filmus2023randomized,ben1997online}.

\subsection{Randomized Littlestone Dimension of Warmuth's $C_5$ Class}\label{val:randomized_littlestone_dimension_warmuth_c5}
\noindent\textbf{Statement.} The Randomized Littlestone Dimension of Warmuth's $C_5$ class is $13/8$.
\begin{proof}
Enumerate the ten cyclic $1001$ concepts and apply the finite randomized-tree recurrence $R(V)=\max_x(1+R(V_{x=0})+R(V_{x=1}))/2$ to every version space V. The exact dynamic program gives $R(C_5)=13/8$; the repository's C5 verifier checks this with rational arithmetic.
\end{proof}
\noindent\textbf{Reference.} \cite{filmus2023randomized,palvolgyi2021unlabeled}.

\subsection{Occasional One-Sided Getback of Full Cube}\label{val:occasional_onesided_getback_full_cube}
\noindent\textbf{Statement.} The Occasional One-Sided Getback of Full Cube is $n$.
\begin{proof}
The established O1G/compression lemma gives $\mathrm{M}_{\mathrm{o1g}}\geq\mathrm{USC}$, and the full cube has established unlabeled-compression value n. The same lemma gives $\mathrm{M}_{\mathrm{o1g}}\leq\mathrm L$, while the full cube has Littlestone dimension n. Thus both bounds coincide at n.
\end{proof}
\noindent\textbf{Reference.} \cite{ben1997online,chalopin2022unlabeled,palvolgyi2021unlabeled}.

\subsection{Random Equivalence Queries Complexity of Full Cube}\label{val:random_equivalence_queries_complexity_full_cube}
\noindent\textbf{Statement.} The Random Equivalence Queries Complexity of Full Cube is $n$.
\begin{proof}
The full cube has VC dimension $n$, and the random-counterexample lower bound gives $\mathrm Q_{\mathrm{req}}\geq\mathrm{VC}$. Conversely, every equivalence-query strategy is valid against random counterexamples, so $\mathrm Q_{\mathrm{req}}\leq\mathrm L$; the full cube has Littlestone dimension $n$. Thus both bounds equal $n$.
\end{proof}
\noindent\textbf{Reference.} \cite{angluin2017random,littlestone1988learning}.

\subsection{Random Proper Equivalence Queries Complexity of Full Cube}\label{val:random_proper_equivalence_queries_complexity_full_cube}
\noindent\textbf{Statement.} The Random Proper Equivalence Queries Complexity of Full Cube is $n$.
\begin{proof}
Proper random equivalence queries are a restriction of random equivalence queries, so $\mathrm Q_{\mathrm{rpeq}}\geq\mathrm Q_{\mathrm{req}}=n$ on the full cube. The established proper random-equivalence upper bound is $\mathrm Q_{\mathrm{rpeq}}\leq\log_2|\mathcal H|$, and the full cube has size $2^n$. Hence $\mathrm Q_{\mathrm{rpeq}}=n$.
\end{proof}
\noindent\textbf{Reference.} \cite{angluin2017random}.

\subsection{Self-Directed Queries Complexity of Full Cube}\label{val:selfdirected_queries_complexity_full_cube}
\noindent\textbf{Statement.} The Self-Directed Queries Complexity of Full Cube is $n$.
\begin{proof}
Query each of the $n$ coordinates once, making an arbitrary prediction at each query. This gives at most $n$ mistakes on the full cube. Conversely, self-directed query complexity is at least recursive teaching dimension, whose established full-cube value is $n$. Therefore the value is exactly $n$.
\end{proof}
\noindent\textbf{Reference.} \cite{doliwa2014recursive}.

\subsection{Order Sample Compression of Full Cube}\label{val:order_sample_compression_full_cube}
\noindent\textbf{Statement.} The Order Sample Compression of Full Cube is $n$.
\begin{proof}
Order sample compression lower-bounds monotonic projected minimum teaching-set size, whose established full-cube value is $n$. Conversely, its established effective-range upper bound gives $\mathrm{OSC}\leq\mathrm E$, and the full cube has effective range $n$. Hence $\mathrm{OSC}=n$.
\end{proof}
\noindent\textbf{Reference.} \cite{darnstadt2016order}.

\subsection{Proper Ordered Sample Compression of Full Cube}\label{val:proper_ordered_sample_compression_full_cube}
\noindent\textbf{Statement.} The Proper Ordered Sample Compression of Full Cube is $n$.
\begin{proof}
Proper ordered sample compression is a restriction of ordinary order sample compression, whose established full-cube value is $n$. Its established effective-range upper bound gives $\mathrm{POSC}\leq\mathrm E$, while the full cube has effective range $n$. Thus $\mathrm{POSC}=n$.
\end{proof}
\noindent\textbf{Reference.} \cite{darnstadt2016order}.

\subsection{Proper Labeled Sample Compression of Full Cube}\label{val:proper_labeled_sample_compression_full_cube}
\noindent\textbf{Statement.} The Proper Labeled Sample Compression of Full Cube is $\Theta(n)$.
\begin{proof}
Write $Q_n=2^{[n]}$. We first give a proper width-two decoder for $Q_6$ on coordinates $0,\ldots,5$. Encode a word by $\sum_{i=0}^5 2^i h(i)$ and write $i^b$ for the labeled coordinate $(i,b)$. Put $\rho(\varnothing)=0$, $\rho(i^0)=63-2^i$, and $\rho(i^1)=63$ for even $i$. The three remaining singleton outputs are $\rho(1^1)=42$, $\rho(3^1)=44$, and $\rho(5^1)=52$. For distinct $i,j$, set $\rho(\{i^1,j^1\})=2^i+2^j$ and $\rho(\{i^0,j^0\})=63-2^i-2^j$. Indices in the mixed-key rules are taken modulo six. If $j=i+1$ or $j=i+2$, put $\rho(\{i^1,j^0\})=2^i$. The other mixed-key outputs are specified by the following table, whose columns correspond to $j=i+3,i+4,i+5$: $$ \begin{array}{c|rrr} i& i+3&i+4&i+5\\\hline 0&7&41&11\\ 1&14&19&50\\ 2&28&38&21\\ 3&56&13&26\\ 4&49&25&22\\ 5&35&37&32 \end{array} $$ These rules specify all 73 keys and each output extends its key. The three odd singleton outputs and the seventeen three-bit outputs in the mixed table are distinct and exhaust all twenty words of weight three; the remaining table entry, 32, has weight one.

To verify decoding, let $P,N$ be the positive and negative observed coordinates. If $N$ is empty and $|P|\le2$, retain $P$ (or its single labeled coordinate); the output agrees with every observation. If $N$ is empty and $|P|\ge3$, either $P$ contains an even coordinate, whose positive singleton key reconstructs all ones, or $P=\{1,3,5\}$, reconstructed by $1^1$. If $1\le|N|\le2$, retaining the negative coordinates reconstructs their complement and fits the sample. It remains that $|N|\ge3$. When $P$ is empty, use the empty key; when $|P|=2$, use the positive pair. When $P=\{i\}$ and $N$ meets $\{i+1,i+2\}$, the corresponding mixed key reconstructs the singleton $i$. Otherwise $N=\{i+3,i+4,i+5\}$ and the first entry of row $i$ reconstructs exactly $\{i,i+1,i+2\}$. Finally, $|P|=3$ forces $|N|=3$, so the sample is a full word of weight three and occurs as one of the decoder outputs listed above. Its key is contained in that word because every output extends its key. These cases exhaust all partial samples. The encoder can choose the first successful key in a fixed order; no label, order, message, seed or repetition beyond the retained labeled subset is transmitted. Properness is automatic on a full cube. One example cannot encode all 64 full concepts using only thirteen keys, so both ordinary and proper labeled compression of $Q_6$ equal two.

For $n=6q+r$, $0\le r<6$, apply this decoder on $q$ fixed six-coordinate blocks. The residual costs are $(a_0,\ldots,a_5)=(0,1,1,2,2,2)$. The first two are immediate. On a two-coordinate block, outputs $0,2,3,1,3$ at keys $\varnothing,0^0,0^1,1^0,1^1$ give a one-example scheme. Restricting the six-cube decoder to samples on any fixed three, four or five coordinates, and restricting its output back to those coordinates, gives cost two. The union of block keys has size at most $2q+a_r$, and the decoder recovers each part by intersection with the fixed block domains. This is $n/3+O(1)$. No stability is claimed. The older four-case five-cube scheme of \cite{darnstadt2016order}, Example 2, remains a valid independently checked comparison, but is not the source of the six-cube table.

For the lower bound, the $2^n$ full inputs need distinct retained keys, of which there are $\sum_{i=0}^k2^i\binom ni$. Define $b_n=\min\{k\ge0:\sum_{i=0}^k2^i\binom ni\ge2^n\}$. Thus $b_n\le\mathrm{LSC}(Q_n)=\mathrm{pLSC}(Q_n)$, and direct integer substitution matches the upper bound for $n=1,\ldots,8,11,12$. No exact general formula is claimed. For all $n$, $\sum_{i=0}^k2^i\binom ni\le4^k(3/2)^n$ gives $k\ge(n/2)\log_2(4/3)>n/5$, proving linear asymptotic order. The solver-free checker \texttt{sync/verify\_six\_cube\_compression.py} reads this table from its sole database source, checks all 729 partial samples in two representations, and verifies restrictions, block composition and counting bounds.
\end{proof}
\noindent\textbf{Reference.} Explicit six-cube construction and complete finite proof in this database record, independently replayed by sync/verify\_six\_cube\_compression.py. For the older four-case five-cube scheme and block-composition method: \cite{darnstadt2016order}, Example 2, Corollary 11 and Lemma 12 in the cached August 2014 preprint..

\subsection{Proper Stable Labeled Sample Compression of Full Cube}\label{val:proper_stable_labeled_sample_compression_full_cube}
\noindent\textbf{Statement.} The Proper Stable Labeled Sample Compression of Full Cube is $\lceil n/2\rceil$.
\begin{proof}
Here is a proper stable width-one scheme on two coordinates. In the following complete table, * means an unobserved coordinate: $$\begin{array}{c|c|c}s&K(s)&\rho(K(s))\\\hline **&**&00\\ 0*&0*&00\\ 1*&1*&11\\ *0&*0&10\\ *1&*1&01\\ 00&0*&00\\ 01&*1&01\\ 10&*0&10\\ 11&1*&11\end{array}$$ Each retained sample is contained in its input, and every decoded word agrees with that input. Deleting an unretained coordinate from a full two-coordinate sample leaves its selected singleton, whose compressed sample and reconstruction are unchanged. Other nonempty inputs select their only coordinate. Thus both compression-map stability and reconstruction stability hold. Partition $[n]$ into fixed disjoint pairs and, if necessary, one singleton. Apply this table independently on each pair; on the singleton retain its label when observed and decode an unobserved singleton as zero. The union of the retained samples has size at most $\lceil n/2\rceil$. The decoder identifies blocks from their fixed coordinate locations; no side message, order or repeated example is transmitted. Consistency and stability hold separately in each block, hence globally. Every resulting binary word belongs to the full cube, so the scheme is proper. Conversely the stable-compression lower bound gives at least projected NCTD, whose full-cube value is $\lceil n/2\rceil$. The two bounds coincide. The table is the sole block-certificate source read by sync/verify\_full\_cube\_labeled\_compression.py.
\end{proof}
\noindent\textbf{Reference.} Explicit block construction above; the matching lower bound uses the survey theorem thm:stable-comp-low-NCTD-star and the cube NCTD calculation in \cite{KSZ2019}..

\subsection{Occasional One-Sided Getback of Singletons}\label{val:occasional_onesided_getback_singletons}
\noindent\textbf{Statement.} The Occasional One-Sided Getback of Singletons is $1$.
\begin{proof}
O1G lower-bounds unlabeled sample compression, whose established singleton-class value is $1$. It is upper-bounded by Littlestone dimension, also $1$ for singletons. Thus the O1G value is $1$.
\end{proof}
\noindent\textbf{Reference.} \cite{ben1997online,chalopin2022unlabeled}.

\subsection{Random Equivalence Queries Complexity of Singletons}\label{val:random_equivalence_queries_complexity_singletons}
\noindent\textbf{Statement.} The Random Equivalence Queries Complexity of Singletons is $1$.
\begin{proof}
Random equivalence-query complexity is at least VC dimension and at most Littlestone dimension. Both established singleton-class values equal $1$, hence the random equivalence-query value is $1$.
\end{proof}
\noindent\textbf{Reference.} \cite{angluin2017random,littlestone1988learning}.

\subsection{Self-Directed Queries Complexity of Singletons}\label{val:selfdirected_queries_complexity_singletons}
\noindent\textbf{Statement.} The Self-Directed Queries Complexity of Singletons is $1$.
\begin{proof}
Self-directed query complexity lower-bounds recursive teaching dimension, whose singleton-class value is $1$. It is at most partial equivalence-query complexity, also $1$ on this class. Hence its value is $1$.
\end{proof}
\noindent\textbf{Reference.} \cite{doliwa2014recursive}.

\subsection{Monotonic Projected Minimum Teaching Set Size of Half-Intervals}\label{val:monotonic_projected_minimum_teaching_set_size_halfintervals}
\noindent\textbf{Statement.} The Monotonic Projected Minimum Teaching Set Size of Half-intervals is $1$.
\begin{proof}
The monotonic projected parameter dominates monotonic minimum teaching-set size, whose established half-interval value is $1$. It is at most order sample compression, also $1$ for half-intervals. Therefore the value is $1$.
\end{proof}
\noindent\textbf{Reference.} \cite{darnstadt2016order}.

\subsection{Occasional One-Sided Getback of Warmuth C5}\label{val:occasional_onesided_getback_warmuth_c5}
\noindent\textbf{Statement.} The Occasional One-Sided Getback of Warmuth's $C_5$ class is $3$.
\begin{proof}
O1G lower-bounds unlabeled sample compression, whose established value on Warmuth's C5 class is $3$. It is upper-bounded by Littlestone dimension, also $3$ on C5. Thus the O1G value is $3$.
\end{proof}
\noindent\textbf{Reference.} \cite{ben1997online,chalopin2022unlabeled}.

\subsection{Stable Unlabeled Complexity of Warmuth C5}\label{val:stable_unlabeled_complexity_warmuth_c5}
\noindent\textbf{Statement.} The Stable Unlabeled Complexity of Warmuth's $C_5$ class is $3$.
\begin{proof}
Stable unlabeled compression is at least unlabeled compression, whose established C5 value is $3$. It is at most Littlestone dimension, also $3$ on C5. Hence stable unlabeled complexity equals $3$.
\end{proof}
\noindent\textbf{Reference.} \cite{hanneke2021stable,bousquet2020proper,chalopin2022unlabeled}.

\subsection{Monotonic Projected Minimum Teaching Set Size of Singletons}\label{val:monotonic_projected_minimum_teaching_set_size_singletons}
\noindent\textbf{Statement.} The Monotonic Projected Minimum Teaching Set Size of Singletons is $n-1$.
\begin{proof}
The class itself is an admissible subfamily, and its minimum projected teaching-set size is $n-1$, so the monotonic projected value is at least $n-1$. Conversely, every subfamily of the singleton class has at most $n$ concepts. For a subfamily with $m\geq2$ singleton traces, projecting away the target coordinate shows that its minimum projected teaching-set size is at most $m-1\leq n-1$; one-concept subfamilies have value zero. Taking the maximum over subfamilies gives exactly $n-1$.
\end{proof}

\subsection{Monotonic Projected Minimum Teaching Set Size of Singletons Plus Empty Set}\label{val:monotonic_projected_minimum_teaching_set_size_singletons_plus_empty_set}
\noindent\textbf{Statement.} The Monotonic Projected Minimum Teaching Set Size of Singletons plus empty set is $n-1$.
\begin{proof}
The singleton subfamily is admissible and has minimum projected teaching-set size $n-1$, giving the lower bound. Every subfamily consists of the empty trace together with at most $m\leq n$ singleton traces. If it contains a singleton, that target has worst projected teaching-set size at most $m-1\leq n-1$; if it contains only the empty trace, its value is zero. Thus every subfamily has minimum projected teaching-set size at most $n-1$, proving equality.
\end{proof}

\subsection{Occasional One-Sided Getback of Singletons Plus Empty Set}\label{val:occasional_onesided_getback_singletons_plus_empty_set}
\noindent\textbf{Statement.} The Occasional One-Sided Getback of Singletons plus empty set is $1$.
\begin{proof}
O1G lower-bounds unlabeled sample compression, whose established value is $1$ for this class. It is upper-bounded by Littlestone dimension, also $1$. Therefore O1G equals $1$.
\end{proof}
\noindent\textbf{Reference.} \cite{ben1997online,chalopin2022unlabeled}.

\subsection{Self-Directed Queries Complexity of Singletons Plus Empty Set}\label{val:selfdirected_queries_complexity_singletons_plus_empty_set}
\noindent\textbf{Statement.} The Self-Directed Queries Complexity of Singletons plus empty set is $1$.
\begin{proof}
Self-directed query complexity lower-bounds recursive teaching dimension, whose established value is $1$ here. It is at most partial equivalence-query complexity, also $1$. Thus self-directed query complexity equals $1$.
\end{proof}
\noindent\textbf{Reference.} \cite{doliwa2014recursive}.

\subsection{Minimum Largest Order Shattered Set of Tagged Singletons}\label{val:minimum_largest_order_shattered_set_tagged_singletons}
\noindent\textbf{Statement.} The Minimum Largest Order Shattered Set of Tagged Singletons is $1$.
\begin{proof}
Minimum largest order-shattered-set size is at least largest strongly shattered-set size and at most largest shattered-set size. Both established tagged-singleton values are $1$, so the minimum order value is $1$.
\end{proof}
\noindent\textbf{Reference.} \cite{bollobas1995defect}.

\subsection{Random Equivalence Queries Complexity of Universal Class}\label{val:random_equivalence_queries_complexity_universal_class}
\noindent\textbf{Statement.} The Random Equivalence Queries Complexity of Universal Class is $\lfloor\log_2 n\rfloor$.
\begin{proof}
Random equivalence-query complexity is at least VC dimension and at most Littlestone dimension. Both established universal-class values equal $\lfloor\log_2 n\rfloor$, proving equality.
\end{proof}
\noindent\textbf{Reference.} \cite{angluin2017random,littlestone1988learning}.

\subsection{Worst Mistakes of Universal Class}\label{val:worst_mistakes_universal_class}
\noindent\textbf{Statement.} The Worst Mistakes of Universal Class is $\lfloor\log_2 n\rfloor$.
\begin{proof}
Worst mistakes is bounded below by VC dimension and above by Littlestone dimension. On the universal class the two established values are both $\lfloor\log_2 n\rfloor$, so worst mistakes has that value.
\end{proof}
\noindent\textbf{Reference.} \cite{ben1997online}.

\subsection{Random Equivalence Queries Complexity of Addressing}\label{val:random_equivalence_queries_complexity_addressing}
\noindent\textbf{Statement.} The Random Equivalence Queries Complexity of Addressing is $\lfloor\log n\rfloor\quad(n=2^r,\ r\ge1)$.
\begin{proof}
Restrict here to $n=2^r$ with all r-bit words present, so $\lfloor\log_2n\rfloor=\log_2n=r$. Random equivalence-query complexity lies between VC dimension and Littlestone dimension. The established addressing-class values of both bounds are $\lfloor\log n\rfloor$, hence the random equivalence-query value is the same.
\end{proof}
\noindent\textbf{Reference.} \cite{angluin2017random,littlestone1988learning}.

\subsection{Worst Mistakes of Addressing}\label{val:worst_mistakes_addressing}
\noindent\textbf{Statement.} The Worst Mistakes of Addressing is $\lfloor\log n\rfloor\quad(n=2^r,\ r\ge1)$.
\begin{proof}
Restrict here to $n=2^r$ with all r-bit words present, so $\lfloor\log_2n\rfloor=\log_2n=r$. Worst mistakes is at least VC dimension and at most Littlestone dimension. Both established addressing-class values are $\lfloor\log n\rfloor$, giving the claimed equality.
\end{proof}
\noindent\textbf{Reference.} \cite{ben1997online}.

\subsection{Minimum Projected Teaching Set Size of Half-Intervals}\label{val:minimum_projected_teaching_set_size_halfintervals}
\noindent\textbf{Statement.} The Minimum Projected Teaching Set Size of Half-intervals is $1$.
\begin{proof}
Minimum projected teaching-set size dominates minimum teaching-set size, whose half-interval value is $1$. It is at most monotonic projected minimum teaching-set size, also $1$ on this class. Hence the value is $1$.
\end{proof}
\noindent\textbf{Reference.} \cite{doliwa2014recursive}.

\subsection{Minimum Projected Teaching Set Size of Warmuth C5}\label{val:minimum_projected_teaching_set_size_warmuth_c5}
\noindent\textbf{Statement.} The Minimum Projected Teaching Set Size of Warmuth's $C_5$ class is $3$.
\begin{proof}
Minimum projected teaching-set size is at least minimum teaching-set size, whose established C5 value is $3$. It is at most maximum projected teaching-set size, also $3$ on C5. Therefore its value is $3$.
\end{proof}
\noindent\textbf{Reference.} \cite{doliwa2014recursive}.

\subsection{Monotonic Projected Minimum Teaching Set Size of Warmuth C5}\label{val:monotonic_projected_minimum_teaching_set_size_warmuth_c5}
\noindent\textbf{Statement.} The Monotonic Projected Minimum Teaching Set Size of Warmuth's $C_5$ class is $3$.
\begin{proof}
The full class has projected minimum teaching-set size three, giving the lower bound. Exhaustive enumeration of all 1,023 nonempty subfamilies and all 32 coordinate projections finds no projected minimum teaching-set size above three. Thus the subfamily maximum is exactly three; this finite calculation is checked by \texttt{sync/verify\_projected\_teaching\_values.py}.
\end{proof}

\subsection{Majority Complexity of Addressing}\label{val:majority_complexity_addressing}
\noindent\textbf{Statement.} The Majority Complexity of Addressing is $\lfloor\log n\rfloor\quad(n=2^r,\ r\ge1)$.
\begin{proof}
Restrict here to $n=2^r$ with all r-bit words present, so $\lfloor\log_2n\rfloor=\log_2n=r$. Majority complexity is at least Littlestone dimension and at most log size. Both established addressing values equal $\lfloor\log n\rfloor$, proving the equality.
\end{proof}
\noindent\textbf{Reference.} \cite{littlestone1988learning,maass1992lower}.

\subsection{Majority Complexity of Blockwise Addressing}\label{val:majority_complexity_blockwiseaddressing}
\noindent\textbf{Statement.} The Majority Complexity of Blockwise-Addressing is $d+1$.
\begin{proof}
Majority complexity is bounded below by Littlestone dimension and above by log size. The two established blockwise-addressing values are both $d+1$, so majority complexity is $d+1$.
\end{proof}
\noindent\textbf{Reference.} \cite{littlestone1988learning,maass1992lower}.

\subsection{Majority Complexity of Half-Intervals}\label{val:majority_complexity_halfintervals}
\noindent\textbf{Statement.} The Majority Complexity of Half-intervals is $\log n$.
\begin{proof}
Majority complexity is at least Littlestone dimension and at most log size. Both established half-interval values equal $\log n$, so majority complexity equals $\log n$.
\end{proof}
\noindent\textbf{Reference.} \cite{littlestone1988learning,maass1992lower}.

\subsection{Random Equivalence Queries Complexity of Majority}\label{val:random_equivalence_queries_complexity_majority}
\noindent\textbf{Statement.} The Random Equivalence Queries Complexity of Majority is $n$.
\begin{proof}
Random equivalence-query complexity lies between VC dimension and Littlestone dimension. Both established majority-class values are $n$, hence the value is $n$.
\end{proof}
\noindent\textbf{Reference.} \cite{angluin2017random,littlestone1988learning}.

\subsection{Worst Mistakes of Majority}\label{val:worst_mistakes_majority}
\noindent\textbf{Statement.} The Worst Mistakes of Majority is $n$.
\begin{proof}
Worst mistakes is bounded below by VC dimension and above by Littlestone dimension. Both established majority-class values equal $n$, so worst mistakes equals $n$.
\end{proof}
\noindent\textbf{Reference.} \cite{ben1997online}.

\subsection{Maximum Teaching Set Size of Majority}\label{val:maximum_teaching_set_size_majority}
\noindent\textbf{Statement.} The Maximum Teaching Set Size of Majority is $n$.
\begin{proof}
Maximum teaching-set size is at least maximum degree, whose established majority-class value is $n$. It is at most maximum projected teaching-set size, also $n$ on this class. Therefore it equals $n$.
\end{proof}
\noindent\textbf{Reference.} \cite{shinohara1991teachability,goldman1995complexity}.

\subsection{Projected Minimum Degree of Majority}\label{val:projected_minimum_degree_majority}
\noindent\textbf{Statement.} The Projected Minimum Degree of Majority is $n$.
\begin{proof}
Projected minimum degree is at least VC dimension, whose established majority-class value is $n$. It is at most projected maximal degree, also $n$ on this class. Hence projected minimum degree equals $n$.
\end{proof}
\noindent\textbf{Reference.} \cite{rubinstein2007shifting}.

\subsection{Degeneracy of Addressing}\label{val:degeneracy_addressing}
\noindent\textbf{Statement.} The Degeneracy of Addressing is $0$.
\begin{proof}
Degeneracy is at least minimum degree and at most densest-subgraph twice. Both established addressing-class values are $0$, so the degeneracy is $0$.
\end{proof}
\noindent\textbf{Reference.} \cite{haussler1988predicting,alon1992colorings,haussler1994predicting}.

\subsection{Random Proper Equivalence Queries Complexity of Addressing}\label{val:random_proper_equivalence_queries_complexity_addressing}
\noindent\textbf{Statement.} The Random Proper Equivalence Queries Complexity of Addressing is $\lfloor\log n\rfloor\quad(n=2^r,\ r\ge1)$.
\begin{proof}
Restrict here to $n=2^r$ with all r-bit words present, so $\lfloor\log_2n\rfloor=\log_2n=r$. Random proper equivalence-query complexity is at least random equivalence-query complexity, whose established addressing value is $\lfloor\log n\rfloor$. It is at most log size, also $\lfloor\log n\rfloor$. Hence equality holds.
\end{proof}
\noindent\textbf{Reference.} \cite{angluin2017random}.

\subsection{Majority Complexity of Majority}\label{val:majority_complexity_majority}
\noindent\textbf{Statement.} The Majority Complexity of Majority is $n$.
\begin{proof}
Majority complexity is at least Littlestone dimension and at most log size. Both established majority-class values equal $n$, so majority complexity is $n$.
\end{proof}
\noindent\textbf{Reference.} \cite{littlestone1988learning,maass1992lower}.

\subsection{Projected Double Density of Majority}\label{val:projected_double_density_or_projected_average_degree_majority}
\noindent\textbf{Statement.} The Projected Double Density or Projected Average Degree of Majority is $n$.
\begin{proof}
Projected double density is at least projected minimum degree and at most projected maximal degree. The established majority-class values of both bounds are $n$, proving the value is $n$.
\end{proof}
\noindent\textbf{Reference.} \cite{haussler1994predicting}.

\subsection{Random Proper Equivalence Queries Complexity of Majority}\label{val:random_proper_equivalence_queries_complexity_majority}
\noindent\textbf{Statement.} The Random Proper Equivalence Queries Complexity of Majority is $n$.
\begin{proof}
Random proper equivalence-query complexity is at least random equivalence-query complexity, whose established majority value is $n$. The established log-size upper bound is also $n$ for this class. Hence the value is $n$.
\end{proof}
\noindent\textbf{Reference.} \cite{angluin2017random}.

\subsection{No-Clashing Teaching Dimension of Warmuth's $C_5$ Class}\label{val:noclashing_teaching_dimension_warmuth_c5}
\noindent\textbf{Statement.} The No-Clashing Teaching Dimension of Warmuth's $C_5$ class is $2$.
\begin{proof}
The repository's exact C5 verifier constructs the ten cyclic $1001$ concepts, enumerates every target-consistent labeled sample of size at most one, and exhaustively rejects every no-clashing teacher map of width one. It then returns an explicit no-clashing map of width two. Hence $\mathrm{NCTD}(C_5)=2$.
\end{proof}
\noindent\textbf{Reference.} \cite{KSZ2019,palvolgyi2021unlabeled}.

\subsection{Positive No-Clashing Teaching Dimension of Warmuth's $C_5$ Class}\label{val:positive_noclashing_teaching_dimension_warmuth_c5}
\noindent\textbf{Statement.} The Positive No-Clashing Teaching Dimension of Warmuth's $C_5$ class is $2$.
\begin{proof}
The exact C5 verifier enumerates every target-consistent positive labeled sample of size at most one and rejects every positive no-clashing map of width one. It returns an explicit positive no-clashing map of width two, so $\mathrm{NCTD}^{+}(C_5)=2$.
\end{proof}
\noindent\textbf{Reference.} \cite{KSZ2019,palvolgyi2021unlabeled}.

\subsection{Unlabeled Sample Compression of the Three-Code Class}\label{val:unlabeled_sample_compression_three_code_class}
\noindent\textbf{Statement.} The Unlabeled sample compression of Three-Code Class is $1$.
\begin{proof}
A size-one improper scheme reconstructs $011$ from the empty code, $000$ from $\{x_0\}$, $001$ from $\{x_1\}$, and $101$ from $\{x_2\}$. For a realizable partial sample, use the empty code if it is consistent with $011$; otherwise use $\{x_0\}$ when it contains $x_0=0$, use $\{x_1\}$ when it contains $x_1=0$, and use $\{x_2\}$ in the remaining case. Each stated reconstruction is consistent with its assigned sample. Size zero is impossible because the full samples $000$ and $011$ disagree, whereas a zero-size scheme has one fixed reconstruction. The exhaustive check in $\texttt{sync/verify\_properness\_separation.py}$ independently verifies this value.
\end{proof}

\subsection{Proper Unlabeled Sample Compression of the Three-Code Class}\label{val:proper_unlabeled_sample_compression_three_code_class}
\noindent\textbf{Statement.} The Proper Unlabeled Sample Compression of Three-Code Class is $2$.
\begin{proof}
A proper size-two scheme exists, as verified by exhaustive enumeration of all realizable partial samples and all unlabeled codes of size at most two. No proper size-one scheme exists: the same finite enumeration checks every possible compressor assignment and every reconstruction in $\mathcal C_{\triangle}$. The check is recorded in $\texttt{sync/verify\_properness\_separation.py}$; it also verifies the matching improper size-one scheme.
\end{proof}

\subsection{Teaching Dimension of Majority}\label{val:teaching_dimension_majority}
\noindent\textbf{Statement.} The Teaching Dimension of Majority is $n$.
\begin{proof}
Teaching dimension is exactly maximum teaching-set size. The latter has established value n on the majority class, so teaching dimension is n as well.
\end{proof}
\noindent\textbf{Reference.} \cite{goldman1995complexity,shinohara1991teachability}.

\subsection{Relative Hitting Size of Majority}\label{val:relative_hitting_size_majority}
\noindent\textbf{Statement.} The Relative Hitting Size of Majority is $n$.
\begin{proof}
Relative hitting size is a definitional reformulation of maximum teaching-set size. The latter has established value n on the majority class, so relative hitting size is n as well.
\end{proof}
\noindent\textbf{Reference.} \cite{goldman1995complexity,shinohara1991teachability}.

\subsection{Stable Unlabeled Complexity of the Three-Code Class}\label{val:stable_unlabeled_complexity_three_code_class}
\noindent\textbf{Statement.} The Stable Unlabeled Complexity of Three-Code Class is $1$.
\begin{proof}
The size-one improper unlabeled scheme recorded for $\mathrm{USC}(\mathcal C_{\triangle})$ can be chosen stable. Exhaustive enumeration of every compressor assignment on the realizable partial samples verifies both consistency and the identity that deleting an unselected example leaves its code unchanged. Size zero is impossible because two full samples disagree and would have the same empty code. The check is recorded in $\texttt{sync/verify\_properness\_separation.py}$.
\end{proof}
\noindent\textbf{Reference.} \cite{hanneke2021stable}.

\subsection{Proper Stable Sample Compression of the Three-Code Class}\label{val:proper_stable_sample_compression_three_code_class}
\noindent\textbf{Statement.} The Proper Stable Sample Compression of Three-Code Class is $2$.
\begin{proof}
Exhaustive enumeration of all proper stable unlabeled schemes on the realizable partial samples finds a size-two scheme and rules out every size-one scheme. The enumeration imposes sample consistency, reconstruction in $\mathcal C_{\triangle}$, and stability under deletion of an unselected example; it is recorded in $\texttt{sync/verify\_properness\_separation.py}$.
\end{proof}
\noindent\textbf{Reference.} \cite{hanneke2021stable}.

\subsection{Labeled Sample Compression of the Three-Code Class}\label{val:labeled_sample_compression_three_code_class}
\noindent\textbf{Statement.} The Labeled Sample Compression of Three-Code Class is $1$.
\begin{proof}
The explicit proper stable labeled size-one scheme checked in $\texttt{sync/verify\_properness\_separation.py}$ is in particular a labeled scheme. Size zero is impossible for a class containing two distinct concepts, because the empty code has one fixed reconstruction.
\end{proof}
\noindent\textbf{Reference.} \cite{FW1995}.

\subsection{Stable Labeled Sample Compression of the Three-Code Class}\label{val:stable_labeled_sample_compression_three_code_class}
\noindent\textbf{Statement.} The Stable Labeled Sample Compression of Three-Code Class is $1$.
\begin{proof}
The explicit proper stable labeled size-one scheme checked in $\texttt{sync/verify\_properness\_separation.py}$ remains valid after dropping properness. Size zero is impossible for a non-singleton class.
\end{proof}
\noindent\textbf{Reference.} \cite{hanneke2021stable}.

\subsection{Proper Labeled Sample Compression of the Three-Code Class}\label{val:proper_labeled_sample_compression_three_code_class}
\noindent\textbf{Statement.} The Proper Labeled Sample Compression of Three-Code Class is $1$.
\begin{proof}
The explicit proper stable labeled size-one scheme checked in $\texttt{sync/verify\_properness\_separation.py}$ is in particular proper labeled. Size zero is impossible for a non-singleton class.
\end{proof}
\noindent\textbf{Reference.} \cite{FW1995}.

\subsection{Proper Stable Labeled Sample Compression of the Three-Code Class}\label{val:proper_stable_labeled_sample_compression_three_code_class}
\noindent\textbf{Statement.} The Proper Stable Labeled Sample Compression of Three-Code Class is $1$.
\begin{proof}
Flip coordinate $x_0$ in every concept, converting $\{000,011,101\}$ into the three-singleton class. Identify its coordinates with $\mathbb Z/3\mathbb Z$. Decode the empty key as $\{0\}$, a positive example at $i$ as $\{i\}$, and a negative example at $i$ as $\{i+1\}$. Retain the positive example whenever present. For a nonempty all-negative sample retain an observed $i$ whose successor $i+1$ is unobserved, using a fixed tie rule. Such an $i$ exists because any realizable all-negative sample is a proper subset of the three-cycle. Reconstruction is a consistent singleton. Deleting an unretained example leaves the same key: a two-negative sample becomes its retained singleton, and a sample with a positive still retains that positive. The empty sample retains nothing. Undo the fixed label flip on samples and reconstructions to obtain a proper, selected-key-stable labeled scheme for the Three-Code Class. Size zero cannot reconstruct distinct full concepts. The cyclic scheme is checked on all 19 realizable partial samples by sync/verify\_proper\_compression\_covc.py; the earlier finite checker supplies an independent scheme.
\end{proof}
\noindent\textbf{Reference.} \cite{hanneke2021stable}.

\subsection{Equivalence Queries Complexity of the Three-Code Class}\label{val:equivalence_queries_complexity_three_code_class}
\noindent\textbf{Statement.} The Equivalence Queries Complexity of Three-Code Class is $1$.
\begin{proof}
Use the improper proposal $001$. Against the three targets $000,011,101$, its disagreement coordinate is respectively $x_0,x_1,x_2$, and each counterexample identifies the target uniquely. The verifier $\texttt{sync/verify\_properness\_separation.py}$ checks these three response classes. Size zero is impossible for a non-singleton class.
\end{proof}
\noindent\textbf{Reference.} \cite{angluin1988queries}.

\subsection{Proper Equivalence Queries Complexity of the Three-Code Class}\label{val:proper_equivalence_queries_complexity_three_code_class}
\noindent\textbf{Statement.} The Proper Equivalence Queries Complexity of Three-Code Class is $2$.
\begin{proof}
No proper first proposal identifies every target after one counterexample: the exhaustive response check in $\texttt{sync/verify\_properness\_separation.py}$ finds, for each proposal in $\mathcal C_{\triangle}$, a disagreement coordinate compatible with at least two targets. Conversely, query any two class members; if both are rejected, the third member is forced. Thus the exact proper-query complexity is two.
\end{proof}
\noindent\textbf{Reference.} \cite{angluin1988queries}.

\subsection{Membership Query Complexity of the Three-Code Class}\label{val:membership_query_complexity_three_code_class}
\noindent\textbf{Statement.} The Membership Query Complexity of Three-Code Class is $2$.
\begin{proof}
Every coordinate has one label value shared by two of $000,011,101$, so one membership query cannot always identify the target; this is checked in $\texttt{sync/verify\_properness\_separation.py}$. Query $x_0$ first. If its label is zero the target is $000$; otherwise query $x_1$, which distinguishes $011$ from $101$. Thus two queries are necessary and sufficient.
\end{proof}
\noindent\textbf{Reference.} \cite{angluin1988queries}.

\subsection{Membership and Equivalence Queries Complexity of the Three-Code Class}\label{val:membership_and_equivalence_queries_complexity_three_code_class}
\noindent\textbf{Statement.} The Membership and Equivalence Queries Complexity of Three-Code Class is $1$.
\begin{proof}
The one improper equivalence proposal $001$ identifies the target after its counterexample, as in the established equivalence-query value. Hence one mixed query suffices; size zero is impossible for a non-singleton class.
\end{proof}
\noindent\textbf{Reference.} \cite{angluin1988queries}.

\subsection{Membership and Proper Equivalence Queries Complexity of the Three-Code Class}\label{val:membership_and_proper_equivalence_queries_complexity_three_code_class}
\noindent\textbf{Statement.} The Membership and Proper Equivalence Queries Complexity of Three-Code Class is $2$.
\begin{proof}
A single allowed query cannot identify every target: every membership coordinate leaves two possible targets for one answer, and every proper equivalence proposal has a counterexample response compatible with two targets; both checks are in $\texttt{sync/verify\_properness\_separation.py}$. Two membership queries suffice by the established membership-query strategy, so the exact mixed proper-query complexity is two.
\end{proof}
\noindent\textbf{Reference.} \cite{angluin1988queries}.

\subsection{Random Equivalence Queries Complexity of the Three-Code Class}\label{val:random_equivalence_queries_complexity_three_code_class}
\noindent\textbf{Statement.} The Random Equivalence Queries Complexity of Three-Code Class is $1$.
\begin{proof}
For every full-support counterexample distribution $\mu$, use the improper proposal $001$. For each target in $\mathcal C_{\triangle}$ its disagreement set with this proposal is a distinct singleton coordinate, so the random counterexample is deterministic and identifies the target. The verifier $\texttt{sync/verify\_properness\_separation.py}$ checks these response classes. Zero queries are impossible for a non-singleton class.
\end{proof}
\noindent\textbf{Reference.} \cite{angluin2017random}.

\subsection{Partial Equivalence Queries Complexity of the Three-Code Class}\label{val:partial_equivalence_queries_complexity_three_code_class}
\noindent\textbf{Statement.} The Partial Equivalence Queries Complexity of Three-Code Class is $1$.
\begin{proof}
A full Boolean hypothesis is an allowed partial-equivalence query. Use the full improper proposal $001$; its three possible counterexamples identify respectively $000,011,101$, as checked in $\texttt{sync/verify\_properness\_separation.py}$. Thus one query suffices, while zero queries cannot identify a non-singleton class.
\end{proof}
\noindent\textbf{Reference.} \cite{maass1992lower}.

\subsection{Size of the Three-Code Class}\label{val:size_three_code_class}
\noindent\textbf{Statement.} The Size of Three-Code Class is $3$.
\begin{proof}
The defining list is $\mathcal C_{\triangle}=\{000,011,101\}$, whose three words are distinct; see the exhaustive certificate in $\texttt{sync/verify\_properness\_separation.py}$.
\end{proof}
\noindent\textbf{Reference.} Direct calculation from the class definition..

\subsection{Effective Range of the Three-Code Class}\label{val:effective_range_three_code_class}
\noindent\textbf{Statement.} The Effective Range of Three-Code Class is $3$.
\begin{proof}
At each of the three coordinates, the defining words take both labels. Thus every coordinate is active, and the effective range is three. This is also checked directly in $\texttt{sync/verify\_properness\_separation.py}$.
\end{proof}
\noindent\textbf{Reference.} Direct calculation from the class definition..

\subsection{VC Dimension of the Three-Code Class}\label{val:vc_dimension_three_code_class}
\noindent\textbf{Statement.} The VC Dimension of Three-Code Class is $1$.
\begin{proof}
Every individual coordinate realizes both labels, so $\mathrm{VC}\geq1$. Each two-coordinate projection has exactly three traces, hence is not shattered; larger sets cannot be shattered either. Therefore $\mathrm{VC}=1$. The exhaustive check is in $\texttt{sync/verify\_properness\_separation.py}$.
\end{proof}
\noindent\textbf{Reference.} Direct calculation from the class definition..

\subsection{Littlestone Dimension of the Three-Code Class}\label{val:littlestone_dimension_three_code_class}
\noindent\textbf{Statement.} The Littlestone Dimension of Three-Code Class is $1$.
\begin{proof}
An active coordinate gives a complete Littlestone tree of depth one. At every coordinate, one label is taken by exactly one of the three concepts, so one branch is a singleton and cannot support another split. No depth-two complete tree exists, hence $\operatorname{Ldim}=1$.
\end{proof}
\noindent\textbf{Reference.} Direct calculation from the class definition..

\subsection{Teaching Dimension of the Three-Code Class}\label{val:teaching_dimension_three_code_class}
\noindent\textbf{Statement.} The Teaching Dimension of Three-Code Class is $1$.
\begin{proof}
For $000,011,101$, respectively, the labels on $x_0,x_1,x_2$ uniquely identify the target among the three concepts. No non-singleton class has teaching dimension zero. Thus $\mathrm{TD}=1$; the finite certificate is checked in $\texttt{sync/verify\_properness\_separation.py}$.
\end{proof}
\noindent\textbf{Reference.} Direct calculation from the class definition..

\subsection{Star Number of the Three-Code Class}\label{val:star_number_three_code_class}
\noindent\textbf{Statement.} The Star number of Three-Code Class is $2$.
\begin{proof}
Projecting to $\{x_0,x_1\}$ gives $\{00,01,11\}$, with $01$ adjacent to both other traces, so the star number is at least two. No projection of a three-concept class has a vertex of degree above two. Hence $\mathrm S=2$; this is exhaustively checked in $\texttt{sync/verify\_properness\_separation.py}$.
\end{proof}
\noindent\textbf{Reference.} Direct calculation from the class definition; compare the projected-star formulation in \cite{hanneke2015minimax}..

\subsection{co-VC Dimension of the Three-Code Class}\label{val:covc_dimension_three_code_class}
\noindent\textbf{Statement.} The co-VC dimension of Three-Code Class is $3$.
\begin{proof}
In the full three-coordinate projection, the word $001$ is absent, while its three one-bit neighbours $000,011,101$ are precisely the class members. Hence it is a hollow three-star. No hollow star can have more than three leaves because its distinct realized neighbours would require more than the three available concepts. Thus $\mathrm{coVC}=3$; the finite enumeration is in $\texttt{sync/verify\_properness\_separation.py}$.
\end{proof}
\noindent\textbf{Reference.} Direct calculation from the class definition; compare the hollow-star formulation in \cite{bousquet2020proper}..

\subsection{Extended Teaching Dimension of the Three-Code Class}\label{val:extended_teaching_dimension_three_code_class}
\noindent\textbf{Statement.} The Extended Teaching Dimension of Three-Code Class is $2$.
\begin{proof}
The absent labeling $001$ needs two coordinates to leave at most one class member consistent, so $\mathrm{XTD}\geq2$. Every labeling is specified by some two-coordinate restriction, so $\mathrm{XTD}\leq2$. Equivalently, $\mathrm{XTD}=\mathrm S=2$ by Theorem 13 of the cited paper. The direct finite check is in $\texttt{sync/verify\_properness\_separation.py}$.
\end{proof}
\noindent\textbf{Reference.} \cite{hanneke2015minimax}.

\subsection{Minimum Star Number of the Three-Code Class}\label{val:minimum_star_number_three_code_class}
\noindent\textbf{Statement.} The Minimum Star Number of Three-Code Class is $1$.
\begin{proof}
Every coordinate is shattered, so every centre labeling has a centred one-star and $\mathrm S_{\min}\geq1$. At the centre $000$, no two-star exists because the other two class members differ from it on two coordinates. Hence $\mathrm S_{\min}=1$. The enumeration over all eight centre labelings is checked in $\texttt{sync/verify\_properness\_separation.py}$.
\end{proof}
\noindent\textbf{Reference.} \cite{hanneke2024star_eluder}.

\subsection{Maximum Projected Teaching-Set Size of the Three-Code Class}\label{val:maximum_projected_teaching_set_size_three_code_class}
\noindent\textbf{Statement.} The Maximum Projected Teaching Set Size of Three-Code Class is $2$.
\begin{proof}
The maximum projected teaching-set size equals the star number. The projection $\{00,01,11\}$ has teaching set size two for its middle trace, and no three-concept projection can require more. The exhaustive check is in $\texttt{sync/verify\_properness\_separation.py}$.
\end{proof}
\noindent\textbf{Reference.} \cite{hanneke2015minimax}.

\subsection{Projected Maximal Degree of the Three-Code Class}\label{val:projected_maximal_degree_three_code_class}
\noindent\textbf{Statement.} The Projected Maximal Degree of Three-Code Class is $2$.
\begin{proof}
On the projection $\{x_0,x_1\}$, the trace $01$ is adjacent to both $00$ and $11$, so projected maximal degree is at least two. A three-concept projection has no degree above two. Thus the value is two, as checked in $\texttt{sync/verify\_properness\_separation.py}$.
\end{proof}
\noindent\textbf{Reference.} \cite{hanneke2015minimax}.

\subsection{Maximum Teaching-Set Size of the Three-Code Class}\label{val:maximum_teaching_set_size_three_code_class}
\noindent\textbf{Statement.} The Maximum Teaching Set Size of Three-Code Class is $1$.
\begin{proof}
This is the classical teaching dimension. Each of $000,011,101$ is uniquely identified by its label respectively at $x_0,x_1,x_2$, so the maximum teaching-set size is one. The direct certificate is in $\texttt{sync/verify\_properness\_separation.py}$.
\end{proof}
\noindent\textbf{Reference.} Direct calculation from the class definition..

\subsection{Recursive Teaching Dimension of the Three-Code Class}\label{val:recursive_teaching_dimension_three_code_class}
\noindent\textbf{Statement.} The Recursive Teaching Dimension of Three-Code Class is $1$.
\begin{proof}
In the initial class, every concept already has a one-coordinate teaching set, so they may all form the first recursive batch. A non-singleton class cannot have recursive teaching dimension zero. Hence $\mathrm{RTD}=1$, verified by the residual-class recursion in $\texttt{sync/verify\_properness\_separation.py}$.
\end{proof}
\noindent\textbf{Reference.} Direct calculation from the class definition; see the teaching-model definition in \cite{zilles2011models}..

\subsection{Projected Recursive Teaching Dimension of the Three-Code Class}\label{val:monotone_recursive_teaching_dimension_three_code_class}
\noindent\textbf{Statement.} The Projected Recursive Teaching Dimension of Three-Code Class is $1$.
\begin{proof}
The full class has recursive teaching dimension one. Each proper coordinate projection is either a singleton/two-concept class or a three-trace path, and its endpoints give a recursive teaching batch of size one. Thus every projection has RTD at most one, and the maximum is one. All projections are enumerated in $\texttt{sync/verify\_properness\_separation.py}$.
\end{proof}
\noindent\textbf{Reference.} Direct calculation from the class definition; compare the projected variant in \cite{doliwa2014recursive}..

\subsection{No-Clashing Teaching Dimension of the Three-Code Class}\label{val:noclashing_teaching_dimension_three_code_class}
\noindent\textbf{Statement.} The No-Clashing Teaching Dimension of Three-Code Class is $1$.
\begin{proof}
Assign respectively the one-point samples $(x_0,0),(x_1,1),(x_2,1)$ to $000,011,101$. Each sample is inconsistent with both other concepts, so no pair clashes. A zero-size teacher map on a non-singleton class makes every pair clash. Hence $\mathrm{NCTD}=1$; the map is checked in $\texttt{sync/verify\_properness\_separation.py}$.
\end{proof}
\noindent\textbf{Reference.} Direct calculation from the class definition; see \cite{KSZ2019} for NCTD..

\subsection{Positive No-Clashing Teaching Dimension of the Three-Code Class}\label{val:positive_noclashing_teaching_dimension_three_code_class}
\noindent\textbf{Statement.} The Positive No-Clashing Teaching Dimension of Three-Code Class is $1$.
\begin{proof}
Give $000$ the empty sample and give $011,101$ respectively their positive examples $(x_1,1),(x_2,1)$. The latter two samples are inconsistent with $000$ and with each other, so this is a positive no-clashing map of width one. Width zero is impossible for a non-singleton class. Hence $\mathrm{NCTD}^{+}=1$; the map is checked in $\texttt{sync/verify\_properness\_separation.py}$.
\end{proof}
\noindent\textbf{Reference.} Direct calculation from the class definition; see \cite{KSZ2019} for positive NCTD..

\subsection{Preference-Based Teaching Dimension of the Three-Code Class}\label{val:preferencebased_teaching_dimension_three_code_class}
\noindent\textbf{Statement.} The Preference-Based Teaching Dimension of Three-Code Class is $1$.
\begin{proof}
Use preference order $000\succ011\succ101$. The empty sample teaches the most preferred target; $(x_1,1)$ teaches $011$, and $(x_2,1)$ teaches $101$. Thus the value is at most one, and zero is impossible for a non-singleton class. Hence $\mathrm{PBTD}=1$.
\end{proof}
\noindent\textbf{Reference.} Direct calculation from the class definition; compare \cite{gao2017preference}..

\subsection{Positive Recursive Teaching Dimension of the Three-Code Class}\label{val:positive_recursive_teaching_dimension_three_code_class}
\noindent\textbf{Statement.} The Positive Recursive Teaching Dimension of Three-Code Class is $1$.
\begin{proof}
Remove $011$ and $101$ first, using their positive one-point teaching samples $(x_1,1)$ and $(x_2,1)$, each unique among the current class. The remaining $000$ is removed with the empty sample. Thus positive RTD is at most one; it cannot be zero on a non-singleton class. Hence $\mathrm{RTD}^{+}=1$.
\end{proof}
\noindent\textbf{Reference.} Direct calculation from the class definition; see \cite{fallat2023batch}..

\subsection{Teaching Dimension of Blockwise-Addressing}\label{val:teaching_dimension_blockwiseaddressing}
\noindent\textbf{Statement.} The Teaching Dimension of Blockwise-Addressing is $2$.
\begin{proof}
For $f_i$, the two labeled examples $(z,0),(x_i,1)$ exclude every $g_j$ and every other $f_j$; symmetrically $(z,1),(y_i,1)$ teaches $g_i$. Thus $\mathrm{TD}\le2$. No one-coordinate sample teaches $f_i$: the $z$-label leaves all $f_j$; an $x$-coordinate either leaves another $f_j$ or agrees with some $g_j$; and a coordinate in a $y$-block is shared by $2^{d-1}$ of the $f_j$. Hence $\mathrm{TD}=2$.
\end{proof}
\noindent\textbf{Reference.} Direct calculation from the Blockwise-Addressing construction in \cite{ben1997online}..

\subsection{Maximum Teaching-Set Size of Blockwise-Addressing}\label{val:maximum_teaching_set_size_blockwiseaddressing}
\noindent\textbf{Statement.} The Maximum Teaching Set Size of Blockwise-Addressing is $2$.
\begin{proof}
Maximum teaching-set size is the classical teaching dimension. The direct Blockwise-Addressing teaching calculation gives $\mathrm{TD}=2$, hence $\mathrm{TS}_{\max}=2$.
\end{proof}
\noindent\textbf{Reference.} Direct calculation from the Blockwise-Addressing construction in \cite{ben1997online}..

\subsection{Star Number of Blockwise-Addressing}\label{val:star_number_blockwiseaddressing}
\noindent\textbf{Statement.} The Star number of Blockwise-Addressing is $2^d\leq\mathrm S\leq2^{d+1}-1$.
\begin{proof}
The established Blockwise-Addressing bound gives $\mathrm{coVC}\ge2^d+1$. The direct affine relation $\mathrm S\ge\mathrm{coVC}-1$ therefore gives $\mathrm S\ge2^d$. Conversely any $r$-star uses its centre and $r$ distinct leaves, so $r\le|\mathcal H|-1=2^{d+1}-1$. This proves the displayed interval.
\end{proof}
\noindent\textbf{Reference.} Derived from the established Blockwise-Addressing co-VC bound and the direct star/co-VC and size/star relationships; see \cite{ben1997online,bousquet2020proper}..

\subsection{Spherical Dimension of the Three-Code Class}\label{val:spherical_dimension_three_code_class}
\noindent\textbf{Statement.} The Spherical Dimension of Three-Code Class is $1$.
\begin{proof}
Write the three concepts as $a=000$, $b=011$, and $c=101$.  Across these concepts, the three coordinates have minority singleton supports $\{c\}$, $\{b\}$, and $\{a\}$, respectively; coordinate bit flips do not change those minority singleton supports.  In a threshold subclass, every nonconstant coordinate is an initial or final segment of one common linear order of the concepts, so its minority singleton must be an endpoint.  A three-element order has only two endpoints, so this class is not, even after bit flips, a threshold subclass.  Its established VC dimension is one.  Lemma 31 of the cited paper classifies precisely this case as spherical dimension one.
\end{proof}
\noindent\textbf{Reference.} \cite{chornomaz2025spherical}.

\subsection{Sign Rank of the Three-Code Class}\label{val:sign_rank_three_code_class}
\noindent\textbf{Statement.} The Sign Rank of Three-Code Class is $2$.
\begin{proof}
In the concept order $(000,011,101)$, the coordinate sign columns are $(--+),(-+-),(-++)$.  They have a rank-two realization: take row vectors $p=(-1,0)$, $q=(1,1)$, $r=(1,-1)$, and column vectors $s=(1/2,-1)$, $t=(1/2,1)$, $w=(1,0)$, respectively.  The signs of their three corresponding dot-product columns are exactly the displayed columns.  Rank one is impossible because all nonzero columns of a rank-one real matrix have sign vectors that are equal or opposite, whereas $(--+)$ and $(-+-)$ are neither.  Hence the sign rank is two.
\end{proof}
\noindent\textbf{Reference.} \cite{chornomaz2025spherical}.

\subsection{SCO Dimension of the Three-Code Class}\label{val:sco_dimension_three_code_class}
\noindent\textbf{Statement.} The SCO Dimension of Three-Code Class is $2$.
\begin{proof}
The established spherical dimension of this class is one.  Theorem 12 of Chornomaz--Moran--Waknine gives $\mathrm{sd}(\mathcal H)+1\leq\mathrm{SCO}(\mathcal H)\leq\mathrm{SR}(\mathcal H)$.  The directly verified sign-rank value here is two, so the two bounds coincide at $\mathrm{SCO}=2$.
\end{proof}
\noindent\textbf{Reference.} \cite{chornomaz2025spherical}.

\subsection{Dual Sign Rank of the Three-Code Class}\label{val:dual_sign_rank_three_code_class}
\noindent\textbf{Statement.} The Dual Sign Rank of Three-Code Class is $2$.
\begin{proof}
The coordinate sign columns are $(--+),(-+-),(-++)$.  Any two are neither equal nor opposite.  Hence in every realizing real matrix, any selected pair is linearly independent: dependent nonzero columns would have identical signs (positive scalar multiple) or opposite signs (negative scalar multiple).  Thus dual sign rank is at least two.  It is at most sign rank, whose directly verified value for this class is two.  Therefore the dual sign rank is exactly two.
\end{proof}
\noindent\textbf{Reference.} \cite{chornomaz2025spherical}.

\subsection{Maximum Degree of the Three-Code Class}\label{val:maximum_degree_three_code_class}
\noindent\textbf{Statement.} The Maximum Degree of Three-Code Class is $0$.
\begin{proof}
The three strings $000,011,101$ are pairwise at Hamming distance two.  Thus the $1$-inclusion graph has no edges and its maximum degree is zero.
\end{proof}

\subsection{Minimum Degree of the Three-Code Class}\label{val:minimum_degree_three_code_class}
\noindent\textbf{Statement.} The Minimum Degree of Three-Code Class is $0$.
\begin{proof}
The three strings $000,011,101$ are pairwise at Hamming distance two.  Hence every vertex of the $1$-inclusion graph is isolated, so its minimum degree is zero.
\end{proof}

\subsection{Degeneracy of the Three-Code Class}\label{val:degeneracy_three_code_class}
\noindent\textbf{Statement.} The Degeneracy of Three-Code Class is $0$.
\begin{proof}
Its $1$-inclusion graph is edgeless because its three concepts are pairwise at Hamming distance two.  Every nonempty induced subgraph therefore has minimum degree zero, which is exactly the degeneracy.
\end{proof}

\subsection{Double Density or Average Degree of the Three-Code Class}\label{val:double_density_or_average_degree_three_code_class}
\noindent\textbf{Statement.} The Double Density or Average Degree of Three-Code Class is $0$.
\begin{proof}
Its $1$-inclusion graph has no edges: $000,011,101$ are pairwise at Hamming distance two.  Its average degree, equivalently its double density, is therefore zero.
\end{proof}

\subsection{Densest Subgraph (twice) of the Three-Code Class}\label{val:densest_subgraph_twice_three_code_class}
\noindent\textbf{Statement.} The Densest Subgraph (twice) of Three-Code Class is $0$.
\begin{proof}
The $1$-inclusion graph is edgeless, and so is every one of its subgraphs.  Every subgraph consequently has average degree zero, proving that its densest-subgraph double density is zero.
\end{proof}

\subsection{Maximum Positive Degree of the Three-Code Class}\label{val:maximum_positive_degree_three_code_class}
\noindent\textbf{Statement.} The Maximum Positive Degree of Three-Code Class is $0$.
\begin{proof}
There are no edges in the $1$-inclusion graph, since every pair of distinct concepts differs on two coordinates.  Thus no concept has a smaller adjacent neighbour and every positive degree is zero.
\end{proof}

\subsection{Projected Minimum Degree of the Three-Code Class}\label{val:projected_minimum_degree_three_code_class}
\noindent\textbf{Statement.} The Projected Minimum Degree of Three-Code Class is $1$.
\begin{proof}
On any one coordinate, the projection is the full one-cube, whose minimum degree is one.  On two coordinates, every projection is a three-vertex path and has minimum degree one; on all three coordinates it is edgeless, and the empty projection has degree zero.  Hence the maximum projected minimum degree is one.
\end{proof}

\subsection{Projected Double Density or Projected Average Degree of the Three-Code Class}\label{val:projected_double_density_or_projected_average_degree_three_code_class}
\noindent\textbf{Statement.} The Projected Double Density or Projected Average Degree of Three-Code Class is $\frac{4}{3}$.
\begin{proof}
A one-coordinate projection is one edge and has average degree one.  Every two-coordinate projection is the three-vertex path with two edges, hence average degree $4/3$.  The full projection is edgeless and the empty projection has degree zero.  Thus the maximum projected average degree is $4/3$.
\end{proof}

\subsection{Effective VC Radius of the Three-Code Class}\label{val:effective_vc_radius_three_code_class}
\noindent\textbf{Statement.} The Effective VC Radius of Three-Code Class is $1$.
\begin{proof}
Every coordinate is nonconstant on the class and is therefore shattered.  On every pair of coordinates the class realizes exactly three of the four binary traces, so no pair is shattered.  Since the effective range is all three coordinates, the effective VC radius is exactly one.
\end{proof}

\subsection{Largest Strongly Shattered Set of the Three-Code Class}\label{val:largest_strongly_shattered_set_three_code_class}
\noindent\textbf{Statement.} The Largest Strongly Shattered Set of Three-Code Class is $0$.
\begin{proof}
A nonempty strongly shattered set contains a one-dimensional subcube and hence gives an edge of the $1$-inclusion graph.  The three concepts are pairwise at Hamming distance two, so that graph is edgeless.  The empty set is strongly shattered, proving that the maximum size is zero.
\end{proof}

\subsection{Diameter of the Three-Code Class}\label{val:diameter_three_code_class}
\noindent\textbf{Statement.} The Diameter of Three-Code Class is $2$.
\begin{proof}
Each pair among $000,011,101$ differs on exactly two coordinates.  Thus the largest Hamming distance between concepts in the class is two.
\end{proof}

\subsection{Hamming Radius of the Three-Code Class}\label{val:hamming_radius_three_code_class}
\noindent\textbf{Statement.} The Hamming Radius of Three-Code Class is $1$.
\begin{proof}
The center $001$ is at Hamming distance one from each of $000,011,101$, so radius one suffices.  Radius zero contains only one binary string, whereas the class has three distinct concepts.  Hence the minimum enclosing radius is one.
\end{proof}

\subsection{Oriented Diameter of the Three-Code Class}\label{val:oriented_diameter_three_code_class}
\noindent\textbf{Statement.} The Oriented Diameter of Three-Code Class is $2$.
\begin{proof}
Restricting $000$ and $011$ to the two coordinates where they differ gives the oriented pair $\emptyset,\{1,2\}$, so the value is at least two.  An oriented pair on three coordinates would require the all-zero and all-one traces on the full domain, but neither $111$ nor any other three-one concept is in the class.  Thus the value is exactly two.
\end{proof}

\subsection{Distinguishing Range of the Three-Code Class}\label{val:distinguishing_range_three_code_class}
\noindent\textbf{Statement.} The Distinguishing Range of Three-Code Class is $2$.
\begin{proof}
One binary coordinate has only two possible traces and cannot distinguish all three concepts.  On the first two coordinates, the traces are $00,01,10$, which are distinct.  Hence exactly two coordinates are necessary and sufficient.
\end{proof}

\subsection{Threshold Dimension of the Three-Code Class}\label{val:threshold_dimension_three_code_class}
\noindent\textbf{Statement.} The Threshold Dimension of Three-Code Class is $2$.
\begin{proof}
On the first and third coordinates, the traces are $00,01,11$, which form a threshold chain of length two.  A length-three threshold chain requires four distinct traces on three coordinates, but the class has only three concepts.  Therefore the threshold dimension is two.
\end{proof}

\subsection{Path Dimension of the Three-Code Class}\label{val:path_dimension_three_code_class}
\noindent\textbf{Statement.} The Path Dimension of Three-Code Class is $2$.
\begin{proof}
On the first two coordinates, the traces $01,00,10$ form a full path of length two: flip the second coordinate and then the first.  A full path of length three has four distinct successive traces, whereas this class has only three concepts.  Hence the path dimension is two.
\end{proof}

\subsection{Largest Shattered Set of the Three-Code Class}\label{val:largest_shattered_set_three_code_class}
\noindent\textbf{Statement.} The Largest Shattered Set of Three-Code Class is $1$.
\begin{proof}
Largest Shattered Set is the VC dimension by definition.  The established VC dimension of this class is one, so this value is one as well.
\end{proof}

\subsection{Projected Strongly Shattered Dimension of the Three-Code Class}\label{val:projected_strongly_shattered_dimension_three_code_class}
\noindent\textbf{Statement.} The Projected Strongly Shattered Dimension of Three-Code Class is $1$.
\begin{proof}
Projected strongly shattered dimension equals VC dimension: a set shattered by the class becomes a contained cube after projecting to that set.  The established VC dimension of the three-code class is one, proving this value.
\end{proof}

\subsection{Log Size of the Three-Code Class}\label{val:log_size_three_code_class}
\noindent\textbf{Statement.} The Log Size of Three-Code Class is $\log_2 3$.
\begin{proof}
The class has exactly the three distinct concepts $000,011,101$.  Its binary logarithmic size is therefore $\log_2 3$.
\end{proof}

\subsection{Log Shattering of the Three-Code Class}\label{val:log_shattering_three_code_class}
\noindent\textbf{Statement.} The Log Shattering of Three-Code Class is $2$.
\begin{proof}
The empty set and all three singleton coordinate sets are shattered.  No coordinate pair is shattered because every two-coordinate projection has only three traces.  Hence there are $1+3=4$ shattered sets and their binary logarithm is two.
\end{proof}

\subsection{Log Size Ratio of the Three-Code Class}\label{val:log_size_ratio_three_code_class}
\noindent\textbf{Statement.} The Log Size Ratio of Three-Code Class is $\frac12$.
\begin{proof}
The class has three concepts on a three-point domain.  Substitution into $\log_2(|\mathcal H|-1)/\log_2(n+1)$ gives $\log_2 2/\log_2 4=1/2$.
\end{proof}

\subsection{Maximum Largest Order Shattered Set of the Three-Code Class}\label{val:maximum_largest_order_shattered_set_three_code_class}
\noindent\textbf{Statement.} The Maximum Largest Order Shattered Set of Three-Code Class is $1$.
\begin{proof}
Maximum Largest Order Shattered Set equals VC dimension by definition.  The established VC dimension of the three-code class is one.
\end{proof}

\subsection{Projected VC Radius of the Three-Code Class}\label{val:projected_vc_radius_three_code_class}
\noindent\textbf{Statement.} The Projected VC Radius of Three-Code Class is $1$.
\begin{proof}
Projected VC radius equals VC dimension.  The established VC dimension of this class is one.
\end{proof}

\subsection{Monotonic Minimum Teaching Set Size of the Three-Code Class}\label{val:monotonic_minimum_teaching_set_size_three_code_class}
\noindent\textbf{Statement.} The Monotonic Minimum Teaching Set Size of Three-Code Class is $1$.
\begin{proof}
For finite classes, monotonic minimum teaching-set size equals recursive teaching dimension.  The established recursive teaching dimension of this class is one.
\end{proof}

\subsection{Minimum Teaching Set Size of the Three-Code Class}\label{val:minimum_teaching_set_size_three_code_class}
\noindent\textbf{Statement.} The Minimum Teaching Set Size of Three-Code Class is $1$.
\begin{proof}
The third coordinate with label zero identifies $000$ among the three concepts.  The empty sample is consistent with every concept, so teaching-set size zero is impossible.  Thus the minimum teaching-set size is one.
\end{proof}

\subsection{Hitting Size of the Three-Code Class}\label{val:hitting_size_three_code_class}
\noindent\textbf{Statement.} The Hitting Size of Three-Code Class is $1$.
\begin{proof}
Interpreting strings as supports, the two nonempty concepts are $\{2,3\}$ and $\{1,3\}$, which share their third coordinate.  That coordinate is a hitting set, while the empty set cannot hit either nonempty concept.  Hence the hitting size is one.
\end{proof}

\subsection{Relative Hitting Size of the Three-Code Class}\label{val:relative_hitting_size_three_code_class}
\noindent\textbf{Statement.} The Relative Hitting Size of Three-Code Class is $1$.
\begin{proof}
Relative hitting size equals the maximum teaching-set size.  The established teaching dimension of this class is one, so its relative hitting size is one.
\end{proof}

\subsection{Best Mistakes of the Three-Code Class}\label{val:best_mistakes_three_code_class}
\noindent\textbf{Statement.} The Best Mistakes of Three-Code Class is $1$.
\begin{proof}
Present the third coordinate first and predict one.  A negative answer makes one mistake and identifies $000$.  After a positive answer, query the first coordinate and predict zero: it is correct for $011$, and its one possible mistake identifies $101$.  Thus at most one mistake occurs.  Zero mistakes are impossible for a nonconstant class, so the value is one.
\end{proof}

\subsection{Self-Directed Queries Complexity of the Three-Code Class}\label{val:selfdirected_queries_complexity_three_code_class}
\noindent\textbf{Statement.} The Self-Directed Queries Complexity of Three-Code Class is $1$.
\begin{proof}
The established chain $\mathrm{RTD}\leq\mathrm{SDC}\leq\mathrm{LC}_{\mathrm{partial}}$ has both endpoint values equal to one for this class.  Therefore self-directed query complexity is exactly one.
\end{proof}

\subsection{Worst Mistakes of the Three-Code Class}\label{val:worst_mistakes_three_code_class}
\noindent\textbf{Statement.} The Worst Mistakes of Three-Code Class is $1$.
\begin{proof}
Worst mistakes is at least VC dimension, whose established value is one.  It is at most Littlestone dimension because the fixed-sequence class trees are among all Littlestone trees; the established Littlestone dimension is also one.  Hence worst mistakes equals one.
\end{proof}

\subsection{Majority Complexity of the Three-Code Class}\label{val:majority_complexity_three_code_class}
\noindent\textbf{Statement.} The Majority Complexity of Three-Code Class is $1$.
\begin{proof}
The pointwise majority hypothesis is $001$.  It differs from each of $000,011,101$ on exactly one coordinate, so a counterexample to this first majority query identifies the target.  One query therefore suffices, and zero cannot identify a three-concept class.
\end{proof}

\subsection{Combinatorial Eluder Dimension of the Three-Code Class}\label{val:combinatorial_eluder_dimension_three_code_class}
\noindent\textbf{Statement.} The Combinatorial Eluder Dimension of Three-Code Class is $2$.
\begin{proof}
Use $000$ as base.  First query the first coordinate with witness $101$; then query the second coordinate with witness $011$, which agrees with the base at the first coordinate.  This gives a length-two eluder sequence.  Such witnesses are distinct and cannot equal the base, so a class of three concepts admits no sequence longer than two.  Thus the value is two.
\end{proof}

\subsection{Dual VC Dimension of the Three-Code Class}\label{val:dual_vc_dimension_three_code_class}
\noindent\textbf{Statement.} The Dual VC Dimension of Three-Code Class is $1$.
\begin{proof}
The three dual concepts, viewed on the three original hypotheses, are the coordinate columns $001,010,011$.  Each of the two nonconstant original hypotheses is shattered as a one-point dual domain.  Any dual two-point restriction realizes only three traces, so no pair is shattered.  Hence dual VC dimension is one.
\end{proof}

\subsection{Dual Antipodal VC Dimension of the Three-Code Class}\label{val:dual_antipodal_vc_dimension_three_code_class}
\noindent\textbf{Statement.} The Dual Antipodal VC Dimension of Three-Code Class is $2$.
\begin{proof}
For the two hypotheses $011,101$, the coordinate traces include $01$, $10$, and $11$; the last also represents its complement $00$.  Thus this pair is dually antipodally shattered.  On all three hypotheses, the realized traces and their complements are $000,011,101,111,100,010$, missing the complementary pair $001,110$.  Hence no triple is dually antipodally shattered and the value is two.
\end{proof}

\subsection{2-Block Littlestone Dimension of the Three-Code Class}\label{val:2block_littlestone_dimension_three_code_class}
\noindent\textbf{Statement.} The 2-Block Littlestone Dimension of Three-Code Class is $1$.
\begin{proof}
The block-Littlestone sandwich gives $\mathrm{VC}\leq\mathrm B_2\leq\mathrm L$.  Both established endpoint values are one for this class, so $\mathrm B_2=1$.
\end{proof}

\subsection{$k$-Block Littlestone Dimension of the Three-Code Class}\label{val:kblock_littlestone_dimension_three_code_class}
\noindent\textbf{Statement.} The $k$-Block Littlestone Dimension ($k>2$) of Three-Code Class is $1$.
\begin{proof}
For every fixed $k>2$, the block-Littlestone sandwich gives $\mathrm{VC}\leq\mathrm B_k\leq\mathrm L$. Both established endpoint values are one for this class, so $\mathrm B_k=1$.
\end{proof}

\subsection{Prefix-$k$ Littlestone Dimension of the Three-Code Class}\label{val:prefix_littlestone_dimension_three_code_class}
\noindent\textbf{Statement.} The Prefix-$k$ Littlestone Dimension ($k>1$) of Three-Code Class is $1$.
\begin{proof}
For every fixed $k>1$, a one-point VC witness is a prefix-$k$ tree of depth one, while every prefix-$k$ tree is an ordinary Littlestone tree. The established VC and Littlestone dimensions are both one, forcing $\mathrm{PL}_k=1$.
\end{proof}

\subsection{Positive Teaching Dimension of the Three-Code Class}\label{val:positive_teaching_dimension_three_code_class}
\noindent\textbf{Statement.} The Positive Teaching Dimension of Three-Code Class is $\infty$.
\begin{proof}
The concept $000$ has no positive examples.  Its only positive sample is therefore empty, but that sample is also consistent with $011$ and $101$.  Hence $000$ has no finite positive teaching set, so the positive teaching dimension is infinite.
\end{proof}

\subsection{Order Sample Compression of the Three-Code Class}\label{val:order_sample_compression_three_code_class}
\noindent\textbf{Statement.} The Order Sample Compression of Three-Code Class is $1$.
\begin{proof}
Use the decoder order $001<000<011<101$, allowing the missing hollow-star center $001$ as an improper reconstruction. A sample consistent with $001$ retains nothing. Otherwise retain one example disagreeing with $001$. Exactly one class member disagrees with $001$ at that coordinate, so realizability forces that member to realize the whole sample. The retained singleton therefore has the same first consistent hypothesis as the original sample. Width zero cannot reconstruct both distinct full samples $000$ and $011$, proving equality. The former proof used only the proper order and incorrectly assumed its proposed distinguishing coordinate was present in every partial sample. The corrected scheme is independently replayed by sync/verify\_proper\_compression\_covc.py.
\end{proof}
\noindent\textbf{Reference.} \cite{darnstadt2016order}.

\subsection{Proper Ordered Sample Compression of the Three-Code Class}\label{val:proper_ordered_sample_compression_three_code_class}
\noindent\textbf{Statement.} The Proper Ordered Sample Compression of Three-Code Class is $2$.
\begin{proof}
The co-VC value is three, so Proposition \ref{prop:proper-compression-covc} gives proper ordered compression at least two. Conversely, fix any order of the three hypotheses. For a sample whose first consistent hypothesis is $h$, choose from the sample one disagreement with each earlier hypothesis. There are at most two earlier hypotheses, and the retained examples preserve $h$ as the first consistent reconstruction; remove redundant examples if necessary. Thus width two suffices. The former width-one proof was invalid on partial samples: in the order $000<011<101$, the sample $x_0=1,x_1=0$ selects $101$ but does not contain the proposed key $x_2=1$; either of its one-example subsets selects an earlier hypothesis. All six proper orders are independently checked by sync/verify\_proper\_compression\_covc.py.
\end{proof}
\noindent\textbf{Reference.} \cite{darnstadt2016order}.

\subsection{Minimum Projected Teaching-Set Size of the Three-Code Class}\label{val:minimum_projected_teaching_set_size_three_code_class}
\noindent\textbf{Statement.} The Minimum Projected Teaching Set Size of Three-Code Class is $2$.
\begin{proof}
Each target requires two examples in a suitable two-coordinate projection.  Omitting $x_0$ gives the traces $000\mapsto00$, $011\mapsto01$, $101\mapsto10$, so $000$ is the middle corner and both coordinates are needed to teach it.  Omitting $x_1$ similarly makes $011$ the middle corner, and omitting $x_2$ makes $101$ the middle corner.  Thus every target has worst projected teaching-set size at least two.  With only three distinct traces, two coordinates always suffice to distinguish any target in every projection, so the minimum is two.
\end{proof}

\subsection{Monotonic Projected Minimum Teaching-Set Size of the Three-Code Class}\label{val:monotonic_projected_minimum_teaching_set_size_three_code_class}
\noindent\textbf{Statement.} The Monotonic Projected Minimum Teaching Set Size of Three-Code Class is $2$.
\begin{proof}
The three-concept subfamily itself has minimum projected teaching-set size two, so the monotonic maximum is at least two.  No subfamily has more than three distinct traces, and a target trace among at most three traces is always distinguished by at most two coordinates in any projection.  Hence every subfamily has minimum projected teaching-set size at most two, and the maximum is exactly two.
\end{proof}

\subsection{Minimum Largest Order Shattered Set of the Three-Code Class}\label{val:minimum_largest_order_shattered_set_three_code_class}
\noindent\textbf{Statement.} The Minimum Largest Order Shattered Set of Three-Code Class is $1$.
\begin{proof}
For every coordinate order, downshifting preserves the class cardinality three and yields a nontrivial downset.  Such a downset contains a one-cube, so its largest order-shattered set has size at least one.  A two-cube has four traces, which cannot occur in a three-trace shifted class.  Thus the largest order-shattered set has size exactly one for every order, and so does the minimum over orders.
\end{proof}

\subsection{Balbach Teaching Dimension of the Three-Code Class}\label{val:balbach_teaching_dimension_three_code_class}
\noindent\textbf{Statement.} The Balbach Teaching Dimension of Three-Code Class is $1$.
\begin{proof}
For this class, ordinary teaching dimension and recursive teaching dimension are both one.  The established finite-class sandwich $\mathrm{TD}\geq\mathrm{BTD}\geq\mathrm{RTD}$ therefore gives $\mathrm{BTD}=1$.
\end{proof}

\subsection{Occasional One-Sided Getback of the Three-Code Class}\label{val:occasional_onesided_getback_three_code_class}
\noindent\textbf{Statement.} The Occasional One-Sided Getback of Three-Code Class is $1$.
\begin{proof}
The established O1G bounds put unlabeled sample compression below O1G and Littlestone dimension above it.  Both endpoint values are one for the three-code class, so $\mathrm{M}_{\mathrm{o1g}}=1$.
\end{proof}

\subsection{Randomized Littlestone Dimension of the Three-Code Class}\label{val:randomized_littlestone_dimension_three_code_class}
\noindent\textbf{Statement.} The Randomized Littlestone Dimension of Three-Code Class is $\frac34$.
\begin{proof}
Query $x_0$ first: its zero branch identifies $000$, while its one branch contains $011,101$ and is separated by $x_1$.  This shattered tree has expected random-branch depth $3/2$, giving value $3/4$.  Conversely, every shattered tree has at most three leaves, and a binary tree with at most three leaves has expected random-branch depth at most $3/2$.  Hence $\mathrm{RL}=3/4$.
\end{proof}

\subsection{Subset Teaching Dimension of the Three-Code Class}\label{val:subset_teaching_dimension_three_code_class}
\noindent\textbf{Statement.} The Subset Teaching Dimension of Three-Code Class is $1$.
\begin{proof}
The ordinary one-point teaching samples $x_0=0$ for $000$, $x_1=1$ for $011$, and $x_2=1$ for $101$ are each contained in no ordinary minimum teaching sample of either other target.  They therefore remain subset-teaching samples at every iteration.  An empty sample identifies no target, so the stabilized maximum is one.
\end{proof}

\subsection{Interpolation Degree of the Three-Code Class}\label{val:interpolation_degree_three_code_class}
\noindent\textbf{Statement.} The Interpolation degree of Three-Code Class is $1$.
\begin{proof}
In coordinate order $(x_0,x_1,x_2)$ the codewords are $000,110,101$.  The evaluation vectors of $1,x_1,x_2$ on these three points are $(1,1,1),(0,1,0),(0,0,1)$, which are linearly independent.  Thus degree-one polynomials interpolate every real-valued function on the class.  Degree zero only gives constants, so the interpolation degree is one.
\end{proof}

\subsection{Extremal-Extension VC Dimension of the Three-Code Class}\label{val:extremal_extension_vc_dimension_three_code_class}
\noindent\textbf{Statement.} The Extremal-Extension VC Dimension of Three-Code Class is $1$.
\begin{proof}
The class has VC dimension one, giving the lower bound.  Add the codeword $001$ to obtain $\mathcal E=\{000,011,101,001\}$.  Every singleton coordinate is shattered, while each two-coordinate projection has only three traces: the new trace repeats an existing trace in every such projection.  Hence the shattered sets of $\mathcal E$ are precisely the empty set and its three singletons, four in total, equal to $|\mathcal E|$.  Thus $\mathcal E$ is extremal with VC dimension one and proves the matching upper bound.
\end{proof}

\subsection{List Replicability Number of the Three-Code Class}\label{val:list_replicability_three_code_class}
\noindent\textbf{Statement.} The List Replicability Number of Three-Code Class is $2$.
\begin{proof}
The three-code class has spherical dimension one.  The spherical-dimension list-replicability theorem therefore gives $\mathrm{LR}\geq\lceil(1+3)/2\rceil=2$.  Its sign rank is two, and the established sign-rank upper bound gives $\mathrm{LR}\leq2$.  Hence $\mathrm{LR}=2$.
\end{proof}
\noindent\textbf{Reference.} \cite{chornomaz2025spherical,blondal2026simplicialcovering}.

\subsection{Projected Membership Queries Complexity of the Three-Code Class}\label{val:projected_membership_queries_complexity_three_code_class}
\noindent\textbf{Statement.} The Projected Membership Queries Complexity of Three-Code Class is $2$.
\begin{proof}
The star number is two, so the established star lower bound gives $\mathrm Q_{\mathrm m}^p\geq2$.  Every projection has at most three distinct traces.  If it has three, query a coordinate that separates one trace from the other two and then, only on the two-trace branch, a coordinate that separates those; if it has fewer, one query suffices.  Thus every projection is learned with at most two membership queries, proving equality.
\end{proof}

\subsection{Radon Number of the Three-Code Class}\label{val:radon_number_three_code_class}
\noindent\textbf{Statement.} The Radon Number of Three-Code Class is $3$.
\begin{proof}
The whole three-concept class is Radon independent.  For each partition into a singleton and a pair, the pair agrees on one coordinate on which the singleton disagrees: $011,101$ agree at $x_0=1$; $000,101$ agree at $x_1=0$; and $000,011$ agree at $x_2=0$.  The pair's convex hull is consequently disjoint from the singleton's hull.  This checks every nontrivial partition, and no Radon-independent subset can have more than the three available concepts.
\end{proof}

\subsection{Self-Directed Queries Complexity of Half-Intervals}\label{val:selfdirected_queries_complexity_halfintervals}
\noindent\textbf{Statement.} The Self-Directed Queries Complexity of Half-intervals is $1$.
\begin{proof}
The half-interval class has VC dimension one.  Corollary 9(1) of Doliwa--Fan--Simon--Zilles states that every nontrivial VC-one class has self-directed learning complexity one.  The class is nontrivial, so the value is exactly one.
\end{proof}
\noindent\textbf{Reference.} \cite{doliwa2014recursive}.

\subsection{Self-Directed Queries Complexity of Tagged Singletons}\label{val:selfdirected_queries_complexity_tagged_singletons}
\noindent\textbf{Statement.} The Self-Directed Queries Complexity of Tagged Singletons is $1$.
\begin{proof}
Tagged singletons have VC dimension one.  Corollary 9(1) of Doliwa--Fan--Simon--Zilles gives self-directed learning complexity one for every nontrivial VC-one class.  Thus the value is one.
\end{proof}
\noindent\textbf{Reference.} \cite{doliwa2014recursive}.

\subsection{Order Sample Compression of Tagged Singletons}\label{val:order_sample_compression_tagged_singletons}
\noindent\textbf{Statement.} The Order Sample Compression of Tagged Singletons is $1$.
\begin{proof}
Tagged singletons have VC dimension one.  Corollary 4 of Darnstädt--Kiss--Simon--Zilles gives an order-compression scheme of size one for every VC-one class.  A nontrivial class cannot have order-compression size zero, because one fixed reconstruction cannot be consistent with two distinct full samples.
\end{proof}
\noindent\textbf{Reference.} \cite{darnstadt2016order}.

\subsection{Recursive Teaching Dimension of the Repetition-Free Teaching Separation Class}\label{val:recursive_teaching_dimension_repetitionfree_teaching_separation}
\noindent\textbf{Statement.} The Recursive Teaching Dimension of Repetition-Free Teaching Separation Class is $2$.
\begin{proof}
The source gives a recursive teaching plan of order two for this thirteen-concept class, and proves that its recursive teaching dimension is exactly two.
\end{proof}
\noindent\textbf{Reference.} \cite{doliwa2014recursive}.

\subsection{Repetition-Free Recursive Teaching Dimension of the Repetition-Free Teaching Separation Class}\label{val:repetitionfree_recursive_teaching_dimension_repetitionfree_teaching_separation}
\noindent\textbf{Statement.} The Repetition-Free Recursive Teaching Dimension of Repetition-Free Teaching Separation Class is $3$.
\begin{proof}
Doliwa--Fan--Simon--Zilles prove that every order-two teaching plan repeats the coordinate set $\{x_1,x_2\}$ for its first two concepts, while an order-three repetition-free plan exists.  Thus $\mathrm{rfRTD}=3$.
\end{proof}
\noindent\textbf{Reference.} \cite{doliwa2014recursive}.

\subsection{Average Teaching Dimension of the Full Cube}\label{val:average_teaching_dimension_full_cube}
\noindent\textbf{Statement.} The Average Teaching Dimension of Full Cube is $n$.
\begin{proof}
Every full-cube concept needs all $n$ coordinates in a teaching set, since changing an omitted coordinate gives another consistent cube concept. Thus every individual value, and hence their average, is $n$.
\end{proof}
\noindent\textbf{Reference.} \cite{doliwa2014recursive}.

\subsection{Average Teaching Dimension of Trivial Concepts}\label{val:average_teaching_dimension_trivial_concepts}
\noindent\textbf{Statement.} The Average Teaching Dimension of Trivial concepts is $1$.
\begin{proof}
One coordinate with the target's label distinguishes either of the two constant concepts, while the empty sample does not. Both individual minimum teaching-set sizes are one, so their average is one.
\end{proof}

\subsection{Average Teaching Dimension of Singletons}\label{val:average_teaching_dimension_singletons}
\noindent\textbf{Statement.} The Average Teaching Dimension of Singletons is $1$.
\begin{proof}
For every singleton, its unique positive coordinate is a teaching set of size one, and the empty sample identifies no target in the nontrivial class. Thus every individual value and their average equal one.
\end{proof}

\subsection{Average Teaching Dimension of Singletons plus Empty Set}\label{val:average_teaching_dimension_singletons_plus_empty_set}
\noindent\textbf{Statement.} The Average Teaching Dimension of Singletons plus empty set is $\frac{2n}{n+1}$.
\begin{proof}
Each of the $n$ singleton concepts is taught by its positive coordinate. The empty concept requires all $n$ negative coordinates to exclude every singleton. Averaging the total $n+n$ over the $n+1$ concepts gives $2n/(n+1)$.
\end{proof}

\subsection{Average Teaching Dimension of Half-Intervals}\label{val:average_teaching_dimension_halfintervals}
\noindent\textbf{Statement.} The Average Teaching Dimension of Half-intervals is $\begin{cases}0&n=1,\\2-\frac2n&n\ge2.\end{cases}$.
\begin{proof}
For $n\ge2$, the two endpoint initial intervals have one-coordinate teaching sets, while every one of the $n-2$ interior intervals needs and has the adjacent positive/negative pair. Their total is $2+2(n-2)=2n-2$, whose average is $2-2/n$. For $n=1$, the unique concept has the empty teaching set, so the value is zero.
\end{proof}

\subsection{Average Teaching Dimension of Addressing}\label{val:average_teaching_dimension_addressing}
\noindent\textbf{Statement.} The Average Teaching Dimension of Addressing is $1\quad(n\ge2)$.
\begin{proof}
Both the established minimum teaching-set size and teaching dimension of Addressing are one. Every individual minimum teaching-set size is therefore one, so their average is one.
\end{proof}
\noindent\textbf{Reference.} \cite{maass1992lower}.

\subsection{Average Teaching Dimension of the Three-Code Class}\label{val:average_teaching_dimension_three_code_class}
\noindent\textbf{Statement.} The Average Teaching Dimension of Three-Code Class is $1$.
\begin{proof}
The three-code class has both minimum teaching-set size and teaching dimension equal to one. Thus every concept has minimum teaching-set size one, and the average is one.
\end{proof}

\subsection{Average Teaching Dimension of the Universal Class}\label{val:average_teaching_dimension_universal_class}
\noindent\textbf{Statement.} The Average Teaching Dimension of Universal Class is $1$.
\begin{proof}
The established minimum teaching-set size and teaching dimension of the universal class are both one. Hence every individual minimum teaching-set size equals one, giving average teaching dimension one.
\end{proof}
\noindent\textbf{Reference.} \cite{shinohara1991teachability,goldman1995complexity}.

\subsection{Average Teaching Dimension of Warmuth's $C_5$ Class}\label{val:average_teaching_dimension_warmuth_c5}
\noindent\textbf{Statement.} The Average Teaching Dimension of Warmuth's $C_5$ class is $3$.
\begin{proof}
Warmuth's $C_5$ class has both established minimum teaching-set size and teaching dimension equal to three. Therefore every individual minimum teaching-set size is three, and so is their average.
\end{proof}
\noindent\textbf{Reference.} \cite{palvolgyi2021unlabeled}.

\subsection{Extended Teaching Dimension of Blockwise-Addressing}\label{val:extended_teaching_dimension_blockwiseaddressing}
\noindent\textbf{Statement.} The Extended Teaching Dimension of Blockwise-Addressing is $2^d\leq\mathrm{XTD}\leq2^{d+1}-1$.
\begin{proof}
For a finite binary class, $\mathrm{XTD}=\mathrm S$ by Hanneke--Yang's Theorem 13.  Apply this equality to the established Blockwise-Addressing star-number interval $2^d\leq\mathrm S\leq2^{d+1}-1$.
\end{proof}
\noindent\textbf{Reference.} \cite{hanneke2015minimax,ben1997online,bousquet2020proper}.

\subsection{Projected Maximal Degree of Blockwise-Addressing}\label{val:projected_maximal_degree_blockwiseaddressing}
\noindent\textbf{Statement.} The Projected Maximal Degree of Blockwise-Addressing is $2^d\leq\mathrm d^p_{\max}\leq2^{d+1}-1$.
\begin{proof}
The projected maximal degree equals the star number by Lemma~\ref{lem:star-projected-teaching}.  Substitute the established Blockwise-Addressing star-number interval $2^d\leq\mathrm S\leq2^{d+1}-1$.
\end{proof}
\noindent\textbf{Reference.} \cite{hanneke2015minimax,ben1997online,bousquet2020proper}.

\subsection{Maximum Projected Teaching-Set Size of Blockwise-Addressing}\label{val:maximum_projected_teaching_set_size_blockwiseaddressing}
\noindent\textbf{Statement.} The Maximum Projected Teaching Set Size of Blockwise-Addressing is $2^d\leq\mathrm{TS}_{\max}^p\leq2^{d+1}-1$.
\begin{proof}
The maximum projected teaching-set size equals the star number by Lemma~\ref{lem:star-projected-teaching}.  Substitute the established Blockwise-Addressing star-number interval $2^d\leq\mathrm S\leq2^{d+1}-1$.
\end{proof}
\noindent\textbf{Reference.} \cite{hanneke2015minimax,ben1997online,bousquet2020proper}.

\subsection{Relative Hitting Size of Blockwise-Addressing}\label{val:relative_hitting_size_blockwiseaddressing}
\noindent\textbf{Statement.} The Relative Hitting Size of Blockwise-Addressing is $2$.
\begin{proof}
Relative hitting size is exactly teaching dimension by the standard teaching-set / relative-hitting-set equivalence.  The direct Blockwise-Addressing teaching calculation gives $\mathrm{TD}=2$.
\end{proof}
\noindent\textbf{Reference.} \cite{goldman1995complexity,shinohara1991teachability}.

\subsection{Preference-Based Teaching Dimension of the Repetition-Free Teaching Separation Class}\label{val:preferencebased_teaching_dimension_repetitionfree_teaching_separation}
\noindent\textbf{Statement.} The Preference-Based Teaching Dimension of Repetition-Free Teaching Separation Class is $2$.
\begin{proof}
For finite classes, preference-based teaching dimension equals recursive teaching dimension.  The cited separation class has $\mathrm{RTD}=2$.
\end{proof}
\noindent\textbf{Reference.} \cite{gao2017preference,doliwa2014recursive}.

\subsection{Monotonic Minimum Teaching-Set Size of the Repetition-Free Teaching Separation Class}\label{val:monotonic_minimum_teaching_set_size_repetitionfree_teaching_separation}
\noindent\textbf{Statement.} The Monotonic Minimum Teaching Set Size of Repetition-Free Teaching Separation Class is $2$.
\begin{proof}
For finite classes, the max-over-subclasses minimum teaching-set size is exactly recursive teaching dimension.  The cited separation class has $\mathrm{RTD}=2$.
\end{proof}
\noindent\textbf{Reference.} \cite{doliwa2014recursive}.

\subsection{Average Teaching Dimension of Blockwise-Addressing}\label{val:average_teaching_dimension_blockwiseaddressing}
\noindent\textbf{Statement.} The Average Teaching Dimension of Blockwise-Addressing is $2$.
\begin{proof}
Each $f_i$ is taught by $(z,0),(x_i,1)$ and each $g_i$ by $(z,1),(y_i,1)$.  No one-coordinate sample identifies any member of either family: at $z$ it leaves all members of that family, at a private coordinate it agrees with a member of the opposite family, and at a block coordinate its label occurs for multiple members of the same family.  Thus every individual minimum teaching-set size is two, and their average is two.
\end{proof}
\noindent\textbf{Reference.} \cite{ben1997online,kushilevitz1996witness}.

\subsection{Average Teaching Dimension of the Repetition-Free Teaching Separation Class}\label{val:average_teaching_dimension_repetitionfree_teaching_separation}
\noindent\textbf{Statement.} The Average Teaching Dimension of Repetition-Free Teaching Separation Class is $\frac{40}{13}$.
\begin{proof}
For the thirteen displayed strings in the class definition, exhaustive inspection of coordinate subsets gives the individual minimum teaching-set sizes $2,2,4,4,3,3,3,3,3,4,3,3,3$.  Their sum is $40$, so the average is $40/13$; the enumeration is checked by $\texttt{sync/verify\_repetitionfree\_teaching.py}$.
\end{proof}
\noindent\textbf{Reference.} \cite{doliwa2014recursive,kushilevitz1996witness}.

\subsection{Star Number of the Repetition-Free Teaching Separation Class}\label{val:star_number_repetitionfree_teaching_separation}
\noindent\textbf{Statement.} The Star number of Repetition-Free Teaching Separation Class is $4$.
\begin{proof}
Enumerating every coordinate projection and every realizable centre trace gives a largest centred star of dimension four; no five-star is possible on the five-coordinate domain. The exhaustive check is in $\texttt{sync/verify\_repetitionfree\_teaching.py}$.
\end{proof}
\noindent\textbf{Reference.} \cite{hanneke2015minimax,doliwa2014recursive}.

\subsection{co-VC Dimension of the Repetition-Free Teaching Separation Class}\label{val:covc_dimension_repetitionfree_teaching_separation}
\noindent\textbf{Statement.} The co-VC dimension of Repetition-Free Teaching Separation Class is $4$.
\begin{proof}
Enumerating coordinate projections and unrealizable centre traces gives a hollow star of dimension four, and no larger one on the five-coordinate domain. The exhaustive check is in $\texttt{sync/verify\_repetitionfree\_teaching.py}$.
\end{proof}
\noindent\textbf{Reference.} \cite{bousquet2020proper,doliwa2014recursive}.

\subsection{Minimum Star Number of the Repetition-Free Teaching Separation Class}\label{val:minimum_star_number_repetitionfree_teaching_separation}
\noindent\textbf{Statement.} The Minimum Star Number of Repetition-Free Teaching Separation Class is $3$.
\begin{proof}
For every one of the $32$ possible centre labelings, enumerate the coordinate projections on which that centre is realized and the corresponding stars. The least resulting maximum is three, while the enumeration supplies a four-star for some centres. The exhaustive check is in $\texttt{sync/verify\_repetitionfree\_teaching.py}$.
\end{proof}
\noindent\textbf{Reference.} \cite{hanneke2024star_eluder,doliwa2014recursive}.

\subsection{Extended Teaching Dimension of the Repetition-Free Teaching Separation Class}\label{val:extended_teaching_dimension_repetitionfree_teaching_separation}
\noindent\textbf{Statement.} The Extended Teaching Dimension of Repetition-Free Teaching Separation Class is $4$.
\begin{proof}
For finite binary classes, extended teaching dimension equals the star number. The directly enumerated star number of this class is four, so $\mathrm{XTD}=4$.
\end{proof}
\noindent\textbf{Reference.} \cite{hanneke2015minimax,doliwa2014recursive}.

\subsection{Projected Maximal Degree of the Repetition-Free Teaching Separation Class}\label{val:projected_maximal_degree_repetitionfree_teaching_separation}
\noindent\textbf{Statement.} The Projected Maximal Degree of Repetition-Free Teaching Separation Class is $4$.
\begin{proof}
Projected maximal degree equals the star number. The direct exhaustive calculation for this class gives star number four, hence projected maximal degree four.
\end{proof}
\noindent\textbf{Reference.} \cite{hanneke2015minimax,doliwa2014recursive}.

\subsection{Maximum Projected Teaching-Set Size of the Repetition-Free Teaching Separation Class}\label{val:maximum_projected_teaching_set_size_repetitionfree_teaching_separation}
\noindent\textbf{Statement.} The Maximum Projected Teaching Set Size of Repetition-Free Teaching Separation Class is $4$.
\begin{proof}
For finite concept classes, maximum projected teaching-set size equals the star number. The direct exhaustive calculation gives star number four for this class.
\end{proof}
\noindent\textbf{Reference.} \cite{hanneke2015minimax,doliwa2014recursive}.

\subsection{Size of the Repetition-Free Teaching Separation Class}\label{val:size_repetitionfree_teaching_separation}
\noindent\textbf{Statement.} The Size of Repetition-Free Teaching Separation Class is $13$.
\begin{proof}
The displayed definition lists thirteen distinct binary strings; this is also checked in $\texttt{sync/verify\_repetitionfree\_teaching.py}$.
\end{proof}
\noindent\textbf{Reference.} \cite{doliwa2014recursive}.

\subsection{Effective Range of the Repetition-Free Teaching Separation Class}\label{val:effective_range_repetitionfree_teaching_separation}
\noindent\textbf{Statement.} The Effective Range of Repetition-Free Teaching Separation Class is $5$.
\begin{proof}
Each of the five displayed coordinates takes both labels among the thirteen strings, so all five are effective. The direct check is in $\texttt{sync/verify\_repetitionfree\_teaching.py}$.
\end{proof}
\noindent\textbf{Reference.} \cite{doliwa2014recursive}.

\subsection{Log Size of the Repetition-Free Teaching Separation Class}\label{val:log_size_repetitionfree_teaching_separation}
\noindent\textbf{Statement.} The Log Size of Repetition-Free Teaching Separation Class is $\log_2 13$.
\begin{proof}
The class has thirteen concepts, so by definition its binary log-size is $\log_2 13$.
\end{proof}
\noindent\textbf{Reference.} \cite{doliwa2014recursive}.

\subsection{VC Dimension of the Repetition-Free Teaching Separation Class}\label{val:vc_dimension_repetitionfree_teaching_separation}
\noindent\textbf{Statement.} The VC Dimension of Repetition-Free Teaching Separation Class is $3$.
\begin{proof}
Enumerating coordinate subsets finds a shattered three-set and no shattered four-set. The exhaustive check is in $\texttt{sync/verify\_repetitionfree\_teaching.py}$.
\end{proof}
\noindent\textbf{Reference.} \cite{doliwa2014recursive}.

\subsection{Littlestone Dimension of the Repetition-Free Teaching Separation Class}\label{val:littlestone_dimension_repetitionfree_teaching_separation}
\noindent\textbf{Statement.} The Littlestone Dimension of Repetition-Free Teaching Separation Class is $3$.
\begin{proof}
The standard restriction recurrence $L(\mathcal H)=\max_x(1+\min\{L(\mathcal H_{x=0}),L(\mathcal H_{x=1})\})$ evaluates to three on the displayed finite class. The recursive computation is checked in $\texttt{sync/verify\_repetitionfree\_teaching.py}$.
\end{proof}
\noindent\textbf{Reference.} \cite{littlestone1988learning,doliwa2014recursive}.

\subsection{Diameter of the Repetition-Free Teaching Separation Class}\label{val:diameter_repetitionfree_teaching_separation}
\noindent\textbf{Statement.} The Diameter of Repetition-Free Teaching Separation Class is $5$.
\begin{proof}
The strings $10101$ and $01010$ differ in all five coordinates, and no distance can exceed the five-coordinate domain size. The enumeration is checked in $\texttt{sync/verify\_repetitionfree\_teaching.py}$.
\end{proof}
\noindent\textbf{Reference.} \cite{doliwa2014recursive}.

\subsection{Hamming Radius of the Repetition-Free Teaching Separation Class}\label{val:hamming_radius_repetitionfree_teaching_separation}
\noindent\textbf{Statement.} The Hamming Radius of Repetition-Free Teaching Separation Class is $4$.
\begin{proof}
Enumerating all $32$ possible binary centres gives minimum enclosing radius four. The exhaustive check is in $\texttt{sync/verify\_repetitionfree\_teaching.py}$.
\end{proof}
\noindent\textbf{Reference.} \cite{doliwa2014recursive}.

\subsection{Minimum Teaching-Set Size of the Repetition-Free Teaching Separation Class}\label{val:minimum_teaching_set_size_repetitionfree_teaching_separation}
\noindent\textbf{Statement.} The Minimum Teaching Set Size of Repetition-Free Teaching Separation Class is $2$.
\begin{proof}
Exhaustive inspection of coordinate subsets gives individual minimum teaching-set sizes $2,2,4,4,3,3,3,3,3,4,3,3,3$. Their minimum is two; the calculation is checked in $\texttt{sync/verify\_repetitionfree\_teaching.py}$.
\end{proof}
\noindent\textbf{Reference.} \cite{doliwa2014recursive,kushilevitz1996witness}.

\subsection{Teaching Dimension of the Repetition-Free Teaching Separation Class}\label{val:teaching_dimension_repetitionfree_teaching_separation}
\noindent\textbf{Statement.} The Teaching Dimension of Repetition-Free Teaching Separation Class is $4$.
\begin{proof}
The individual minimum teaching-set sizes are $2,2,4,4,3,3,3,3,3,4,3,3,3$. Teaching dimension is their maximum, hence four; the exhaustive calculation is checked in $\texttt{sync/verify\_repetitionfree\_teaching.py}$.
\end{proof}
\noindent\textbf{Reference.} \cite{doliwa2014recursive,kushilevitz1996witness}.

\subsection{Maximum Teaching-Set Size of the Repetition-Free Teaching Separation Class}\label{val:maximum_teaching_set_size_repetitionfree_teaching_separation}
\noindent\textbf{Statement.} The Maximum Teaching Set Size of Repetition-Free Teaching Separation Class is $4$.
\begin{proof}
Maximum teaching-set size is the teaching dimension. The direct enumeration of this class's individual minimum teaching-set sizes has maximum four; it is checked in $\texttt{sync/verify\_repetitionfree\_teaching.py}$.
\end{proof}
\noindent\textbf{Reference.} \cite{doliwa2014recursive,kushilevitz1996witness}.

\subsection{Minimum Degree of the Repetition-Free Teaching Separation Class}\label{val:minimum_degree_repetitionfree_teaching_separation}
\noindent\textbf{Statement.} The Minimum Degree of Repetition-Free Teaching Separation Class is $1$.
\begin{proof}
The one-inclusion graph has degree sequence $1,1,3,3,2,2,1,3,2,2,1,2,1$, so its minimum degree is one. The enumeration is checked in $\texttt{sync/verify\_repetitionfree\_teaching.py}$.
\end{proof}
\noindent\textbf{Reference.} \cite{doliwa2014recursive,haussler1988predicting}.

\subsection{Maximum Degree of the Repetition-Free Teaching Separation Class}\label{val:maximum_degree_repetitionfree_teaching_separation}
\noindent\textbf{Statement.} The Maximum Degree of Repetition-Free Teaching Separation Class is $3$.
\begin{proof}
The one-inclusion graph has degree sequence $1,1,3,3,2,2,1,3,2,2,1,2,1$, so its maximum degree is three. The enumeration is checked in $\texttt{sync/verify\_repetitionfree\_teaching.py}$.
\end{proof}
\noindent\textbf{Reference.} \cite{doliwa2014recursive,haussler1988predicting}.

\subsection{Double Density of the Repetition-Free Teaching Separation Class}\label{val:double_density_or_average_degree_repetitionfree_teaching_separation}
\noindent\textbf{Statement.} The Double Density or Average Degree of Repetition-Free Teaching Separation Class is $\frac{24}{13}$.
\begin{proof}
The one-inclusion graph has thirteen vertices and twelve edges, hence average degree $2\cdot12/13=24/13$. The edge enumeration is checked in $\texttt{sync/verify\_repetitionfree\_teaching.py}$.
\end{proof}
\noindent\textbf{Reference.} \cite{doliwa2014recursive,haussler1988predicting}.

\subsection{Densest Subgraph (twice) of the Repetition-Free Teaching Separation Class}\label{val:densest_subgraph_twice_repetitionfree_teaching_separation}
\noindent\textbf{Statement.} The Densest Subgraph (twice) of Repetition-Free Teaching Separation Class is $2$.
\begin{proof}
The four concepts $01010,01011,01110,01111$ form a two-cube and thus have average degree two. Exhaustive inspection of all induced subgraphs finds none with larger average degree; this is checked in $\texttt{sync/verify\_repetitionfree\_teaching.py}$.
\end{proof}
\noindent\textbf{Reference.} \cite{doliwa2014recursive,haussler1988predicting}.

\subsection{Degeneracy of the Repetition-Free Teaching Separation Class}\label{val:degeneracy_repetitionfree_teaching_separation}
\noindent\textbf{Statement.} The Degeneracy of Repetition-Free Teaching Separation Class is $2$.
\begin{proof}
The displayed two-cube has minimum degree two. Exhaustive inspection of every induced subgraph finds no subgraph with minimum degree above two, so the degeneracy is two; this is checked in $\texttt{sync/verify\_repetitionfree\_teaching.py}$.
\end{proof}
\noindent\textbf{Reference.} \cite{doliwa2014recursive,haussler1988predicting}.

\subsection{Largest Strongly Shattered Set of the Repetition-Free Teaching Separation Class}\label{val:largest_strongly_shattered_set_repetitionfree_teaching_separation}
\noindent\textbf{Statement.} The Largest Strongly Shattered Set of Repetition-Free Teaching Separation Class is $2$.
\begin{proof}
Fixing coordinates $0,1,3$ to $0,1,1$ leaves all four labelings on coordinates $2,4$, giving a contained two-cube. Exhaustive inspection finds no contained three-cube; this is checked in $\texttt{sync/verify\_repetitionfree\_teaching.py}$.
\end{proof}
\noindent\textbf{Reference.} \cite{doliwa2014recursive,bollobas1995defect}.

\subsection{Projected Minimum Degree of the Repetition-Free Teaching Separation Class}\label{val:projected_minimum_degree_repetitionfree_teaching_separation}
\noindent\textbf{Statement.} The Projected Minimum Degree of Repetition-Free Teaching Separation Class is $3$.
\begin{proof}
Projecting to coordinates $0,2,3$ realizes the full three-cube, of minimum degree three. Exhaustive inspection of all nonempty coordinate projections finds none with larger minimum degree; this is checked in $\texttt{sync/verify\_repetitionfree\_teaching.py}$.
\end{proof}
\noindent\textbf{Reference.} \cite{doliwa2014recursive,rubinstein2007shifting}.

\subsection{Projected Double Density of the Repetition-Free Teaching Separation Class}\label{val:projected_double_density_or_projected_average_degree_repetitionfree_teaching_separation}
\noindent\textbf{Statement.} The Projected Double Density or Projected Average Degree of Repetition-Free Teaching Separation Class is $3$.
\begin{proof}
The projection to coordinates $0,2,3$ is the full three-cube and has average degree three. Exhaustive inspection of all nonempty coordinate projections finds no larger average degree; this is checked in $\texttt{sync/verify\_repetitionfree\_teaching.py}$.
\end{proof}
\noindent\textbf{Reference.} \cite{doliwa2014recursive,haussler1988predicting}.

\subsection{Largest Shattered Set of the Repetition-Free Teaching Separation Class}\label{val:largest_shattered_set_repetitionfree_teaching_separation}
\noindent\textbf{Statement.} The Largest Shattered Set of Repetition-Free Teaching Separation Class is $3$.
\begin{proof}
Largest shattered set is the VC dimension. The directly verified VC dimension of this class is three.
\end{proof}
\noindent\textbf{Reference.} \cite{vapnik1971uniform,doliwa2014recursive}.

\subsection{Projected VC Radius of the Repetition-Free Teaching Separation Class}\label{val:projected_vc_radius_repetitionfree_teaching_separation}
\noindent\textbf{Statement.} The Projected VC Radius of Repetition-Free Teaching Separation Class is $3$.
\begin{proof}
Projected VC radius equals VC dimension. The directly verified VC dimension of this class is three.
\end{proof}
\noindent\textbf{Reference.} \cite{bollobas1995defect,doliwa2014recursive}.

\subsection{Projected Strongly Shattered Dimension of the Repetition-Free Teaching Separation Class}\label{val:projected_strongly_shattered_dimension_repetitionfree_teaching_separation}
\noindent\textbf{Statement.} The Projected Strongly Shattered Dimension of Repetition-Free Teaching Separation Class is $3$.
\begin{proof}
Projected strongly shattered dimension equals VC dimension. The directly verified VC dimension of this class is three.
\end{proof}
\noindent\textbf{Reference.} \cite{bollobas1995defect,doliwa2014recursive}.

\subsection{Maximum Largest Order Shattered Set of the Repetition-Free Teaching Separation Class}\label{val:maximum_largest_order_shattered_set_repetitionfree_teaching_separation}
\noindent\textbf{Statement.} The Maximum Largest Order Shattered Set of Repetition-Free Teaching Separation Class is $3$.
\begin{proof}
Maximum largest order-shattered-set size equals VC dimension. The directly verified VC dimension of this class is three.
\end{proof}
\noindent\textbf{Reference.} \cite{bollobas1995defect,doliwa2014recursive}.

\subsection{Effective VC Radius of the Repetition-Free Teaching Separation Class}\label{val:effective_vc_radius_repetitionfree_teaching_separation}
\noindent\textbf{Statement.} The Effective VC Radius of Repetition-Free Teaching Separation Class is $2$.
\begin{proof}
Every pair of the five effective coordinates is shattered. The three-coordinate set $\{0,1,2\}$ realizes only six traces, so it is not shattered. Thus the all-effective-subsets shattering radius is two; this is checked in $\texttt{sync/verify\_repetitionfree\_teaching.py}$.
\end{proof}
\noindent\textbf{Reference.} \cite{doliwa2014recursive,vapnik1971uniform}.

\subsection{Minimum Largest Order Shattered Set of the Repetition-Free Teaching Separation Class}\label{val:minimum_largest_order_shattered_set_repetitionfree_teaching_separation}
\noindent\textbf{Statement.} The Minimum Largest Order Shattered Set of Repetition-Free Teaching Separation Class is $2$.
\begin{proof}
The effective VC radius is two, so every coordinate order has an order-shattered two-set. Standard downshifting in the coordinate order $(4,2,0,1,3)$ produces a downset whose largest face has size two. Thus the minimum over orders is exactly two; all 120 orders are checked in $\texttt{sync/verify\_repetitionfree\_teaching.py}$.
\end{proof}
\noindent\textbf{Reference.} \cite{bollobas1995defect,doliwa2014recursive}.

\subsection{Interpolation Degree of the Repetition-Free Teaching Separation Class}\label{val:interpolation_degree_repetitionfree_teaching_separation}
\noindent\textbf{Statement.} The Interpolation degree of Repetition-Free Teaching Separation Class is $2$.
\begin{proof}
The class contains a two-cube, so interpolation degree is at least its largest strongly shattered-set size, namely two. Its minimum largest order-shattered-set size is two, and this upper-bounds interpolation degree. Hence interpolation degree is exactly two.
\end{proof}
\noindent\textbf{Reference.} \cite{moran2015shattered,bollobas1995defect,doliwa2014recursive}.

\subsection{Interpolation Degree of the Interpolation-Conjecture Counterexample}\label{val:interpolation_degree_interpolation_conjecture_counterexample}
\noindent\textbf{Statement.} The Interpolation degree of Interpolation-Conjecture Counterexample is $1$.
\begin{proof}
For the concepts $001,010,011,100$, the evaluation vectors of $1,x_1,x_2,x_3$ are linearly independent, so degree-one polynomials interpolate every function. Degree zero gives only constants, hence the interpolation degree is one. This is checked by $\texttt{sync/verify\_interpolation\_conjecture.py}$.
\end{proof}

\subsection{Shattering-Cover Projection VC Dimension of the Interpolation-Conjecture Counterexample}\label{val:shattering_cover_projection_vc_dimension_interpolation_conjecture_counterexample}
\noindent\textbf{Statement.} The Shattering-Cover Projection VC Dimension of Interpolation-Conjecture Counterexample is $2$.
\begin{proof}
The projection to $\{x_2,x_3\}$ is the full two-cube, with four shattered subsets, and is therefore admissible. It has VC dimension two. Every projection to at most one coordinate has at most two shattered subsets and fails the admissibility condition $|\operatorname{sh}(\mathcal H_{|S})|\ge4$. Thus the minimum admissible VC dimension is two. This is checked by $\texttt{sync/verify\_interpolation\_conjecture.py}$.
\end{proof}

\subsection{Shattering-Cover Projection VC Dimension of Trivial Concepts}\label{val:shattering_cover_projection_vc_dimension_trivial_concepts}
\noindent\textbf{Statement.} The Shattering-Cover Projection VC Dimension of Trivial concepts is $1$.
\begin{proof}
A one-coordinate projection is the full one-cube, so it is admissible and has VC dimension one. The empty projection has one shattered subset, fewer than the two concepts in the class; thus no admissible projection has VC dimension zero. This is checked by $\texttt{sync/verify\_interpolation\_conjecture.py}$.
\end{proof}

\subsection{Shattering-Cover Projection VC Dimension of Singletons}\label{val:shattering_cover_projection_vc_dimension_singletons}
\noindent\textbf{Statement.} The Shattering-Cover Projection VC Dimension of Singletons is $1$.
\begin{proof}
Project onto any $n-1$ coordinates. The traces are the empty set and their $n-1$ singletons, so all $n$ concepts remain distinguished and exactly the empty set and the singletons are shattered. This projection is admissible and has VC dimension one. A VC-zero projection has one trace and only one shattered subset, hence cannot be admissible for a class of size $n\geq2$. This is checked by $\texttt{sync/verify\_interpolation\_conjecture.py}$.
\end{proof}

\subsection{Shattering-Cover Projection VC Dimension of Full Cube}\label{val:shattering_cover_projection_vc_dimension_full_cube}
\noindent\textbf{Statement.} The Shattering-Cover Projection VC Dimension of Full Cube is $n$.
\begin{proof}
A projection to $m$ coordinates is the full $m$-cube, with $2^m$ shattered subsets. To meet the admissibility requirement $2^m\geq|\mathcal H|=2^n$ requires $m=n$. The only admissible projection is therefore the full cube, whose VC dimension is $n$. This is checked by $\texttt{sync/verify\_interpolation\_conjecture.py}$.
\end{proof}

\subsection{Shattering-Cover Projection VC Dimension of Singletons plus Empty Set}\label{val:shattering_cover_projection_vc_dimension_singletons_plus_empty_set}
\noindent\textbf{Statement.} The Shattering-Cover Projection VC Dimension of Singletons plus empty set is $1$.
\begin{proof}
The full class consists of the empty set and all $n$ singletons. Its shattered subsets are exactly the empty set and the $n$ singleton coordinates, so it is itself admissible and has VC dimension one. A VC-zero projection has only one shattered subset and cannot be admissible because the class has at least two concepts. This is checked by $\texttt{sync/verify\_interpolation\_conjecture.py}$.
\end{proof}

\subsection{Shattering-Cover Projection VC Dimension of Half-intervals}\label{val:shattering_cover_projection_vc_dimension_halfintervals}
\noindent\textbf{Statement.} The Shattering-Cover Projection VC Dimension of Half-intervals is $1$.
\begin{proof}
The full threshold chain has exactly the empty set and its $n$ singleton coordinates shattered, so it is admissible and has VC dimension one. No projection of a threshold chain has VC dimension above one. A VC-zero projection has one shattered subset and is not admissible for this $n$-concept class. Thus the minimum admissible VC dimension is one. This is checked by $\texttt{sync/verify\_interpolation\_conjecture.py}$.
\end{proof}

\subsection{VC Dimension of the SCVC/VC Separation Class}\label{val:vc_dimension_scvc_vc_separation}
\noindent\textbf{Statement.} The VC Dimension of SCVC/VC Separation Class is $2$.
\begin{proof}
On $\{x_2,z\}$, the four concepts give the traces $00,01,10,11$, so VC dimension is at least two. Every three-coordinate projection has at most four traces, whereas shattering a three-set requires eight; hence VC dimension is exactly two. This is checked by $\texttt{sync/verify\_interpolation\_conjecture.py}$.
\end{proof}

\subsection{Shattering-Cover Projection VC Dimension of the SCVC/VC Separation Class}\label{val:shattering_cover_projection_vc_dimension_scvc_vc_separation}
\noindent\textbf{Statement.} The Shattering-Cover Projection VC Dimension of SCVC/VC Separation Class is $1$.
\begin{proof}
Projection to $\{x_1,x_2,x_3\}$ gives the four-element threshold chain $000,100,110,111$. It has exactly the empty set and its three singleton coordinates shattered, hence four shattered subsets and is admissible for this four-concept class; its VC dimension is one. An admissible projection has at least two traces and therefore cannot have VC dimension zero. Thus $\mathrm{SCVC}=1$. This is checked by $\texttt{sync/verify\_interpolation\_conjecture.py}$.
\end{proof}

\subsection{Interpolation Degree of the SCVC/Interpolation Reverse Separation Class}\label{val:interpolation_degree_scvc_interpolation_reverse_separation}
\noindent\textbf{Statement.} The Interpolation degree of SCVC/Interpolation Reverse Separation Class is $2$.
\begin{proof}
The evaluation matrix of $1,x_1,x_2,x_3,z$ on the four concepts has rank three, so degree one cannot interpolate every function. Including degree-two monomials has rank four, so degree two does suffice. This exact rank calculation is checked by $\texttt{sync/verify\_interpolation\_conjecture.py}$.
\end{proof}

\subsection{Shattering-Cover Projection VC Dimension of the SCVC/Interpolation Reverse Separation Class}\label{val:shattering_cover_projection_vc_dimension_scvc_interpolation_reverse_separation}
\noindent\textbf{Statement.} The Shattering-Cover Projection VC Dimension of SCVC/Interpolation Reverse Separation Class is $1$.
\begin{proof}
Project to $\{x_1,x_2,x_3\}$; the trace is $\{011,100\}$, which has VC dimension one and exactly its empty set and three singleton coordinates shattered. Thus the projection has four shattered subsets, matching the four original concepts. A VC-dimension-zero projection has only one shattered subset, so $\mathrm{SCVC}=1$. This is checked by $\texttt{sync/verify\_interpolation\_conjecture.py}$.
\end{proof}

\subsection{Interpolation Degree of the Interpolation/Effective-VC-Radius Counterexample}\label{val:interpolation_degree_interpolation_effective_vc_radius_counterexample}
\noindent\textbf{Statement.} The Interpolation degree of Interpolation/Effective-VC-Radius Counterexample is $1$.
\begin{proof}
On $000,011,101,110$, the four evaluation vectors of $1,x_1,x_2,x_3$ are linearly independent, so affine polynomials interpolate every function. Degree zero gives only constants. Hence $\mathrm{intdeg}=1$, as checked by $\texttt{sync/verify\_interpolation\_conjecture.py}$.
\end{proof}

\subsection{Effective VC Radius of the Interpolation/Effective-VC-Radius Counterexample}\label{val:effective_vc_radius_interpolation_effective_vc_radius_counterexample}
\noindent\textbf{Statement.} The Effective VC Radius of Interpolation/Effective-VC-Radius Counterexample is $2$.
\begin{proof}
All three coordinates are effective, and every two-coordinate projection is the full two-cube, so every pair is shattered. The full three-coordinate trace has only four concepts and is not shattered. Thus $\mathrm{VCR}=2$, as checked by $\texttt{sync/verify\_interpolation\_conjecture.py}$.
\end{proof}

\subsection{Minimum Largest Order-Shattered Set of the Order-Shattering/Effective-VC-Radius Separation Class}\label{val:minimum_largest_order_shattered_set_order_shattering_effective_vc_radius_separation}
\noindent\textbf{Statement.} The Minimum Largest Order Shattered Set of Order-Shattering/Effective-VC-Radius Separation Class is $2$.
\begin{proof}
For each of the six coordinate orders, standard downshifting produces a downset with a two-element face; hence every order has an order-shattered two-set. The class has five concepts, so no shifted class can contain a three-cube. Thus the minimum largest order-shattered-set size is two. All orders are checked by $\texttt{sync/verify\_order\_shattering\_separation.py}$.
\end{proof}

\subsection{Effective VC Radius of the Order-Shattering/Effective-VC-Radius Separation Class}\label{val:effective_vc_radius_order_shattering_effective_vc_radius_separation}
\noindent\textbf{Statement.} The Effective VC Radius of Order-Shattering/Effective-VC-Radius Separation Class is $1$.
\begin{proof}
All three coordinates are effective and hence individually shattered. On $\{x_1,x_2\}$ the class realizes only $00,01,10$, so not every pair is shattered. Therefore the effective VC radius is one; this is checked by $\texttt{sync/verify\_order\_shattering\_separation.py}$.
\end{proof}

\subsection{Repetition-Free Recursive Teaching Dimension of Trivial Concepts}\label{val:repetitionfree_recursive_teaching_dimension_trivial_concepts}
\noindent\textbf{Statement.} The Repetition-Free Recursive Teaching Dimension of Trivial concepts is $1$.
\begin{proof}
Choose either constant concept first and teach it at one coordinate.  The remaining concept is then the sole candidate and uses the empty teaching set.  These coordinate sets are distinct, so this is a repetition-free plan of order one.  No nontrivial two-concept class has a teaching plan of order zero.
\end{proof}
\noindent\textbf{Reference.} \cite{doliwa2014recursive}.

\subsection{Repetition-Free Recursive Teaching Dimension of Singletons}\label{val:repetitionfree_recursive_teaching_dimension_singletons}
\noindent\textbf{Statement.} The Repetition-Free Recursive Teaching Dimension of Singletons is $1$.
\begin{proof}
Remove the singleton concepts in any order.  The positive example at its unique coordinate teaches each singleton, and these one-coordinate teaching sets are pairwise distinct.  Thus this is a repetition-free recursive teaching plan of order one; order zero cannot identify a non-singleton class.
\end{proof}
\noindent\textbf{Reference.} \cite{doliwa2014recursive}.

\subsection{Repetition-Free Recursive Teaching Dimension of Singletons plus Empty Set}\label{val:repetitionfree_recursive_teaching_dimension_singletons_plus_empty_set}
\noindent\textbf{Statement.} The Repetition-Free Recursive Teaching Dimension of Singletons plus empty set is $1$.
\begin{proof}
Remove every singleton using its unique positive coordinate, in any order, and then remove the empty concept with the empty teaching set.  The singleton coordinate sets are pairwise distinct and differ from the final empty set.  This is a repetition-free plan of order one, which is optimal for a non-singleton class.
\end{proof}
\noindent\textbf{Reference.} \cite{doliwa2014recursive}.

\subsection{Repetition-Free Recursive Teaching Dimension of the Three-Code Class}\label{val:repetitionfree_recursive_teaching_dimension_three_code_class}
\noindent\textbf{Statement.} The Repetition-Free Recursive Teaching Dimension of Three-Code Class is $1$.
\begin{proof}
For the three concepts $000,011,101$, use respectively the one-coordinate teaching samples at coordinates $3,2,1$ (with their displayed labels).  Each sample identifies its target among all three concepts, and the coordinate sets $\{3\},\{2\},\{1\}$ are distinct.  Thus these removals give a repetition-free recursive plan of order one, which is optimal for a non-singleton class.
\end{proof}
\noindent\textbf{Reference.} \cite{doliwa2014recursive}.

\subsection{Repetition-Free Recursive Teaching Dimension of Half-intervals}\label{val:repetitionfree_recursive_teaching_dimension_halfintervals}
\noindent\textbf{Statement.} The Repetition-Free Recursive Teaching Dimension of Half-intervals is $1$.
\begin{proof}
Repeatedly remove an extreme remaining initial segment.  If the remaining endpoints form an interval $[a,b]$, remove its left endpoint using the negative label at coordinate $a+1$, or its right endpoint using the positive label at coordinate $b$.  Alternating sides uses each singleton coordinate set at most once; the final sole concept uses the empty set.  Hence this is a repetition-free plan of order one, and order zero is impossible for a non-singleton class.
\end{proof}
\noindent\textbf{Reference.} \cite{doliwa2014recursive}.

\subsection{Self-Directed Dimension of Trivial Concepts}\label{val:selfdirected_dimension_trivial_concepts}
\noindent\textbf{Statement.} The Self-Directed Dimension of Trivial concepts is $1$.
\begin{proof}
The established equality $\mathrm{SDdim}=\mathrm M_{\mathrm{sd}}$ transfers the recorded self-directed-queries-complexity value $1$ on trivial concepts.
\end{proof}
\noindent\textbf{Reference.} \cite{devulapalli2024selfdirected}; derived from the established self-directed-queries-complexity benchmark value.

\subsection{Self-Directed Dimension of Blockwise Addressing}\label{val:selfdirected_dimension_blockwiseaddressing}
\noindent\textbf{Statement.} The Self-Directed Dimension of Blockwise-Addressing is $\leq 2$.
\begin{proof}
The established equality $\mathrm{SDdim}=\mathrm M_{\mathrm{sd}}$ transfers the recorded upper bound $\mathrm M_{\mathrm{sd}}\leq2$ for the blockwise-addressing class.
\end{proof}
\noindent\textbf{Reference.} \cite{devulapalli2024selfdirected}; derived from the established self-directed-queries-complexity benchmark value.

\subsection{Self-Directed Dimension of Full Cube}\label{val:selfdirected_dimension_full_cube}
\noindent\textbf{Statement.} The Self-Directed Dimension of Full Cube is $n$.
\begin{proof}
The established equality $\mathrm{SDdim}=\mathrm M_{\mathrm{sd}}$ transfers the recorded self-directed-queries-complexity value $n$ on the full cube.
\end{proof}
\noindent\textbf{Reference.} \cite{devulapalli2024selfdirected}; derived from the established self-directed-queries-complexity benchmark value.

\subsection{Self-Directed Dimension of Singletons}\label{val:selfdirected_dimension_singletons}
\noindent\textbf{Statement.} The Self-Directed Dimension of Singletons is $1$.
\begin{proof}
The established equality $\mathrm{SDdim}=\mathrm M_{\mathrm{sd}}$ transfers the recorded self-directed-queries-complexity value $1$ on singletons.
\end{proof}
\noindent\textbf{Reference.} \cite{devulapalli2024selfdirected}; derived from the established self-directed-queries-complexity benchmark value.

\subsection{Self-Directed Dimension of Singletons plus Empty Set}\label{val:selfdirected_dimension_singletons_plus_empty_set}
\noindent\textbf{Statement.} The Self-Directed Dimension of Singletons plus empty set is $1$.
\begin{proof}
The established equality $\mathrm{SDdim}=\mathrm M_{\mathrm{sd}}$ transfers the recorded self-directed-queries-complexity value $1$ on singletons plus the empty set.
\end{proof}
\noindent\textbf{Reference.} \cite{devulapalli2024selfdirected}; derived from the established self-directed-queries-complexity benchmark value.

\subsection{Self-Directed Dimension of the Three-Code Class}\label{val:selfdirected_dimension_three_code_class}
\noindent\textbf{Statement.} The Self-Directed Dimension of Three-Code Class is $1$.
\begin{proof}
The established equality $\mathrm{SDdim}=\mathrm M_{\mathrm{sd}}$ transfers the recorded self-directed-queries-complexity value $1$ on the three-code class.
\end{proof}
\noindent\textbf{Reference.} \cite{devulapalli2024selfdirected}; derived from the established self-directed-queries-complexity benchmark value.

\subsection{Self-Directed Dimension of Halfintervals}\label{val:selfdirected_dimension_halfintervals}
\noindent\textbf{Statement.} The Self-Directed Dimension of Half-intervals is $1$.
\begin{proof}
The established equality $\mathrm{SDdim}=\mathrm M_{\mathrm{sd}}$ transfers the recorded self-directed-queries-complexity value $1$ on halfintervals.
\end{proof}
\noindent\textbf{Reference.} \cite{devulapalli2024selfdirected}; derived from the established self-directed-queries-complexity benchmark value.

\subsection{Self-Directed Dimension of Tagged Singletons}\label{val:selfdirected_dimension_tagged_singletons}
\noindent\textbf{Statement.} The Self-Directed Dimension of Tagged Singletons is $1$.
\begin{proof}
The established equality $\mathrm{SDdim}=\mathrm M_{\mathrm{sd}}$ transfers the recorded self-directed-queries-complexity value $1$ on tagged singletons.
\end{proof}
\noindent\textbf{Reference.} \cite{devulapalli2024selfdirected}; derived from the established self-directed-queries-complexity benchmark value.

\subsection{Monotonic co-VC Dimension of Trivial Concepts}\label{val:monotonic_covc_dimension_trivial_concepts}
\noindent\textbf{Statement.} The Monotonic co-VC Dimension of Trivial concepts is $2$.
\begin{proof}
The class itself has co-VC dimension two, giving the lower bound. Its Star Number is one, and the established inequality $\mathrm{coVC}^*\le\mathrm S+1$ gives the matching upper bound.
\end{proof}

\subsection{Monotonic co-VC Dimension of Singletons}\label{val:monotonic_covc_dimension_singletons}
\noindent\textbf{Statement.} The Monotonic co-VC Dimension of Singletons is $n$.
\begin{proof}
The singleton class itself has co-VC dimension $n$, so the envelope is at least $n$. Its effective range has $n$ coordinates, which upper-bounds every co-VC value of every subfamily.
\end{proof}

\subsection{Monotonic co-VC Dimension of Full Cube}\label{val:monotonic_covc_dimension_full_cube}
\noindent\textbf{Statement.} The Monotonic co-VC Dimension of Full Cube is $n$.
\begin{proof}
The full cube has Star Number $n$, so the star lower bound gives $\mathrm{coVC}^*\ge n$. Its effective range has $n$ coordinates, giving the matching upper bound. This contrasts with its ordinary co-VC dimension, which is zero.
\end{proof}

\subsection{Monotonic co-VC Dimension of Singletons plus Empty Set}\label{val:monotonic_covc_dimension_singletons_plus_empty_set}
\noindent\textbf{Statement.} The Monotonic co-VC Dimension of Singletons plus empty set is $n$.
\begin{proof}
The singleton subfamily has co-VC dimension $n$, giving the lower bound. The effective range has $n$ coordinates, so no subfamily can have larger co-VC dimension.
\end{proof}

\subsection{Projected Interpolation Degree of Trivial Concepts}\label{val:projected_interpolation_degree_trivial_concepts}
\noindent\textbf{Statement.} The Projected Interpolation Degree of Trivial concepts is $1$.
\begin{proof}
Every nonempty projection remains a two-point class of opposite constants and has interpolation degree one; the empty projection has degree zero. Hence the maximum is one.
\end{proof}

\subsection{Projected Interpolation Degree of Singletons}\label{val:projected_interpolation_degree_singletons}
\noindent\textbf{Statement.} The Projected Interpolation Degree of Singletons is $1$.
\begin{proof}
Every nonempty coordinate projection is the empty concept together with singleton concepts on the retained coordinates. Affine-linear polynomials interpolate arbitrary values on that class, while a nontrivial projection is not interpolated by constants.
\end{proof}

\subsection{Projected Interpolation Degree of Full Cube}\label{val:projected_interpolation_degree_full_cube}
\noindent\textbf{Statement.} The Projected Interpolation Degree of Full Cube is $n$.
\begin{proof}
The full-domain projection is the $n$-cube and has interpolation degree $n$. No projection has more than $n$ active coordinates, so degree $n$ is also an upper bound.
\end{proof}

\subsection{Projected Interpolation Degree of the Interpolation-Conjecture Counterexample}\label{val:projected_interpolation_degree_interpolation_conjecture_counterexample}
\noindent\textbf{Statement.} The Projected Interpolation Degree of Interpolation-Conjecture Counterexample is $2$.
\begin{proof}
Projecting $\{001,010,011,100\}$ onto its last two coordinates gives the full two-cube, so the value is at least two. The other two-coordinate projections have at most three distinct traces and interpolation degree one, while the original class and all one-coordinate projections have degree at most one. Thus the maximum is exactly two.
\end{proof}

\subsection{Monotonic co-VC Dimension of Warmuth's $C_5$}\label{val:monotonic_covc_dimension_warmuth_c5}
\noindent\textbf{Statement.} The Monotonic co-VC Dimension of Warmuth's $C_5$ class is $4$.
\begin{proof}
The class itself has co-VC dimension four, giving $\mathrm{coVC}^*\ge4$. Its Star Number is three, and the established affine inequality $\mathrm{coVC}^*\le\mathrm S+1$ gives the matching upper bound.
\end{proof}

\subsection{Effective Range of the Interpolation-Conjecture Counterexample}\label{val:effective_range_interpolation_conjecture_counterexample}
\noindent\textbf{Statement.} The Effective Range of Interpolation-Conjecture Counterexample is $3$.
\begin{proof}
In $\{001,010,011,100\}$, each of the three coordinates takes both labels: compare $001$ with $100$ for the first coordinate, $001$ with $010$ for the second, and $010$ with $011$ for the third.
\end{proof}

\subsection{Minimum Teaching-Set Size of the Minimum Teaching-Set Separation Class}\label{val:minimum_teaching_set_size_minimum_teaching_separation}
\noindent\textbf{Statement.} The Minimum Teaching Set Size of Minimum Teaching-Set Separation Class is $1$.
\begin{proof}
The exceptional concept $001$ is uniquely specified by its positive label on $x_3$, so the minimum teaching-set size is at most one. The class is non-singleton, so size zero is impossible.
\end{proof}
\noindent\textbf{Reference.} Direct calculation; verified by sync/verify\_monotonic\_minimum\_teaching.py..

\subsection{Monotonic Minimum Teaching-Set Size of the Minimum Teaching-Set Separation Class}\label{val:monotonic_minimum_teaching_set_size_minimum_teaching_separation}
\noindent\textbf{Statement.} The Monotonic Minimum Teaching Set Size of Minimum Teaching-Set Separation Class is $2$.
\begin{proof}
The four-concept subfamily 000, 100, 010, 110 is a two-cube, in which every concept needs both x1 and x2, giving the lower bound two. Every subfamily containing 001 has a one-point teacher for 001 at x3; every remaining subfamily is contained in that two-cube and has a teacher of size at most two. Hence the maximum over subfamilies is exactly two.
\end{proof}
\noindent\textbf{Reference.} Direct calculation; verified by sync/verify\_monotonic\_minimum\_teaching.py..

\subsection{Membership Query Complexity of Coded Singletons}\label{val:membership_query_complexity_coded_singletons}
\noindent\textbf{Statement.} The Membership Query Complexity of Coded Singletons is $\lceil\log_2(n+1)\rceil$.
\begin{proof}
Querying all r tag coordinates recovers the distinct code $c_i$ and hence identifies the target. Conversely, any deterministic membership-query tree of depth q has at most $2^q$ leaves, but there are n+1 targets; therefore q is at least ceil($log_2$(n+1)).
\end{proof}
\noindent\textbf{Reference.} Direct calculation; verified by sync/verify\_coded\_singletons.py..

\subsection{Projected Membership Query Complexity of Coded Singletons}\label{val:projected_membership_queries_complexity_coded_singletons}
\noindent\textbf{Statement.} The Projected Membership Queries Complexity of Coded Singletons is $n$.
\begin{proof}
Projecting away the tag coordinates gives the n singleton concepts together with the empty concept, whose membership-query complexity is n. Conversely, every projection has at most n+1 distinct traces, and any finite binary class with m traces can be identified by at most m-1 coordinate queries. Thus no projection needs more than n queries.
\end{proof}
\noindent\textbf{Reference.} Direct calculation; verified by sync/verify\_coded\_singletons.py..

\subsection{Degeneracy of the Middle Two Layers}\label{val:degeneracy_middle_layers}
\noindent\textbf{Statement.} The Degeneracy of Middle Two Layers is $m+1$.
\begin{proof}
Every m-set has 2m+1-m=m+1 one-inclusion neighbors in the upper layer, and every (m+1)-set has m+1 neighbors in the lower layer. Thus the full one-inclusion graph is (m+1)-regular, so its degeneracy is at least m+1. Its maximum degree is also m+1, so every induced subgraph has minimum degree at most m+1. Hence the degeneracy is exactly m+1.
\end{proof}
\noindent\textbf{Reference.} Direct calculation; verified by sync/verify\_middle\_layers.py..

\subsection{Largest Strongly Shattered Set of the Middle Two Layers}\label{val:largest_strongly_shattered_set_middle_layers}
\noindent\textbf{Statement.} The Largest Strongly Shattered Set of Middle Two Layers is $1$.
\begin{proof}
The class has one-inclusion edges, so it contains a one-cube. Every two-cube contains vertices in three consecutive cardinality layers. Since $ML_m$ uses only the two layers m and m+1, it contains no two-cube; hence its largest contained cube has dimension one.
\end{proof}
\noindent\textbf{Reference.} Direct calculation; verified by sync/verify\_middle\_layers.py..

\subsection{Littlestone Dimension of the Majority-Query Separation Class}\label{val:littlestone_dimension_majority_gap}
\noindent\textbf{Statement.} The Littlestone Dimension of Majority-Query Separation Class is $2$.
\begin{proof}
The finite Littlestone recurrence gives value two: each coordinate split has a branch of rank at most one, while a depth-two tree is obtained by first querying x1 and then a separating coordinate in either branch. The exhaustive recurrence certificate is in sync/verify\_majority\_gap.py.
\end{proof}
\noindent\textbf{Reference.} Direct calculation; verified by sync/verify\_majority\_gap.py..

\subsection{Majority Complexity of the Majority-Query Separation Class}\label{val:majority_complexity_majority_gap}
\noindent\textbf{Statement.} The Majority Complexity of Majority-Query Separation Class is $3$.
\begin{proof}
At each version space the prescribed query is its coordinatewise majority, with zero tie-breaking. The adversary chooses a counterexample coordinate whose minority branch maximizes the remaining recursive cost. Exhaustive evaluation of this recurrence gives exactly three; see sync/verify\_majority\_gap.py.
\end{proof}
\noindent\textbf{Reference.} Direct calculation; verified by sync/verify\_majority\_gap.py..

\subsection{Littlestone Dimension of Majority-Query Separation Products}\label{val:littlestone_dimension_majority_gap_products}
\noindent\textbf{Statement.} The Littlestone Dimension of Products of the Majority-Query Separation Class is $2t$.
\begin{proof}
Littlestone dimension is additive on Cartesian products over disjoint coordinate blocks: after a coordinate query in one factor, both child version spaces remain products, and the standard rank recurrence gives the sum of the factor ranks by induction. Each $C_maj$ factor has rank 2, so the t-fold product has rank 2t.
\end{proof}
\noindent\textbf{Reference.} Direct recurrence argument; checked for one and two factors by sync/verify\_majority\_gap\_products.py..

\subsection{Majority Complexity of Majority-Query Separation Products}\label{val:majority_complexity_majority_gap_products}
\noindent\textbf{Statement.} The Majority Complexity of Products of the Majority-Query Separation Class is $3t$.
\begin{proof}
Coordinatewise majority factors across disjoint Cartesian factors. After an adversarial counterexample in one block, the version space remains a product with only that factor restricted. The majority-query recurrence is therefore additive by induction. Each $C_maj$ factor costs 3, giving total cost 3t.
\end{proof}
\noindent\textbf{Reference.} Direct recurrence argument; checked for one and two factors by sync/verify\_majority\_gap\_products.py..

\subsection{Recursive Teaching Dimension of the Projected RTD Separation Class}\label{val:recursive_teaching_dimension_projected_rtd_gap}
\noindent\textbf{Statement.} The Recursive Teaching Dimension of Projected Recursive-Teaching Separation Class is $1$.
\begin{proof}
Remove 100 using its first coordinate, then 001 using its second coordinate. Of the two remaining concepts 010 and 011, remove either with the third coordinate and finally the last one with the empty set. Thus RTD is at most one; it cannot be zero for a non-singleton class.
\end{proof}
\noindent\textbf{Reference.} Direct calculation; verified by sync/verify\_projected\_rtd\_gap.py..

\subsection{Projected Recursive Teaching Dimension of the Projected RTD Separation Class}\label{val:monotone_recursive_teaching_dimension_projected_rtd_gap}
\noindent\textbf{Statement.} The Projected Recursive Teaching Dimension of Projected Recursive-Teaching Separation Class is $2$.
\begin{proof}
Projecting away the first coordinate produces all four strings on the remaining two coordinates, i.e. the full two-cube, whose recursive teaching dimension is two. No projection has more than three coordinates and the original class has RTD one; exhaustive evaluation shows the maximum is exactly two.
\end{proof}
\noindent\textbf{Reference.} Direct calculation; verified by sync/verify\_projected\_rtd\_gap.py..

\subsection{VC Dimension of the Rubinstein Minimum-Degree Class}\label{val:vc_dimension_rubinstein_minimum_degree}
\noindent\textbf{Statement.} The VC Dimension of Rubinstein Minimum-Degree Class is $d$.
\begin{proof}
Corollary 44 of Rubinstein--Bartlett--Rubinstein states that their class $V_d$ has VC dimension exactly d for every d at least 10.
\end{proof}
\noindent\textbf{Reference.} \cite{rubinstein2007shifting}, Corollary 44..

\subsection{Projected Minimum Degree of the Rubinstein Minimum-Degree Class}\label{val:projected_minimum_degree_rubinstein_minimum_degree}
\noindent\textbf{Statement.} The Projected Minimum Degree of Rubinstein Minimum-Degree Class is $\ge d+1$.
\begin{proof}
Corollary 44 gives ordinary one-inclusion minimum degree d+1 for $V_d$. The full-coordinate projection is among those maximized by projected minimum degree, so the projected parameter is at least d+1.
\end{proof}
\noindent\textbf{Reference.} \cite{rubinstein2007shifting}, Corollary 44..

\subsection{Littlestone Dimension of the Littlestone/Worst-Order Separation Class}\label{val:littlestone_dimension_littlestone_worst_gap}
\noindent\textbf{Statement.} The Littlestone Dimension of Littlestone/Worst-Order Separation Class is $3$.
\begin{proof}
The standard finite Littlestone rank recurrence evaluates to three. The exhaustive recurrence calculation is recorded in sync/verify\_littlestone\_worst\_gap.py.
\end{proof}
\noindent\textbf{Reference.} Direct calculation; verified by sync/verify\_littlestone\_worst\_gap.py..

\subsection{Worst Mistakes of the Littlestone/Worst-Order Separation Class}\label{val:worst_mistakes_littlestone_worst_gap}
\noindent\textbf{Statement.} The Worst Mistakes of Littlestone/Worst-Order Separation Class is $2$.
\begin{proof}
For each of the 24 coordinate orders, solve the transductive minimax prediction recurrence from the final coordinate backwards. Every order has value two, and no smaller uniform bound is possible. Repeated coordinates do not help an adversary because their labels are known after first occurrence. See sync/verify\_littlestone\_worst\_gap.py.
\end{proof}
\noindent\textbf{Reference.} Direct calculation; verified by sync/verify\_littlestone\_worst\_gap.py..

\subsection{Best Mistakes of the Best/Worst Mistakes Separation Class}\label{val:best_mistakes_best_worst_mistakes_gap}
\noindent\textbf{Statement.} The Best Mistakes of Best/Worst Mistakes Separation Class is $1$.
\begin{proof}
The favourable orders $x_2$,$x_1$,$x_3$ and $x_3$,$x_1$,$x_2$ have transductive minimax value one. Zero mistakes are impossible for a nonconstant class. The exhaustive calculation over all six coordinate orders is recorded in sync/verify\_best\_worst\_mistakes\_gap.py.
\end{proof}
\noindent\textbf{Reference.} Direct calculation; verified by sync/verify\_best\_worst\_mistakes\_gap.py..

\subsection{Worst Mistakes of the Best/Worst Mistakes Separation Class}\label{val:worst_mistakes_best_worst_mistakes_gap}
\noindent\textbf{Statement.} The Worst Mistakes of Best/Worst Mistakes Separation Class is $2$.
\begin{proof}
The orders $x_1$,$x_2$,$x_3$ and $x_1$,$x_3$,$x_2$ each have transductive minimax value two, and no order has larger value. Repeated coordinates do not help after their labels have been revealed. See sync/verify\_best\_worst\_mistakes\_gap.py.
\end{proof}
\noindent\textbf{Reference.} Direct calculation; verified by sync/verify\_best\_worst\_mistakes\_gap.py..

\subsection{Minimum Largest Order-Shattered Set of the Interpolation/Effective-VC-Radius Counterexample}\label{val:minimum_largest_order_shattered_set_interpolation_effective_vc_radius_counterexample}
\noindent\textbf{Statement.} The Minimum Largest Order Shattered Set of Interpolation/Effective-VC-Radius Counterexample is $2$.
\begin{proof}
For each of the six coordinate orders, apply the three binary downshifts and inspect the resulting down-closed family. Its largest face has size two, so the minimum over orders is two. The exhaustive finite calculation is recorded in sync/verify\_interpolation\_conjecture.py.
\end{proof}
\noindent\textbf{Reference.} Direct calculation; verified by sync/verify\_interpolation\_conjecture.py..

\subsection{Maximum Positive Degree of the Positive-Degree/Strong-Shattering Separation Class}\label{val:maximum_positive_degree_positive_degree_gap}
\noindent\textbf{Statement.} The Maximum Positive Degree of Positive-Degree/Strong-Shattering Separation Class is $3$.
\begin{proof}
The top vertex 111 has the three one-inclusion neighbours 110, 101, and 011, all smaller under inclusion. Thus its positive degree is three, which is maximal on a three-coordinate domain. This is checked in sync/verify\_positive\_degree\_gap.py.
\end{proof}
\noindent\textbf{Reference.} Direct calculation; verified by sync/verify\_positive\_degree\_gap.py..

\subsection{Largest Strongly Shattered Set of the Positive-Degree/Strong-Shattering Separation Class}\label{val:largest_strongly_shattered_set_positive_degree_gap}
\noindent\textbf{Statement.} The Largest Strongly Shattered Set of Positive-Degree/Strong-Shattering Separation Class is $1$.
\begin{proof}
Every coordinate is strongly shattered by an edge from 111 to the corresponding two-element set. No two-dimensional cube is contained in the four-concept class: its members have Hamming weights two or three, whereas every two-cube has two vertices whose weights differ by two. Hence the largest strongly shattered set has size one. See sync/verify\_positive\_degree\_gap.py.
\end{proof}
\noindent\textbf{Reference.} Direct calculation; verified by sync/verify\_positive\_degree\_gap.py..

\subsection{Best Mistakes of the Best-Mistakes/Order-Shattering Separation Class}\label{val:best_mistakes_best_mistakes_order_shattering_gap}
\noindent\textbf{Statement.} The Best Mistakes of Best-Mistakes/Order-Shattering Separation Class is $2$.
\begin{proof}
Solve the transductive minimax prediction recurrence from the last coordinate backwards for every one of the 24 coordinate orders. Every order has value at least two, and some order has value two. Repeated coordinates cannot help after their label has been seen. The exact computation is in sync/verify\_best\_mistakes\_order\_shattering\_gap.py.
\end{proof}
\noindent\textbf{Reference.} Direct calculation; verified by sync/verify\_best\_mistakes\_order\_shattering\_gap.py..

\subsection{Minimum Largest Order-Shattered Set of the Best-Mistakes/Order-Shattering Separation Class}\label{val:minimum_largest_order_shattered_set_best_mistakes_order_shattering_gap}
\noindent\textbf{Statement.} The Minimum Largest Order Shattered Set of Best-Mistakes/Order-Shattering Separation Class is $1$.
\begin{proof}
For example, applying the four binary downshifts in the order $x_1$,$x_2$,$x_3$,$x_4$ produces a down-closed family with largest support size one. No nonconstant class has order-shattering dimension zero. Thus the minimum over coordinate orders is one. The exhaustive calculation is in sync/verify\_best\_mistakes\_order\_shattering\_gap.py.
\end{proof}
\noindent\textbf{Reference.} Direct calculation; verified by sync/verify\_best\_mistakes\_order\_shattering\_gap.py..

\subsection{Minimum Projected Teaching Set Size of the Projected-Teaching Monotonicity Separation Class}\label{val:minimum_projected_teaching_set_size_projected_teaching_monotonicity_gap}
\noindent\textbf{Statement.} The Minimum Projected Teaching Set Size of Projected-Teaching Monotonicity Separation Class is $2$.
\begin{proof}
For each target and each of the 16 coordinate projections, compute its minimum distinguishing set in the distinct trace class. The minimum, over targets, of the resulting worst projected teaching size is two. The exact calculation is in sync/verify\_projected\_teaching\_values.py.
\end{proof}
\noindent\textbf{Reference.} Direct calculation; verified by sync/verify\_projected\_teaching\_values.py..

\subsection{Monotonic Projected Minimum Teaching Set Size of the Projected-Teaching Monotonicity Separation Class}\label{val:monotonic_projected_minimum_teaching_set_size_projected_teaching_monotonicity_gap}
\noindent\textbf{Statement.} The Monotonic Projected Minimum Teaching Set Size of Projected-Teaching Monotonicity Separation Class is $3$.
\begin{proof}
The four-concept subfamily $\{1000,0100,0010,0001\}$ has projected minimum teaching-set size three. Exhaustive enumeration of all 31 nonempty subfamilies shows none has larger value. Therefore the subclass maximum is three; see sync/verify\_projected\_teaching\_values.py.
\end{proof}
\noindent\textbf{Reference.} Direct calculation; verified by sync/verify\_projected\_teaching\_values.py..

\subsection{Minimum Largest Order-Shattered Set of Addressing}\label{val:minimum_largest_order_shattered_set_addressing}
\noindent\textbf{Statement.} The Minimum Largest Order Shattered Set of Addressing is $1\quad(n\ge2)$.
\begin{proof}
Order the bit coordinates before the address coordinates. Each bit downshift replaces every concept's address code by zero, since no other concept has its address coordinate. The subsequent address downshifts leave an empty concept and address singletons, whose largest support has size one. Thus $OSH_min$ is at most one; it is at least one because the addressing class is nonconstant. See sync/verify\_addressing\_order\_shattering.py.
\end{proof}
\noindent\textbf{Reference.} Direct calculation; verified by sync/verify\_addressing\_order\_shattering.py..

\subsection{Positive No-Clashing Teaching Dimension of the Positive-Teaching Triangle}\label{val:positive_noclashing_teaching_dimension_positive_teaching_triangle}
\noindent\textbf{Statement.} The Positive No-Clashing Teaching Dimension of Positive-Teaching Triangle is $1$.
\begin{proof}
Assign the one-coordinate positive samples $x_1$ to 110, $x_3$ to 101, and $x_2$ to 011. For every pair, at least one assigned sample is negative on the other target, so no pair clashes. Width zero is impossible in a non-singleton class. See sync/verify\_positive\_teaching\_triangle.py.
\end{proof}
\noindent\textbf{Reference.} Direct calculation; verified by sync/verify\_positive\_teaching\_triangle.py..
